% Généré depuis les contenus pédagogiques du manuel complet.
\documentclass[12pt,twoside,openany]{book}
\usepackage{optitab-livre}
\usepackage{pc-livre}
\usepackage{adjustbox}
% Adaptation collège de la charte éditoriale utilisée en Terminale.
\renewcommand{\otNiveau}{Cinquième · Collège}
\renewcommand{\otNbChapitres}{20}
\renewcommand{\otNbExosChapitres}{300}
\renewcommand{\otNbDefis}{20}
\renewcommand{\otNbExosTotal}{320}
\renewcommand{\otSousTitre}{Cours et exercices corrigés}
\renewcommand{\otCleardoublepage}{\clearpage}
\setlength{\parskip}{2.2mm}
\setlength{\headheight}{15pt}
\usepackage{xurl}
\urlstyle{same}
\linespread{1.05}
\tcbset{otEncadre/.append style={fontupper=\fontsize{11.5}{15}\selectfont,
 before skip=3mm,after skip=3mm,top=2.8mm,bottom=2.8mm}}
\newenvironment{pcEncadre}[2]{%
 \def\pcCol{otNavy}\def\pcFond{otNavyTint}%
 \ifstrequal{#1}{method}{\def\pcCol{otGreenDark}\def\pcFond{otGreenTint}}{}%
 \ifstrequal{#1}{attention}{\def\pcCol{otOrange!80!black}\def\pcFond{otOrange!10}}{}%
 \begin{tcolorbox}[otEncadre,colback=\pcFond,borderline west={1.1mm}{0pt}{\pcCol}]
 \otEtiq{\pcCol}{#2}\par\vspace{1mm}}{\end{tcolorbox}}

\newcommand{\pcPage}[2]{%
 \clearpage\markright{#1}\otPageTitre{#1}{#2}\vspace{-2mm}}
\newcommand{\pcSousTitre}[1]{%
 \par\addvspace{3.5mm}\noindent{\otSemi\fontsize{12}{15}\selectfont\color{otNavy}#1}\par\nopagebreak\vspace{1mm}}
\newcommand{\pcRetour}[2]{%
 \par{\otMedium\fontsize{8.5}{11}\selectfont\color{otGreenDark}\hyperref[#1]{#2\enspace·\enspace p.~\pageref*{#1}}}\par}
\newcommand{\pcFigure}[2]{%
 \par\vspace{2mm}\noindent\begin{minipage}{\linewidth}\centering
 \adjincludegraphics[max width=\linewidth]{#1}\par
 \ifstrempty{#2}{}{\vspace{1.3mm}{\otLight\fontsize{8.7}{11.5}\selectfont\color{otMuted}\raggedright #2\par}}
 \end{minipage}\par\vspace{1.5mm}}
% La largeur maximale est prise en charge par adjustbox, sans déformer l'image.

\newsavebox{\pcFicheBox}
\newenvironment{pcFiche}[1]{%
 \def\pcFicheId{#1}\par\noindent
 \begin{lrbox}{\pcFicheBox}\begin{minipage}{\linewidth}
}{%
 \end{minipage}\end{lrbox}%
 \dimen0=\dimexpr\pagegoal-\pagetotal-4mm\relax
 \ifdim\dimen0>\textheight\dimen0=\textheight\fi
 \typeout{PCFICHE|\pcFicheId|\the\dimexpr\ht\pcFicheBox+\dp\pcFicheBox\relax|\the\dimen0}%
 \ifdim\dimexpr\ht\pcFicheBox+\dp\pcFicheBox\relax>\dimen0
   \PackageError{cinquieme}{PCFICHE-TROP|\pcFicheId}{Repartir le contenu sans reduire le texte.}%
 \else\noindent\usebox{\pcFicheBox}\fi\par}

\newcommand{\pcSomPartie}[2]{%
 \par\addvspace{3.2mm}{\otSemi\fontsize{8}{10}\selectfont\color{otGreenDark}\otCaps[12]{Partie #1}}\enspace
 {\otSemi\fontsize{10}{12}\selectfont\color{otNavy}#2}\par\vspace{1.8mm}}
\newcommand{\pcSomLigne}[3]{%
 \noindent\makebox[9mm][l]{\otSemi\fontsize{9}{11}\selectfont\color{otNavy}#1}%
 \hyperref[#3]{{\otRegular\fontsize{9.4}{12}\selectfont #2}%
 \otPoints{\otSemi\fontsize{9}{11}\selectfont\color{otNavy}\pageref*{#3}}}\par\vspace{0.5mm}}

\newcommand{\pcMission}[3]{%
 \phantomsection\label{mission-#1-#2}%
 \otEntete{otNavy}{#1.#2}{#3}{0}{\otRenvoi{Corrigé}{cor-mission-#1-#2}}}
\newcommand{\pcSolutionMission}[3]{%
 \phantomsection\label{cor-mission-#1-#2}%
 \otEntete{otGreen}{#1.#2}{#3}{0}{\otRenvoi{Énoncé}{mission-#1-#2}}}
\newenvironment{pcLexique}[2]{%
 \par\addvspace{3.8mm}\Needspace{5\baselineskip}\noindent
 {\otSemi\fontsize{11}{14}\selectfont\color{otNavy}#1}\hfill
 {\otMedium\fontsize{8}{11}\selectfont\color{otGreenDark}\hyperref[chap-#2]{Chapitre #2\enspace·\enspace p.~\pageref*{chap-#2}}}\par\nopagebreak\vspace{1mm}
 \fontsize{11.5}{15}\selectfont}{\par}

% Les ouvertures conservent le bandeau bleu, le grand numéro, les deux
% colonnes d'objectifs et de repères, ainsi que les onglets de Terminale.
\renewcommand{\otQRCours}{%
 \begin{tikzpicture}\node[fill=otNavyTint,rounded corners=2mm,inner sep=3.5mm,text width=48mm]
 {{\otSemi\fontsize{8.5}{11}\selectfont\color{otNavy}Mon carnet de sciences}\par\vspace{1mm}
  {\otLight\fontsize{8}{11}\selectfont Une question, un schéma légendé et une conclusion à garder.}};\end{tikzpicture}}
\renewcommand{\otLegende}{%
 {\otLight\fontsize{7.8}{10}\selectfont\color{otMuted}Difficulté\hspace{3mm}%
 \otDots{1}\enspace débuter\hspace{4mm}\otDots{2}\enspace s’entraîner\hspace{4mm}\otDots{3}\enspace raisonner}}
\renewcommand{\otDansCe}{%
 \pcOuvertureLigne{L’essentiel}{ess-\otChapNum}%
 \pcOuvertureLigne{Je démarre}{rub-\otChapNum-1}%
 \pcOuvertureLigne{Je m’entraîne}{rub-\otChapNum-2}%
 \pcOuvertureLigne{Je raisonne}{rub-\otChapNum-3}%
 \pcOuvertureLigne{Défi}{defi-ex-\otChapNum}%
 {\color{otRule}\rule{\linewidth}{0.4pt}}\par\vspace{1mm}%
 \pcOuvertureLigne{Corrigés}{corchap-\otChapNum}}
\newcommand{\pcOuvertureLigne}[2]{%
 \noindent{\otSemi\fontsize{9}{12}\selectfont\color{otNavy}#1}\hfill
 {\otMedium\fontsize{8}{11}\selectfont\hyperref[#2]{p.~\pageref*{#2}}}\par\vspace{2mm}}

% Décors propres aux quatre thèmes du collège.
\expandafter\newcommand\csname otDecocollege1\endcsname{\otDecoGrille
 \foreach \x/\y in {151/70,160/72,169/69,154/80,164/82,174/79,157/90,167/93,178/90}{%
 \fill[white,opacity=0.26] (\x mm,\y mm) circle (1.3mm);}
 \draw[white,opacity=0.4,line width=0.8pt] (145mm,99mm) -- (145mm,61mm) -- (185mm,61mm) -- (185mm,99mm);}
\expandafter\newcommand\csname otDecocollege2\endcsname{\otDecoGrille
 \draw[white,opacity=0.3,line width=1pt] (138mm,70mm) .. controls (159mm,103mm) and (172mm,54mm) .. (196mm,84mm);
 \foreach \x/\y in {138/70,151/80,164/81,177/78,190/79}{\fill[otGreen] (\x mm,\y mm) circle (1.8mm);}}
\expandafter\newcommand\csname otDecocollege3\endcsname{\otDecoGrille
 \draw[white,opacity=0.4,line width=1pt] (141mm,67mm) -- (141mm,91mm) -- (189mm,91mm) -- (189mm,67mm) -- cycle;
 \draw[white,opacity=0.6,line width=1pt] (141mm,77mm) -- (141mm,81mm);
 \fill[otNavy] (162mm,86mm) rectangle (172mm,96mm);
 \draw[otGreen,line width=1pt] (167mm,91mm) circle (4.5mm);
 \draw[otGreen,line width=0.7pt] (164mm,88mm) -- (170mm,94mm) (164mm,94mm) -- (170mm,88mm);}
\expandafter\newcommand\csname otDecocollege4\endcsname{\otDecoGrille
 \draw[white,opacity=0.45,line width=1pt] plot[domain=0:6,samples=110] ({135mm+\x*11mm},{79mm+sin(\x*3 r)*9mm});}

\hypersetup{pdftitle={Physique-Chimie - Cinquième - Manuel complet},
 pdfsubject={Cours illustrés, activités, exercices progressifs et corrigés - Cycle 4 - 2026-2027},pdfauthor={OptiTAB}}
\begin{document}
\pagenumbering{alph}
%% =====================================================================
%%  1re de couverture
%% =====================================================================
\thispagestyle{empty}
\begin{tikzpicture}[remember picture,overlay]
\begin{scope}[shift={(current page.south west)}]
  % Fond
  \fill[otNavy] (0,0) rectangle (210mm,297mm);
  \shade[bottom color=otNavyDark,top color=otNavy] (0,0) rectangle (210mm,160mm);
  \draw[white,opacity=0.06,line width=0.3pt,step=7mm] (0,0) grid (210mm,297mm);
  % Ondes circulaires issues de deux sources (interférences)
  \begin{scope}
    \clip (0,0) rectangle (210mm,297mm);
    \foreach \r in {5,10,...,40}{%
      \draw[white,opacity=0.08,line width=0.6pt] (183mm,270mm) circle (\r mm);
      \draw[white,opacity=0.08,line width=0.6pt] (205mm,270mm) circle (\r mm);}
    \fill[white,opacity=0.5] (183mm,270mm) circle (1.1mm);
    \fill[white,opacity=0.5] (205mm,270mm) circle (1.1mm);
  \end{scope}
  % Trajectoire lumineuse et mesures de distance : décor de collège.
  \draw[white,opacity=0.2,line width=0.6pt] (105mm,18mm) -- (105mm,90mm) -- (210mm,90mm);
  \foreach \x in {117,137,157,177,197}{
    \draw[white,opacity=0.14] (\x mm,18mm) -- (\x mm,90mm);}
  \draw[otGreen,line width=2.2pt,line cap=round] (105mm,20mm) -- (205mm,83mm);
  \foreach \x/\y in {105/20,125/32.6,145/45.2,165/57.8,185/70.4,205/83}{
    \fill[white] (\x mm,\y mm) circle (1.5mm);}
  % Logo
  \node[anchor=north west,inner sep=0pt] at (22mm,279mm)
    {\includegraphics[width=60mm]{logo-blanc}};

  % Titre
  \node[anchor=base west,inner sep=0pt,text=otGreen,
    font=\otSemi\fontsize{11}{13}\selectfont] at (24mm,236mm) {\otCaps[22]{\otNiveau}};
  \node[anchor=base west,inner sep=0pt,text=white,
    font=\otBold\fontsize{50}{54}\selectfont] at (21.5mm,212mm) {\otTitreLivre};
  \node[anchor=base west,inner sep=0pt,text=white,
    font=\otLight\fontsize{32}{36}\selectfont] at (23mm,194mm) {\otSousTitre};
  \fill[otGreen] (24mm,181mm) rectangle (54mm,182.4mm);
  \node[anchor=base west,inner sep=0pt,text=white,opacity=0.85,
    font=\otRegular\fontsize{11}{14}\selectfont] at (24mm,168mm)
    {Année complète\enspace·\enspace Cours illustré\enspace·\enspace Exercices progressifs\enspace·\enspace Corrigés détaillés};

  % Points forts
  \node[anchor=north west,inner sep=0pt,text width=100mm,text=white,
    font=\otRegular\fontsize{10.5}{14}\selectfont] at (24mm,146mm)
    {\otCheck\hspace{3mm}Exercices classés par difficulté\par\vspace{3.2mm}
     \otCheck\hspace{3mm}Schémas clairs et activités guidées\par\vspace{3.2mm}
     \otCheck\hspace{3mm}\otNbDefis\ défis avec indications progressives\par\vspace{3.2mm}
     \otCheck\hspace{3mm}Quatre bilans pour faire le point};

  % Pastille
  \fill[otGreen] (150mm,120mm) circle (18mm);
  \node[text=white,align=center,inner sep=0pt] at (150mm,120mm)
    {{\otMedium\fontsize{7}{9}\selectfont Année complète}\\[0.6mm]
     {\otBold\fontsize{30}{30}\selectfont \otNbExosTotal}\\[1.4mm]
     {\otSemi\fontsize{7.5}{9}\selectfont\otCaps[10]{exercices et}}\\
     {\otSemi\fontsize{7.5}{9}\selectfont\otCaps[10]{défis corrigés}}};

  % Pied
  \node[anchor=base west,inner sep=0pt,text=white,opacity=0.75,
    font=\otMedium\fontsize{10}{12}\selectfont] at (24mm,16mm) {\otSite};
\end{scope}
\end{tikzpicture}
\clearpage

\pagenumbering{arabic}
\pagestyle{otFront}
\pcPage{Sommaire}{Cinquième · Physique-chimie · Programme français de cycle 4}
\phantomsection\label{sommaire}
\pdfbookmark[0]{Sommaire}{bm-sommaire}
\begingroup\setlength{\parskip}{0pt}
\pcSomPartie{1}{Organisation et transformations de la matière}
\pcSomLigne{01}{La matière et ses états}{chap-1}
\pcSomLigne{02}{Mesurer une masse et un volume}{chap-2}
\pcSomLigne{03}{Les changements d’état de l’eau}{chap-3}
\pcSomLigne{04}{Corps purs, mélanges et composition de l’air}{chap-4}
\pcSomLigne{05}{Dissolution et miscibilité}{chap-5}
\pcSomLigne{06}{Séparer les constituants d’un mélange}{chap-6}
\pcSomLigne{07}{Reconnaître une transformation chimique}{chap-7}
\pcSomLigne{08}{La Terre dans le Système solaire}{chap-8}
\pcSomPartie{2}{Mouvement et interactions}
\pcSomLigne{09}{Décrire un mouvement}{chap-9}
\pcSomLigne{10}{Mesurer et calculer une vitesse}{chap-10}
\pcSomLigne{11}{Décrire des actions et des interactions}{chap-11}
\pcSomPartie{3}{L’énergie, ses transferts et ses conversions}
\pcSomLigne{12}{Sources, formes et conversions d’énergie}{chap-12}
\pcSomLigne{13}{Réaliser et schématiser un circuit}{chap-13}
\pcSomLigne{14}{Conducteurs, isolants et sécurité électrique}{chap-14}
\pcSomLigne{15}{Circuits en série et en dérivation}{chap-15}
\pcSomLigne{16}{Énergie et choix responsables}{chap-16}
\pcSomPartie{4}{Des signaux pour observer et communiquer}
\pcSomLigne{17}{Sources de lumière et vision}{chap-17}
\pcSomLigne{18}{Propagation de la lumière et ombres}{chap-18}
\pcSomLigne{19}{Produire et transmettre un son}{chap-19}
\pcSomLigne{20}{Un signal pour transmettre une information}{chap-20}
\vspace{3mm}\pcSomLigne{B}{Bilans des quatre thèmes}{bilans-debut}
\pcSomLigne{C}{Corrigés détaillés}{corriges-debut}
\pcSomLigne{i}{Méthodes, unités, lexique et références}{annexes-debut}
\endgroup
\pcPage{Mode d’emploi}{Une démarche pour comprendre, puis s’entraîner}
\phantomsection\label{mode}
\pdfbookmark[0]{Mode d’emploi}{bm-mode}
\begin{pcEncadre}{definition}{1 · Observer}
Décrire les faits, préciser les conditions et poser une question.\par
\end{pcEncadre}
\begin{pcEncadre}{method}{2 · Chercher}
Formuler une hypothèse et proposer un test qui pourrait la mettre à l’épreuve.\par
\end{pcEncadre}
\begin{pcEncadre}{definition}{3 · Mesurer}
Choisir l’instrument, noter les résultats avec leurs unités et contrôler le protocole.\par
\end{pcEncadre}
\begin{pcEncadre}{method}{4 · Expliquer}
Relier les observations aux connaissances et rédiger une conclusion argumentée.\par
\end{pcEncadre}
\pcSousTitre{Trois niveaux pour progresser}
\noindent\otDots{1}\quad\textbf{Je démarre}\enspace Reconnaître, nommer et appliquer.\par\vspace{2mm}
\noindent\otDots{2}\quad\textbf{Je m’entraîne}\enspace Lire un schéma, mesurer, calculer et expliquer.\par\vspace{2mm}
\noindent\otDots{3}\quad\textbf{Je raisonne}\enspace Croiser des données, critiquer une méthode et argumenter.\par\vspace{2mm}
\pcPage{Parcours et ressources}{Retrouver une méthode et comprendre ses erreurs}
\phantomsection\label{ressources}
\pdfbookmark[0]{Parcours et ressources}{bm-ressources}
\otLigneRessource{securite}{Travailler en sécurité}
\otLigneRessource{corriges-debut}{Corrigés des exercices}
\otLigneRessource{bilans-debut}{Bilans des quatre thèmes}
\otLigneRessource{methode-mesure}{Mesurer et contrôler}
\otLigneRessource{methode-graphique}{Construire et lire un graphique}
\otLigneRessource{methode-calcul}{Calculer avec les unités}
\otLigneRessource{unites}{Grandeurs et conversions}
\otLigneRessource{lexique}{Lexique scientifique}
\otLigneRessource{programme}{Programme et progression}
\otLigneRessource{references}{Références}
\begin{pcEncadre}{method}{Une méthode qui fonctionne}
Cherche d’abord seul, écris tes essais et explique ta réponse. Utilise le corrigé pour comprendre l’erreur, puis refais la question sans regarder la solution.\par
\end{pcEncadre}
\begin{pcEncadre}{definition}{Le cadre du manuel}
Programme français du cycle 4 : BO du 30 juillet 2020, référence confirmée pour 2026-2027. La progression de 5e est un choix pédagogique. Manuel OptiTAB indépendant, sans validation institutionnelle revendiquée.\par
\end{pcEncadre}
\pcPage{Travailler avec soin}{Démarche expérimentale}
\phantomsection\label{securite}
\pdfbookmark[0]{Travailler avec soin}{bm-securite}
\begin{pcEncadre}{definition}{AVANT DE COMMENCER}
Lire la consigne entière, identifier la question et faire vérifier le montage par le professeur. Les manipulations de chimie, de chauffage et d’électricité sont réalisées en classe selon ses consignes.\par
\end{pcEncadre}
\begin{questions}
\item Ne jamais goûter un échantillon ni sentir directement un produit. Lire les étiquettes et utiliser les protections demandées.
\item Porter les lunettes et autres protections prévues par le protocole. Organiser une paillasse dégagée et attacher les cheveux si nécessaire.
\item Pour l’électricité : utiliser exclusivement le matériel scolaire de très basse tension ; ne jamais manipuler une prise du secteur.
\item Ne pas regarder le Soleil, un laser ou une source intense. Les activités de lumière utilisent des lampes adaptées.
\item Limiter le niveau sonore ; s’éloigner des sources fortes et interrompre une exposition gênante.
\item Trier les déchets selon la consigne ; demander avant de jeter une solution. Se laver les mains après les manipulations.
\end{questions}
\begin{pcEncadre}{method}{UNE TRACE SCIENTIFIQUE}
Question \(\rightarrow\) hypothèse \(\rightarrow\) protocole \(\rightarrow\) observations et mesures \(\rightarrow\) interprétation \(\rightarrow\) conclusion. Une conclusion peut rester limitée : le préciser fait partie du travail scientifique.\par
\end{pcEncadre}
\clearpage\pagestyle{otMain}
\begin{ouverture}[college1]{1}{La matière et ses états}{Partie 1 · Organisation et transformations de la matière}{Un verre « vide » est-il vraiment vide ?}
\objectif{Distinguer solide, liquide et gaz à partir de propriétés observables.}
\objectif{Expliquer que l’air occupe un volume et peut être comprimé.}
\objectif{Utiliser un modèle simple et distinguer observation et interprétation.}
\end{ouverture}
\pcPage{L’essentiel}{La matière et ses états}
\phantomsection\label{ch-1-etats}
\phantomsection\label{ess-1}
\begin{pcFiche}{ch-1-etats}
\essSection{1}{Décrire avant de classer}
Une même matière peut exister sous différents états physiques. On précise les conditions de l’étude, notamment la température. Pour comparer, on observe la forme, le volume occupé et l’effet d’une compression douce.\par
\pcFigure{figures/42d1614f4ae878e0.pdf}{}
\begin{center}
{\renewcommand{\arraystretch}{1.35}\fontsize{11}{14}\selectfont
\setlength{\tabcolsep}{2mm}\begin{tabular}{@{}>{\raggedright\arraybackslash}p{42.80mm}>{\raggedright\arraybackslash}p{34.77mm}>{\raggedright\arraybackslash}p{34.77mm}>{\raggedright\arraybackslash}p{35.66mm}@{}}\toprule
\otTete{Propriété} & \otTete{Solide compact} & \otTete{Liquide} & \otTete{Gaz}\\
\midrule
Forme propre & Oui & Non & Non\\
Volume propre & Oui & Oui & Non\\
Compression douce & Très faible effet & Très faible effet & Volume réduit\\
\bottomrule\end{tabular}}\end{center}
\begin{pcEncadre}{attention}{POUR BIEN COMPARER}
On étudie ici un solide compact, un liquide et un gaz dans les conditions usuelles. Les déformations, les changements de température et les fortes pressions peuvent modifier certaines propriétés.\par
\end{pcEncadre}
\end{pcFiche}
\pcPage{Un liquide change de forme}{La matière et ses états}
\phantomsection\label{ch-1-liquide}
\begin{pcFiche}{ch-1-liquide}
\essSection{2}{Transvaser ne fait pas disparaître la matière}
Sans perte de liquide et sans changement de température, transvaser une même quantité d’eau modifie sa forme et sa hauteur dans le récipient, mais ne modifie pas son volume.\par
\pcFigure{figures/34126ecab693bc6e.pdf}{Le récipient large donne une hauteur d’eau plus faible, pour la même quantité d’eau.}
\pcSousTitre{La surface libre d’un liquide au repos}
\pcFigure{figures/1b86ec55cb3a425a.pdf}{}
Loin des parois, une surface libre suffisamment étendue est plane et horizontale. Elle reste horizontale quand le récipient immobile est incliné.\par
Près des parois, la surface peut se courber : nous apprendrons à lire ce ménisque au chapitre 2.\par
\end{pcFiche}
\pcPage{L’air est une matière}{La matière et ses états}
\phantomsection\label{ch-1-gaz}
\begin{pcFiche}{ch-1-gaz}
\essSection{3}{Un gaz occupe l’espace accessible}
L’air est invisible, mais il occupe une place et possède une masse. Dans un récipient fermé, il se répartit dans l’espace accessible. Son volume peut diminuer lorsqu’on le comprime, ou augmenter si on lui offre davantage d’espace.\par
\pcFigure{figures/41411f5a625ff530.pdf}{Une seringue fermée permet d’observer la compressibilité de l’air. La quantité d’air reste la même si rien ne fuit.}
\begin{pcEncadre}{definition}{UNE OBSERVATION, UNE INTERPRÉTATION}
Observation : le piston avance et le volume occupé diminue.\par Interprétation : l’air est compressible. Il ne disparaît pas.\par
\end{pcEncadre}
\begin{pcEncadre}{attention}{NE PAS CONFONDRE VAPEUR ET BROUILLARD}
La vapeur d’eau est un gaz invisible. Le brouillard et les nuages visibles contiennent de fines gouttelettes d’eau liquide, parfois aussi des cristaux de glace. Une « fumée blanche » ne prouve donc pas que l’on voit un gaz.\par
\end{pcEncadre}
\end{pcFiche}
\pcPage{Le verre protège-t-il le papier ?}{La matière et ses états}
\phantomsection\label{ch-1-activite}
\begin{pcFiche}{ch-1-activite}
\pcSousTitre{Question : que contient un verre que l’on pense vide ?}
Matériel : bac d’eau, gobelet transparent rigide et papier. Activité en classe avec le professeur. Utiliser un récipient adapté et ne pas forcer.\par
\pcFigure{figures/0c6231a1d9e60a7e.pdf}{Gobelet entièrement sous l’eau, ouverture vers le bas. Garder le gobelet vertical pendant la première observation. Schéma non à l’échelle.}
\begin{questions}
\item \textbf{Avant :} formule une hypothèse sur l’état du papier après immersion.
\item \textbf{Tester :} coince le papier au fond, retourne le gobelet, puis descends-le verticalement. Fais valider le montage.
\item \textbf{Observer :} relève l’état du papier. Incline ensuite le gobelet et décris les bulles ainsi que l’entrée de l’eau.
\end{questions}
\begin{pcEncadre}{method}{POUR CONCLURE}
Écris deux observations et une explication. Si tu ne réalises pas la manipulation, exploite le document : papier sec au départ ; des bulles s’échappent et l’eau entre quand on incline le gobelet.\par
\end{pcEncadre}
\end{pcFiche}
\pcPage{Un premier modèle de la matière}{La matière et ses états}
\phantomsection\label{ch-1-modele}
\begin{pcFiche}{ch-1-modele}
\begin{pcEncadre}{definition}{MODÉLISER}
Un modèle est une représentation simplifiée. Ici, chaque point représente une particule. Il aide à relier l’organisation de la matière à ses propriétés, mais ce n’est pas une photographie.\par
\end{pcEncadre}
\pcFigure{figures/281e473284f774b5.pdf}{Même nombre de points et même taille des points dans les trois cases. Le modèle dessiné pour le solide est un modèle ordonné simplifié.}
\begin{questions}
\item \textbf{Solide :} les particules restent proches et vibrent autour de positions fixes dans ce modèle. Cela explique une forme propre.
\item \textbf{Liquide :} les particules sont proches et peuvent se déplacer les unes par rapport aux autres.
\item \textbf{Gaz :} les particules sont espacées et se déplacent dans tout l’espace accessible. Les espaces entre elles peuvent diminuer lors d’une compression.
\end{questions}
Les points ne sont ni des gouttes ni des grains visibles. Ils ne grossissent pas simplement parce que le volume du gaz augmente.\par
\end{pcFiche}
\pcPage{L’essentiel à retenir}{La matière et ses états}
\phantomsection\label{ch-1-bilan}
\begin{pcFiche}{ch-1-bilan}
\begin{pcEncadre}{definition}{LES QUATRE IDÉES CLÉS}
Un solide compact possède une forme et un volume propres.\par Un liquide possède un volume propre et prend la forme du récipient.\par Un gaz occupe l’espace accessible et peut être comprimé.\par Invisible ne signifie pas absent : l’air est une matière.\par
\end{pcEncadre}
\essSection{\(\rightarrow\)}{Expliquer avec une preuve}
\begin{questions}
\item Je nomme \textbf{ce qui a été observé}.
\item Je cite la \textbf{propriété} que cette observation permet d’utiliser.
\item Je relie les deux avec une phrase : « J’observe…, donc… ».
\end{questions}
\begin{pcEncadre}{method}{UN EXEMPLE DE RÉDACTION}
« Dans la seringue fermée, le volume de l’air diminue lorsque le piston avance. Comme l’air ne s’échappe pas, cette observation montre qu’un gaz est compressible. »\par
\end{pcEncadre}
Avant les exercices : sais-tu expliquer pourquoi du sable qui coule peut rester solide et pourquoi on ne voit pas la vapeur d’eau ?\par
\end{pcFiche}
\clearpage\markright{Exercices}
\begin{blocLivre}{exo-1-1}
\rubrique{Je démarre}
\begin{exercice}[1]{Associer une matière à son état}
\begin{questions}
\item Indique l’état physique de chaque matière à température ambiante : acier d’une règle, lait, air, huile.
\item Indique l’état de l’eau dans un glaçon et celui de la vapeur d’eau présente dans l’air.
\end{questions}
\end{exercice}\end{blocLivre}
\begin{blocLivre}{exo-1-2}
\begin{exercice}[1]{Les bons mots}
\begin{questions}
\item Complète : un solide compact conserve sa … et son … lorsqu’on le change de récipient.
\item Complète : un liquide possède un … propre, mais il prend la … du récipient.
\item Complète : un gaz occupe tout l’espace qui lui est … .
\end{questions}
\end{exercice}\end{blocLivre}
\begin{blocLivre}{exo-1-3}
\begin{exercice}[1]{Dessiner une surface libre}
\pcFigure{figures/eba7680d12a4107b.pdf}{Les deux récipients sont immobiles. Dessine dans ton cahier.}
\begin{questions}
\item Recopie les deux récipients et dessine la surface libre de l’eau, au repos et loin des parois.
\item Ajoute une légende qui explique l’orientation de cette surface.
\end{questions}
\end{exercice}\end{blocLivre}
\begin{blocLivre}{exo-1-4}
\begin{exercice}[1]{De l’air dans une seringue}
\begin{donnees}La seringue ne contient que de l’air. Son embout est fermé de façon étanche.\par
\end{donnees}
\pcFigure{figures/423ab0f5bf636d6b.pdf}{}
\begin{questions}
\item Décris ce qui change lorsqu’on pousse doucement le piston.
\item L’air s’est-il échappé ? Quelle propriété du gaz cette observation met-elle en évidence ?
\end{questions}
\end{exercice}\end{blocLivre}
\begin{blocLivre}{exo-1-5}
\begin{exercice}[1]{Vrai ou faux ? Justifier}
\begin{questions}
\item « Un gaz est toujours visible. »
\item « La vapeur d’eau est constituée de petites gouttes liquides. »
\item « Un liquide au repos peut avoir une surface libre horizontale. »
\end{questions}
\end{exercice}\end{blocLivre}
\begin{blocLivre}{exo-1-6}
\rubrique{Je m’entraîne}
\begin{exercice}[2]{Peut-on reconnaître un liquide à sa forme ?}
\pcFigure{figures/b742bf2c6f9c8dd4.pdf}{}
\begin{questions}
\item Le sable et l’eau prennent tous les deux la forme du récipient. Ce critère suffit-il pour les classer comme liquides ?
\item Observe un grain de sable : quelle propriété permet de préciser son état ?
\end{questions}
\end{exercice}\end{blocLivre}
\begin{blocLivre}{exo-1-7}
\begin{exercice}[2]{Le papier reste sec}
\begin{donnees}Observe de nouveau le schéma de l’activité « Le verre protège-t-il le papier ? ».\par
\end{donnees}
\pcFigure{figures/992a247d3beea186.pdf}{Document de travail · modèle simplifié.}
\begin{questions}
\item Nomme la matière présente entre le papier et l’eau à l’intérieur du verre.
\item Explique pourquoi l’eau n’atteint pas immédiatement le papier.
\item Que peut-on observer si l’on incline le verre sous l’eau ?
\end{questions}
\end{exercice}\end{blocLivre}
\begin{blocLivre}{exo-1-8}
\begin{exercice}[2]{Changer de récipient}
\begin{questions}
\item On transvase 120 mL d’eau d’un récipient étroit dans un récipient large, sans perte. Quel est le volume final ?
\item Explique pourquoi la hauteur d’eau peut diminuer sans que le volume diminue.
\end{questions}
\end{exercice}\end{blocLivre}
\begin{blocLivre}{exo-1-9}
\begin{exercice}[2]{Observation ou interprétation ?}
\begin{questions}
\item Classe ces phrases : « Des bulles quittent le verre » ; « L’air occupait une partie du verre » ; « Le papier est sec ».
\item Rédige une courte conclusion à partir de ces observations.
\end{questions}
\end{exercice}\end{blocLivre}
\begin{blocLivre}{exo-1-10}
\begin{exercice}[2]{Lire un modèle}
\pcFigure{figures/7b263cdf44f47936.pdf}{}
\begin{questions}
\item Quel état correspond à des particules proches les unes des autres et capables de se déplacer les unes par rapport aux autres ?
\item Dans le modèle du gaz, les particules sont-elles plus grosses ou plus espacées ?
\item Ce schéma est-il une photographie de la matière ?
\end{questions}
\end{exercice}\end{blocLivre}
\begin{blocLivre}{exo-1-11}
\begin{exercice}[2]{Tester une idée correctement}
\begin{questions}
\item Lina veut comparer la compressibilité de l’air et de l’eau. Elle utilise une seringue ouverte pour l’air et une seringue fermée pour l’eau. Son test est-il convaincant ?
\item Propose les conditions à garder identiques et une précaution pour la seringue contenant de l’eau.
\end{questions}
\end{exercice}\end{blocLivre}
\begin{blocLivre}{exo-1-12}
\rubrique{Je raisonne}
\begin{exercice}[3]{L’air possède-t-il une masse ?}
\begin{donnees}Un dispositif de laboratoire permet de peser le même récipient rigide fermé. Avec très peu d’air : 238,4 g. Après admission d’air, sans autre modification : 239,6 g. Données pédagogiques.\par
\end{donnees}
\begin{questions}
\item Calcule l’augmentation de masse du récipient.
\item Que permettent de conclure ces mesures à propos de l’air ?
\end{questions}
\end{exercice}\end{blocLivre}
\begin{blocLivre}{exo-1-13}
\begin{exercice}[3]{Une légende scientifique}
\begin{questions}
\item Une image de brouillard porte la légende : « On voit la vapeur d’eau dans l’air ». Explique ce qui pose problème.
\item Propose une légende scientifiquement correcte et compréhensible.
\end{questions}
\end{exercice}\end{blocLivre}
\begin{blocLivre}{exo-1-14}
\begin{exercice}[3]{Le sable met la règle à l’épreuve}
\begin{questions}
\item Noé affirme : « Tout ce qui s’écoule est liquide. » Utilise un contre-exemple pour discuter son affirmation.
\item Écris une description qui distingue le comportement de l’ensemble et l’état de chaque grain.
\end{questions}
\end{exercice}\end{blocLivre}
\begin{blocLivre}{exo-1-15}
\begin{exercice}[3]{Identifier trois échantillons}
\begin{donnees}A : forme conservée lors d’un transfert, volume conservé. B : forme modifiée lors d’un transfert, volume conservé, très peu compressible. C : tout l’espace accessible est occupé ; une compression réduit son volume.\par
\end{donnees}
\begin{questions}
\item Associe A, B et C à un état physique en citant à chaque fois deux propriétés.
\item Explique pourquoi une seule propriété peut parfois être insuffisante.
\end{questions}
\end{exercice}\end{blocLivre}
\pcPage{Défi · Prendre une initiative}{Transférer de l’air sous l’eau}
\phantomsection\label{defi-ex-1}
\begin{pcFiche}{defi-1}
Deux gobelets transparents rigides sont retournés dans un bac d’eau. A contient de l’air ; B est entièrement rempli d’eau. Comment transférer une partie de l’air de A vers B, sans les sortir de l’eau ?\par
\pcFigure{figures/d349fd055e5e8aa3.pdf}{Les ouvertures des gobelets sont dirigées vers le bas. Situation schématisée.}
\begin{questions}
\item Propose une disposition et un geste. Fais un schéma légendé.
\item Prévois le trajet de l’air et ce qui arrive à l’eau dans B.
\item Justifie ta proposition avec une propriété étudiée dans le chapitre.
\end{questions}
Indice progressif : les bulles d’air remontent dans l’eau. Le test éventuel se réalise avec le professeur et du matériel adapté.\par
\end{pcFiche}
\clearpage\pagestyle{otMain}
\begin{ouverture}[college1]{2}{Mesurer une masse et un volume}{Partie 1 · Organisation et transformations de la matière}{Un même volume implique-t-il une même masse ?}
\objectif{Mesurer une masse et lire un volume avec les instruments adaptés.}
\objectif{Utiliser la tare, les conversions et le déplacement d’eau.}
\objectif{Distinguer les deux grandeurs et exploiter une proportionnalité simple.}
\end{ouverture}
\pcPage{L’essentiel}{Mesurer une masse et un volume}
\phantomsection\label{ch-2-masse}
\phantomsection\label{ess-2}
\begin{pcFiche}{ch-2-masse}
\begin{pcEncadre}{definition}{GRANDEUR ET UNITÉ}
La masse, notée m, est une grandeur que l’on mesure avec une balance. L’unité du Système international est le kilogramme (kg). Au laboratoire, on utilise souvent le gramme (g).\par
\end{pcEncadre}
\pcSousTitre{Peser un liquide avec la tare}
\pcFigure{figures/b39a2e119fd58ff3.pdf}{« Vide » signifie ici sans liquide : le récipient contient de l’air. Après la tare, garder le même récipient ; l’affichage donne la masse de liquide ajoutée.}
\begin{pcEncadre}{method}{SI L’ON N’A PAS FAIT LA TARE}
Masse du liquide = masse du récipient plein - masse du récipient vide.\par Exemple : 123,1 g - 48,6 g = 74,5 g.\par
\end{pcEncadre}
Poser la balance à plat et attendre que l’affichage se stabilise. « Masse » et « poids » désignent des grandeurs différentes.\par
\end{pcFiche}
\pcPage{Lire un volume au bon endroit}{Mesurer une masse et un volume}
\phantomsection\label{ch-2-volume}
\begin{pcFiche}{ch-2-volume}
Le volume, noté V, mesure l’espace occupé. Son unité SI est le mètre cube (m\textsuperscript{3}). Pour un liquide au laboratoire, on utilise souvent une éprouvette graduée en millilitres (mL).\par
\pcFigure{figures/116d73ab8d2f43be.pdf}{Pour l’eau dans une éprouvette usuelle, lire le bas du ménisque. Ici V = 36 mL.}
\begin{questions}
\item Poser l’éprouvette sur un support \textbf{horizontal}.
\item Repérer l’\textbf{unité}, puis compter les intervalles entre deux valeurs connues.
\item Placer l’\textbf{œil à la hauteur} du bas du ménisque et noter valeur + unité.
\end{questions}
Une division vaut ici 10 mL \(\div\) 5 = 2 mL. Compter les intervalles, pas le nombre de traits.\par
\end{pcFiche}
\pcPage{Convertir en gardant la grandeur}{Mesurer une masse et un volume}
\phantomsection\label{ch-2-unites}
\begin{pcFiche}{ch-2-unites}
\begin{center}
{\renewcommand{\arraystretch}{1.35}\fontsize{11}{14}\selectfont
\setlength{\tabcolsep}{2mm}\begin{tabular}{@{}>{\raggedright\arraybackslash}p{78.00mm}>{\raggedright\arraybackslash}p{78.00mm}@{}}\toprule
\otTete{Masse} & \otTete{Volume}\\
\midrule
1 kg = 1 000 g & 1 L = 1 000 mL\\
1 g = 1 000 mg & 1 L = 1 dm\textsuperscript{3}\\
0,250 kg = 250 g & 1 mL = 1 cm\textsuperscript{3}\\
\bottomrule\end{tabular}}\end{center}
\pcFigure{figures/0b02dc097ca50c7e.pdf}{Le cm\textsuperscript{3} est le volume d’un cube d’arête 1 cm. Le dm\textsuperscript{3} est le volume d’un cube d’arête 1 dm. Représentation non à l’échelle.}
\begin{pcEncadre}{definition}{UNE MÉTHODE SIMPLE}
Repérer la grandeur, choisir l’équivalence puis calculer.\par 0,35 L = 0,35 \(\times\) 1 000 mL = 350 mL.\par 750 g = 750 \(\div\) 1 000 kg = 0,750 kg.\par
\end{pcEncadre}
\end{pcFiche}
\pcPage{Mesurer un solide par déplacement d’eau}{Mesurer une masse et un volume}
\phantomsection\label{ch-2-deplacement}
\begin{pcFiche}{ch-2-deplacement}
Une éprouvette permet aussi de mesurer le volume d’un solide de forme irrégulière. On compare le volume lu avant et après immersion complète.\par
\pcFigure{figures/78910d2a01de2bf4.pdf}{Le solide est entièrement dans l’eau. Les graduations sont identiques sur les deux éprouvettes.}
\begin{pcEncadre}{method}{LIRE, SOUSTRAIRE ET EXPRIMER}
Volume du solide = volume final - volume initial.\par V = 58 mL - 40 mL = 18 mL = 18 cm\textsuperscript{3}.\par
\end{pcEncadre}
\begin{pcEncadre}{attention}{LES CONDITIONS DU PROTOCOLE}
Le solide ne doit pas se dissoudre, absorber l’eau ni réagir avec elle. Éviter les bulles. L’objet doit être entièrement immergé et l’éprouvette assez grande pour ne pas déborder.\par
\end{pcEncadre}
\end{pcFiche}
\pcPage{Masse et volume : deux grandeurs}{Mesurer une masse et un volume}
\phantomsection\label{ch-2-comparer}
\begin{pcFiche}{ch-2-comparer}
À volume égal, deux matières différentes peuvent avoir des masses différentes. On ne convertit donc pas des mL en g sans une information sur la matière étudiée.\par
\pcFigure{figures/38a6414a2b00e488.pdf}{Résultats d’une expérience donnée : 20 mL d’eau, environ 20 g ; 20 mL de cette huile, 18 g. Données pédagogiques arrondies.}
\begin{center}
{\renewcommand{\arraystretch}{1.35}\fontsize{11}{14}\selectfont
\setlength{\tabcolsep}{2mm}\begin{tabular}{@{}>{\raggedright\arraybackslash}p{46.84mm}>{\raggedright\arraybackslash}p{24.29mm}>{\raggedright\arraybackslash}p{24.29mm}>{\raggedright\arraybackslash}p{24.29mm}>{\raggedright\arraybackslash}p{24.29mm}@{}}\toprule
\otTete{Volume d’eau} & \otTete{20 mL} & \otTete{40 mL} & \otTete{60 mL} & \otTete{80 mL}\\
\midrule
Masse approximative & 20 g & 40 g & 60 g & 80 g\\
\bottomrule\end{tabular}}\end{center}
\begin{pcEncadre}{method}{POUR UNE MÊME SUBSTANCE}
À température fixée, la masse est proportionnelle au volume. Pour l’eau usuelle à température ambiante, 1 mL a une masse d’environ 1 g. Ce repère n’est ni une égalité entre unités ni une règle pour tous les liquides.\par
\end{pcEncadre}
\end{pcFiche}
\pcPage{Préparer et vérifier une mesure}{Mesurer une masse et un volume}
\phantomsection\label{ch-2-activite}
\begin{pcFiche}{ch-2-activite}
\pcSousTitre{Mission : préparer 50 mL d’eau et déterminer leur masse.}
\begin{questions}
\item \textbf{Choisir :} éprouvette adaptée, balance, récipient sec et eau. Vérifier la capacité et la graduation.
\item \textbf{Mesurer :} préparer 50 mL d’eau en lisant à hauteur du ménisque.
\item \textbf{Peser :} faire la tare avec le récipient sec, transférer l’eau sans perte puis relever la masse.
\end{questions}
\begin{center}
{\renewcommand{\arraystretch}{1.35}\fontsize{11}{14}\selectfont
\setlength{\tabcolsep}{2mm}\begin{tabular}{@{}>{\raggedright\arraybackslash}p{24.72mm}>{\raggedright\arraybackslash}p{63.18mm}>{\raggedright\arraybackslash}p{64.10mm}@{}}\toprule
\otTete{Essai} & \otTete{Volume préparé} & \otTete{Masse mesurée}\\
\midrule
1 & ........................ & ........................\\
2 & ........................ & ........................\\
3 & ........................ & ........................\\
\bottomrule\end{tabular}}\end{center}
\begin{pcEncadre}{definition}{ANALYSER LES RÉSULTATS}
Compare les essais. Une masse voisine de 50 g est attendue pour 50 mL d’eau à température ambiante. Si l’écart est important, vérifie la tare, la lecture et les pertes au transfert avant de conclure.\par
\end{pcEncadre}
Activité en classe avec le professeur. Relever les résultats réellement obtenus, même s’ils diffèrent de la valeur attendue.\par
\end{pcFiche}
\pcPage{La méthode en quatre questions}{Mesurer une masse et un volume}
\phantomsection\label{ch-2-bilan}
\begin{pcFiche}{ch-2-bilan}
\begin{pcEncadre}{definition}{1 · Quelle grandeur ?}
Masse m ou volume V ? Je ne mélange pas les deux.\par
\end{pcEncadre}
\begin{pcEncadre}{definition}{2 · Quel instrument et quelle unité ?}
Balance et g ou kg ; éprouvette et mL ou L. Un résultat comporte une valeur et une unité.\par
\end{pcEncadre}
\begin{pcEncadre}{definition}{3 · Quel protocole ?}
Tare ou différence pour une masse ; bas du ménisque pour un liquide ; différence avant/après pour un solide immergé.\par
\end{pcEncadre}
\begin{pcEncadre}{definition}{4 · Mon résultat est-il cohérent ?}
Je contrôle l’unité, l’ordre de grandeur et les conditions de mesure. Une différence de volumes doit rester positive dans l’expérience d’immersion.\par
\end{pcEncadre}
La physique-chimie demande autant de soin dans le choix de la méthode que dans le calcul.\par
\end{pcFiche}
\clearpage\markright{Exercices}
\begin{blocLivre}{exo-2-1}
\rubrique{Je démarre}
\begin{exercice}[1]{Choisir la grandeur et l’instrument}
\begin{questions}
\item Pour connaître la masse d’une trousse, choisis : balance ou éprouvette graduée ? Propose une unité.
\item Pour connaître le volume d’un peu d’eau, choisis un instrument et une unité adaptés.
\item Que mesure une valeur exprimée en cm\textsuperscript{3} ?
\end{questions}
\end{exercice}\end{blocLivre}
\begin{blocLivre}{exo-2-2}
\begin{exercice}[1]{Lire une pesée}
\begin{questions}
\item Après la tare avec le récipient vide, la balance affiche 37,6 g quand on ajoute le liquide. Quelle est la masse du liquide ?
\item Pourquoi ne faut-il pas ajouter à nouveau la masse du récipient ?
\end{questions}
\end{exercice}\end{blocLivre}
\begin{blocLivre}{exo-2-3}
\begin{exercice}[1]{Lire une éprouvette}
\pcFigure{figures/4f957faa1b510e2e.pdf}{}
\begin{questions}
\item Détermine le volume correspondant à une petite division.
\item Lis le volume au bas du ménisque et écris le résultat avec son unité.
\end{questions}
\end{exercice}\end{blocLivre}
\begin{blocLivre}{exo-2-4}
\begin{exercice}[1]{Convertir sans confondre}
\begin{questions}
\item Convertis : 1,25 kg en g ; 750 g en kg.
\item Convertis : 2,4 L en mL ; 350 mL en L.
\item Complète : 18 mL = … cm\textsuperscript{3} ; 3 L = … dm\textsuperscript{3}.
\end{questions}
\end{exercice}\end{blocLivre}
\begin{blocLivre}{exo-2-5}
\begin{exercice}[1]{Corriger l’unité}
\begin{questions}
\item Lina lit 35 mL dans une éprouvette et écrit : « La masse du liquide est 35 g ». Explique son erreur.
\item Écris ce que cette mesure permet réellement de connaître.
\end{questions}
\end{exercice}\end{blocLivre}
\begin{blocLivre}{exo-2-6}
\rubrique{Je m’entraîne}
\begin{exercice}[2]{Peser sans utiliser la tare}
\begin{donnees}Récipient vide : 62,8 g. Même récipient contenant un liquide : 187,3 g.\par
\end{donnees}
\begin{questions}
\item Calcule la masse du liquide.
\item Explique l’opération à l’aide des grandeurs qui interviennent.
\end{questions}
\end{exercice}\end{blocLivre}
\begin{blocLivre}{exo-2-7}
\begin{exercice}[2]{Deux méthodes, un même résultat}
\begin{questions}
\item Méthode A : récipient vide 45,2 g ; récipient et liquide 118,7 g. Détermine la masse du liquide.
\item Méthode B : après tare avec le même récipient vide, on verse le même liquide. Quelle valeur devrait s’afficher ?
\end{questions}
\end{exercice}\end{blocLivre}
\begin{blocLivre}{exo-2-8}
\begin{exercice}[2]{Mesurer un petit solide}
\begin{donnees}Solide entièrement immergé, insoluble et non absorbant. Aucune bulle n’est accrochée.\par
\end{donnees}
\pcFigure{figures/2c98165e80ece330.pdf}{}
\begin{questions}
\item Lis le volume avant et après immersion.
\item Détermine le volume du solide et exprime-le en cm\textsuperscript{3}.
\end{questions}
\end{exercice}\end{blocLivre}
\begin{blocLivre}{exo-2-9}
\begin{exercice}[2]{Un repère pour l’eau}
\begin{questions}
\item À température ambiante, quelle masse approximative correspond à 250 mL d’eau ?
\item Une huile de même volume a-t-elle nécessairement cette masse ? Justifie.
\end{questions}
\end{exercice}\end{blocLivre}
\begin{blocLivre}{exo-2-10}
\begin{exercice}[2]{Exploiter une proportionnalité}
\begin{donnees}Eau à température ambiante, données arrondies : 20 mL \(\rightarrow\) 20 g ; 40 mL \(\rightarrow\) 40 g.\par
\end{donnees}
\begin{questions}
\item Complète les masses correspondant à 60 mL et 80 mL.
\item Explique comment tu passes de 20 mL à 80 mL.
\item Quelle condition faut-il conserver pour comparer ces mesures ?
\end{questions}
\end{exercice}\end{blocLivre}
\begin{blocLivre}{exo-2-11}
\begin{exercice}[2]{Choisir une verrerie adaptée}
\begin{questions}
\item Pour mesurer environ 38 mL, choisis entre une éprouvette de 50 mL graduée tous les 1 mL et une éprouvette de 1 L graduée tous les 20 mL.
\item Justifie ton choix avec la capacité et la graduation.
\end{questions}
\end{exercice}\end{blocLivre}
\begin{blocLivre}{exo-2-12}
\rubrique{Je raisonne}
\begin{exercice}[3]{Comparer sans mélanger les grandeurs}
\begin{donnees}Échantillon A : volume 45 mL, masse 40 g. Échantillon B, constitué d’une autre matière : volume 30 mL, masse 60 g.\par
\end{donnees}
\begin{questions}
\item Quel échantillon possède le plus grand volume ? Lequel possède la plus grande masse ?
\item Explique pourquoi « le plus volumineux est toujours le plus lourd » n’est pas une conclusion générale.
\end{questions}
\end{exercice}\end{blocLivre}
\begin{blocLivre}{exo-2-13}
\begin{exercice}[3]{Un protocole à discuter}
\pcFigure{figures/0d2f1afa8f5842f8.pdf}{}
\begin{questions}
\item Le morceau de bois flotte et dépasse de l’eau. La différence des deux lectures donne-t-elle son volume total ?
\item Indique une condition qui n’est pas respectée et propose une façon de poursuivre l’étude.
\end{questions}
\end{exercice}\end{blocLivre}
\begin{blocLivre}{exo-2-14}
\begin{exercice}[3]{Une identification prudente}
\begin{donnees}Pour un liquide de référence, dans les conditions de l’étude : 25 mL ont une masse de 22,5 g. Pour l’échantillon inconnu : 50 mL ont une masse de 45 g.\par
\end{donnees}
\begin{questions}
\item Un liquide inconnu a une masse de 45 g pour 50 mL. Est-ce compatible avec le liquide de référence ?
\item Ces seules mesures permettent-elles de prouver l’identité du liquide ?
\end{questions}
\end{exercice}\end{blocLivre}
\begin{blocLivre}{exo-2-15}
\begin{exercice}[3]{Réunir deux mesures indépendantes}
\begin{donnees}Coupelle vide : 91,4 g ; avec le solide : 125,6 g. L’expérience ci-dessous concerne ce même solide, entièrement immergé, non absorbant, sans bulle et sans réaction avec l’eau.\par
\end{donnees}
\pcFigure{figures/9f2e051f0fa80f66.pdf}{}
\begin{questions}
\item Calcule la masse du solide.
\item Calcule son volume en mL, puis en cm\textsuperscript{3}.
\item L’affirmation « sa masse en g est forcément égale à son volume en mL » est-elle vérifiée ? Rédige une conclusion.
\end{questions}
\end{exercice}\end{blocLivre}
\pcPage{Défi · Prendre une initiative}{Une bille trop petite pour être lue seule}
\phantomsection\label{defi-ex-2}
\begin{pcFiche}{defi-2}
Le volume d’une seule bille est trop petit pour être facilement lu avec l’éprouvette disponible. On dispose de vingt billes identiques et d’une balance.\par
\pcFigure{figures/1484120210509d15.pdf}{}
\begin{pcEncadre}{definition}{DONNÉES À EXPLOITER}
Avec les 20 billes entièrement immergées, le volume passe de 40 mL à 50 mL. La masse totale des 20 billes est de 78,0 g. On suppose les billes identiques, insolubles et sans bulle accrochée.\par
\end{pcEncadre}
\begin{questions}
\item Explique pourquoi mesurer plusieurs billes est utile.
\item Détermine le volume moyen d’une bille, en cm\textsuperscript{3}.
\item Détermine sa masse moyenne et explique les limites de la mesure.
\end{questions}
\end{pcFiche}
\clearpage\pagestyle{otMain}
\begin{ouverture}[college1]{3}{Les changements d’état de l’eau}{Partie 1 · Organisation et transformations de la matière}{Pourquoi la glace fond-elle sans que l’eau disparaisse ?}
\objectif{Nommer les changements d’état et distinguer état et changement d’état.}
\objectif{Mesurer une température et lire une courbe de changement d’état.}
\objectif{Distinguer conservation de la masse et variation du volume.}
\end{ouverture}
\pcPage{L’essentiel}{Chapitre 03 · Les changements d’état de l’eau}
\phantomsection\label{ch-3-cours-1}
\phantomsection\label{ess-3}
\begin{pcFiche}{cours-3-1}
\essSection{1}{Passer d’un état à un autre}
Un changement d’état est une transformation physique : la substance reste la même. La fusion transforme un solide en liquide ; la solidification réalise le passage inverse. La vaporisation transforme un liquide en gaz ; la liquéfaction transforme un gaz en liquide. L’évaporation se produit à la surface du liquide, même sans ébullition. L’ébullition se produit dans tout le liquide avec formation de bulles de vapeur.\par
\pcFigure{figures/3832df9e2dc735b0.pdf}{Les flèches indiquent les changements, et non le déplacement d’une matière entre deux récipients.}
\begin{pcEncadre}{definition}{À retenir}
Un glaçon qui fond et de l’eau qui bout restent constitués d’eau. Chauffer ne provoque donc pas toujours une transformation chimique.\par
\end{pcEncadre}
\begin{pcEncadre}{method}{Une méthode concrète}
Décrire d’abord l’état initial et l’état final. Écrire ensuite le nom du changement : eau liquide \(\rightarrow\) eau gazeuse : vaporisation.\par
\end{pcEncadre}
\end{pcFiche}
\pcPage{L’essentiel}{Chapitre 03 · Les changements d’état de l’eau}
\phantomsection\label{ch-3-cours-2}
\begin{pcFiche}{cours-3-2}
\essSection{2}{Température et changement d’état}
La température se mesure avec un thermomètre et s’exprime ici en degré Celsius (°C). Lorsqu’on chauffe de la glace d’eau pure sous la pression atmosphérique usuelle, la fusion se déroule autour de 0 °C. Tant que glace et eau liquide coexistent, l’énergie reçue sert au changement d’état et la température reste voisine de 0 °C. Ce palier est un repère expérimental, avec les limites de précision des mesures.\par
\pcFigure{figures/d8e8069526b50879.pdf}{Exemple pédagogique de chauffage de glace d’eau pure ; pression atmosphérique usuelle.}
\begin{pcEncadre}{definition}{À retenir}
Un palier correspond à plusieurs mesures successives de température presque égale. Un simple point ne constitue pas un palier. L’eau pure bout vers 100 °C à la pression atmosphérique normale.\par
\end{pcEncadre}
\begin{pcEncadre}{method}{Une méthode concrète}
Lire le titre des axes, leurs unités et leurs graduations. Décrire séparément les portions qui montent et la portion horizontale.\par
\end{pcEncadre}
\end{pcFiche}
\pcPage{L’essentiel}{Chapitre 03 · Les changements d’état de l’eau}
\phantomsection\label{ch-3-cours-3}
\begin{pcFiche}{cours-3-3}
\essSection{3}{Même masse, volume parfois différent}
Lors d’un changement d’état, la masse se conserve si aucune matière n’entre ni ne sort du système étudié. Le volume peut changer. Pour l’eau, la glace occupe davantage de place que la même masse d’eau liquide. Une bouteille pleine et fermée peut donc être endommagée par le gel. Le modèle particulaire conserve les mêmes particules, tout en modifiant leur organisation et leurs mouvements.\par
\pcFigure{figures/cac35b4bcd88d707.pdf}{La masse affichée inclut le récipient fermé. Schéma qualitatif : le volume augmente lors du gel de l’eau.}
\begin{pcEncadre}{definition}{À retenir}
Conserver la masse ne signifie pas conserver le volume. Si de la vapeur s’échappe d’un récipient ouvert, la masse de ce qui reste dans le récipient diminue.\par
\end{pcEncadre}
\begin{pcEncadre}{method}{Une méthode concrète}
Préciser les limites du système pesé : eau seule, ou eau et récipient. Comparer les mesures sur le même ensemble et dans les mêmes conditions.\par
\end{pcEncadre}
\end{pcFiche}
\pcPage{Activité expérimentale}{Suivre la fusion de la glace}
\phantomsection\label{ch-3-activite}
\begin{pcFiche}{activite-3}
\begin{pcEncadre}{method}{Mission}
Exploiter des mesures de température pour repérer la fusion de l’eau pure.\par
\end{pcEncadre}
\begin{questions}
\item En classe : placer un thermomètre adapté dans un mélange de glace pilée et d’un peu d’eau. Éviter le contact du capteur avec la paroi.
\item Relever la température à intervalles réguliers pendant un réchauffement doux choisi par le professeur. Noter si de la glace reste visible.
\item Tracer la courbe avec les valeurs réellement mesurées. Pour s’entraîner, utiliser la série pédagogique du schéma.
\end{questions}
\pcFigure{figures/aac75681dcc7e81c.pdf}{Données d’entraînement : temps en min ; température en °C. Le réchauffement et la durée du palier dépendent du montage.}
\begin{questions}
\item Entre quelles dates observe-t-on un palier ? Quel changement d’état est compatible avec ce palier ?
\item Pourquoi faut-il aussi observer le contenu du récipient ?
\item Rédige une conclusion qui associe les mesures et l’observation de la glace.
\end{questions}
\end{pcFiche}
\pcPage{L’essentiel à retenir}{Chapitre 03 · Synthèse et autoévaluation}
\phantomsection\label{ch-3-bilan}
\begin{pcFiche}{synthese-3}
\begin{pcEncadre}{definition}{Repère 1}
Un changement d’état conserve la nature de la substance.\par
\end{pcEncadre}
\begin{pcEncadre}{method}{Repère 2}
Fusion, solidification, vaporisation et liquéfaction décrivent des passages précis.\par
\end{pcEncadre}
\begin{pcEncadre}{definition}{Repère 3}
Une température se mesure et se donne avec son unité.\par
\end{pcEncadre}
\begin{pcEncadre}{method}{Repère 4}
Un corps pur peut présenter un palier lors d’un changement d’état à pression fixée.\par
\end{pcEncadre}
\begin{pcEncadre}{definition}{Repère 5}
Dans un système fermé, la masse se conserve ; le volume peut varier.\par
\end{pcEncadre}
\begin{pcEncadre}{attention}{Une erreur à éviter}
« L’eau s’évapore, donc elle n’existe plus. » Elle devient gazeuse et peut quitter le récipient ; elle ne disparaît pas.\par
\end{pcEncadre}
\end{pcFiche}
\clearpage\markright{Exercices}
\begin{blocLivre}{exo-3-1}
\rubrique{Je démarre}
\begin{exercice}[1]{Nommer les passages}
\begin{questions}
\item Un glaçon devient liquide. Quel changement d’état se produit ?
\item L’eau liquide devient de la glace. Quel est le changement inverse ?
\end{questions}
\end{exercice}\end{blocLivre}
\begin{blocLivre}{exo-3-2}
\begin{exercice}[1]{La buée sur une vitre}
\begin{questions}
\item De la vapeur d’eau forme des gouttelettes sur une vitre froide. Nomme le changement.
\item Quel est l’état physique de ces gouttelettes ?
\end{questions}
\end{exercice}\end{blocLivre}
\begin{blocLivre}{exo-3-3}
\begin{exercice}[1]{Lire une température}
\pcFigure{figures/fde5db9303e00188.pdf}{}
\begin{questions}
\item Quelle grandeur porte l’axe vertical ? Quelle est son unité ?
\item À quelle température se situe le palier ?
\end{questions}
\end{exercice}\end{blocLivre}
\begin{blocLivre}{exo-3-4}
\begin{exercice}[1]{Évaporation ou ébullition ?}
\begin{questions}
\item Du linge sèche à 20 °C. L’eau a-t-elle nécessairement bouilli ?
\item Une casserole d’eau forme des bulles de vapeur dans tout le liquide. Comment nommer cette vaporisation ?
\end{questions}
\end{exercice}\end{blocLivre}
\begin{blocLivre}{exo-3-5}
\begin{exercice}[1]{Choisir la bonne unité}
\begin{questions}
\item Complète : la température du mélange vaut 0 … .
\item Pourquoi « 0 min » ne convient-il pas ?
\end{questions}
\end{exercice}\end{blocLivre}
\begin{blocLivre}{exo-3-6}
\rubrique{Je m’entraîne}
\begin{exercice}[2]{Peser pendant la fusion}
\begin{donnees}Un récipient fermé contenant de la glace a une masse totale de 180,0 g. Il reste fermé pendant la fusion.\par
\end{donnees}
\begin{questions}
\item Quelle masse doit-on retrouver après fusion complète ?
\item Cite la condition indispensable à cette prévision.
\end{questions}
\end{exercice}\end{blocLivre}
\begin{blocLivre}{exo-3-7}
\begin{exercice}[2]{Deux grandeurs à distinguer}
\pcFigure{figures/4bf3cc5c97b47074.pdf}{}
\begin{questions}
\item Le niveau augmente lors du gel. Peut-on conclure que la masse d’eau augmente ?
\item Explique en distinguant masse et volume.
\end{questions}
\end{exercice}\end{blocLivre}
\begin{blocLivre}{exo-3-8}
\begin{exercice}[2]{Durée du palier}
\begin{questions}
\item Le palier commence à 3 min et finit à 8 min. Quelle est sa durée ?
\item Que peut-on observer dans le récipient pendant ce palier de fusion ?
\end{questions}
\end{exercice}\end{blocLivre}
\begin{blocLivre}{exo-3-9}
\begin{exercice}[2]{Dessiner une courbe}
\begin{donnees}Eau pure, pression usuelle. Temps : 0, 1, 2, 3, 4 min. Température : -5, 0, 0, 0, 5 °C.\par
\end{donnees}
\begin{questions}
\item Trace la courbe des températures en fonction du temps.
\item Repère la portion correspondant à la fusion.
\end{questions}
\end{exercice}\end{blocLivre}
\begin{blocLivre}{exo-3-10}
\begin{exercice}[2]{Le récipient ouvert}
\begin{donnees}On chauffe doucement de l’eau dans un récipient ouvert. Aucune éclaboussure n’est observée.\par
\end{donnees}
\begin{questions}
\item La masse mesurée passe de 250 g à 242 g. Quelle masse a quitté le récipient ?
\item Pourquoi cela ne contredit-il pas la conservation de la masse ?
\end{questions}
\end{exercice}\end{blocLivre}
\begin{blocLivre}{exo-3-11}
\begin{exercice}[2]{Le bon placement du thermomètre}
\begin{questions}
\item Pourquoi éviter d’appuyer le capteur sur le fond chauffé ?
\item Quelle information supplémentaire faut-il noter à chaque mesure ?
\end{questions}
\end{exercice}\end{blocLivre}
\begin{blocLivre}{exo-3-12}
\rubrique{Je raisonne}
\begin{exercice}[3]{Un glaçon dans un sachet}
\begin{questions}
\item Propose une expérience pour tester la conservation de la masse pendant la fusion.
\item Indique une observation qui invaliderait le protocole.
\end{questions}
\end{exercice}\end{blocLivre}
\begin{blocLivre}{exo-3-13}
\begin{exercice}[3]{Une eau qui bout à 94 °C}
\begin{questions}
\item Un document donne une pression atmosphérique plus faible qu’au niveau de la mer. Le thermomètre est-il forcément défectueux ?
\item Quelle condition faut-il préciser lorsqu’on annonce « l’eau bout à 100 °C » ?
\end{questions}
\end{exercice}\end{blocLivre}
\begin{blocLivre}{exo-3-14}
\begin{exercice}[3]{Un palier trop vite annoncé}
\begin{questions}
\item Une élève mesure une seule fois 0 °C puis écrit « il y a un palier ». Pourquoi sa conclusion est-elle insuffisante ?
\item Quelles mesures et observations compléter ?
\end{questions}
\end{exercice}\end{blocLivre}
\begin{blocLivre}{exo-3-15}
\begin{exercice}[3]{Interpréter le modèle}
\begin{questions}
\item Les particules d’eau deviennent-elles plus grosses quand la glace fond ?
\item Que conserve-t-on et que modifie-t-on dans un modèle simple ?
\end{questions}
\end{exercice}\end{blocLivre}
\pcPage{Défi · Prendre une initiative}{Enquêter sur une bouteille déformée}
\phantomsection\label{defi-ex-3}
\begin{pcFiche}{defi-3}
Une bouteille presque pleine d’eau a été oubliée au congélateur. Après gel, elle est déformée. Un camarade affirme : « La glace pèse plus lourd, donc elle pousse la bouteille. » Construis une explication scientifique sans refaire cette manipulation.\par
\pcFigure{figures/7bb00c91ea07e1d4.pdf}{}
\begin{questions}
\item Identifie les deux grandeurs confondues.
\item Propose une vérification sûre de la conservation de la masse.
\item Explique la déformation avec la propriété de l’eau lors du gel.
\end{questions}
\begin{pcEncadre}{method}{Indice progressif}
Comparer une même masse d’eau sous deux états et distinguer l’espace occupé de la masse.\par
\end{pcEncadre}
\end{pcFiche}
\clearpage\pagestyle{otMain}
\begin{ouverture}[college1]{4}{Corps purs, mélanges et composition de l’air}{Partie 1 · Organisation et transformations de la matière}{Un liquide transparent est-il forcément de l’eau pure ?}
\objectif{Distinguer corps pur, mélange homogène et mélange hétérogène.}
\objectif{Relier un modèle particulaire à la présence d’une ou plusieurs espèces.}
\objectif{Décrire la composition approximative de l’air sec.}
\end{ouverture}
\pcPage{L’essentiel}{Chapitre 04 · Corps purs, mélanges et composition de l’air}
\phantomsection\label{ch-4-cours-1}
\phantomsection\label{ess-4}
\begin{pcFiche}{cours-4-1}
\essSection{1}{Une seule espèce ou plusieurs ?}
Une espèce chimique désigne une catégorie de constituants identiques du point de vue chimique. Un corps pur ne contient qu’une seule espèce ; un mélange en contient plusieurs. Une eau minérale contient de l’eau et des substances dissoutes : c’est un mélange. L’eau distillée est considérée comme un corps pur dans le modèle scolaire si les impuretés sont négligées. L’aspect seul ne permet pas de compter les espèces présentes.\par
\pcFigure{figures/eb162b7ceff9cbaa.pdf}{Modèles simplifiés : les couleurs représentent des espèces différentes, sans représenter leur couleur réelle.}
\begin{pcEncadre}{definition}{À retenir}
« Pur » a un sens scientifique précis. Une eau potable n’est pas nécessairement un corps pur ; une eau chimiquement pure n’est pas automatiquement garantie potable.\par
\end{pcEncadre}
\begin{pcEncadre}{method}{Une méthode concrète}
Pour classer un échantillon, rechercher sa composition. Une couleur, une étiquette publicitaire ou la transparence ne remplacent pas cette information.\par
\end{pcEncadre}
\end{pcFiche}
\pcPage{L’essentiel}{Chapitre 04 · Corps purs, mélanges et composition de l’air}
\phantomsection\label{ch-4-cours-2}
\begin{pcFiche}{cours-4-2}
\essSection{2}{Homogène ou hétérogène ?}
À l’œil nu, on distingue plusieurs constituants ou plusieurs phases d’un mélange hétérogène. Dans un mélange homogène, on ne distingue qu’une seule phase à cette échelle d’observation. L’eau salée est homogène lorsque le sel est entièrement dissous. L’eau et l’huile forment après repos deux phases visibles. Le classement dépend des conditions et de l’échelle d’observation ; une agitation peut produire une dispersion temporaire.\par
\pcFigure{figures/801d683a4dd8d300.pdf}{À gauche, une solution homogène ; à droite, deux phases liquides après repos.}
\begin{pcEncadre}{definition}{À retenir}
Une seule phase visible n’implique pas une seule espèce chimique. Un mélange homogène peut être incolore, transparent ou coloré.\par
\end{pcEncadre}
\begin{pcEncadre}{method}{Une méthode concrète}
Observer après agitation puis après repos. Écrire « je distingue… » avant de conclure « le mélange est… ». Ne jamais goûter un produit pour l’identifier.\par
\end{pcEncadre}
\end{pcFiche}
\pcPage{L’essentiel}{Chapitre 04 · Corps purs, mélanges et composition de l’air}
\phantomsection\label{ch-4-cours-3}
\begin{pcFiche}{cours-4-3}
\essSection{3}{L’air, un mélange de gaz}
L’air qui nous entoure est un mélange. Pour l’air sec, on retient environ 78 \% de diazote et 21 \% de dioxygène en volume. Le pour cent restant rassemble d’autres gaz, notamment l’argon et le dioxyde de carbone. L’air ambiant contient aussi une quantité variable de vapeur d’eau. Ces proportions ne signifient pas qu’un petit récipient contient des couches de gaz séparées.\par
\pcFigure{figures/85ec3b167784b239.pdf}{Cent points représentent cent parts égales de volume d’air sec ; valeurs arrondies.}
\begin{pcEncadre}{definition}{À retenir}
L’air n’est pas du dioxygène pur. Les mots « gaz », « air » et « dioxygène » n’ont pas le même sens.\par
\end{pcEncadre}
\begin{pcEncadre}{method}{Une méthode concrète}
Sur 100 parts d’air sec, environ 21 correspondent au dioxygène. Pour 200 parts, toutes les valeurs sont multipliées par 2 ; la proportion reste la même.\par
\end{pcEncadre}
\end{pcFiche}
\pcPage{Activité expérimentale}{Classer sans se fier à la couleur}
\phantomsection\label{ch-4-activite}
\begin{pcFiche}{activite-4}
\begin{pcEncadre}{method}{Mission}
Construire un classement à partir d’observations et de compositions fournies.\par
\end{pcEncadre}
\begin{questions}
\item Préparer avec le professeur de l’eau salée sans cristaux visibles, de l’eau et de l’huile, puis de l’eau distillée de laboratoire.
\item Observer les récipients, les laisser reposer et noter le nombre de phases visibles. Ne goûter aucun échantillon.
\item Comparer l’observation aux compositions connues : eau seule dans le modèle ; eau et sel ; eau et huile.
\end{questions}
\pcFigure{figures/9a252fdceeeb3df9.pdf}{Deux aspects possibles ne suffisent pas à distinguer un corps pur d’un mélange homogène.}
\begin{questions}
\item Quel échantillon est un corps pur dans le modèle ? Lesquels sont des mélanges ?
\item Pourquoi l’œil ne permet-il pas de distinguer eau distillée et eau salée ?
\item Rédige une phrase qui utilise correctement « transparent » et « homogène ».
\end{questions}
\end{pcFiche}
\pcPage{L’essentiel à retenir}{Chapitre 04 · Synthèse et autoévaluation}
\phantomsection\label{ch-4-bilan}
\begin{pcFiche}{synthese-4}
\begin{pcEncadre}{definition}{Repère 1}
Un corps pur contient une seule espèce chimique ; un mélange en contient plusieurs.\par
\end{pcEncadre}
\begin{pcEncadre}{method}{Repère 2}
Homogène : une seule phase visible à l’échelle considérée.\par
\end{pcEncadre}
\begin{pcEncadre}{definition}{Repère 3}
Hétérogène : plusieurs phases ou constituants visibles.\par
\end{pcEncadre}
\begin{pcEncadre}{method}{Repère 4}
L’aspect transparent ne prouve pas la pureté.\par
\end{pcEncadre}
\begin{pcEncadre}{definition}{Repère 5}
L’air sec contient environ 78 \% de diazote et 21 \% de dioxygène en volume.\par
\end{pcEncadre}
\begin{pcEncadre}{attention}{Une erreur à éviter}
« Je ne vois pas le sel, donc il n’y en a plus. » Une substance dissoute reste présente même si elle n’est plus visible.\par
\end{pcEncadre}
\end{pcFiche}
\clearpage\markright{Exercices}
\begin{blocLivre}{exo-4-1}
\rubrique{Je démarre}
\begin{exercice}[1]{Corps pur ou mélange ?}
\begin{questions}
\item Classe : eau seule, eau salée, air sec.
\item Quel critère utilises-tu ?
\end{questions}
\end{exercice}\end{blocLivre}
\begin{blocLivre}{exo-4-2}
\begin{exercice}[1]{Nommer l’aspect}
\begin{questions}
\item On distingue de la pulpe dans un jus. Le mélange est-il homogène ou hétérogène ?
\item Une eau salée sans cristaux visibles présente une seule phase. Comment la classer ?
\end{questions}
\end{exercice}\end{blocLivre}
\begin{blocLivre}{exo-4-3}
\begin{exercice}[1]{Les deux gaz principaux}
\begin{questions}
\item Nomme les deux principaux gaz de l’air sec.
\item Donne leurs proportions approximatives en volume.
\end{questions}
\end{exercice}\end{blocLivre}
\begin{blocLivre}{exo-4-4}
\begin{exercice}[1]{Lire un modèle}
\pcFigure{figures/2b735c49a64f52ea.pdf}{}
\begin{questions}
\item Quel dessin représente un corps pur ?
\item Que représente le changement de couleur des points ?
\end{questions}
\end{exercice}\end{blocLivre}
\begin{blocLivre}{exo-4-5}
\begin{exercice}[1]{Déjouer une affirmation}
\begin{questions}
\item « Un liquide incolore est pur. » Vrai ou faux ? Donne un contre-exemple.
\end{questions}
\end{exercice}\end{blocLivre}
\begin{blocLivre}{exo-4-6}
\rubrique{Je m’entraîne}
\begin{exercice}[2]{Lire une étiquette}
\begin{questions}
\item Une eau contient des ions calcium, magnésium et des hydrogénocarbonates en plus de l’eau. Est-ce un mélange ?
\item Peut-elle être potable tout en étant un mélange ?
\end{questions}
\end{exercice}\end{blocLivre}
\begin{blocLivre}{exo-4-7}
\begin{exercice}[2]{Air sec dans une seringue}
\begin{questions}
\item Dans 100 mL d’air sec, quel volume correspond approximativement au dioxygène ?
\item Quel volume correspond au diazote ?
\end{questions}
\end{exercice}\end{blocLivre}
\begin{blocLivre}{exo-4-8}
\begin{exercice}[2]{Changer de quantité}
\begin{questions}
\item Calcule les parts de dioxygène et de diazote correspondant à 200 mL d’air sec.
\item Les pourcentages ont-ils doublé ?
\end{questions}
\end{exercice}\end{blocLivre}
\begin{blocLivre}{exo-4-9}
\begin{exercice}[2]{Agiter puis laisser reposer}
\begin{questions}
\item Après agitation, eau et huile semblent dispersées. Après repos, deux couches apparaissent. Quelle observation est décisive ?
\item Comment classer le mélange après repos ?
\end{questions}
\end{exercice}\end{blocLivre}
\begin{blocLivre}{exo-4-10}
\begin{exercice}[2]{Faire le total}
\begin{donnees}Air sec : 78 \% de diazote et 21 \% de dioxygène, valeurs arrondies.\par
\end{donnees}
\begin{questions}
\item Calcule la part des autres gaz.
\item Pourquoi ne peut-on pas attribuer tout ce reste au dioxyde de carbone ?
\end{questions}
\end{exercice}\end{blocLivre}
\begin{blocLivre}{exo-4-11}
\begin{exercice}[2]{Préciser l’échelle}
\begin{questions}
\item Un lait paraît uniforme à l’œil nu, mais contient de minuscules gouttes de matière grasse dispersées. Que faut-il préciser dans sa description ?
\item Pourquoi les observations au microscope et à l’œil nu peuvent-elles différer ?
\end{questions}
\end{exercice}\end{blocLivre}
\begin{blocLivre}{exo-4-12}
\rubrique{Je raisonne}
\begin{exercice}[3]{Deux flacons identiques}
\begin{questions}
\item Un flacon contient de l’eau pure, l’autre de l’eau salée. Tous deux sont limpides. Propose une méthode pour les distinguer sans goûter.
\item Quelle limite faut-il garder en tête ?
\end{questions}
\end{exercice}\end{blocLivre}
\begin{blocLivre}{exo-4-13}
\begin{exercice}[3]{Un modèle à corriger}
\pcFigure{figures/25edd406a9df0934.pdf}{Document de travail · modèle simplifié.}
\begin{questions}
\item Un dessin représente l’air par une seule sorte de points. Est-ce cohérent avec sa composition ?
\item Comment améliorer la légende et le dessin ?
\end{questions}
\end{exercice}\end{blocLivre}
\begin{blocLivre}{exo-4-14}
\begin{exercice}[3]{Passer à un litre}
\begin{questions}
\item Détermine les volumes correspondant aux deux gaz principaux dans 1 L d’air sec.
\item Exprime les résultats en mL.
\end{questions}
\end{exercice}\end{blocLivre}
\begin{blocLivre}{exo-4-15}
\begin{exercice}[3]{Argumenter avec deux indices}
\begin{questions}
\item Un liquide transparent laisse un dépôt après évaporation. Que peut-on en conclure ?
\item Peut-on identifier précisément la substance du dépôt par cette seule observation ?
\end{questions}
\end{exercice}\end{blocLivre}
\pcPage{Défi · Prendre une initiative}{Préparer un modèle de l’air}
\phantomsection\label{defi-ex-4}
\begin{pcFiche}{defi-4}
Tu disposes de 100 jetons de même taille pour représenter l’air sec. Construis une affiche qui explique à un élève de 6e ce que représente chaque jeton et ce que le modèle ne montre pas.\par
\pcFigure{figures/4c54cd7d3062b491.pdf}{}
\begin{questions}
\item Répartis les jetons et prépare une légende.
\item Explique pourquoi ce modèle ne représente pas des couches de gaz.
\item Adapte la répartition à 200 jetons.
\end{questions}
\begin{pcEncadre}{method}{Indice progressif}
Garder les proportions, et préciser que l’air réel contient des constituants en mouvement.\par
\end{pcEncadre}
\end{pcFiche}
\clearpage\pagestyle{otMain}
\begin{ouverture}[college1]{5}{Dissolution et miscibilité}{Partie 1 · Organisation et transformations de la matière}{Le sucre a-t-il disparu lorsqu’on ne le voit plus ?}
\objectif{Identifier solvant, soluté et solution.}
\objectif{Tester une dissolution et reconnaître une saturation.}
\objectif{Distinguer dissolution d’un solide et miscibilité de deux liquides.}
\end{ouverture}
\pcPage{L’essentiel}{Chapitre 05 · Dissolution et miscibilité}
\phantomsection\label{ch-5-cours-1}
\phantomsection\label{ess-5}
\begin{pcFiche}{cours-5-1}
\essSection{1}{Dissoudre sans faire disparaître}
Lorsqu’un solide se dissout dans un liquide, ses constituants se dispersent dans le liquide. Le liquide qui dissout est le solvant ; la substance dissoute est le soluté. Le mélange homogène obtenu est une solution. Avec de l’eau comme solvant, on parle de solution aqueuse. Le soluté ne devient pas nécessairement liquide et ne disparaît pas : une dissolution n’est pas une fusion.\par
\pcFigure{figures/7860c21c7ed00911.pdf}{Les points colorés symbolisent le soluté réparti dans l’eau ; ils ne représentent pas des grains visibles.}
\begin{pcEncadre}{definition}{À retenir}
Si tout le sel est dissous, il reste présent dans la solution. Dans un système sans perte, masse de solution = masse de solvant + masse de soluté ajouté.\par
\end{pcEncadre}
\begin{pcEncadre}{method}{Une méthode concrète}
Pour préparer une solution, peser le soluté et le solvant séparément, puis mélanger sans perte. Ne pas additionner automatiquement les volumes.\par
\end{pcEncadre}
\end{pcFiche}
\pcPage{L’essentiel}{Chapitre 05 · Dissolution et miscibilité}
\phantomsection\label{ch-5-cours-2}
\begin{pcFiche}{cours-5-2}
\essSection{2}{Il existe une limite à la dissolution}
Un solide est soluble dans un solvant s’il peut s’y dissoudre dans les conditions étudiées. Si l’on ajoute progressivement trop de soluté, un dépôt peut persister malgré une agitation suffisante : la solution est saturée. La quantité maximale dissoute dépend notamment du soluté, du solvant et de la température. Agiter facilite le contact et accélère souvent la dissolution, sans garantir que toute quantité pourra se dissoudre.\par
\pcFigure{figures/7860c21c7ed00911.pdf}{Dans le dernier récipient, une partie du solide reste au fond : la solution au-dessus est saturée.}
\begin{pcEncadre}{definition}{À retenir}
Un dépôt juste après l’ajout n’est pas à lui seul la preuve d’une saturation. Laisser assez de temps et agiter selon le protocole avant de conclure.\par
\end{pcEncadre}
\begin{pcEncadre}{method}{Une méthode concrète}
Comparer des essais avec la même quantité de solvant à la même température. Modifier une seule variable à la fois.\par
\end{pcEncadre}
\end{pcFiche}
\pcPage{L’essentiel}{Chapitre 05 · Dissolution et miscibilité}
\phantomsection\label{ch-5-cours-3}
\begin{pcFiche}{cours-5-3}
\essSection{3}{Mélanger deux liquides}
Deux liquides miscibles forment un mélange homogène lorsqu’on les mélange dans les proportions étudiées. L’eau et l’éthanol sont miscibles ; leur manipulation éventuelle reste encadrée en classe. L’eau et l’huile sont non miscibles : après agitation puis repos, elles forment des phases distinctes. Une dispersion temporaire de gouttelettes ne constitue pas une preuve de miscibilité.\par
\pcFigure{figures/801d683a4dd8d300.pdf}{Une phase visible après mélange ; deux phases après repos dans le cas de liquides non miscibles.}
\begin{pcEncadre}{definition}{À retenir}
On parle ici de dissolution pour un solide dans un liquide, et de miscibilité pour comparer deux liquides. Toujours nommer les substances et les conditions.\par
\end{pcEncadre}
\begin{pcEncadre}{method}{Une méthode concrète}
Agiter doucement les deux liquides dans un récipient adapté, puis observer après repos. Décrire les phases avant d’utiliser le mot miscible.\par
\end{pcEncadre}
\end{pcFiche}
\pcPage{Activité expérimentale}{Chercher une limite expérimentale}
\phantomsection\label{ch-5-activite}
\begin{pcFiche}{activite-5}
\begin{pcEncadre}{method}{Mission}
Déterminer entre quelles masses se situe une limite de dissolution.\par
\end{pcEncadre}
\begin{questions}
\item Le professeur fournit trois essais réalisés avec 50 g d’eau à la même température et un solide X : 10 g entièrement dissous ; 15 g entièrement dissous ; 20 g avec dépôt persistant.
\item Comparer les essais après une même durée d’agitation et de repos. Les données sont pédagogiques et ne décrivent pas une substance à identifier.
\item Proposer une nouvelle masse à tester entre les deux essais qui encadrent la limite.
\end{questions}
\pcFigure{figures/42794c855063a88c.pdf}{L’encadrement suppose une agitation suffisante et une température constante.}
\begin{questions}
\item Que peut-on affirmer de la masse maximale qui se dissout dans 50 g d’eau ?
\item Quel nouvel essai permettrait de préciser l’encadrement ?
\item Pourquoi ne faut-il pas changer la température pendant la comparaison ?
\end{questions}
\end{pcFiche}
\pcPage{L’essentiel à retenir}{Chapitre 05 · Synthèse et autoévaluation}
\phantomsection\label{ch-5-bilan}
\begin{pcFiche}{synthese-5}
\begin{pcEncadre}{definition}{Repère 1}
Une solution contient un solvant et un ou plusieurs solutés.\par
\end{pcEncadre}
\begin{pcEncadre}{method}{Repère 2}
Dissolution et fusion sont deux phénomènes différents.\par
\end{pcEncadre}
\begin{pcEncadre}{definition}{Repère 3}
La masse se conserve pendant une dissolution sans perte de matière.\par
\end{pcEncadre}
\begin{pcEncadre}{method}{Repère 4}
Une solution saturée ne peut plus dissoudre davantage de soluté dans les conditions données.\par
\end{pcEncadre}
\begin{pcEncadre}{definition}{Repère 5}
Deux liquides non miscibles forment des phases distinctes après repos.\par
\end{pcEncadre}
\begin{pcEncadre}{attention}{Une erreur à éviter}
« Agiter plus fort permet de dissoudre une quantité illimitée. » L’agitation peut accélérer la dissolution, mais ne supprime pas la saturation.\par
\end{pcEncadre}
\end{pcFiche}
\clearpage\markright{Exercices}
\begin{blocLivre}{exo-5-1}
\rubrique{Je démarre}
\begin{exercice}[1]{Le vocabulaire de la solution}
\begin{questions}
\item On dissout du sucre dans l’eau. Nomme le solvant et le soluté.
\item Comment nomme-t-on le mélange homogène obtenu ?
\end{questions}
\end{exercice}\end{blocLivre}
\begin{blocLivre}{exo-5-2}
\begin{exercice}[1]{Dissoudre ou fondre ?}
\begin{questions}
\item Un glaçon devient liquide au-dessus de 0 °C. Est-ce une dissolution ?
\item Du sel se disperse dans l’eau sans rester visible. Quel phénomène se produit ?
\end{questions}
\end{exercice}\end{blocLivre}
\begin{blocLivre}{exo-5-3}
\begin{exercice}[1]{Lire les récipients}
\pcFigure{figures/42794c855063a88c.pdf}{}
\begin{questions}
\item Quel récipient montre un dépôt persistant ?
\item Que signifie « solution saturée » ?
\end{questions}
\end{exercice}\end{blocLivre}
\begin{blocLivre}{exo-5-4}
\begin{exercice}[1]{Deux liquides au repos}
\begin{questions}
\item Eau et huile donnent deux couches après repos. Sont-elles miscibles ?
\item Le mélange est-il homogène ou hétérogène à l’œil nu ?
\end{questions}
\end{exercice}\end{blocLivre}
\begin{blocLivre}{exo-5-5}
\begin{exercice}[1]{La substance invisible}
\begin{questions}
\item « Le sel n’est plus visible, donc la solution ne contient que de l’eau. » Corrige cette phrase.
\end{questions}
\end{exercice}\end{blocLivre}
\begin{blocLivre}{exo-5-6}
\rubrique{Je m’entraîne}
\begin{exercice}[2]{Conserver la masse}
\begin{questions}
\item On mélange 80 g d’eau et 12 g de sucre sans perte. Calcule la masse de solution.
\item Quel type de grandeur additionnes-tu ?
\end{questions}
\end{exercice}\end{blocLivre}
\begin{blocLivre}{exo-5-7}
\begin{exercice}[2]{Retrouver la masse du soluté}
\begin{questions}
\item Une solution a une masse de 125 g. Elle a été préparée avec 110 g d’eau. Quelle masse de soluté a été ajoutée ?
\item Quelle hypothèse utilises-tu ?
\end{questions}
\end{exercice}\end{blocLivre}
\begin{blocLivre}{exo-5-8}
\begin{exercice}[2]{Ne changer qu’un facteur}
\begin{questions}
\item Pour comparer l’effet de l’agitation, Léa utilise 50 mL d’eau froide et 100 mL d’eau chaude. Son protocole convient-il ?
\item Propose une correction.
\end{questions}
\end{exercice}\end{blocLivre}
\begin{blocLivre}{exo-5-9}
\begin{exercice}[2]{Encadrer une limite}
\begin{donnees}Pour 40 g d’eau à température fixée : 8 g dissous ; 12 g dissous ; 16 g avec un dépôt persistant après agitation suffisante.\par
\end{donnees}
\begin{questions}
\item Quelle masse maximale dissoute peut-on encadrer ?
\item Propose un quatrième essai utile.
\end{questions}
\end{exercice}\end{blocLivre}
\begin{blocLivre}{exo-5-10}
\begin{exercice}[2]{Une masse totale avec un dépôt}
\begin{questions}
\item On ajoute 25 g d’un solide à 60 g d’eau. Une partie reste au fond. Quelle masse totale contient le récipient, sans perte ?
\item Cette masse est-elle celle de la seule phase liquide ?
\end{questions}
\end{exercice}\end{blocLivre}
\begin{blocLivre}{exo-5-11}
\begin{exercice}[2]{Comparer deux volumes de solvant}
\begin{questions}
\item Une substance peut se dissoudre jusqu’à 18 g dans 50 g d’eau. Dans les mêmes conditions, quelle masse correspond à 100 g d’eau ?
\item Quelle hypothèse permet ce calcul ?
\end{questions}
\end{exercice}\end{blocLivre}
\begin{blocLivre}{exo-5-12}
\rubrique{Je raisonne}
\begin{exercice}[3]{Le dépôt disparaît lentement}
\begin{questions}
\item Juste après ajout du sucre, un dépôt est visible. Dix minutes plus tard, il a disparu. Pouvait-on conclure immédiatement à la saturation ?
\item Quelle précaution de langage adopter ?
\end{questions}
\end{exercice}\end{blocLivre}
\begin{blocLivre}{exo-5-13}
\begin{exercice}[3]{Choisir le mot précis}
\begin{questions}
\item Une élève écrit « l’huile est insoluble, donc l’eau est un mauvais liquide ». Reformule scientifiquement.
\item Pourquoi doit-on nommer les deux liquides ?
\end{questions}
\end{exercice}\end{blocLivre}
\begin{blocLivre}{exo-5-14}
\begin{exercice}[3]{Calculer la masse dissoute}
\begin{questions}
\item On ajoute 22 g de X à 50 g d’eau. Après saturation, 4 g de solide sec restent non dissous. Quelle masse a été dissoute ?
\item Quelle est la masse de la solution séparée du dépôt, sans perte ?
\end{questions}
\end{exercice}\end{blocLivre}
\begin{blocLivre}{exo-5-15}
\begin{exercice}[3]{Le volume ne s’additionne pas toujours}
\begin{questions}
\item Deux liquides miscibles de volumes 50 mL et 50 mL donnent 97 mL de mélange. Y a-t-il forcément eu une perte de matière ?
\item Quelle mesure permettrait de contrôler une éventuelle perte ?
\end{questions}
\end{exercice}\end{blocLivre}
\pcPage{Défi · Prendre une initiative}{Obtenir une solution sans dépôt}
\phantomsection\label{defi-ex-5}
\begin{pcFiche}{defi-5}
À température fixée, au maximum 12 g d’un solide X se dissolvent dans 40 g d’eau. On a versé 18 g de X dans 40 g d’eau. Sans chauffer ni retirer de solide, quelle quantité minimale d’eau faut-il ajouter pour pouvoir tout dissoudre ?\par
\pcFigure{figures/8769e03894a55a95.pdf}{}
\begin{questions}
\item Calcule la masse d’eau nécessaire pour 18 g de X.
\item Déduis la masse d’eau à ajouter.
\item Explique comment vérifier le résultat expérimentalement.
\end{questions}
\begin{pcEncadre}{method}{Indice progressif}
18 g représentent une fois et demie 12 g ; garder la même proportion entre soluté et eau.\par
\end{pcEncadre}
\end{pcFiche}
\clearpage\pagestyle{otMain}
\begin{ouverture}[college1]{6}{Séparer les constituants d’un mélange}{Partie 1 · Organisation et transformations de la matière}{Comment récupérer du sel et de l’eau séparément ?}
\objectif{Choisir décantation ou filtration pour un mélange hétérogène.}
\objectif{Distinguer évaporation et distillation selon le constituant recherché.}
\objectif{Décrire un protocole et ses limites, notamment pour la potabilité.}
\end{ouverture}
\pcPage{L’essentiel}{Chapitre 06 · Séparer les constituants d’un mélange}
\phantomsection\label{ch-6-cours-1}
\phantomsection\label{ess-6}
\begin{pcFiche}{cours-6-1}
\essSection{1}{Laisser reposer puis séparer}
La décantation consiste à laisser un mélange hétérogène au repos pour que les phases se séparent. Des particules solides peuvent se déposer au fond ; des liquides non miscibles peuvent former des couches. On peut ensuite prélever une phase avec précaution. La décantation ne retire pas les substances dissoutes dans l’eau et ne garantit pas l’élimination des particules les plus fines.\par
\pcFigure{figures/27b8d8da40c1f5b3.pdf}{Une partie des particules se dépose pendant le repos ; l’eau surnageante peut rester imparfaitement limpide.}
\begin{pcEncadre}{definition}{À retenir}
Un liquide plus limpide n’est pas nécessairement pur ni potable. La séparation doit répondre à un objectif annoncé.\par
\end{pcEncadre}
\begin{pcEncadre}{method}{Une méthode concrète}
Avant de choisir la méthode, repérer le nombre de phases et l’état des constituants. Laisser le dépôt au fond lorsque l’on prélève le liquide.\par
\end{pcEncadre}
\end{pcFiche}
\pcPage{L’essentiel}{Chapitre 06 · Séparer les constituants d’un mélange}
\phantomsection\label{ch-6-cours-2}
\begin{pcFiche}{cours-6-2}
\essSection{2}{Filtrer un mélange hétérogène}
Une filtration utilise une barrière poreuse, par exemple un papier-filtre. Les particules suffisamment grosses restent sur le filtre : elles constituent le résidu. Le liquide qui traverse est le filtrat. Le sel dissous traverse un papier-filtre ordinaire avec l’eau ; il ne forme pas des cristaux bloqués dans le filtre. La filtration sépare selon des propriétés et des tailles, pas selon la dangerosité des substances.\par
\pcFigure{figures/879d24dea896097b.pdf}{Le liquide recueilli est le filtrat ; les particules retenues forment le résidu.}
\begin{pcEncadre}{definition}{À retenir}
Filtrer de l’eau salée ne permet pas, à lui seul, de retirer le sel dissous. Un filtre scolaire n’est pas un dispositif garantissant la potabilité.\par
\end{pcEncadre}
\begin{pcEncadre}{method}{Une méthode concrète}
Nommer les récipients et les produits recueillis. Verser doucement avec un montage stable, sans remplir le filtre au-delà de son bord.\par
\end{pcEncadre}
\end{pcFiche}
\pcPage{L’essentiel}{Chapitre 06 · Séparer les constituants d’un mélange}
\phantomsection\label{ch-6-cours-3}
\begin{pcFiche}{cours-6-3}
\essSection{3}{Récupérer le soluté ou le solvant}
Pour récupérer un solide dissous non volatil, on peut laisser évaporer le solvant ; le solide demeure. Pour récupérer aussi le solvant, on utilise une distillation : une partie du liquide se vaporise, puis sa vapeur se liquéfie dans un réfrigérant. Le liquide recueilli est le distillat. Au collège, les montages avec chauffage sont réalisés uniquement selon les consignes du professeur.\par
\pcFigure{figures/34523d21d3212538.pdf}{Distillation simplifiée d’eau salée : le sel reste dans le récipient chauffé ; l’eau se vaporise puis se liquéfie.}
\begin{pcEncadre}{definition}{À retenir}
Évaporer permet de récupérer le sel, mais laisse partir l’eau dans l’air. Distiller organise une récupération de l’eau après deux changements d’état.\par
\end{pcEncadre}
\begin{pcEncadre}{method}{Une méthode concrète}
Préciser ce que l’on cherche à récupérer. Une méthode de séparation ne prouve pas que l’eau obtenue répond à tous les critères sanitaires.\par
\end{pcEncadre}
\end{pcFiche}
\pcPage{Activité expérimentale}{Séparer un mélange à trois constituants}
\phantomsection\label{ch-6-activite}
\begin{pcFiche}{activite-6}
\begin{pcEncadre}{method}{Mission}
Proposer un enchaînement pour un mélange d’eau, de sel dissous et de sable.\par
\end{pcEncadre}
\begin{questions}
\item Identifier les constituants visibles et ceux qui restent dissous. Préparer un schéma de chaque étape.
\item Avec le professeur, filtrer pour séparer le sable. Étiqueter résidu et filtrat.
\item Sur une petite partie du filtrat, observer l’évaporation de l’eau selon le protocole de classe. Ne jamais goûter les produits.
\end{questions}
\pcFigure{figures/3bc7ac9680ed49fa.pdf}{Après filtration, le filtrat contient encore de l’eau et du sel dissous.}
\begin{questions}
\item Quel constituant se trouve sur le filtre ?
\item Quel constituant solide récupère-t-on après évaporation du filtrat ?
\item Quelle étape remplacer pour récupérer de l’eau liquide au lieu de la laisser partir ?
\end{questions}
\end{pcFiche}
\pcPage{L’essentiel à retenir}{Chapitre 06 · Synthèse et autoévaluation}
\phantomsection\label{ch-6-bilan}
\begin{pcFiche}{synthese-6}
\begin{pcEncadre}{definition}{Repère 1}
La décantation sépare des phases après repos.\par
\end{pcEncadre}
\begin{pcEncadre}{method}{Repère 2}
La filtration donne un résidu et un filtrat.\par
\end{pcEncadre}
\begin{pcEncadre}{definition}{Repère 3}
Un soluté dissous traverse habituellement le papier-filtre avec le solvant.\par
\end{pcEncadre}
\begin{pcEncadre}{method}{Repère 4}
L’évaporation laisse un solide dissous non volatil.\par
\end{pcEncadre}
\begin{pcEncadre}{definition}{Repère 5}
La distillation permet de recueillir un liquide après vaporisation puis liquéfaction.\par
\end{pcEncadre}
\begin{pcEncadre}{attention}{Une erreur à éviter}
« L’eau est transparente après filtration, donc on peut la boire. » La transparence ne garantit ni la pureté chimique ni la potabilité.\par
\end{pcEncadre}
\end{pcFiche}
\clearpage\markright{Exercices}
\begin{blocLivre}{exo-6-1}
\rubrique{Je démarre}
\begin{exercice}[1]{Choisir une technique}
\begin{questions}
\item Pour séparer du sable et de l’eau, propose une méthode utilisant un filtre.
\item Comment appelle-t-on le liquide recueilli ?
\end{questions}
\end{exercice}\end{blocLivre}
\begin{blocLivre}{exo-6-2}
\begin{exercice}[1]{Nommer les produits}
\pcFigure{figures/3bc7ac9680ed49fa.pdf}{}
\begin{questions}
\item Associe résidu et filtrat aux éléments du schéma.
\item Quel rôle joue le papier-filtre ?
\end{questions}
\end{exercice}\end{blocLivre}
\begin{blocLivre}{exo-6-3}
\begin{exercice}[1]{Laisser reposer}
\begin{questions}
\item Une eau boueuse devient plus limpide au-dessus d’un dépôt. Quelle méthode commence-t-on à utiliser ?
\item Le dépôt s’est-il dissous ?
\end{questions}
\end{exercice}\end{blocLivre}
\begin{blocLivre}{exo-6-4}
\begin{exercice}[1]{Récupérer du sel}
\begin{questions}
\item Quelle opération peut laisser le sel après départ de l’eau d’une solution salée ?
\item L’eau partie dans l’air a-t-elle disparu ?
\end{questions}
\end{exercice}\end{blocLivre}
\begin{blocLivre}{exo-6-5}
\begin{exercice}[1]{Deux changements successifs}
\begin{questions}
\item Dans une distillation, nomme le passage du liquide à la vapeur puis de la vapeur au liquide.
\end{questions}
\end{exercice}\end{blocLivre}
\begin{blocLivre}{exo-6-6}
\rubrique{Je m’entraîne}
\begin{exercice}[2]{Un filtre face au sel}
\begin{questions}
\item Un élève filtre une eau salée limpide. Le filtrat contient-il encore du sel ?
\item Explique pourquoi.
\end{questions}
\end{exercice}\end{blocLivre}
\begin{blocLivre}{exo-6-7}
\begin{exercice}[2]{Comparer deux objectifs}
\begin{questions}
\item Tu veux récupérer le sel d’une eau salée. Quelle méthode simple choisir ?
\item Tu veux récupérer l’eau sous forme liquide. Quelle méthode choisir ?
\end{questions}
\end{exercice}\end{blocLivre}
\begin{blocLivre}{exo-6-8}
\begin{exercice}[2]{Lire un montage}
\pcFigure{figures/9a0a5539496a7e8b.pdf}{}
\begin{questions}
\item Où le liquide se vaporise-t-il ?
\item Quel est le rôle du réfrigérant ?
\end{questions}
\end{exercice}\end{blocLivre}
\begin{blocLivre}{exo-6-9}
\begin{exercice}[2]{Une séparation en deux étapes}
\begin{questions}
\item Pour séparer eau, sable et sel dissous, quelle méthode utiliser d’abord ?
\item Que contient encore le filtrat ?
\end{questions}
\end{exercice}\end{blocLivre}
\begin{blocLivre}{exo-6-10}
\begin{exercice}[2]{Le résultat de l’évaporation}
\begin{questions}
\item Après évaporation complète, on récupère 7,5 g de solide. Quelle masse minimale de substance non volatile était présente dans l’échantillon ?
\item Quelle hypothèse faut-il faire pour parler de toute la masse initialement dissoute ?
\end{questions}
\end{exercice}\end{blocLivre}
\begin{blocLivre}{exo-6-11}
\begin{exercice}[2]{Des volumes de distillat}
\begin{questions}
\item Un montage recueille 45 mL d’eau sur les 60 mL initialement présents. Combien de millilitres n’ont pas été recueillis ?
\item Peut-on affirmer qu’ils ont été détruits ?
\end{questions}
\end{exercice}\end{blocLivre}
\begin{blocLivre}{exo-6-12}
\rubrique{Je raisonne}
\begin{exercice}[3]{Filtrer avant ou après ?}
\begin{questions}
\item Un filtre est bouché rapidement par beaucoup de boue. Quelle étape préalable peut aider ?
\item Pourquoi le résultat doit-il encore être contrôlé ?
\end{questions}
\end{exercice}\end{blocLivre}
\begin{blocLivre}{exo-6-13}
\begin{exercice}[3]{Un mot trompeur}
\begin{questions}
\item Une publicité nomme « eau pure » une eau simplement filtrée. Quelle question scientifique poser ?
\item Propose une formulation plus précise.
\end{questions}
\end{exercice}\end{blocLivre}
\begin{blocLivre}{exo-6-14}
\begin{exercice}[3]{Éviter une mauvaise conclusion}
\begin{questions}
\item Un filtrat est incolore. Donne deux types de constituants qui peuvent encore être présents.
\item Pourquoi ne doit-on pas le goûter ?
\end{questions}
\end{exercice}\end{blocLivre}
\begin{blocLivre}{exo-6-15}
\begin{exercice}[3]{Construire un protocole complet}
\begin{questions}
\item Décris une suite d’opérations pour obtenir du sable sec et du sel à partir d’eau salée contenant du sable.
\item Cite une limite de récupération.
\end{questions}
\end{exercice}\end{blocLivre}
\pcPage{Défi · Prendre une initiative}{Une mini-chaîne de séparation}
\phantomsection\label{defi-ex-6}
\begin{pcFiche}{defi-6}
Un mélange contient de l’eau, du sel dissous, du sable et quelques morceaux de liège qui flottent. Tu dois proposer une suite d’étapes pour récupérer séparément le liège, le sable et le sel. Le matériel et le chauffage sont gérés par le professeur.\par
\pcFigure{figures/5c5e5c92cd065e20.pdf}{}
\begin{questions}
\item Ordonne les opérations en expliquant le rôle de chacune.
\item Précise ce qui se trouve dans le filtrat.
\item Propose une adaptation pour récupérer aussi l’eau.
\end{questions}
\begin{pcEncadre}{method}{Indice progressif}
Commencer par retirer les morceaux faciles à prélever ; séparer ensuite les particules, puis les espèces dissoutes.\par
\end{pcEncadre}
\end{pcFiche}
\clearpage\pagestyle{otMain}
\begin{ouverture}[college1]{7}{Reconnaître une transformation chimique}{Partie 1 · Organisation et transformations de la matière}{Toute apparition de bulles révèle-t-elle une réaction chimique ?}
\objectif{Distinguer mélange, changement d’état et transformation chimique.}
\objectif{Exploiter un test caractéristique fourni avec ses limites.}
\objectif{Raisonner sur la conservation de la masse dans un ensemble fermé.}
\end{ouverture}
\pcPage{L’essentiel}{Chapitre 07 · Reconnaître une transformation chimique}
\phantomsection\label{ch-7-cours-1}
\phantomsection\label{ess-7}
\begin{pcFiche}{cours-7-1}
\essSection{1}{Reconnaître une nouvelle espèce}
Une transformation chimique consomme certaines espèces, les réactifs, et en forme d’autres, les produits. Un changement d’état conserve la nature des espèces ; un mélange n’implique pas à lui seul la formation de nouvelles espèces. Une couleur qui change, une apparition de gaz ou un dépôt constituent des indices à examiner. Ils ne suffisent pas toujours : l’eau qui bout forme des bulles sans devenir une autre substance.\par
\pcFigure{figures/d6c7a20f68a11aef.pdf}{Un gaz se forme et peut être recueilli ; son identification apporte une information supplémentaire.}
\begin{pcEncadre}{definition}{À retenir}
La conclusion doit s’appuyer sur la présence de nouvelles espèces, avec des observations et des tests adaptés. On ne décide pas uniquement à partir d’un effet spectaculaire.\par
\end{pcEncadre}
\begin{pcEncadre}{method}{Une méthode concrète}
Décrire l’état initial et l’état final, puis comparer les espèces présentes. Écrire « cet indice est compatible avec… » si l’identification reste incomplète.\par
\end{pcEncadre}
\end{pcFiche}
\pcPage{L’essentiel}{Chapitre 07 · Reconnaître une transformation chimique}
\phantomsection\label{ch-7-cours-2}
\begin{pcFiche}{cours-7-2}
\essSection{2}{Utiliser une banque de tests}
Un test caractéristique relie une observation à la présence d’une espèce dans des conditions précisées. Dans les situations de ce livre, de l’eau de chaux initialement limpide se trouble en présence de dioxyde de carbone. Une banque peut aussi indiquer que du sulfate de cuivre anhydre blanc devient bleu en présence d’eau. Ce second test est présenté comme document et manipulé uniquement par le professeur avec les précautions adaptées.\par
\pcFigure{figures/d9368b17583b9996.pdf}{Test du dioxyde de carbone dans les conditions scolaires fournies : apparition d’un trouble blanc.}
\begin{pcEncadre}{definition}{À retenir}
Un test positif indique une présence ; il ne démontre pas que l’échantillon est pur. Un résultat négatif s’interprète avec la sensibilité et les conditions du test.\par
\end{pcEncadre}
\begin{pcEncadre}{method}{Une méthode concrète}
Rédiger trois étapes : j’observe le résultat ; je cite le test de la banque ; j’en déduis la présence de l’espèce recherchée.\par
\end{pcEncadre}
\end{pcFiche}
\pcPage{L’essentiel}{Chapitre 07 · Reconnaître une transformation chimique}
\phantomsection\label{ch-7-cours-3}
\begin{pcFiche}{cours-7-3}
\essSection{3}{Compter toute la matière}
Lors d’une transformation chimique dans un ensemble fermé, la masse totale se conserve. Un gaz formé fait partie des produits et doit être compté. Dans un récipient ouvert, il peut partir dans l’air : la masse du contenu restant diminue sans contredire la conservation. Une combustion est un exemple de transformation chimique qui consomme un combustible et du dioxygène ; certains combustibles produisent notamment du dioxyde de carbone.\par
\pcFigure{figures/d6c7a20f68a11aef.pdf}{Le gaz est retenu par un ballon souple adapté, sous le contrôle du professeur. Ne jamais enfermer une production de gaz dans un récipient rigide bouché.}
\begin{pcEncadre}{definition}{À retenir}
Une conservation de masse ne signifie pas que les espèces restent les mêmes. Une nouvelle espèce peut apparaître tout en conservant la masse totale.\par
\end{pcEncadre}
\begin{pcEncadre}{method}{Une méthode concrète}
Définir le système pesé et rechercher toute entrée ou sortie de matière. Une combustion utilise aussi un réactif provenant de l’air, le dioxygène.\par
\end{pcEncadre}
\end{pcFiche}
\pcPage{Activité expérimentale}{Une enquête à partir d’un gaz}
\phantomsection\label{ch-7-activite}
\begin{pcFiche}{activite-7}
\begin{pcEncadre}{method}{Mission}
Comparer deux documents pour identifier une transformation.\par
\end{pcEncadre}
\begin{questions}
\item Document A : de l’eau bout ; la vapeur refroidie redonne de l’eau liquide. Document B : du vinaigre réagit avec de l’hydrogénocarbonate de sodium et un gaz se forme.
\item Dans l’essai B réalisé par le professeur, le gaz mis en contact avec de l’eau de chaux initialement limpide produit un trouble blanc.
\item Utiliser la banque de tests pour distinguer ce qui est observé, ce qui est identifié et ce qui peut être conclu.
\end{questions}
\pcFigure{figures/b594f708d644c622.pdf}{Documents d’expériences encadrées ; aucun test avec une flamme n’est demandé.}
\begin{questions}
\item Quel gaz est mis en évidence dans le document B ?
\item Pourquoi les bulles du document A ne prouvent-elles pas une transformation chimique ?
\item Rédige une conclusion argumentée sur les deux situations.
\end{questions}
\end{pcFiche}
\pcPage{L’essentiel à retenir}{Chapitre 07 · Synthèse et autoévaluation}
\phantomsection\label{ch-7-bilan}
\begin{pcFiche}{synthese-7}
\begin{pcEncadre}{definition}{Repère 1}
Une transformation chimique forme de nouvelles espèces.\par
\end{pcEncadre}
\begin{pcEncadre}{method}{Repère 2}
Les espèces consommées sont les réactifs ; les espèces formées sont les produits.\par
\end{pcEncadre}
\begin{pcEncadre}{definition}{Repère 3}
Un test caractéristique s’utilise avec une banque et des conditions données.\par
\end{pcEncadre}
\begin{pcEncadre}{method}{Repère 4}
Un test positif ne prouve pas la pureté d’un échantillon.\par
\end{pcEncadre}
\begin{pcEncadre}{definition}{Repère 5}
Dans un ensemble fermé, la masse totale se conserve, gaz compris.\par
\end{pcEncadre}
\begin{pcEncadre}{attention}{Une erreur à éviter}
« La masse baisse, donc la matière est détruite. » Vérifier d’abord si un gaz ou un liquide a quitté le système pesé.\par
\end{pcEncadre}
\end{pcFiche}
\clearpage\markright{Exercices}
\begin{blocLivre}{exo-7-1}
\rubrique{Je démarre}
\begin{exercice}[1]{Nommer réactif et produit}
\begin{questions}
\item Dans une transformation, comment nomme-t-on une espèce consommée ?
\item Comment nomme-t-on une espèce formée ?
\end{questions}
\end{exercice}\end{blocLivre}
\begin{blocLivre}{exo-7-2}
\begin{exercice}[1]{Trois situations}
\begin{questions}
\item Classe : fusion d’un glaçon ; mélange de sable et d’eau ; formation de rouille sur du fer.
\end{questions}
\end{exercice}\end{blocLivre}
\begin{blocLivre}{exo-7-3}
\begin{exercice}[1]{Un test fourni}
\begin{questions}
\item L’eau de chaux se trouble au contact d’un gaz. Quelle présence indique la banque de tests du cours ?
\item Le gaz est-il forcément pur ?
\end{questions}
\end{exercice}\end{blocLivre}
\begin{blocLivre}{exo-7-4}
\begin{exercice}[1]{Des bulles sans réaction}
\begin{questions}
\item Donne un exemple de formation de bulles qui ne constitue pas une transformation chimique.
\item Explique ton choix.
\end{questions}
\end{exercice}\end{blocLivre}
\begin{blocLivre}{exo-7-5}
\begin{exercice}[1]{Rendre compte d’un test}
\pcFigure{figures/b594f708d644c622.pdf}{}
\begin{questions}
\item Écris une phrase complète d’identification à partir du schéma.
\end{questions}
\end{exercice}\end{blocLivre}
\begin{blocLivre}{exo-7-6}
\rubrique{Je m’entraîne}
\begin{exercice}[2]{Peser un ensemble fermé}
\begin{questions}
\item Avant mélange, l’ensemble fermé pèse 152,0 g. Un gaz se forme sans fuite. Quelle masse attend-on après ?
\item Que doit-on inclure dans la pesée ?
\end{questions}
\end{exercice}\end{blocLivre}
\begin{blocLivre}{exo-7-7}
\begin{exercice}[2]{Un gaz qui part}
\begin{questions}
\item Un récipient ouvert passe de 98,4 g à 97,8 g pendant une réaction. Quelle masse a quitté le système si seul le gaz s’échappe ?
\item Pourquoi cette condition est-elle importante ?
\end{questions}
\end{exercice}\end{blocLivre}
\begin{blocLivre}{exo-7-8}
\begin{exercice}[2]{Interpréter un test de l’eau}
\pcFigure{figures/ceaa2beaf10e86f3.pdf}{Document de travail · modèle simplifié.}
\begin{questions}
\item Un document indique : poudre anhydre blanche avant le test, bleue après contact avec un liquide. Quelle présence cela met-il en évidence selon la banque ?
\item Peut-on conclure « ce liquide est uniquement de l’eau » ?
\end{questions}
\end{exercice}\end{blocLivre}
\begin{blocLivre}{exo-7-9}
\begin{exercice}[2]{Un ballon qui gonfle}
\pcFigure{figures/420b2fff6466bf8c.pdf}{Document de travail · modèle simplifié.}
\begin{questions}
\item Pourquoi placer un ballon souple adapté sur le récipient dans une expérience encadrée ?
\item Quel mauvais montage faut-il éviter ?
\end{questions}
\end{exercice}\end{blocLivre}
\begin{blocLivre}{exo-7-10}
\begin{exercice}[2]{Le rôle de l’air}
\begin{questions}
\item Une combustion utilise un combustible. Quel autre réactif de l’air est habituellement nécessaire ?
\item Pourquoi faut-il en tenir compte dans un bilan de masse ?
\end{questions}
\end{exercice}\end{blocLivre}
\begin{blocLivre}{exo-7-11}
\begin{exercice}[2]{Présence ou absence ?}
\begin{questions}
\item Un test ne montre aucun changement visible. Peut-on conclure à une absence absolue de l’espèce ?
\item Cite une vérification utile.
\end{questions}
\end{exercice}\end{blocLivre}
\begin{blocLivre}{exo-7-12}
\rubrique{Je raisonne}
\begin{exercice}[3]{Un changement de couleur}
\begin{questions}
\item Un colorant rouge se disperse dans l’eau et tout le liquide devient rose. La nouvelle couleur suffit-elle à prouver une réaction ?
\item Quelle autre interprétation proposer ?
\end{questions}
\end{exercice}\end{blocLivre}
\begin{blocLivre}{exo-7-13}
\begin{exercice}[3]{Une conclusion trop large}
\begin{questions}
\item « L’eau de chaux est trouble, donc le flacon ne contient que du dioxyde de carbone. » Corrige la conclusion.
\item Quelle observation doit précéder la conclusion ?
\end{questions}
\end{exercice}\end{blocLivre}
\begin{blocLivre}{exo-7-14}
\begin{exercice}[3]{Comparer ouvert et fermé}
\begin{questions}
\item Dans A, le gaz s’échappe ; dans B, il reste retenu. Quelle masse est conservée dans chaque situation ?
\item Quelle masse la balance compare-t-elle réellement ?
\end{questions}
\end{exercice}\end{blocLivre}
\begin{blocLivre}{exo-7-15}
\begin{exercice}[3]{Une preuve construite}
\begin{questions}
\item Propose un raisonnement pour distinguer l’ébullition de l’eau et la production de dioxyde de carbone par une réaction.
\item Pourquoi faut-il éviter le seul critère « je vois des bulles » ?
\end{questions}
\end{exercice}\end{blocLivre}
\pcPage{Défi · Prendre une initiative}{Le mystère des 1,2 gramme}
\phantomsection\label{defi-ex-7}
\begin{pcFiche}{defi-7}
Un document annonce une baisse de 1,2 g dans un récipient ouvert pendant une réaction qui dégage un gaz. Un autre groupe obtient la même masse avant et après avec un dispositif fermé adapté. Les deux groupes s’accusent d’avoir faux.\par
\pcFigure{figures/2ee068bf8af22dde.pdf}{}
\begin{questions}
\item Explique comment les deux résultats peuvent être compatibles.
\item Précise la masse du gaz parti si aucune autre perte n’a eu lieu.
\item Propose deux contrôles pour améliorer la comparaison.
\end{questions}
\begin{pcEncadre}{method}{Indice progressif}
Dessiner la frontière du système réellement posé sur la balance.\par
\end{pcEncadre}
\end{pcFiche}
\clearpage\pagestyle{otMain}
\begin{ouverture}[college1]{8}{La Terre dans le Système solaire}{Partie 1 · Organisation et transformations de la matière}{Comment représenter des distances aussi grandes ?}
\objectif{Distinguer étoile, planète et satellite naturel.}
\objectif{Se repérer dans le Système solaire et distinguer rotation et révolution.}
\objectif{Lire des distances astronomiques et construire une échelle simple.}
\end{ouverture}
\pcPage{L’essentiel}{Chapitre 08 · La Terre dans le Système solaire}
\phantomsection\label{ch-8-cours-1}
\phantomsection\label{ess-8}
\begin{pcFiche}{cours-8-1}
\essSection{1}{Notre place dans le Système solaire}
Le Soleil est une étoile : il émet sa propre lumière. Les huit planètes tournent autour de lui et diffusent la lumière qu’elles reçoivent. Dans l’ordre d’éloignement moyen : Mercure, Vénus, Terre, Mars, Jupiter, Saturne, Uranus et Neptune. La Lune est un satellite naturel de la Terre. Le Système solaire comprend aussi des planètes naines, des astéroïdes et des comètes ; il appartient à la galaxie appelée Voie lactée.\par
\pcFigure{figures/72652ef79f2a6bc0.pdf}{Représentation d’ordre : les tailles, les intervalles et les positions ne respectent pas une échelle commune.}
\begin{pcEncadre}{definition}{À retenir}
Le Soleil n’est pas une planète. La Lune n’est pas une étoile et ne produit pas la lumière visible qu’elle nous renvoie.\par
\end{pcEncadre}
\begin{pcEncadre}{method}{Une méthode concrète}
Distinguer le nom de l’objet, sa catégorie et l’astre autour duquel il se déplace. Ne pas lire une distance sur une illustration annoncée sans échelle.\par
\end{pcEncadre}
\end{pcFiche}
\pcPage{L’essentiel}{Chapitre 08 · La Terre dans le Système solaire}
\phantomsection\label{ch-8-cours-2}
\begin{pcFiche}{cours-8-2}
\essSection{2}{Rotation, révolution et alternance jour-nuit}
La Terre tourne sur elle-même : c’est sa rotation. Elle effectue aussi une révolution autour du Soleil en environ une année. L’alternance du jour et de la nuit est liée à la rotation terrestre : à un instant donné, le Soleil éclaire approximativement une moitié de la Terre. Une durée proche de 24 h décrit un jour solaire ; ces ordres de grandeur suffisent ici pour distinguer les mouvements.\par
\pcFigure{figures/2d6fcc9ad0ced6f0.pdf}{La moitié tournée vers le Soleil est éclairée ; l’autre est dans la nuit. Les tailles et distances ne sont pas à l’échelle.}
\begin{pcEncadre}{definition}{À retenir}
Le Soleil ne s’éteint pas chaque soir. Un lieu terrestre passe progressivement de la partie éclairée à la partie non éclairée.\par
\end{pcEncadre}
\begin{pcEncadre}{method}{Une méthode concrète}
Pour expliquer le jour et la nuit, faire tourner un globe devant une lampe fixe. Ne pas confondre ce mouvement avec la révolution annuelle.\par
\end{pcEncadre}
\end{pcFiche}
\pcPage{L’essentiel}{Chapitre 08 · La Terre dans le Système solaire}
\phantomsection\label{ch-8-cours-3}
\begin{pcFiche}{cours-8-3}
\essSection{3}{Choisir une échelle de distance}
Les distances du Système solaire s’expriment souvent en millions de kilomètres. La distance moyenne Terre-Soleil est voisine de 150 millions de kilomètres ; elle donne le repère d’environ une unité astronomique (ua). Une représentation à l’échelle utilise le même rapport pour toutes les distances représentées. Une échelle pratique pour les planètes proches devient vite trop grande pour inclure Neptune sur une feuille.\par
\pcFigure{figures/23e8d10958f0f3b7.pdf}{Distances moyennes arrondies depuis le Soleil : Mercure 58, Vénus 108, Terre 150, Mars 228 millions de km.}
\begin{pcEncadre}{definition}{À retenir}
Une distance moyenne au Soleil n’est pas une distance fixe entre deux planètes. Leurs positions varient au cours de leurs révolutions.\par
\end{pcEncadre}
\begin{pcEncadre}{method}{Une méthode concrète}
Si 1 cm représente 50 millions de km, la Terre est placée à 150 \(\div\) 50 = 3 cm du Soleil. Appliquer le même coefficient à tous les objets du modèle.\par
\end{pcEncadre}
\end{pcFiche}
\pcPage{Activité expérimentale}{Une maquette sur une bande de papier}
\phantomsection\label{ch-8-activite}
\begin{pcFiche}{activite-8}
\begin{pcEncadre}{method}{Mission}
Placer les quatre planètes internes à une même échelle de distance.\par
\end{pcEncadre}
\begin{questions}
\item Utiliser les distances moyennes arrondies : 58, 108, 150 et 228 millions de km pour Mercure, Vénus, Terre et Mars.
\item Choisir 1 cm pour 50 millions de km et calculer la position de chaque planète depuis le Soleil.
\item Tracer une ligne, placer le Soleil à l’origine et reporter les distances calculées. Les disques des planètes servent seulement de repères.
\end{questions}
\pcFigure{figures/0e02bcf845d3f6a7.pdf}{Le schéma du livre emploie une échelle graphique ; ta maquette utilise celle annoncée dans l’activité.}
\begin{questions}
\item À quelles distances du Soleil places-tu les quatre repères ?
\item Jupiter est à environ 780 millions de km du Soleil. Tiendrait-il sur une bande de 20 cm ?
\item Pourquoi faut-il préciser que les diamètres dessinés ne sont pas à cette échelle ?
\end{questions}
\end{pcFiche}
\pcPage{L’essentiel à retenir}{Chapitre 08 · Synthèse et autoévaluation}
\phantomsection\label{ch-8-bilan}
\begin{pcFiche}{synthese-8}
\begin{pcEncadre}{definition}{Repère 1}
Le Soleil est l’étoile du Système solaire.\par
\end{pcEncadre}
\begin{pcEncadre}{method}{Repère 2}
Huit planètes gravitent autour de lui ; la Lune est un satellite de la Terre.\par
\end{pcEncadre}
\begin{pcEncadre}{definition}{Repère 3}
La rotation de la Terre explique l’alternance du jour et de la nuit.\par
\end{pcEncadre}
\begin{pcEncadre}{method}{Repère 4}
La révolution de la Terre autour du Soleil dure environ une année.\par
\end{pcEncadre}
\begin{pcEncadre}{definition}{Repère 5}
Une échelle doit être annoncée et appliquée de manière identique aux distances.\par
\end{pcEncadre}
\begin{pcEncadre}{attention}{Une erreur à éviter}
« Mars est à 228 millions de km de la Terre. » Cette valeur arrondie concerne sa distance moyenne au Soleil, pas sa distance à la Terre.\par
\end{pcEncadre}
\end{pcFiche}
\clearpage\markright{Exercices}
\begin{blocLivre}{exo-8-1}
\rubrique{Je démarre}
\begin{exercice}[1]{Classer les astres}
\begin{questions}
\item Classe le Soleil, la Terre et la Lune : étoile, planète ou satellite naturel.
\end{questions}
\end{exercice}\end{blocLivre}
\begin{blocLivre}{exo-8-2}
\begin{exercice}[1]{Retrouver l’ordre}
\begin{questions}
\item Nomme les quatre planètes les plus proches du Soleil dans l’ordre.
\item Quelle planète vient après Mars ?
\end{questions}
\end{exercice}\end{blocLivre}
\begin{blocLivre}{exo-8-3}
\begin{exercice}[1]{Identifier la source de lumière}
\begin{questions}
\item Pourquoi voit-on la Lune lorsqu’elle est éclairée ?
\item Est-elle une source primaire de lumière visible ?
\end{questions}
\end{exercice}\end{blocLivre}
\begin{blocLivre}{exo-8-4}
\begin{exercice}[1]{Rotation ou révolution ?}
\begin{questions}
\item Quel mouvement explique l’alternance jour-nuit ?
\item Quel mouvement terrestre dure environ une année ?
\end{questions}
\end{exercice}\end{blocLivre}
\begin{blocLivre}{exo-8-5}
\begin{exercice}[1]{Lire une représentation}
\pcFigure{figures/d6ce308df0a77bcf.pdf}{}
\begin{questions}
\item Peut-on mesurer une distance réelle avec une règle sur ce dessin ?
\item Quelle indication l’empêche ?
\end{questions}
\end{exercice}\end{blocLivre}
\begin{blocLivre}{exo-8-6}
\rubrique{Je m’entraîne}
\begin{exercice}[2]{Une unité de distance}
\begin{questions}
\item Écris 150 millions de km en kilomètres.
\item Cette valeur mesure-t-elle une durée ?
\end{questions}
\end{exercice}\end{blocLivre}
\begin{blocLivre}{exo-8-7}
\begin{exercice}[2]{Placer la Terre}
\begin{questions}
\item À l’échelle 1 cm pour 50 millions de km, à quelle distance du Soleil placer la Terre ?
\item Précise l’unité de la mesure sur la feuille.
\end{questions}
\end{exercice}\end{blocLivre}
\begin{blocLivre}{exo-8-8}
\begin{exercice}[2]{Placer Mars}
\begin{questions}
\item À la même échelle, calcule la position de Mars, située à 228 millions de km du Soleil.
\item Est-elle plus de quatre fois plus loin que le repère d’un centimètre ?
\end{questions}
\end{exercice}\end{blocLivre}
\begin{blocLivre}{exo-8-9}
\begin{exercice}[2]{Jour et nuit}
\pcFigure{figures/756e7a4628904c01.pdf}{}
\begin{questions}
\item Explique pourquoi deux lieux très éloignés peuvent être au jour et dans la nuit au même instant.
\item Faut-il que le Soleil s’éteigne pour produire la nuit ?
\end{questions}
\end{exercice}\end{blocLivre}
\begin{blocLivre}{exo-8-10}
\begin{exercice}[2]{Comparer des distances}
\begin{questions}
\item La Terre est à environ 150 millions de km du Soleil, Jupiter à 780 millions. Combien de fois la seconde distance représente-t-elle la première ?
\item Cette comparaison donne-t-elle la distance Terre-Jupiter ?
\end{questions}
\end{exercice}\end{blocLivre}
\begin{blocLivre}{exo-8-11}
\begin{exercice}[2]{Lire un repère en ua}
\begin{questions}
\item On donne 1 ua \(\approx\) 150 millions de km. Une distance de 3 ua correspond à combien de millions de km ?
\item Pourquoi utiliser ce repère ?
\end{questions}
\end{exercice}\end{blocLivre}
\begin{blocLivre}{exo-8-12}
\rubrique{Je raisonne}
\begin{exercice}[3]{Le modèle trop petit}
\begin{questions}
\item Neptune est à environ 4 500 millions de km du Soleil. Quelle longueur faut-il à l’échelle 1 cm pour 50 millions de km ?
\item Une feuille de 21 cm de large suffit-elle ?
\end{questions}
\end{exercice}\end{blocLivre}
\begin{blocLivre}{exo-8-13}
\begin{exercice}[3]{Une échelle incohérente}
\begin{questions}
\item Un élève place la Terre à 3 cm pour 150 millions de km et Mars à 3 cm pour 228 millions de km. Pourquoi la maquette est-elle fausse ?
\item Comment corriger Mars ?
\end{questions}
\end{exercice}\end{blocLivre}
\begin{blocLivre}{exo-8-14}
\begin{exercice}[3]{Relier plusieurs mouvements}
\begin{questions}
\item Une personne assise dans son jardin semble immobile. Cite deux mouvements qu’elle accompagne à l’échelle astronomique.
\item Pourquoi faut-il préciser l’observateur ?
\end{questions}
\end{exercice}\end{blocLivre}
\begin{blocLivre}{exo-8-15}
\begin{exercice}[3]{Évaluer une source}
\begin{questions}
\item Un document mélange planètes et étoiles dans une liste sans date ni auteur. Quels contrôles effectuer avant de l’utiliser ?
\item Cite un organisme pertinent pour vérifier les informations.
\end{questions}
\end{exercice}\end{blocLivre}
\pcPage{Défi · Prendre une initiative}{Le Système solaire dans la cour}
\phantomsection\label{defi-ex-8}
\begin{pcFiche}{defi-8}
Tu disposes d’une longueur de 45 m pour placer Neptune, à environ 4 500 millions de km du Soleil. Tu souhaites une seule échelle pour toutes les distances. Terre : 150 millions de km ; Mars : 228 millions de km ; Jupiter : 780 millions de km.\par
\pcFigure{figures/3516b1ebc9998a8e.pdf}{}
\begin{questions}
\item Quelle distance réelle doit représenter 1 m dans la cour ?
\item Place la Terre, Mars et Jupiter à cette échelle.
\item Explique pourquoi les diamètres des planètes devront être traités séparément.
\end{questions}
\begin{pcEncadre}{method}{Indice progressif}
Diviser la distance de Neptune par 45 ; utiliser ensuite la même proportion pour chaque planète.\par
\end{pcEncadre}
\end{pcFiche}
\clearpage\pagestyle{otMain}
\begin{ouverture}[college2]{9}{Décrire un mouvement}{Partie 2 · Mouvement et interactions}{Peut-on être immobile et en mouvement à la fois ?}
\objectif{Préciser l’objet et l’observateur pour décrire un mouvement.}
\objectif{Reconnaître des trajectoires rectilignes, circulaires et curvilignes.}
\objectif{Exploiter des positions relevées à intervalles réguliers.}
\end{ouverture}
\pcPage{L’essentiel}{Chapitre 09 · Décrire un mouvement}
\phantomsection\label{ch-9-cours-1}
\phantomsection\label{ess-9}
\begin{pcFiche}{cours-9-1}
\essSection{1}{Le mouvement dépend de l’observateur}
Décrire un mouvement demande de choisir l’objet étudié et le repère de l’observateur. Un voyageur assis dans un train est immobile par rapport au siège, mais se déplace par rapport au quai lorsque le train avance. Les deux descriptions sont compatibles. On étudie souvent un point représentatif de l’objet, par exemple le centre d’un ballon ; les autres points peuvent suivre des trajectoires différentes.\par
\pcFigure{figures/d1d94dacbbe77273.pdf}{Le voyageur est au repos par rapport au train et en mouvement par rapport au quai.}
\begin{pcEncadre}{definition}{À retenir}
Il n’existe pas de description complète sans référence à un observateur ou à un repère. Écrire simplement « il bouge » laisse une information essentielle de côté.\par
\end{pcEncadre}
\begin{pcEncadre}{method}{Une méthode concrète}
Commencer la phrase par « Par rapport à… ». Nommer ensuite l’objet et décrire comment sa position évolue.\par
\end{pcEncadre}
\end{pcFiche}
\pcPage{L’essentiel}{Chapitre 09 · Décrire un mouvement}
\phantomsection\label{ch-9-cours-2}
\begin{pcFiche}{cours-9-2}
\essSection{2}{Nommer une trajectoire}
La trajectoire est l’ensemble des positions successives du point étudié dans le repère choisi. Elle est rectiligne si ces positions appartiennent à une droite ; circulaire si elles appartiennent à un cercle ; curviligne si elles suivent une courbe. Une trajectoire renseigne sur la forme du chemin, mais pas à elle seule sur la rapidité du déplacement. Une même trajectoire peut être parcourue dans deux sens.\par
\pcFigure{figures/0dd992971694c19c.pdf}{Trois formes de trajectoire ; la flèche précise le sens du parcours représenté.}
\begin{pcEncadre}{definition}{À retenir}
Rectiligne et uniforme ne sont pas synonymes. Le premier mot décrit la trajectoire ; le second décrit la constance de la valeur de la vitesse.\par
\end{pcEncadre}
\begin{pcEncadre}{method}{Une méthode concrète}
Repérer plusieurs positions dans le même repère. Décrire séparément la forme de la trajectoire et le sens du déplacement.\par
\end{pcEncadre}
\end{pcFiche}
\pcPage{L’essentiel}{Chapitre 09 · Décrire un mouvement}
\phantomsection\label{ch-9-cours-3}
\begin{pcFiche}{cours-9-3}
\essSection{3}{Comparer les positions dans le temps}
Une chronophotographie superpose des positions à des dates connues. Pour comparer la valeur de la vitesse, on peut relever les positions à intervalles de temps égaux. Des espacements égaux indiquent alors un mouvement de valeur de vitesse constante. Des espacements croissants correspondent à une augmentation de la valeur de la vitesse sur les intervalles observés ; des espacements décroissants à une diminution.\par
\pcFigure{figures/d28706f310244d89.pdf}{Pour chaque ligne, les dates 0, 1, 2, 3 et 4 sont séparées par la même durée.}
\begin{pcEncadre}{definition}{À retenir}
L’espacement ne renseigne sur l’évolution de la vitesse que si l’on connaît les durées correspondantes. Un mouvement circulaire peut avoir une valeur de vitesse constante tout en changeant de direction.\par
\end{pcEncadre}
\begin{pcEncadre}{method}{Une méthode concrète}
Vérifier l’égalité des intervalles de temps avant de comparer les distances entre positions successives.\par
\end{pcEncadre}
\end{pcFiche}
\pcPage{Activité expérimentale}{Photographier un petit mobile}
\phantomsection\label{ch-9-activite}
\begin{pcFiche}{activite-9}
\begin{pcEncadre}{method}{Mission}
Décrire un déplacement avec une série de positions datées.\par
\end{pcEncadre}
\begin{questions}
\item En classe, filmer un mobile sur une piste droite avec un appareil fixe. Garder un repère de longueur dans le plan du mouvement.
\item Extraire cinq positions à intervalles de temps égaux. Identifier le même point sur chaque image.
\item Comparer les espacements. Pour un entraînement sans vidéo, analyser les trois lignes du schéma.
\end{questions}
\pcFigure{figures/b9d2213e035c1cd3.pdf}{On compare les espacements à l’intérieur d’une même ligne ; l’origine des positions est arbitraire.}
\begin{questions}
\item Quelle ligne présente des espacements constants ? Lesquels augmentent ou diminuent ?
\item Pourquoi l’appareil doit-il rester fixe par rapport à la piste ?
\item Rédige une description complète d’une ligne avec trajectoire et évolution de la vitesse.
\end{questions}
\end{pcFiche}
\pcPage{L’essentiel à retenir}{Chapitre 09 · Synthèse et autoévaluation}
\phantomsection\label{ch-9-bilan}
\begin{pcFiche}{synthese-9}
\begin{pcEncadre}{definition}{Repère 1}
Décrire un mouvement exige un objet étudié et un repère.\par
\end{pcEncadre}
\begin{pcEncadre}{method}{Repère 2}
La trajectoire est l’ensemble des positions successives du point choisi.\par
\end{pcEncadre}
\begin{pcEncadre}{definition}{Repère 3}
Rectiligne, circulaire et curviligne décrivent la forme de la trajectoire.\par
\end{pcEncadre}
\begin{pcEncadre}{method}{Repère 4}
Le sens du parcours est distinct de la forme du chemin.\par
\end{pcEncadre}
\begin{pcEncadre}{definition}{Repère 5}
À durées égales, les espacements renseignent sur l’évolution de la valeur de la vitesse.\par
\end{pcEncadre}
\begin{pcEncadre}{attention}{Une erreur à éviter}
« Les points sont plus éloignés, donc le mobile va plus vite. » Vérifier aussi les durées séparant les positions.\par
\end{pcEncadre}
\end{pcFiche}
\clearpage\markright{Exercices}
\begin{blocLivre}{exo-9-1}
\rubrique{Je démarre}
\begin{exercice}[1]{Préciser le repère}
\begin{questions}
\item Une élève est assise dans un bus qui roule. Est-elle immobile par rapport à son siège ?
\item Est-elle immobile par rapport à un arrêt de bus ?
\end{questions}
\end{exercice}\end{blocLivre}
\begin{blocLivre}{exo-9-2}
\begin{exercice}[1]{Nommer un chemin}
\begin{questions}
\item Un point suit une droite. Comment nomme-t-on sa trajectoire ?
\item Un point suit un cercle. Comment nomme-t-on sa trajectoire ?
\end{questions}
\end{exercice}\end{blocLivre}
\begin{blocLivre}{exo-9-3}
\begin{exercice}[1]{Lire le sens}
\pcFigure{figures/07bfff6c23b887fa.pdf}{Document de travail · modèle simplifié.}
\begin{questions}
\item Deux cyclistes parcourent la même ligne droite en sens opposés. Leur trajectoire peut-elle avoir la même forme ?
\item Quelle information faut-il ajouter pour les distinguer ?
\end{questions}
\end{exercice}\end{blocLivre}
\begin{blocLivre}{exo-9-4}
\begin{exercice}[1]{Un passager observé}
\pcFigure{figures/5f58c69c5265b1b2.pdf}{}
\begin{questions}
\item Nomme deux repères possibles pour le voyageur du schéma.
\item Décris son mouvement par rapport à chacun.
\end{questions}
\end{exercice}\end{blocLivre}
\begin{blocLivre}{exo-9-5}
\begin{exercice}[1]{Des mots différents}
\begin{questions}
\item « Rectiligne » décrit-il la rapidité ou la forme de la trajectoire ?
\item Que décrit « uniforme » lorsqu’on parle de la valeur de la vitesse ?
\end{questions}
\end{exercice}\end{blocLivre}
\begin{blocLivre}{exo-9-6}
\rubrique{Je m’entraîne}
\begin{exercice}[2]{Des positions régulières}
\pcFigure{figures/2f5fc91e576092f8.pdf}{Document de travail · modèle simplifié.}
\begin{questions}
\item Des positions sont séparées de 2 cm toutes les 0,5 s. Que peux-tu dire de la valeur de la vitesse ?
\item Quelle condition sur les temps est essentielle ?
\end{questions}
\end{exercice}\end{blocLivre}
\begin{blocLivre}{exo-9-7}
\begin{exercice}[2]{Trois lignes à comparer}
\pcFigure{figures/b9d2213e035c1cd3.pdf}{}
\begin{questions}
\item Classe A, B et C selon l’évolution de la valeur de la vitesse.
\item Cite les indices utilisés.
\end{questions}
\end{exercice}\end{blocLivre}
\begin{blocLivre}{exo-9-8}
\begin{exercice}[2]{Un tour de manège}
\begin{questions}
\item Dans le repère du sol, quelle est la trajectoire d’un siège de manège tournant autour d’un axe fixe ?
\item La direction du déplacement reste-t-elle la même tout le tour ?
\end{questions}
\end{exercice}\end{blocLivre}
\begin{blocLivre}{exo-9-9}
\begin{exercice}[2]{Le bon point à suivre}
\begin{questions}
\item Pourquoi faut-il suivre le même point d’un ballon sur toute une série d’images ?
\item Quel point simple peut-on choisir pour décrire le déplacement global ?
\end{questions}
\end{exercice}\end{blocLivre}
\begin{blocLivre}{exo-9-10}
\begin{exercice}[2]{Des durées inégales}
\begin{questions}
\item Un point parcourt 1 m en 1 s, puis 2 m en 2 s. Le second intervalle prouve-t-il une accélération ?
\item Compare les distances parcourues par seconde.
\end{questions}
\end{exercice}\end{blocLivre}
\begin{blocLivre}{exo-9-11}
\begin{exercice}[2]{Une description complète}
\pcFigure{figures/0ee4546ede86e6f7.pdf}{Document de travail · modèle simplifié.}
\begin{questions}
\item Un mobile suit une ligne droite vers la droite ; les espacements augmentent à durées égales. Rédige une description.
\end{questions}
\end{exercice}\end{blocLivre}
\begin{blocLivre}{exo-9-12}
\rubrique{Je raisonne}
\begin{exercice}[3]{L’appareil se déplace}
\begin{questions}
\item Un élève marche en filmant un autre élève immobile au sol. Pourquoi l’image peut-elle montrer un déplacement apparent ?
\item Quel repère permettrait de dire que le second élève est au repos ?
\end{questions}
\end{exercice}\end{blocLivre}
\begin{blocLivre}{exo-9-13}
\begin{exercice}[3]{Une photographie unique}
\begin{questions}
\item Peut-on déterminer l’évolution de la vitesse à partir d’une seule photo d’un mobile sans autre information ?
\item Quelles données supplémentaires demander ?
\end{questions}
\end{exercice}\end{blocLivre}
\begin{blocLivre}{exo-9-14}
\begin{exercice}[3]{Une hypothèse à tester}
\begin{questions}
\item Une balle parcourt des distances égales sur trois images prises à des dates inconnues. Peut-on conclure à un mouvement uniforme ?
\item Quelle vérification faut-il obtenir ?
\end{questions}
\end{exercice}\end{blocLivre}
\begin{blocLivre}{exo-9-15}
\begin{exercice}[3]{Comparer deux observateurs}
\begin{questions}
\item Deux amis restent assis face à face dans un train en mouvement. L’un affirme « nous ne bougeons pas » ; un passant affirme le contraire. Qui a raison ?
\item Formule les deux phrases sans ambiguïté.
\end{questions}
\end{exercice}\end{blocLivre}
\pcPage{Défi · Prendre une initiative}{Le reportage de deux témoins}
\phantomsection\label{defi-ex-9}
\begin{pcFiche}{defi-9}
Une balle est tenue immobile par un voyageur dans un train avançant en ligne droite. Une caméra est fixée dans le train ; une autre est fixée sur le quai. Prépare deux commentaires scientifiques compatibles.\par
\pcFigure{figures/f552d736b6238df0.pdf}{}
\begin{questions}
\item Choisis le point étudié et le repère de chaque caméra.
\item Décris ce que chaque caméra peut relever.
\item Ajoute ce qui changerait si le voyageur lâchait la balle, sans prévoir de calcul.
\end{questions}
\begin{pcEncadre}{method}{Indice progressif}
Séparer la situation balle tenue de la situation après lâcher ; ne pas mélanger les repères.\par
\end{pcEncadre}
\end{pcFiche}
\clearpage\pagestyle{otMain}
\begin{ouverture}[college2]{10}{Mesurer et calculer une vitesse}{Partie 2 · Mouvement et interactions}{Comment comparer deux trajets de durées différentes ?}
\objectif{Calculer une vitesse moyenne à partir d’une distance parcourue et d’une durée.}
\objectif{Utiliser des unités cohérentes et convertir les durées.}
\objectif{Distinguer vitesse moyenne et variation de vitesse pendant un trajet.}
\end{ouverture}
\pcPage{L’essentiel}{Chapitre 10 · Mesurer et calculer une vitesse}
\phantomsection\label{ch-10-cours-1}
\phantomsection\label{ess-10}
\begin{pcFiche}{cours-10-1}
\essSection{1}{Distance parcourue, durée et vitesse}
La vitesse moyenne compare la distance parcourue à la durée du trajet : vitesse moyenne = distance parcourue \(\div\) durée. Avec des mètres et des secondes, elle s’exprime en mètres par seconde (m/s). Avec des kilomètres et des heures, elle s’exprime en kilomètres par heure (km/h). Si la valeur de la vitesse reste constante, ce quotient donne aussi cette valeur constante. Sur un trajet varié, il résume le trajet entier.\par
\pcFigure{figures/dcb60dc104f701dd.pdf}{Exemple de mouvement uniforme : 3 m sont parcourus chaque seconde ; d désigne la distance depuis le départ.}
\begin{pcEncadre}{definition}{À retenir}
Une vitesse de 3 m/s signifie 3 m parcourus par seconde dans un mouvement uniforme. La distance parcourue suit le trajet, pas forcément la ligne droite entre départ et arrivée.\par
\end{pcEncadre}
\begin{pcEncadre}{method}{Une méthode concrète}
Écrire les données, vérifier leurs unités, calculer le quotient puis écrire une phrase avec l’unité : 12 m \(\div\) 4 s = 3 m/s.\par
\end{pcEncadre}
\end{pcFiche}
\pcPage{L’essentiel}{Chapitre 10 · Mesurer et calculer une vitesse}
\phantomsection\label{ch-10-cours-2}
\begin{pcFiche}{cours-10-2}
\essSection{2}{Rendre les unités cohérentes}
Une heure vaut 60 min et 3 600 s. Une minute vaut 60 s. Un kilomètre vaut 1 000 m. Avant un calcul, convertir pour utiliser une même famille d’unités. Pour passer de m/s à km/h, on multiplie par 3,6 : en une heure, 1 m/s correspond à 3 600 m, soit 3,6 km. Dans le sens inverse, on divise par 3,6.\par
\pcFigure{figures/d11f798c6d921968.pdf}{Le trajet mesure 120 m ; la séparation directe des extrémités est différente et ne sert pas à ce calcul.}
\begin{pcEncadre}{definition}{À retenir}
1 min 30 s vaut 90 s, et non 1,30 min ni 130 s. 0,5 h vaut 30 min ; une écriture décimale d’heure ne se lit pas comme des minutes.\par
\end{pcEncadre}
\begin{pcEncadre}{method}{Une méthode concrète}
Choisir d’abord l’unité demandée pour la vitesse. Convertir la distance et la durée avant de diviser, puis contrôler l’ordre de grandeur.\par
\end{pcEncadre}
\end{pcFiche}
\pcPage{L’essentiel}{Chapitre 10 · Mesurer et calculer une vitesse}
\phantomsection\label{ch-10-cours-3}
\begin{pcFiche}{cours-10-3}
\essSection{3}{Une moyenne ne raconte pas tout}
Deux trajets peuvent avoir la même vitesse moyenne sans se dérouler de la même manière. Un cycliste peut s’arrêter puis repartir, alors qu’un autre avance régulièrement. La moyenne sur le trajet entier inclut les arrêts si leur durée fait partie du temps mesuré. Une valeur donnée par un compteur à un instant ne décrit pas nécessairement la moyenne de tout le trajet.\par
\pcFigure{figures/d28706f310244d89.pdf}{Des mouvements différents peuvent conduire à des distances et durées totales identiques.}
\begin{pcEncadre}{definition}{À retenir}
Pour un mouvement uniforme, distance = vitesse \(\times\) durée. Pour un trajet non uniforme, cette relation s’utilise avec la vitesse moyenne sur la même durée totale.\par
\end{pcEncadre}
\begin{pcEncadre}{method}{Une méthode concrète}
Pour une moyenne sur plusieurs étapes, additionner les distances et les durées, puis diviser les totaux. Ne pas faire automatiquement la moyenne des vitesses des étapes.\par
\end{pcEncadre}
\end{pcFiche}
\pcPage{Activité expérimentale}{Mesurer une vitesse de marche}
\phantomsection\label{ch-10-activite}
\begin{pcFiche}{activite-10}
\begin{pcEncadre}{method}{Mission}
Prévoir une mesure reproductible sur une distance connue, sans course.\par
\end{pcEncadre}
\begin{questions}
\item Choisir une zone sûre et une distance de 10 m mesurée. Définir les lignes de départ et d’arrivée et le moment de déclenchement du chronomètre.
\item Effectuer trois essais de marche normale. Relever la durée de chaque essai avec son unité ; ne pas modifier la distance.
\item Calculer les vitesses moyennes. Données d’entraînement possibles : 10 m en 8 s, 10 s et 12 s.
\end{questions}
\pcFigure{figures/a4c8cc8fee0689f9.pdf}{La droite illustre un modèle uniforme de 10 m en 8 s ; une marche réelle présente de petites variations.}
\begin{questions}
\item Calcule les trois vitesses moyennes en m/s ; arrondis la dernière au centième.
\item Quel essai correspond à la plus grande vitesse moyenne ?
\item Cite deux causes possibles de variation des résultats.
\end{questions}
\end{pcFiche}
\pcPage{L’essentiel à retenir}{Chapitre 10 · Synthèse et autoévaluation}
\phantomsection\label{ch-10-bilan}
\begin{pcFiche}{synthese-10}
\begin{pcEncadre}{definition}{Repère 1}
Vitesse moyenne = distance parcourue \(\div\) durée du même trajet.\par
\end{pcEncadre}
\begin{pcEncadre}{method}{Repère 2}
Mètres et secondes conduisent à des m/s ; kilomètres et heures à des km/h.\par
\end{pcEncadre}
\begin{pcEncadre}{definition}{Repère 3}
1 h = 3 600 s ; 1 min = 60 s ; 1 km = 1 000 m.\par
\end{pcEncadre}
\begin{pcEncadre}{method}{Repère 4}
De m/s à km/h : multiplier par 3,6 ; inversement : diviser par 3,6.\par
\end{pcEncadre}
\begin{pcEncadre}{definition}{Repère 5}
Une moyenne peut inclure des arrêts et masquer des variations de vitesse.\par
\end{pcEncadre}
\begin{pcEncadre}{attention}{Une erreur à éviter}
« J’additionne les vitesses puis je divise par deux. » Pour un trajet complet, calculer distance totale \(\div\) durée totale.\par
\end{pcEncadre}
\end{pcFiche}
\clearpage\markright{Exercices}
\begin{blocLivre}{exo-10-1}
\rubrique{Je démarre}
\begin{exercice}[1]{Choisir une unité}
\begin{questions}
\item Quelle unité de vitesse associe-t-on aux mètres et aux secondes ?
\item Quelle unité associe-t-on aux kilomètres et aux heures ?
\end{questions}
\end{exercice}\end{blocLivre}
\begin{blocLivre}{exo-10-2}
\begin{exercice}[1]{Une première vitesse}
\begin{questions}
\item Un mobile parcourt 20 m en 5 s. Calcule sa vitesse moyenne.
\item Que signifie le résultat pour un mouvement uniforme ?
\end{questions}
\end{exercice}\end{blocLivre}
\begin{blocLivre}{exo-10-3}
\begin{exercice}[1]{Convertir une durée}
\begin{questions}
\item Convertis 2 min en secondes.
\item Convertis 1 min 30 s en secondes.
\end{questions}
\end{exercice}\end{blocLivre}
\begin{blocLivre}{exo-10-4}
\begin{exercice}[1]{Lire une droite}
\pcFigure{figures/c12d6c191831498b.pdf}{}
\begin{questions}
\item Quelle distance est parcourue en 4 s ?
\item Quelle est la vitesse dans le modèle représenté ?
\end{questions}
\end{exercice}\end{blocLivre}
\begin{blocLivre}{exo-10-5}
\begin{exercice}[1]{Le bon trajet}
\pcFigure{figures/38c5864ac59a2ed8.pdf}{}
\begin{questions}
\item Pour calculer la vitesse moyenne entre A et B sur ce chemin, faut-il utiliser 120 m ou 80 m ?
\item Pourquoi ?
\end{questions}
\end{exercice}\end{blocLivre}
\begin{blocLivre}{exo-10-6}
\rubrique{Je m’entraîne}
\begin{exercice}[2]{Un trajet à vélo}
\begin{questions}
\item Un cycliste parcourt 12 km en 0,5 h. Calcule sa vitesse moyenne.
\item Exprime la durée en minutes.
\end{questions}
\end{exercice}\end{blocLivre}
\begin{blocLivre}{exo-10-7}
\begin{exercice}[2]{Un petit robot}
\begin{questions}
\item Un robot parcourt 150 cm en 5 s. Calcule sa vitesse en m/s.
\item Pourquoi faut-il convertir avant le calcul ?
\end{questions}
\end{exercice}\end{blocLivre}
\begin{blocLivre}{exo-10-8}
\begin{exercice}[2]{Changer d’unité de vitesse}
\begin{questions}
\item Convertis 5 m/s en km/h.
\item Convertis 36 km/h en m/s.
\end{questions}
\end{exercice}\end{blocLivre}
\begin{blocLivre}{exo-10-9}
\begin{exercice}[2]{Prévoir une distance}
\begin{questions}
\item Un mobile garde une vitesse de 2 m/s pendant 15 s. Quelle distance parcourt-il ?
\item Quelle condition autorise l’utilisation de cette vitesse sur toute la durée ?
\end{questions}
\end{exercice}\end{blocLivre}
\begin{blocLivre}{exo-10-10}
\begin{exercice}[2]{Prévoir une durée}
\begin{questions}
\item À 4 m/s constants, combien de temps faut-il pour parcourir 100 m ?
\item Écris la relation utilisée en mots.
\end{questions}
\end{exercice}\end{blocLivre}
\begin{blocLivre}{exo-10-11}
\begin{exercice}[2]{Inclure un arrêt}
\begin{questions}
\item Une promenade de 3 km dure 1 h, dont 15 min d’arrêt. Quelle vitesse moyenne obtient-on sur la sortie entière ?
\item Pourquoi n’utilise-t-on pas seulement 45 min ici ?
\end{questions}
\end{exercice}\end{blocLivre}
\begin{blocLivre}{exo-10-12}
\rubrique{Je raisonne}
\begin{exercice}[3]{Deux étapes différentes}
\begin{questions}
\item Étape 1 : 2 km en 20 min. Étape 2 : 4 km en 40 min. Calcule la vitesse moyenne totale en km/h.
\item Que remarques-tu sur les deux étapes ?
\end{questions}
\end{exercice}\end{blocLivre}
\begin{blocLivre}{exo-10-13}
\begin{exercice}[3]{Une moyenne trompeuse}
\begin{questions}
\item Un trajet comprend 1 h à 10 km/h puis 2 h à 4 km/h. Calcule distance et durée totales.
\item La vitesse moyenne vaut-elle 7 km/h ?
\end{questions}
\end{exercice}\end{blocLivre}
\begin{blocLivre}{exo-10-14}
\begin{exercice}[3]{Comparer des performances}
\begin{questions}
\item A parcourt 100 m en 20 s ; B parcourt 150 m en 25 s. Qui est le plus rapide en moyenne ?
\item Pourquoi comparer seulement les durées serait-il incorrect ?
\end{questions}
\end{exercice}\end{blocLivre}
\begin{blocLivre}{exo-10-15}
\begin{exercice}[3]{Contrôler un résultat}
\begin{questions}
\item Une élève trouve 500 m/s pour un marcheur ayant parcouru 100 m en 200 s. Repère l’erreur probable.
\item Calcule et interprète la valeur correcte.
\end{questions}
\end{exercice}\end{blocLivre}
\pcPage{Défi · Prendre une initiative}{Le trajet avec une pause}
\phantomsection\label{defi-ex-10}
\begin{pcFiche}{defi-10}
Une élève marche 600 m en 8 min, s’arrête 2 min, puis parcourt encore 400 m en 5 min. Elle affirme : « Ma vitesse moyenne est la moyenne des vitesses des deux parties marchées. »\par
\pcFigure{figures/0b583c8be5e3e8a9.pdf}{}
\begin{questions}
\item Calcule la distance et la durée de toute la sortie.
\item Calcule la vitesse moyenne en m/min puis en km/h.
\item Explique pourquoi la méthode proposée n’est pas adaptée.
\end{questions}
\begin{pcEncadre}{method}{Indice progressif}
L’arrêt compte dans la durée totale demandée. 1 000 m correspondent à 1 km.\par
\end{pcEncadre}
\end{pcFiche}
\clearpage\pagestyle{otMain}
\begin{ouverture}[college2]{11}{Décrire des actions et des interactions}{Partie 2 · Mouvement et interactions}{Un objet immobile subit-il des actions ?}
\objectif{Reconnaître des actions qui déforment ou modifient un mouvement.}
\objectif{Distinguer interactions de contact et interactions à distance.}
\objectif{Identifier les objets en interaction avec un système choisi.}
\end{ouverture}
\pcPage{L’essentiel}{Chapitre 11 · Décrire des actions et des interactions}
\phantomsection\label{ch-11-cours-1}
\phantomsection\label{ess-11}
\begin{pcFiche}{cours-11-1}
\essSection{1}{Observer les effets d’une action}
Une action mécanique peut déformer un objet, le mettre en mouvement ou modifier son mouvement : changer sa direction ou la valeur de sa vitesse. Un objet immobile peut pourtant subir plusieurs actions. Un livre posé sur une table subit notamment l’action de la table et l’attraction de la Terre. L’absence de mouvement visible ne signifie donc pas l’absence d’interactions.\par
\pcFigure{figures/33018c9e5bc31248.pdf}{Quelques effets possibles : déformation, mise en mouvement et modification de trajectoire.}
\begin{pcEncadre}{definition}{À retenir}
Décrire un effet ne suffit pas toujours à identifier toutes les actions. Une action peut exister sans produire une modification visible dans la situation observée.\par
\end{pcEncadre}
\begin{pcEncadre}{method}{Une méthode concrète}
Choisir l’objet étudié, décrire son état ou son évolution, puis rechercher quels autres objets peuvent agir sur lui.\par
\end{pcEncadre}
\end{pcFiche}
\pcPage{L’essentiel}{Chapitre 11 · Décrire des actions et des interactions}
\phantomsection\label{ch-11-cours-2}
\begin{pcFiche}{cours-11-2}
\essSection{2}{Contact ou action à distance ?}
Une action de contact nécessite un contact entre les objets : main et ballon, table et livre, air et voile. Une interaction à distance existe sans contact direct : attraction de la Terre, interaction entre aimants, par exemple. Le mot interaction rappelle que les objets agissent l’un sur l’autre. À ce niveau, on identifie les objets et la nature de l’interaction sans calculer de force.\par
\pcFigure{figures/bb9e14f7cce3917c.pdf}{Un aimant peut agir sur un objet ferromagnétique sans le toucher ; tous les métaux ne sont pas attirés de la même façon.}
\begin{pcEncadre}{definition}{À retenir}
L’air est une matière : son action sur une voile ou un mobile est une action de contact. « Invisible » ne signifie pas « à distance ».\par
\end{pcEncadre}
\begin{pcEncadre}{method}{Une méthode concrète}
Demander si les surfaces sont en contact. Pour une attraction terrestre ou magnétique, identifier les deux objets même si l’un est loin ou peu visible.\par
\end{pcEncadre}
\end{pcFiche}
\pcPage{L’essentiel}{Chapitre 11 · Décrire des actions et des interactions}
\phantomsection\label{ch-11-cours-3}
\begin{pcFiche}{cours-11-3}
\essSection{3}{Construire un diagramme objet-interactions}
Un diagramme objet-interactions place l’objet étudié au centre et les objets de son environnement autour. Un lien indique une interaction retenue dans le modèle ; sa légende distingue contact et distance. Le schéma ne représente ni la trajectoire, ni le sens d’une force, ni une valeur mesurée. Il sert à ne pas oublier des interactions dans une situation clairement définie.\par
\pcFigure{figures/0c93c17d3223e96e.pdf}{Exemple : un livre posé sur une table est poussé par une main. L’interaction avec l’air est négligée dans ce modèle.}
\begin{pcEncadre}{definition}{À retenir}
La Terre reste à considérer même lorsque l’objet est immobile. On annonce les interactions négligées pour simplifier la situation.\par
\end{pcEncadre}
\begin{pcEncadre}{method}{Une méthode concrète}
Nommer le système central, dresser la liste des objets extérieurs, classer les liens et ajouter une légende explicite.\par
\end{pcEncadre}
\end{pcFiche}
\pcPage{Activité expérimentale}{Enquêter autour d’un livre}
\phantomsection\label{ch-11-activite}
\begin{pcFiche}{activite-11}
\begin{pcEncadre}{method}{Mission}
Comparer les interactions dans deux situations proches.\par
\end{pcEncadre}
\begin{questions}
\item Situation A : un livre repose sur une table, sans main. Situation B : une main pousse le livre sur la même table.
\item Choisir le livre comme objet étudié et négliger l’action de l’air pour cette activité.
\item Tracer un diagramme pour chaque situation, en différenciant contact et distance.
\end{questions}
\pcFigure{figures/b0f4a40eb7b49549.pdf}{Le schéma correspond à la situation B ; il ne représente pas des forces orientées.}
\begin{questions}
\item Quel lien doit-on retirer pour représenter la situation A ?
\item Pourquoi conserve-t-on le lien avec la Terre dans les deux situations ?
\item Quelles informations ce diagramme ne permet-il pas de calculer ?
\end{questions}
\end{pcFiche}
\pcPage{L’essentiel à retenir}{Chapitre 11 · Synthèse et autoévaluation}
\phantomsection\label{ch-11-bilan}
\begin{pcFiche}{synthese-11}
\begin{pcEncadre}{definition}{Repère 1}
Une action peut déformer un objet ou modifier son mouvement.\par
\end{pcEncadre}
\begin{pcEncadre}{method}{Repère 2}
Un objet immobile peut subir plusieurs actions.\par
\end{pcEncadre}
\begin{pcEncadre}{definition}{Repère 3}
Une interaction de contact exige un contact entre les objets.\par
\end{pcEncadre}
\begin{pcEncadre}{method}{Repère 4}
L’attraction terrestre et l’interaction magnétique peuvent agir à distance.\par
\end{pcEncadre}
\begin{pcEncadre}{definition}{Repère 5}
Un diagramme objet-interactions est un inventaire avec une légende, pas un calcul de forces.\par
\end{pcEncadre}
\begin{pcEncadre}{attention}{Une erreur à éviter}
« Le livre ne bouge pas, donc la Terre n’agit pas sur lui. » Le repos ne prouve pas l’absence d’interaction.\par
\end{pcEncadre}
\end{pcFiche}
\clearpage\markright{Exercices}
\begin{blocLivre}{exo-11-1}
\rubrique{Je démarre}
\begin{exercice}[1]{Un effet visible}
\begin{questions}
\item Une main étire un élastique. Quel effet observe-t-on ?
\item Nomme les deux objets en interaction.
\end{questions}
\end{exercice}\end{blocLivre}
\begin{blocLivre}{exo-11-2}
\begin{exercice}[1]{Classer les interactions}
\begin{questions}
\item Classe : pied et ballon au moment du tir ; Terre et ballon ; aimant et trombone attiré sans contact.
\end{questions}
\end{exercice}\end{blocLivre}
\begin{blocLivre}{exo-11-3}
\begin{exercice}[1]{Le livre immobile}
\begin{questions}
\item Un livre sur une table subit-il l’attraction de la Terre ?
\item Quelle autre interaction de contact doit-on retenir ?
\end{questions}
\end{exercice}\end{blocLivre}
\begin{blocLivre}{exo-11-4}
\begin{exercice}[1]{L’action de l’air}
\begin{questions}
\item Le vent déforme une voile. L’interaction air-voile est-elle de contact ou à distance ?
\item Pourquoi l’air invisible peut-il agir ?
\end{questions}
\end{exercice}\end{blocLivre}
\begin{blocLivre}{exo-11-5}
\begin{exercice}[1]{Lire un diagramme}
\pcFigure{figures/b0f4a40eb7b49549.pdf}{}
\begin{questions}
\item Quel objet est choisi comme système ?
\item Quel lien représente une interaction à distance ?
\end{questions}
\end{exercice}\end{blocLivre}
\begin{blocLivre}{exo-11-6}
\rubrique{Je m’entraîne}
\begin{exercice}[2]{Après le tir}
\begin{questions}
\item Après que le pied a quitté le ballon, y a-t-il encore un contact pied-ballon ?
\item Quelle interaction à distance reste présente ?
\end{questions}
\end{exercice}\end{blocLivre}
\begin{blocLivre}{exo-11-7}
\begin{exercice}[2]{Une bille déviée}
\begin{questions}
\item Une bille passe près d’un aimant et sa trajectoire change sans contact. Quel type d’interaction est mis en évidence ?
\item Quel effet de l’action est observé ?
\end{questions}
\end{exercice}\end{blocLivre}
\begin{blocLivre}{exo-11-8}
\begin{exercice}[2]{Préciser les objets}
\begin{questions}
\item Pourquoi « la table agit » est-il moins précis que « la table agit sur le livre » ?
\item Quel élément manque dans la première phrase ?
\end{questions}
\end{exercice}\end{blocLivre}
\begin{blocLivre}{exo-11-9}
\begin{exercice}[2]{Changer le système}
\begin{questions}
\item On étudie maintenant la table au lieu du livre. Faut-il placer encore le livre au centre du diagramme ?
\item Quel principe guide la construction ?
\end{questions}
\end{exercice}\end{blocLivre}
\begin{blocLivre}{exo-11-10}
\begin{exercice}[2]{Main retirée}
\pcFigure{figures/262d33b89a64fd99.pdf}{Document de travail · modèle simplifié.}
\begin{questions}
\item Dans le diagramme du livre poussé sur une table, quel lien disparaît lorsque la main est retirée ?
\item Le lien avec la table disparaît-il aussi si le livre reste posé ?
\end{questions}
\end{exercice}\end{blocLivre}
\clearpage
\begin{blocLivre}{exo-11-11}
\begin{exercice}[2]{Ne pas tout déduire}
\begin{questions}
\item Un diagramme donne un lien table-livre. Permet-il de connaître la valeur de la vitesse du livre ?
\item Quelle est sa fonction principale ?
\end{questions}
\end{exercice}\end{blocLivre}
\begin{blocLivre}{exo-11-12}
\rubrique{Je raisonne}
\begin{exercice}[3]{Une chute dans l’air}
\begin{questions}
\item Établis la liste de deux objets pouvant interagir avec une balle en chute dans l’air.
\item Précise la nature de chaque interaction.
\end{questions}
\end{exercice}\end{blocLivre}
\begin{blocLivre}{exo-11-13}
\begin{exercice}[3]{Choisir un modèle simple}
\begin{questions}
\item Pour un livre immobile sur une table, on néglige l’air. Pourquoi cette simplification doit-elle être annoncée ?
\item Quelles interactions garde-t-on ?
\end{questions}
\end{exercice}\end{blocLivre}
\begin{blocLivre}{exo-11-14}
\begin{exercice}[3]{Une généralisation à discuter}
\begin{questions}
\item « Tous les métaux sont attirés par un aimant. » Pourquoi cette phrase est-elle trop générale ?
\item Propose un protocole de comparaison.
\end{questions}
\end{exercice}\end{blocLivre}
\begin{blocLivre}{exo-11-15}
\begin{exercice}[3]{Une action sans mouvement visible}
\begin{questions}
\item Un élève pousse très doucement un meuble qui reste immobile. Peut-on conclure qu’il n’exerce aucune action ?
\item Quelle observation ou information montre le contact ?
\end{questions}
\end{exercice}\end{blocLivre}
\pcPage{Défi · Prendre une initiative}{Trois scènes, trois diagrammes}
\phantomsection\label{defi-ex-11}
\begin{pcFiche}{defi-11}
Décris un ballon dans trois situations : tenu dans la main ; lâché dans l’air ; posé sur le sol. Tu retiendras la Terre et les objets en contact nécessaires, en précisant le rôle de l’air.\par
\pcFigure{figures/d73f50a80031b0d6.pdf}{}
\begin{questions}
\item Choisis le ballon comme système dans les trois scènes.
\item Dresse les listes d’interactions et classe-les.
\item Explique pourquoi une interaction peut rester présente dans les trois scènes.
\end{questions}
\begin{pcEncadre}{method}{Indice progressif}
L’attraction terrestre ne dépend pas de l’existence d’un contact avec une main ou un support.\par
\end{pcEncadre}
\end{pcFiche}
\clearpage\pagestyle{otMain}
\begin{ouverture}[college3]{12}{Sources, formes et conversions d’énergie}{Partie 3 · L’énergie, ses transferts et ses conversions}{D’où vient l’énergie qui fait fonctionner une lampe ?}
\objectif{Distinguer ressource, réserve, forme d’énergie et convertisseur.}
\objectif{Décrire une chaîne de transferts et de conversions.}
\objectif{Distinguer transfert utile et autres transferts sans faire disparaître l’énergie.}
\end{ouverture}
\pcPage{L’essentiel}{Chapitre 12 · Sources, formes et conversions d’énergie}
\phantomsection\label{ch-12-cours-1}
\phantomsection\label{ess-12}
\begin{pcFiche}{cours-12-1}
\essSection{1}{Identifier les sources et les réserves}
Un objet qui fonctionne reçoit de l’énergie ou utilise une réserve. Une pile stocke de l’énergie chimique ; des aliments constituent une réserve chimique pour le corps ; de l’eau retenue en hauteur possède une énergie liée à sa position. Soleil, vent, eau en mouvement et combustibles sont des ressources mobilisées dans différents dispositifs. L’électricité est une façon de transférer de l’énergie, pas une ressource primaire extraite telle quelle.\par
\pcFigure{figures/92cfce9fecfc2df3.pdf}{Quelques ressources utilisées pour l’énergie. La classification concerne leur renouvellement à l’échelle humaine.}
\begin{pcEncadre}{definition}{À retenir}
Distinguer la ressource et l’appareil : le vent est une ressource ; l’éolienne est un convertisseur. Le mot énergie n’est pas synonyme de courant électrique.\par
\end{pcEncadre}
\begin{pcEncadre}{method}{Une méthode concrète}
Demander « d’où vient l’énergie ? » puis « quel dispositif la convertit ? ». Nommer chaque élément au lieu de répondre seulement « c’est électrique ».\par
\end{pcEncadre}
\end{pcFiche}
\pcPage{L’essentiel}{Chapitre 12 · Sources, formes et conversions d’énergie}
\phantomsection\label{ch-12-cours-2}
\begin{pcFiche}{cours-12-2}
\essSection{2}{Suivre une conversion}
Un convertisseur reçoit de l’énergie et la transfère sous d’autres formes ou par d’autres modes de transfert. Une lampe alimentée reçoit de l’énergie par transfert électrique ; elle émet de la lumière et transfère aussi de l’énergie à son environnement sous forme thermique. Un moteur permet notamment de mettre un objet en mouvement. Une chaîne énergétique représente une sélection de ces étapes avec des flèches légendées.\par
\pcFigure{figures/588e0a374579d7d0.pdf}{La pile fournit un transfert électrique à la lampe ; les deux sorties sont des transferts d’énergie vers l’environnement.}
\begin{pcEncadre}{definition}{À retenir}
Une flèche de chaîne énergétique représente un transfert d’énergie, pas le déplacement d’une substance et pas nécessairement le sens du courant dans un circuit.\par
\end{pcEncadre}
\begin{pcEncadre}{method}{Une méthode concrète}
Écrire le nom du convertisseur dans la case et la nature du transfert sur les flèches. Vérifier que les sorties ne se limitent pas au seul effet recherché.\par
\end{pcEncadre}
\end{pcFiche}
\pcPage{L’essentiel}{Chapitre 12 · Sources, formes et conversions d’énergie}
\phantomsection\label{ch-12-cours-3}
\begin{pcFiche}{cours-12-3}
\essSection{3}{L’énergie ne disparaît pas}
Dans un bilan d’énergie, on compte les transferts entrants, les transferts sortants et les variations des réserves. Si un appareil ne stocke pas d’énergie pendant l’intervalle étudié, l’énergie reçue est égale au total de l’énergie transférée. Le joule (J) est l’unité du Système international pour l’énergie. Une énergie dite perdue reste transférée, souvent vers l’environnement sous une forme moins utile pour l’objectif choisi.\par
\pcFigure{figures/4b767bf897db9cdd.pdf}{Bilan pédagogique d’un appareil sans variation de réserve : 100 J reçus, 30 J utiles et 70 J par un autre transfert.}
\begin{pcEncadre}{definition}{À retenir}
Le caractère utile dépend du but : un transfert thermique est recherché pour chauffer, mais peut être secondaire pour éclairer.\par
\end{pcEncadre}
\begin{pcEncadre}{method}{Une méthode concrète}
Avant d’additionner, vérifier que toutes les valeurs concernent la même durée et la même unité, puis rechercher un éventuel stockage.\par
\end{pcEncadre}
\end{pcFiche}
\pcPage{Activité expérimentale}{Construire une chaîne énergétique}
\phantomsection\label{ch-12-activite}
\begin{pcFiche}{activite-12}
\begin{pcEncadre}{method}{Mission}
Comparer les fonctions d’une lampe et d’un petit moteur scolaires.\par
\end{pcEncadre}
\begin{questions}
\item À partir de documents ou de matériels de classe compatibles, identifier le générateur et le récepteur. Les montages sont validés par le professeur.
\item Décrire l’effet recherché : lumière ou mouvement. Rechercher aussi les autres transferts, sans toucher une partie qui pourrait être chaude.
\item Dessiner pour chaque appareil une chaîne comportant la source ou réserve, le convertisseur et les transferts sortants.
\end{questions}
\pcFigure{figures/4884783ed01a165e.pdf}{Le schéma décrit les transferts d’énergie ; il ne remplace pas un schéma électrique.}
\begin{questions}
\item Qu’est-ce qui change entre la chaîne de la lampe et celle du moteur ?
\item Pourquoi une sortie thermique peut-elle apparaître dans les deux chaînes ?
\item Quelle réserve diminue lorsque la pile alimente le montage ?
\end{questions}
\end{pcFiche}
\pcPage{L’essentiel à retenir}{Chapitre 12 · Synthèse et autoévaluation}
\phantomsection\label{ch-12-bilan}
\begin{pcFiche}{synthese-12}
\begin{pcEncadre}{definition}{Repère 1}
Une ressource ou une réserve fournit l’énergie utilisée.\par
\end{pcEncadre}
\begin{pcEncadre}{method}{Repère 2}
Un convertisseur modifie les formes ou les modes de transfert de l’énergie.\par
\end{pcEncadre}
\begin{pcEncadre}{definition}{Repère 3}
Une chaîne énergétique utilise des flèches légendées.\par
\end{pcEncadre}
\begin{pcEncadre}{method}{Repère 4}
L’énergie s’exprime notamment en joules, symbole J.\par
\end{pcEncadre}
\begin{pcEncadre}{definition}{Repère 5}
Une énergie moins utile pour un objectif n’a pas été détruite.\par
\end{pcEncadre}
\begin{pcEncadre}{attention}{Une erreur à éviter}
« La lampe consomme le courant, qui disparaît. » La lampe convertit et transfère de l’énergie ; le courant ne s’épuise pas en traversant le circuit.\par
\end{pcEncadre}
\end{pcFiche}
\clearpage\markright{Exercices}
\begin{blocLivre}{exo-12-1}
\rubrique{Je démarre}
\begin{exercice}[1]{Ressource ou appareil ?}
\begin{questions}
\item Classe : vent, éolienne, Soleil, panneau photovoltaïque.
\end{questions}
\end{exercice}\end{blocLivre}
\begin{blocLivre}{exo-12-2}
\begin{exercice}[1]{Une réserve dans la pile}
\begin{questions}
\item Sous quelle forme l’énergie est-elle stockée dans une pile ?
\item Que devient cette réserve pendant l’alimentation d’une lampe ?
\end{questions}
\end{exercice}\end{blocLivre}
\begin{blocLivre}{exo-12-3}
\begin{exercice}[1]{Nommer une conversion}
\begin{questions}
\item Quel effet utile recherche-t-on avec un moteur de jouet ?
\item Quel effet utile recherche-t-on avec une lampe ?
\end{questions}
\end{exercice}\end{blocLivre}
\begin{blocLivre}{exo-12-4}
\begin{exercice}[1]{L’unité d’énergie}
\begin{questions}
\item Quelle grandeur peut être exprimée en J ?
\item Que signifie ce symbole ?
\end{questions}
\end{exercice}\end{blocLivre}
\begin{blocLivre}{exo-12-5}
\begin{exercice}[1]{Lire une chaîne}
\pcFigure{figures/6d0e39ac014c4fa7.pdf}{}
\begin{questions}
\item Quel est le convertisseur au centre ?
\item Nomme les deux sorties représentées.
\end{questions}
\end{exercice}\end{blocLivre}
\begin{blocLivre}{exo-12-6}
\rubrique{Je m’entraîne}
\begin{exercice}[2]{Compléter un bilan}
\begin{questions}
\item Un appareil reçoit 80 J et en transfère 25 J utilement, sans stockage. Quelle autre énergie doit-on comptabiliser ?
\item Est-elle détruite ?
\end{questions}
\end{exercice}\end{blocLivre}
\begin{blocLivre}{exo-12-7}
\begin{exercice}[2]{Changer d’objectif}
\begin{questions}
\item Un transfert thermique est-il toujours inutile ?
\item Compare une lampe d’éclairage et un appareil destiné à chauffer.
\end{questions}
\end{exercice}\end{blocLivre}
\begin{blocLivre}{exo-12-8}
\begin{exercice}[2]{Une chaîne solaire}
\begin{questions}
\item Un panneau photovoltaïque alimente un moteur. Identifie la ressource, un premier convertisseur et l’effet recherché.
\end{questions}
\end{exercice}\end{blocLivre}
\begin{blocLivre}{exo-12-9}
\begin{exercice}[2]{Un bilan à contrôler}
\pcFigure{figures/1b06dfae6ff2b61d.pdf}{}
\begin{questions}
\item Le schéma reçoit 100 J et indique 30 J utiles, 70 J autrement. Le total est-il cohérent sans stockage ?
\item Quelle information faut-il garder identique pour comparer les valeurs ?
\end{questions}
\end{exercice}\end{blocLivre}
\begin{blocLivre}{exo-12-10}
\begin{exercice}[2]{Une réserve qui augmente}
\begin{questions}
\item Un dispositif reçoit 50 J, en transfère 35 J à l’extérieur et stocke le reste. Quelle réserve gagne-t-il ?
\item Pourquoi les seules sorties ne sont-elles pas égales à l’entrée ?
\end{questions}
\end{exercice}\end{blocLivre}
\begin{blocLivre}{exo-12-11}
\begin{exercice}[2]{Des mots précis}
\begin{questions}
\item « Le vent fabrique de l’énergie dans l’éolienne. » Reformule sans création d’énergie.
\end{questions}
\end{exercice}\end{blocLivre}
\begin{blocLivre}{exo-12-12}
\rubrique{Je raisonne}
\begin{exercice}[3]{Comparer deux appareils}
\begin{questions}
\item A reçoit 100 J pour 20 J utiles ; B reçoit 100 J pour 35 J utiles. Lequel fournit davantage d’énergie utile dans ces conditions ?
\item Calcule les autres transferts si les appareils ne stockent rien.
\end{questions}
\end{exercice}\end{blocLivre}
\begin{blocLivre}{exo-12-13}
\begin{exercice}[3]{Une chaîne incomplète}
\pcFigure{figures/0291720a8f0f0e99.pdf}{Document de travail · modèle simplifié.}
\begin{questions}
\item Une chaîne de moteur ne montre que « pile \(\rightarrow\) mouvement ». Quelles informations manque-t-il ?
\item Propose une amélioration.
\end{questions}
\end{exercice}\end{blocLivre}
\begin{blocLivre}{exo-12-14}
\begin{exercice}[3]{Ne pas additionner des grandeurs différentes}
\begin{questions}
\item Un élève additionne 20 J et 5 s pour trouver « 25 d’énergie ». Pourquoi le calcul est-il invalide ?
\item Quelle règle utiliser ?
\end{questions}
\end{exercice}\end{blocLivre}
\begin{blocLivre}{exo-12-15}
\begin{exercice}[3]{Un effet secondaire observé}
\begin{questions}
\item Un moteur chauffe et fait un peu de bruit. Comment enrichir la chaîne énergétique ?
\item Pourquoi ne faut-il pas écrire « l’énergie restante est perdue pour toujours » ?
\end{questions}
\end{exercice}\end{blocLivre}
\pcPage{Défi · Prendre une initiative}{L’éclairage à améliorer}
\phantomsection\label{defi-ex-12}
\begin{pcFiche}{defi-12}
Deux systèmes reçoivent chacun 200 J pendant la même durée. A transfère 40 J sous forme lumineuse utile ; B en transfère 70 J. Ils ne stockent pas d’énergie pendant l’essai. On recherche le meilleur éclairage pour cette énergie reçue.\par
\pcFigure{figures/0e6d37f79edadf29.pdf}{}
\begin{questions}
\item Complète les deux bilans.
\item Quel système répond mieux au critère annoncé ?
\item Cite une autre information nécessaire pour choisir un éclairage réel.
\end{questions}
\begin{pcEncadre}{method}{Indice progressif}
Le critère porte sur l’énergie lumineuse utile pour une même énergie reçue, pas seulement sur le mot « moderne ».\par
\end{pcEncadre}
\end{pcFiche}
\clearpage\pagestyle{otMain}
\begin{ouverture}[college3]{13}{Réaliser et schématiser un circuit}{Partie 3 · L’énergie, ses transferts et ses conversions}{Pourquoi un seul fil ne suffit-il pas pour allumer une lampe ?}
\objectif{Identifier générateur, récepteur, fils et interrupteur.}
\objectif{Réaliser et lire un circuit simple avec ses symboles.}
\objectif{Distinguer circuit ouvert, circuit fermé et sens conventionnel du courant.}
\end{ouverture}
\pcPage{L’essentiel}{Chapitre 13 · Réaliser et schématiser un circuit}
\phantomsection\label{ch-13-cours-1}
\phantomsection\label{ess-13}
\begin{pcFiche}{cours-13-1}
\essSection{1}{Des composants et des symboles}
Un circuit simple comporte un générateur, des récepteurs et des liaisons conductrices. Une pile peut jouer le rôle de générateur ; une lampe ou un moteur est un récepteur. Un interrupteur permet d’ouvrir ou de fermer un trajet conducteur. Un dipôle possède deux bornes de connexion. Le schéma électrique représente les connexions par des symboles normalisés ; il n’imite pas la forme ou la taille réelle des objets.\par
\pcFigure{figures/ca00ee694b529ba4.pdf}{Exemples de symboles scolaires : pile, lampe, moteur et interrupteur ouvert.}
\begin{pcEncadre}{definition}{À retenir}
Pour la pile, le trait long représente la borne positive, le trait court la borne négative. Les traits du schéma indiquent les liaisons conductrices.\par
\end{pcEncadre}
\begin{pcEncadre}{method}{Une méthode concrète}
Faire l’inventaire des composants puis repérer leurs deux bornes. Dessiner les fils à la règle et placer les symboles sans les remplacer par des dessins décoratifs.\par
\end{pcEncadre}
\end{pcFiche}
\pcPage{L’essentiel}{Chapitre 13 · Réaliser et schématiser un circuit}
\phantomsection\label{ch-13-cours-2}
\begin{pcFiche}{cours-13-2}
\essSection{2}{Une boucle conductrice fermée}
Une lampe adaptée peut fonctionner lorsque ses deux bornes sont reliées aux deux bornes d’un générateur adapté par une boucle conductrice fermée. Si une liaison est interrompue, le circuit est ouvert et le courant ne circule plus dans cette boucle. La présence de la pile ne suffit pas : il faut un trajet complet. Les composants de classe sont choisis compatibles et le montage est vérifié avant mise en fonctionnement.\par
\pcFigure{figures/4ebb9777c7631d28.pdf}{L’interrupteur ouvert interrompt la boucle. La lampe est représentée éteinte.}
\begin{pcEncadre}{definition}{À retenir}
Une lampe éteinte peut avoir plusieurs causes : interrupteur ouvert, mauvais contact, pile inadaptée ou déchargée, lampe défectueuse. Il faut tester méthodiquement.\par
\end{pcEncadre}
\begin{pcEncadre}{method}{Une méthode concrète}
Suivre le trajet depuis une borne de la pile jusqu’à l’autre. Toute interruption ou mauvais contact empêche le fonctionnement de cette boucle.\par
\end{pcEncadre}
\end{pcFiche}
\pcPage{L’essentiel}{Chapitre 13 · Réaliser et schématiser un circuit}
\phantomsection\label{ch-13-cours-3}
\begin{pcFiche}{cours-13-3}
\essSection{3}{Un sens conventionnel pour le courant}
Dans un circuit fermé alimenté par une pile, on représente conventionnellement le courant à l’extérieur du générateur du pôle positif vers le pôle négatif. La lampe ne consomme pas progressivement le courant ; elle reçoit et transfère de l’énergie. Certains récepteurs dépendent du branchement : inverser les bornes d’un petit moteur à courant continu adapté inverse son sens de rotation.\par
\pcFigure{figures/d570d8704be55406.pdf}{Les flèches représentent le sens conventionnel à l’extérieur de la pile, dans la boucle fermée.}
\begin{pcEncadre}{definition}{À retenir}
Le sens de rotation du moteur est distinct du sens conventionnel du courant. Un schéma de circuit et une chaîne énergétique répondent à des questions différentes.\par
\end{pcEncadre}
\begin{pcEncadre}{method}{Une méthode concrète}
Repérer d’abord + et - sur le générateur. Tracer les flèches seulement sur les fils du circuit extérieur ; ne pas confondre ces flèches avec celles d’une chaîne d’énergie.\par
\end{pcEncadre}
\end{pcFiche}
\pcPage{Activité expérimentale}{Allumer, éteindre et dépanner}
\phantomsection\label{ch-13-activite}
\begin{pcFiche}{activite-13}
\begin{pcEncadre}{method}{Mission}
Réaliser un circuit à une lampe et expliquer chaque modification.\par
\end{pcEncadre}
\begin{questions}
\item Avec le matériel scolaire fourni, préparer un générateur protégé, une lampe adaptée, des fils et un interrupteur. Faire valider le montage hors fonctionnement.
\item Fermer puis ouvrir l’interrupteur. Noter le fonctionnement de la lampe et représenter les deux situations.
\item Si la lampe reste éteinte, contrôler une seule cause à la fois avec le professeur : contact, composant ou source.
\end{questions}
\pcFigure{figures/8a70d74bfce1a098.pdf}{Un circuit scolaire simple fermé ; la couleur de remplissage aide ici à repérer la lampe en fonctionnement.}
\begin{questions}
\item Quel trajet complet permet le fonctionnement ?
\item Quelle connexion change lorsqu’on ouvre l’interrupteur ?
\item Pourquoi faut-il éviter de modifier plusieurs éléments à la fois pour dépanner ?
\end{questions}
\end{pcFiche}
\pcPage{L’essentiel à retenir}{Chapitre 13 · Synthèse et autoévaluation}
\phantomsection\label{ch-13-bilan}
\begin{pcFiche}{synthese-13}
\begin{pcEncadre}{definition}{Repère 1}
Un générateur fournit l’énergie transférée aux récepteurs.\par
\end{pcEncadre}
\begin{pcEncadre}{method}{Repère 2}
Une boucle conductrice fermée est nécessaire au fonctionnement du circuit simple.\par
\end{pcEncadre}
\begin{pcEncadre}{definition}{Repère 3}
Un interrupteur ouvre ou ferme une liaison.\par
\end{pcEncadre}
\begin{pcEncadre}{method}{Repère 4}
Le schéma représente les connexions et les symboles des composants.\par
\end{pcEncadre}
\begin{pcEncadre}{definition}{Repère 5}
Dans le circuit extérieur, le sens conventionnel va du + vers le - de la pile.\par
\end{pcEncadre}
\begin{pcEncadre}{attention}{Une erreur à éviter}
« Un seul fil apporte l’électricité à la lampe. » Il faut un trajet de retour vers l’autre borne du générateur pour former une boucle.\par
\end{pcEncadre}
\end{pcFiche}
\clearpage\markright{Exercices}
\begin{blocLivre}{exo-13-1}
\rubrique{Je démarre}
\begin{exercice}[1]{Nommer les composants}
\begin{questions}
\item Quel composant peut jouer le rôle de générateur ?
\item Donne deux exemples de récepteurs.
\end{questions}
\end{exercice}\end{blocLivre}
\begin{blocLivre}{exo-13-2}
\begin{exercice}[1]{Deux bornes}
\begin{questions}
\item Que signifie le mot dipôle ?
\item Combien de bornes comporte le symbole d’une lampe ?
\end{questions}
\end{exercice}\end{blocLivre}
\begin{blocLivre}{exo-13-3}
\begin{exercice}[1]{Ouvrir le circuit}
\begin{questions}
\item Quel est l’effet d’un interrupteur ouvert dans une boucle unique ?
\item La présence d’une pile suffit-elle à allumer la lampe ?
\end{questions}
\end{exercice}\end{blocLivre}
\begin{blocLivre}{exo-13-4}
\begin{exercice}[1]{Lire les symboles}
\pcFigure{figures/d31c9e690caff648.pdf}{}
\begin{questions}
\item Identifie pile, lampe et moteur sur la planche.
\item Quel trait de la pile représente la borne positive ?
\end{questions}
\end{exercice}\end{blocLivre}
\begin{blocLivre}{exo-13-5}
\begin{exercice}[1]{Le sens conventionnel}
\begin{questions}
\item Dans le circuit extérieur d’une pile, de quelle borne vers quelle borne représente-t-on le courant ?
\end{questions}
\end{exercice}\end{blocLivre}
\begin{blocLivre}{exo-13-6}
\rubrique{Je m’entraîne}
\begin{exercice}[2]{Lire une coupure}
\pcFigure{figures/a7491f435a5e5e7a.pdf}{}
\begin{questions}
\item Pourquoi la lampe ne fonctionne-t-elle pas sur ce schéma ?
\item Quelle modification autorisée fermerait la boucle ?
\end{questions}
\end{exercice}\end{blocLivre}
\begin{blocLivre}{exo-13-7}
\begin{exercice}[2]{Un dessin ou un schéma ?}
\begin{questions}
\item Pourquoi un schéma électrique n’a-t-il pas besoin de respecter la taille réelle des composants ?
\item Quelle information doit-il préserver absolument ?
\end{questions}
\end{exercice}\end{blocLivre}
\begin{blocLivre}{exo-13-8}
\begin{exercice}[2]{Une panne à rechercher}
\begin{questions}
\item Une lampe reste éteinte dans un circuit apparemment fermé. Propose trois causes possibles.
\item Pourquoi ne faut-il pas conclure directement que la pile est vide ?
\end{questions}
\end{exercice}\end{blocLivre}
\begin{blocLivre}{exo-13-9}
\begin{exercice}[2]{Le moteur inversé}
\pcFigure{figures/cbf4bc80f2249cf9.pdf}{}
\begin{questions}
\item Un moteur à courant continu adapté tourne dans un sens. On inverse les connexions de la pile. Que peut-on observer ?
\item Faut-il confondre ce sens avec la forme du schéma ?
\end{questions}
\end{exercice}\end{blocLivre}
\begin{blocLivre}{exo-13-10}
\begin{exercice}[2]{Suivre une boucle}
\pcFigure{figures/8a70d74bfce1a098.pdf}{Document de travail · modèle simplifié.}
\begin{questions}
\item Décris un trajet complet depuis la borne positive, en passant par les composants du montage simple.
\item Où faut-il revenir ?
\end{questions}
\end{exercice}\end{blocLivre}
\begin{blocLivre}{exo-13-11}
\begin{exercice}[2]{Courant et énergie}
\begin{questions}
\item Corrige : « le courant entre dans la lampe et disparaît en lumière ».
\end{questions}
\end{exercice}\end{blocLivre}
\clearpage
\begin{blocLivre}{exo-13-12}
\rubrique{Je raisonne}
\begin{exercice}[3]{Le bon témoin}
\begin{questions}
\item Pour tester une lampe, pourquoi faut-il d’abord connaître un montage qui fonctionne ?
\item Quel composant changer ensuite ?
\end{questions}
\end{exercice}\end{blocLivre}
\begin{blocLivre}{exo-13-13}
\begin{exercice}[3]{Une connexion trompeuse}
\begin{questions}
\item Un fil touche l’enveloppe isolante d’une lampe au lieu de sa borne métallique. Pourquoi la boucle peut-elle rester ouverte ?
\item Quelle information matérielle faut-il vérifier ?
\end{questions}
\end{exercice}\end{blocLivre}
\begin{blocLivre}{exo-13-14}
\begin{exercice}[3]{Deux schémas équivalents}
\pcFigure{figures/23cb91a1b45aefe5.pdf}{Document de travail · modèle simplifié.}
\begin{questions}
\item On déplace une lampe sur le dessin sans changer ses connexions. Le circuit électrique change-t-il forcément ?
\item Quel contrôle permet de comparer les dessins ?
\end{questions}
\end{exercice}\end{blocLivre}
\begin{blocLivre}{exo-13-15}
\begin{exercice}[3]{Répondre à un cahier des charges}
\begin{questions}
\item Propose les composants d’un montage qui allume une lampe et permet de l’éteindre sans retirer de fil.
\item Explique le rôle de chaque élément.
\end{questions}
\end{exercice}\end{blocLivre}
\pcPage{Défi · Prendre une initiative}{Le dossier de dépannage}
\phantomsection\label{defi-ex-13}
\begin{pcFiche}{defi-13}
Une lampe ne s’allume pas. On dispose d’une lampe de référence et de fils connus en bon état, compatibles avec le matériel scolaire utilisé. Construis un arbre de recherche qui évite de changer tout le montage à la fois.\par
\pcFigure{figures/e6818b25a3509f4f.pdf}{}
\begin{questions}
\item Indique le premier contrôle visuel hors fonctionnement.
\item Propose deux substitutions successives.
\item Explique comment conclure sans aller au-delà de ce que chaque test montre.
\end{questions}
\begin{pcEncadre}{method}{Indice progressif}
Conserver une source et des composants compatibles ; demander au professeur avant toute modification.\par
\end{pcEncadre}
\end{pcFiche}
\clearpage\pagestyle{otMain}
\begin{ouverture}[college3]{14}{Conducteurs, isolants et sécurité électrique}{Partie 3 · L’énergie, ses transferts et ses conversions}{Comment tester un matériau sans prendre de risque ?}
\objectif{Distinguer conducteur et isolant dans les conditions d’un test.}
\objectif{Construire un test avec des témoins de fonctionnement.}
\objectif{Reconnaître un court-circuit et les règles essentielles de sécurité.}
\end{ouverture}
\pcPage{L’essentiel}{Chapitre 14 · Conducteurs, isolants et sécurité électrique}
\phantomsection\label{ch-14-cours-1}
\phantomsection\label{ess-14}
\begin{pcFiche}{cours-14-1}
\essSection{1}{Un matériau dans un circuit de test}
Un conducteur permet le passage du courant dans les conditions considérées ; un isolant le limite fortement. Dans un circuit scolaire de test, un échantillon complète une boucle comportant une lampe adaptée. Les métaux usuels conduisent généralement bien ; beaucoup de plastiques secs sont isolants. Le résultat dépend aussi des contacts, de la forme de l’échantillon et de la sensibilité du dispositif.\par
\pcFigure{figures/753476b9fd77f40b.pdf}{L’échantillon est placé dans la boucle avec un récepteur ; il ne relie pas directement les bornes du générateur.}
\begin{pcEncadre}{definition}{À retenir}
Une lampe éteinte ne prouve pas à elle seule que l’échantillon est isolant. Il faut vérifier que le montage fonctionne et que les contacts sont corrects.\par
\end{pcEncadre}
\begin{pcEncadre}{method}{Une méthode concrète}
Valider le montage avec un conducteur connu. Tester ensuite l’échantillon sans changer les autres conditions, puis contrôler les contacts.\par
\end{pcEncadre}
\end{pcFiche}
\pcPage{L’essentiel}{Chapitre 14 · Conducteurs, isolants et sécurité électrique}
\phantomsection\label{ch-14-cours-2}
\begin{pcFiche}{cours-14-2}
\essSection{2}{Identifier un court-circuit}
Un fil conducteur qui relie directement les bornes d’un dipôle peut le court-circuiter : le courant dispose d’un chemin de très faible résistance qui le contourne. Si les bornes du générateur sont ainsi reliées presque directement, un courant important peut apparaître, avec échauffement et détérioration. Le schéma sert ici à reconnaître une erreur ; le montage dangereux n’est pas à réaliser.\par
\pcFigure{figures/ef90c747a9d6cc39.pdf}{Situation à repérer sur document seulement : le fil contourne l’unique lampe et crée un trajet presque direct entre les bornes de la pile.}
\begin{pcEncadre}{definition}{À retenir}
Un circuit fermé n’est pas automatiquement correct. Il doit comporter des récepteurs et des liaisons conformes au montage prévu.\par
\end{pcEncadre}
\begin{pcEncadre}{method}{Une méthode concrète}
Suivre les chemins possibles. Rechercher une liaison qui évite tous les récepteurs et rejoint presque directement les deux bornes du générateur.\par
\end{pcEncadre}
\end{pcFiche}
\pcPage{L’essentiel}{Chapitre 14 · Conducteurs, isolants et sécurité électrique}
\phantomsection\label{ch-14-cours-3}
\begin{pcFiche}{cours-14-3}
\essSection{3}{Le corps et l’environnement comptent}
Le corps humain peut conduire le courant. L’humidité peut accroître les risques ; l’eau du robinet n’est pas un isolant fiable. On ne teste jamais la conductivité avec son corps ni avec une prise domestique. Les expériences utilisent exclusivement le matériel scolaire de très basse tension et les protections choisies par le professeur. Une pile peut elle-même chauffer en cas de mauvais branchement.\par
\pcFigure{figures/4ebb9777c7631d28.pdf}{Préparer et modifier un montage hors fonctionnement, puis le faire vérifier avant fermeture.}
\begin{pcEncadre}{definition}{À retenir}
Le secteur domestique et les batteries puissantes ne sont pas des sources pour les montages du livre. L’isolant d’un appareil abîmé ne se répare pas par un simple essai.\par
\end{pcEncadre}
\begin{pcEncadre}{method}{Une méthode concrète}
En cas d’anomalie, arrêter l’activité et prévenir le professeur. Ne pas toucher un composant échauffé et ne pas tenter de dépannage sur une installation domestique.\par
\end{pcEncadre}
\end{pcFiche}
\pcPage{Activité expérimentale}{Comparer trois matériaux}
\phantomsection\label{ch-14-activite}
\begin{pcFiche}{activite-14}
\begin{pcEncadre}{method}{Mission}
Établir un classement argumenté avec un test scolaire et ses limites.\par
\end{pcEncadre}
\begin{questions}
\item Le professeur fournit un montage protégé, un témoin conducteur et des échantillons secs identifiés, par exemple métal, plastique et bois sec.
\item Vérifier le fonctionnement avec le témoin. Remplacer seulement l’échantillon et garder les contacts comparables.
\item Noter l’observation sans forcer les conclusions. Si la lampe reste éteinte, vérifier à nouveau les contacts puis le témoin.
\end{questions}
\pcFigure{figures/4855a50c85e1943a.pdf}{Aucun essai sur le corps, le secteur ou un liquide inconnu n’est demandé.}
\begin{questions}
\item Pourquoi commence-t-on et termine-t-on par un témoin ?
\item Comment noter un résultat négatif sans affirmer une propriété absolue ?
\item Pourquoi le revêtement isolant d’un fil et son cœur métallique ont-ils des rôles différents ?
\end{questions}
\end{pcFiche}
\pcPage{L’essentiel à retenir}{Chapitre 14 · Synthèse et autoévaluation}
\phantomsection\label{ch-14-bilan}
\begin{pcFiche}{synthese-14}
\begin{pcEncadre}{definition}{Repère 1}
Conductivité et isolation s’évaluent dans des conditions données.\par
\end{pcEncadre}
\begin{pcEncadre}{method}{Repère 2}
Un test demande un montage validé et des témoins.\par
\end{pcEncadre}
\begin{pcEncadre}{definition}{Repère 3}
Un mauvais contact peut produire un résultat trompeur.\par
\end{pcEncadre}
\begin{pcEncadre}{method}{Repère 4}
Un court-circuit peut provoquer un échauffement important.\par
\end{pcEncadre}
\begin{pcEncadre}{definition}{Repère 5}
Le corps, le secteur et les sources puissantes sont exclus des essais scolaires du livre.\par
\end{pcEncadre}
\begin{pcEncadre}{attention}{Une erreur à éviter}
« La lampe est éteinte, donc l’échantillon est forcément isolant. » Vérifier le montage, les contacts et la sensibilité du test.\par
\end{pcEncadre}
\end{pcFiche}
\clearpage\markright{Exercices}
\begin{blocLivre}{exo-14-1}
\rubrique{Je démarre}
\begin{exercice}[1]{Conducteur ou isolant ?}
\begin{questions}
\item Dans les conditions usuelles d’un test scolaire, comment classe-t-on un fil de cuivre et une gaine plastique sèche ?
\end{questions}
\end{exercice}\end{blocLivre}
\begin{blocLivre}{exo-14-2}
\begin{exercice}[1]{Le rôle de la gaine}
\begin{questions}
\item Pourquoi un fil électrique possède-t-il souvent un cœur métallique entouré de plastique ?
\end{questions}
\end{exercice}\end{blocLivre}
\begin{blocLivre}{exo-14-3}
\begin{exercice}[1]{Lire le test}
\pcFigure{figures/4855a50c85e1943a.pdf}{}
\begin{questions}
\item Où place-t-on l’échantillon dans le circuit ?
\item Pourquoi conserve-t-on une lampe adaptée dans la boucle ?
\end{questions}
\end{exercice}\end{blocLivre}
\begin{blocLivre}{exo-14-4}
\begin{exercice}[1]{Les expériences autorisées}
\begin{questions}
\item Peut-on tester un matériau avec les deux contacts d’une prise domestique ?
\item Quelle source utilise-t-on ici ?
\end{questions}
\end{exercice}\end{blocLivre}
\begin{blocLivre}{exo-14-5}
\begin{exercice}[1]{Un chemin dangereux}
\begin{questions}
\item Comment nomme-t-on une liaison conductrice presque directe entre les deux bornes d’un générateur ?
\item Quel risque peut-elle présenter ?
\end{questions}
\end{exercice}\end{blocLivre}
\begin{blocLivre}{exo-14-6}
\rubrique{Je m’entraîne}
\begin{exercice}[2]{Un témoin utile}
\begin{questions}
\item La lampe ne fonctionne pas avec un échantillon. Elle ne fonctionne pas non plus avec le témoin conducteur. Peut-on classer l’échantillon ?
\item Quelle étape reprendre ?
\end{questions}
\end{exercice}\end{blocLivre}
\begin{blocLivre}{exo-14-7}
\begin{exercice}[2]{Un métal recouvert}
\begin{questions}
\item Une pièce métallique peinte ne ferme pas le circuit lors du test. Pourquoi cela ne prouve-t-il pas que le métal est isolant ?
\item Que doit-on identifier ?
\end{questions}
\end{exercice}\end{blocLivre}
\begin{blocLivre}{exo-14-8}
\begin{exercice}[2]{Repérer le fil ajouté}
\pcFigure{figures/e82b4cba3fbea2a7.pdf}{}
\begin{questions}
\item Pourquoi le fil supérieur est-il un problème dans ce montage à une seule lampe ?
\item Faut-il réaliser ce montage ?
\end{questions}
\end{exercice}\end{blocLivre}
\begin{blocLivre}{exo-14-9}
\begin{exercice}[2]{Le corps humain}
\begin{questions}
\item Pourquoi ne faut-il pas utiliser ses doigts pour fermer le circuit de test ?
\item L’absence de sensation serait-elle une preuve de sécurité générale ?
\end{questions}
\end{exercice}\end{blocLivre}
\begin{blocLivre}{exo-14-10}
\begin{exercice}[2]{L’eau et les risques}
\begin{questions}
\item Un camarade affirme que toute eau est isolante. Pourquoi cette affirmation est-elle dangereuse ?
\item Quelle règle appliquer près d’appareils électriques ?
\end{questions}
\end{exercice}\end{blocLivre}
\begin{blocLivre}{exo-14-11}
\begin{exercice}[2]{Deux contrôles de contact}
\begin{questions}
\item Donne deux raisons possibles d’un mauvais contact avec un échantillon conducteur.
\item Comment éviter une conclusion trop rapide ?
\end{questions}
\end{exercice}\end{blocLivre}
\begin{blocLivre}{exo-14-12}
\rubrique{Je raisonne}
\begin{exercice}[3]{Rédiger un résultat limité}
\begin{questions}
\item Le bois sec testé n’allume pas la lampe. Propose une formulation scientifique prudente.
\item Pourquoi ne faut-il pas étendre le résultat à tout bois humide ?
\end{questions}
\end{exercice}\end{blocLivre}
\begin{blocLivre}{exo-14-13}
\begin{exercice}[3]{La pile chauffe}
\begin{questions}
\item Pendant un montage, une source devient anormalement chaude. Quelle conduite adopter ?
\item Pourquoi ne faut-il pas poursuivre pour obtenir une mesure ?
\end{questions}
\end{exercice}\end{blocLivre}
\begin{blocLivre}{exo-14-14}
\begin{exercice}[3]{La variable étudiée}
\begin{questions}
\item Un groupe change la pile et l’échantillon entre deux essais. Pourquoi la comparaison est-elle ambiguë ?
\item Comment améliorer le protocole ?
\end{questions}
\end{exercice}\end{blocLivre}
\begin{blocLivre}{exo-14-15}
\begin{exercice}[3]{Concevoir une fiche de test}
\begin{questions}
\item Rédige les rubriques nécessaires pour rendre un essai de conductivité compréhensible.
\item Ajoute une rubrique sur les limites.
\end{questions}
\end{exercice}\end{blocLivre}
\pcPage{Défi · Prendre une initiative}{Le métal déclaré isolant}
\phantomsection\label{defi-ex-14}
\begin{pcFiche}{defi-14}
Deux groupes testent une même tige métallique. L’un allume la lampe et l’autre non. La tige est recouverte d’une peinture isolante sauf sur une petite zone. Prépare une explication et un protocole de vérification encadré.\par
\pcFigure{figures/6c7886c8de396d33.pdf}{}
\begin{questions}
\item Identifie au moins deux causes possibles de la différence.
\item Propose une comparaison avec un seul facteur modifié.
\item Rédige une conclusion qui distingue le matériau et son revêtement.
\end{questions}
\begin{pcEncadre}{method}{Indice progressif}
Le contact avec l’objet n’est pas forcément un contact électrique avec son cœur métallique.\par
\end{pcEncadre}
\end{pcFiche}
\clearpage\pagestyle{otMain}
\begin{ouverture}[college3]{15}{Circuits en série et en dérivation}{Partie 3 · L’énergie, ses transferts et ses conversions}{Comment garder une lampe allumée quand une autre s’arrête ?}
\objectif{Distinguer une association en série et une association en dérivation.}
\objectif{Prévoir les effets d’une ouverture sur les chemins conducteurs.}
\objectif{Placer un interrupteur pour une commande générale ou indépendante.}
\end{ouverture}
\pcPage{L’essentiel}{Chapitre 15 · Circuits en série et en dérivation}
\phantomsection\label{ch-15-cours-1}
\phantomsection\label{ess-15}
\begin{pcFiche}{cours-15-1}
\essSection{1}{Une seule boucle : l’association en série}
Dans un circuit en série simple, les dipôles sont placés les uns à la suite des autres sur une seule boucle. Si une connexion de cette boucle s’ouvre, le courant ne peut plus circuler et tous les récepteurs de la boucle cessent de fonctionner. Modifier l’ordre des composants sans changer les connexions ne crée pas une nouvelle branche. La forme du dessin peut varier tout en représentant le même circuit.\par
\pcFigure{figures/f849864e7c855ef2.pdf}{Deux lampes en série : une seule boucle relie les bornes du générateur.}
\begin{pcEncadre}{definition}{À retenir}
Compter les lampes ne suffit pas pour identifier le montage. Il faut suivre les chemins et les connexions.\par
\end{pcEncadre}
\begin{pcEncadre}{method}{Une méthode concrète}
Partir d’une borne du générateur et suivre l’unique trajet. Si l’on traverse successivement les deux lampes sans choix de chemin, elles sont ici en série.\par
\end{pcEncadre}
\end{pcFiche}
\pcPage{L’essentiel}{Chapitre 15 · Circuits en série et en dérivation}
\phantomsection\label{ch-15-cours-2}
\begin{pcFiche}{cours-15-2}
\essSection{2}{Plusieurs chemins : la dérivation}
Dans une association en dérivation, les récepteurs appartiennent à des branches différentes reliées aux mêmes deux points de connexion. Plusieurs chemins fermés sont possibles avec le générateur. Une branche peut être ouverte tandis qu’une autre reste fermée. Avec une source adaptée et maintenue dans les mêmes conditions, le récepteur de la branche restée fermée peut continuer à fonctionner.\par
\pcFigure{figures/c0d1d729eab36394.pdf}{L₁ et L\textsubscript{2} se trouvent sur des branches différentes ; les points indiquent des jonctions conductrices.}
\begin{pcEncadre}{definition}{À retenir}
Une dérivation ne se reconnaît pas simplement à deux traits parallèles sur la feuille. Les bornes des branches doivent être reliées aux mêmes deux points.\par
\end{pcEncadre}
\begin{pcEncadre}{method}{Une méthode concrète}
Repérer les jonctions puis surligner chaque chemin entre elles. Vérifier que chaque branche de récepteur forme un chemin avec le générateur.\par
\end{pcEncadre}
\end{pcFiche}
\pcPage{L’essentiel}{Chapitre 15 · Circuits en série et en dérivation}
\phantomsection\label{ch-15-cours-3}
\begin{pcFiche}{cours-15-3}
\essSection{3}{Commander ensemble ou séparément}
Un interrupteur placé dans une portion commune à tous les chemins commande l’ensemble des récepteurs. Un interrupteur placé seulement dans une branche commande cette branche, à condition de ne pas interrompre les autres chemins. On choisit donc l’emplacement selon la fonction demandée. Les comparaisons d’éclat nécessitent des lampes identifiées et un générateur adapté ; on ne généralise pas un résultat à n’importe quels composants.\par
\pcFigure{figures/de95f1abc8e204d9.pdf}{L’interrupteur de la branche de L₁ est ouvert ; la branche de L\textsubscript{2} conserve un chemin fermé.}
\begin{pcEncadre}{definition}{À retenir}
Lampe éteinte ne signifie pas absence de générateur. Il faut chercher quel chemin est ouvert et quels autres restent conducteurs.\par
\end{pcEncadre}
\begin{pcEncadre}{method}{Une méthode concrète}
Pour prévoir le fonctionnement, ignorer visuellement la branche ouverte puis chercher s’il reste un chemin complet passant par chaque autre récepteur.\par
\end{pcEncadre}
\end{pcFiche}
\pcPage{Activité expérimentale}{Choisir une commande d’éclairage}
\phantomsection\label{ch-15-activite}
\begin{pcFiche}{activite-15}
\begin{pcEncadre}{method}{Mission}
Comparer la commande de deux lampes dans deux montages scolaires.\par
\end{pcEncadre}
\begin{questions}
\item À partir des schémas, proposer un circuit en série et un circuit en dérivation avec deux lampes adaptées. Faire valider les montages avant fonctionnement.
\item Avec le professeur, ouvrir une liaison prévue dans un chemin de lampe. Noter quelles lampes restent allumées.
\item Représenter une commande indépendante dans une seule branche et une commande générale sur une portion commune.
\end{questions}
\pcFigure{figures/413870db94cfb67c.pdf}{Une ouverture localisée dans une branche ne coupe pas automatiquement toutes les autres.}
\begin{questions}
\item Quel montage permet à L\textsubscript{2} de fonctionner malgré l’ouverture de la branche de L₁ ?
\item Où placer une commande qui éteint les deux lampes ?
\item Pourquoi les deux lampes ne doivent-elles pas être reliées par un fil qui les contourne ?
\end{questions}
\end{pcFiche}
\pcPage{L’essentiel à retenir}{Chapitre 15 · Synthèse et autoévaluation}
\phantomsection\label{ch-15-bilan}
\begin{pcFiche}{synthese-15}
\begin{pcEncadre}{definition}{Repère 1}
En série simple, une seule boucle contient les récepteurs successifs.\par
\end{pcEncadre}
\begin{pcEncadre}{method}{Repère 2}
Une ouverture dans cette boucle interrompt le fonctionnement de tous ses récepteurs.\par
\end{pcEncadre}
\begin{pcEncadre}{definition}{Repère 3}
En dérivation, plusieurs branches sont reliées aux mêmes deux points.\par
\end{pcEncadre}
\begin{pcEncadre}{method}{Repère 4}
Une branche ouverte peut laisser fonctionner les autres branches fermées.\par
\end{pcEncadre}
\begin{pcEncadre}{definition}{Repère 5}
La position d’un interrupteur détermine une commande générale ou indépendante.\par
\end{pcEncadre}
\begin{pcEncadre}{attention}{Une erreur à éviter}
« Deux lampes dessinées côte à côte sont forcément en dérivation. » Les connexions, pas leur position sur la page, déterminent l’association.\par
\end{pcEncadre}
\end{pcFiche}
\clearpage\markright{Exercices}
\begin{blocLivre}{exo-15-1}
\rubrique{Je démarre}
\begin{exercice}[1]{Une boucle unique}
\begin{questions}
\item Dans un circuit simple en série, combien de chemins peut-on suivre dans la boucle ?
\item Traverse-t-on les deux lampes successivement ?
\end{questions}
\end{exercice}\end{blocLivre}
\begin{blocLivre}{exo-15-2}
\begin{exercice}[1]{Nommer l’association}
\pcFigure{figures/1b916fb9a3d436c2.pdf}{}
\begin{questions}
\item Identifie l’association représentée.
\item Quel indice de connexion justifie ta réponse ?
\end{questions}
\end{exercice}\end{blocLivre}
\begin{blocLivre}{exo-15-3}
\begin{exercice}[1]{Une ouverture en série}
\begin{questions}
\item Une des lampes d’un circuit en série devient un circuit ouvert. Que se passe-t-il pour l’autre ?
\item Pourquoi ?
\end{questions}
\end{exercice}\end{blocLivre}
\begin{blocLivre}{exo-15-4}
\begin{exercice}[1]{Une dérivation}
\begin{questions}
\item Quel montage utilise plusieurs branches entre les mêmes deux points ?
\item Faut-il que les fils soient dessinés parallèles pour que la propriété existe ?
\end{questions}
\end{exercice}\end{blocLivre}
\begin{blocLivre}{exo-15-5}
\begin{exercice}[1]{Une commande indépendante}
\begin{questions}
\item Où placer un interrupteur pour commander seulement une branche de lampe ?
\end{questions}
\end{exercice}\end{blocLivre}
\begin{blocLivre}{exo-15-6}
\rubrique{Je m’entraîne}
\begin{exercice}[2]{Suivre les chemins}
\pcFigure{figures/5fb4215f70151542.pdf}{}
\begin{questions}
\item Sur le montage, identifie un chemin passant par L₁ et un autre passant par L\textsubscript{2}.
\item Qu’ont-ils en commun ?
\end{questions}
\end{exercice}\end{blocLivre}
\begin{blocLivre}{exo-15-7}
\begin{exercice}[2]{La branche ouverte}
\pcFigure{figures/413870db94cfb67c.pdf}{}
\begin{questions}
\item Quelle lampe reste sur une boucle fermée ?
\item Peut-elle fonctionner avec le générateur adapté du modèle ?
\end{questions}
\end{exercice}\end{blocLivre}
\begin{blocLivre}{exo-15-8}
\begin{exercice}[2]{Commander les deux lampes}
\pcFigure{figures/c132ea2c6e94fe88.pdf}{Document de travail · modèle simplifié.}
\begin{questions}
\item Un interrupteur est placé sur un fil commun entre la pile et toutes les branches. Quelle est sa fonction ?
\item Que se passe-t-il lorsqu’il est ouvert ?
\end{questions}
\end{exercice}\end{blocLivre}
\begin{blocLivre}{exo-15-9}
\begin{exercice}[2]{Changer l’ordre}
\begin{questions}
\item Dans un circuit en série simple, on échange les positions de deux lampes identiques sans changer les connexions. Crée-t-on une dérivation ?
\item Que reste-t-il identique ?
\end{questions}
\end{exercice}\end{blocLivre}
\begin{blocLivre}{exo-15-10}
\begin{exercice}[2]{Les conditions d’éclat}
\begin{questions}
\item Pourquoi comparer deux lampes inconnues ne suffit-il pas pour établir une règle générale d’éclat en série ?
\item Quelles informations faut-il contrôler ?
\end{questions}
\end{exercice}\end{blocLivre}
\begin{blocLivre}{exo-15-11}
\begin{exercice}[2]{Reconnaître une jonction}
\pcFigure{figures/8787152f2a635390.pdf}{Document de travail · modèle simplifié.}
\begin{questions}
\item À quoi servent les points noirs sur les jonctions du schéma ?
\item Pourquoi un croisement de traits sans convention claire peut-il être ambigu ?
\end{questions}
\end{exercice}\end{blocLivre}
\begin{blocLivre}{exo-15-12}
\rubrique{Je raisonne}
\begin{exercice}[3]{Deux lampes pour un couloir}
\begin{questions}
\item On veut que chaque lampe puisse rester allumée si la branche de l’autre est ouverte. Quelle association choisir ?
\item Justifie par les chemins.
\end{questions}
\end{exercice}\end{blocLivre}
\begin{blocLivre}{exo-15-13}
\begin{exercice}[3]{Une fausse dérivation}
\pcFigure{figures/c3df0e174684c6a7.pdf}{Document de travail · modèle simplifié.}
\begin{questions}
\item Un dessin possède deux branches, mais un fil relie directement les bornes de la pile. Est-ce un montage correct parce qu’il a plusieurs branches ?
\item Quel contrôle prime ?
\end{questions}
\end{exercice}\end{blocLivre}
\begin{blocLivre}{exo-15-14}
\begin{exercice}[3]{Un symptôme, plusieurs causes}
\begin{questions}
\item Dans un montage supposé en dérivation, les deux lampes s’éteignent quand on ouvre une seule branche. Cite deux vérifications.
\end{questions}
\end{exercice}\end{blocLivre}
\begin{blocLivre}{exo-15-15}
\begin{exercice}[3]{Trois fonctions à réaliser}
\begin{questions}
\item Propose l’emplacement de trois interrupteurs : un général, un pour L₁ et un pour L\textsubscript{2}, dans un montage à deux branches.
\item Comment vérifier la commande indépendante ?
\end{questions}
\end{exercice}\end{blocLivre}
\pcPage{Défi · Prendre une initiative}{Le tableau de commande}
\phantomsection\label{defi-ex-15}
\begin{pcFiche}{defi-15}
Un modèle de vitrine scolaire comporte deux lampes. On souhaite quatre situations : aucune allumée, seulement L₁, seulement L\textsubscript{2}, les deux. Une commande générale doit pouvoir tout arrêter. Le montage reste alimenté uniquement par le matériel scolaire adapté.\par
\pcFigure{figures/0c0d980bdeb33bdc.pdf}{}
\begin{questions}
\item Choisis l’association et les emplacements des interrupteurs.
\item Prépare les réglages pour les quatre situations.
\item Explique le rôle de la commande générale.
\end{questions}
\begin{pcEncadre}{method}{Indice progressif}
Chaque lampe doit disposer de sa branche ; ajouter ensuite une commande sur un trajet commun.\par
\end{pcEncadre}
\end{pcFiche}
\clearpage\pagestyle{otMain}
\begin{ouverture}[college3]{16}{Énergie et choix responsables}{Partie 3 · L’énergie, ses transferts et ses conversions}{Comment comparer des choix d’énergie de façon honnête ?}
\objectif{Distinguer ressource renouvelable, ressource non renouvelable et impact environnemental.}
\objectif{Lire des données d’énergie et distinguer énergie et puissance.}
\objectif{Comparer à service rendu équivalent et proposer des économies argumentées.}
\end{ouverture}
\pcPage{L’essentiel}{Chapitre 16 · Énergie et choix responsables}
\phantomsection\label{ch-16-cours-1}
\phantomsection\label{ess-16}
\begin{pcFiche}{cours-16-1}
\essSection{1}{Renouvellement et effets sur l’environnement}
Une ressource renouvelable est disponible sous forme de flux ou se reconstitue à l’échelle humaine dans les conditions de son exploitation. Soleil, vent et cycle de l’eau en fournissent des exemples. Les stocks de charbon, de pétrole et d’uranium ne se reconstituent pas à cette échelle. Renouvelable ne signifie pas sans impact : fabrication, transport, usage et fin de vie des équipements doivent être considérés.\par
\pcFigure{figures/92cfce9fecfc2df3.pdf}{La nature renouvelable d’une ressource et les effets de son exploitation sont deux questions liées mais distinctes.}
\begin{pcEncadre}{definition}{À retenir}
Brûler des combustibles fossiles émet notamment du dioxyde de carbone. L’augmentation des gaz à effet de serre modifie le bilan énergétique du système Terre-atmosphère.\par
\end{pcEncadre}
\begin{pcEncadre}{method}{Une méthode concrète}
Éviter « toujours propre » ou « sans aucune conséquence ». Préciser le critère comparé : ressource, émissions, matériaux, espace occupé ou service rendu.\par
\end{pcEncadre}
\end{pcFiche}
\pcPage{L’essentiel}{Chapitre 16 · Énergie et choix responsables}
\phantomsection\label{ch-16-cours-2}
\begin{pcFiche}{cours-16-2}
\essSection{2}{Lire une consommation d’énergie}
L’énergie s’exprime en joules et, dans de nombreux documents d’usage, en wattheures (Wh) ou kilowattheures (kWh). Le watt (W) est une unité de puissance : il renseigne sur le rythme du transfert d’énergie. Pour un appareil fonctionnant à puissance constante, 100 W pendant une heure correspondent à 100 Wh ; pendant deux heures, à 200 Wh. On utilise ici cette proportionnalité simple sans confondre W et Wh.\par
\pcFigure{figures/cbb217d150348436.pdf}{Consommations annuelles pédagogiques de deux solutions procurant le même service : 120 kWh et 36 kWh.}
\begin{pcEncadre}{definition}{À retenir}
1 kWh = 1 000 Wh. Une valeur de puissance seule ne donne pas l’énergie d’un usage : il faut aussi connaître la durée et les conditions de fonctionnement.\par
\end{pcEncadre}
\begin{pcEncadre}{method}{Une méthode concrète}
Lire l’unité, la période de mesure et le service rendu. Pour comparer des énergies, convertir dans la même unité avant de soustraire.\par
\end{pcEncadre}
\end{pcFiche}
\pcPage{L’essentiel}{Chapitre 16 · Énergie et choix responsables}
\phantomsection\label{ch-16-cours-3}
\begin{pcFiche}{cours-16-3}
\essSection{3}{Comparer à service équivalent}
Un choix d’équipement doit d’abord répondre au même besoin. Comparer deux éclairages demande par exemple un éclairage utile et une durée d’utilisation comparables. Réduire une durée inutile, améliorer l’isolation ou choisir un dispositif moins énergivore peut diminuer l’énergie nécessaire. Il faut néanmoins vérifier le résultat réel et considérer la durée de vie ainsi que les ressources nécessaires à la fabrication.\par
\pcFigure{figures/c0a738efc76f547e.pdf}{Une conversion comporte plusieurs transferts ; un dispositif ne convertit pas toute l’énergie reçue en effet utile.}
\begin{pcEncadre}{definition}{À retenir}
L’effet de serre naturel joue un rôle dans la température terrestre. L’enjeu étudié ici est la modification du bilan par l’augmentation de certains gaz, pas la disparition de tout effet de serre.\par
\end{pcEncadre}
\begin{pcEncadre}{method}{Une méthode concrète}
Construire un tableau : même service, énergie sur la même durée, durée de vie et limites des données. Formuler une conclusion proportionnée aux informations disponibles.\par
\end{pcEncadre}
\end{pcFiche}
\pcPage{Activité expérimentale}{Comparer deux éclairages fictifs}
\phantomsection\label{ch-16-activite}
\begin{pcFiche}{activite-16}
\begin{pcEncadre}{method}{Mission}
Calculer une économie et discuter ce que les chiffres permettent de conclure.\par
\end{pcEncadre}
\begin{questions}
\item Deux systèmes A et B assurent le même éclairage utile sur une année : A utilise 120 kWh, B 36 kWh. Les données sont fictives et servent à l’exercice.
\item Comparer les consommations en conservant le même service et la même période. Calculer la différence et la part économisée par rapport à A.
\item Dresser la liste des autres données utiles pour une décision d’équipement, sans inventer de prix ni de durée de vie.
\end{questions}
\pcFigure{figures/61c62ffc75c6764a.pdf}{Les barres représentent des énergies consommées sur une même année.}
\begin{questions}
\item Quelle énergie est économisée avec B sur un an ?
\item Que représente cette économie sur 120 kWh : plus ou moins de la moitié ?
\item Pourquoi cette seule comparaison ne suffit-elle pas à un bilan environnemental complet ?
\end{questions}
\end{pcFiche}
\pcPage{L’essentiel à retenir}{Chapitre 16 · Synthèse et autoévaluation}
\phantomsection\label{ch-16-bilan}
\begin{pcFiche}{synthese-16}
\begin{pcEncadre}{definition}{Repère 1}
Le renouvellement d’une ressource s’évalue à une échelle de temps.\par
\end{pcEncadre}
\begin{pcEncadre}{method}{Repère 2}
Aucune classification de ressource ne suffit à résumer tous les impacts.\par
\end{pcEncadre}
\begin{pcEncadre}{definition}{Repère 3}
W mesure une puissance ; Wh et kWh mesurent des énergies.\par
\end{pcEncadre}
\begin{pcEncadre}{method}{Repère 4}
1 kWh = 1 000 Wh ; la durée d’usage compte dans l’énergie consommée.\par
\end{pcEncadre}
\begin{pcEncadre}{definition}{Repère 5}
Une comparaison demande le même service et des conditions explicites.\par
\end{pcEncadre}
\begin{pcEncadre}{attention}{Une erreur à éviter}
« L’appareil le moins puissant utilise toujours le moins d’énergie. » La durée et le service rendu doivent aussi être comparés.\par
\end{pcEncadre}
\end{pcFiche}
\clearpage\markright{Exercices}
\begin{blocLivre}{exo-16-1}
\rubrique{Je démarre}
\begin{exercice}[1]{Classer des ressources}
\begin{questions}
\item Classe le vent, le pétrole et le Soleil selon leur renouvellement à l’échelle humaine.
\end{questions}
\end{exercice}\end{blocLivre}
\begin{blocLivre}{exo-16-2}
\begin{exercice}[1]{Lire les unités}
\begin{questions}
\item Associe W, Wh et kWh à puissance ou énergie.
\end{questions}
\end{exercice}\end{blocLivre}
\begin{blocLivre}{exo-16-3}
\begin{exercice}[1]{Une conversion d’énergie}
\begin{questions}
\item Convertis 2 kWh en Wh.
\item Convertis 500 Wh en kWh.
\end{questions}
\end{exercice}\end{blocLivre}
\begin{blocLivre}{exo-16-4}
\begin{exercice}[1]{Le bon critère}
\begin{questions}
\item Deux appareils n’assurent pas le même éclairage. Peut-on les classer uniquement par leur consommation ?
\item Quelle condition demander ?
\end{questions}
\end{exercice}\end{blocLivre}
\begin{blocLivre}{exo-16-5}
\begin{exercice}[1]{Un renouvelable sans impact ?}
\begin{questions}
\item Corrige : « Une éolienne ne nécessite aucune ressource parce que le vent est renouvelable. »
\end{questions}
\end{exercice}\end{blocLivre}
\begin{blocLivre}{exo-16-6}
\rubrique{Je m’entraîne}
\begin{exercice}[2]{Lire deux barres}
\pcFigure{figures/61c62ffc75c6764a.pdf}{}
\begin{questions}
\item Calcule l’écart entre A et B.
\item Quelle période concerne cet écart ?
\end{questions}
\end{exercice}\end{blocLivre}
\begin{blocLivre}{exo-16-7}
\begin{exercice}[2]{La durée de fonctionnement}
\begin{questions}
\item Un appareil maintient 50 W pendant 2 h. Avec le repère du cours, quelle énergie utilise-t-il en Wh ?
\item Que donnerait 1 h ?
\end{questions}
\end{exercice}\end{blocLivre}
\begin{blocLivre}{exo-16-8}
\begin{exercice}[2]{Passer à plusieurs jours}
\begin{questions}
\item Un usage consomme 40 Wh par jour. Quelle énergie correspond à 5 jours identiques ?
\item Exprime-la aussi en kWh.
\end{questions}
\end{exercice}\end{blocLivre}
\begin{blocLivre}{exo-16-9}
\begin{exercice}[2]{Réduire un usage inutile}
\begin{questions}
\item Une lampe reçoit de l’énergie pendant qu’une pièce est vide. Quel changement d’usage peut éviter ce transfert inutile ?
\item Quelle limite faut-il préciser ?
\end{questions}
\end{exercice}\end{blocLivre}
\begin{blocLivre}{exo-16-10}
\begin{exercice}[2]{Un stock non renouvelable}
\begin{questions}
\item L’uranium n’émet pas le même bilan direct qu’un combustible fossile lors de l’usage dans une centrale. Est-il pour autant renouvelable ?
\item Quelles deux questions faut-il distinguer ?
\end{questions}
\end{exercice}\end{blocLivre}
\begin{blocLivre}{exo-16-11}
\begin{exercice}[2]{Un même service sur deux ans}
\begin{questions}
\item A utilise 120 kWh par an et B 36 kWh, avec un usage identique chaque année. Quelle économie de B sur deux ans ?
\item Quelle hypothèse est nécessaire ?
\end{questions}
\end{exercice}\end{blocLivre}
\begin{blocLivre}{exo-16-12}
\rubrique{Je raisonne}
\begin{exercice}[3]{Moins puissant, mais plus longtemps}
\begin{questions}
\item A fonctionne à 20 W pendant 5 h ; B à 40 W pendant 2 h, à puissance constante. Compare les énergies en Wh.
\item Que montre cet exemple ?
\end{questions}
\end{exercice}\end{blocLivre}
\begin{blocLivre}{exo-16-13}
\begin{exercice}[3]{Une donnée publicitaire}
\begin{questions}
\item « 30 \% de moins » est annoncé sans référence. Quelles questions poser ?
\item Pourquoi la durée de mesure compte-t-elle ?
\end{questions}
\end{exercice}\end{blocLivre}
\begin{blocLivre}{exo-16-14}
\begin{exercice}[3]{Calculer une part économisée}
\begin{questions}
\item On passe de 200 kWh à 150 kWh pour le même service. Quelle économie réalise-t-on ?
\item Cette économie représente quelle part de 200 kWh ?
\end{questions}
\end{exercice}\end{blocLivre}
\begin{blocLivre}{exo-16-15}
\begin{exercice}[3]{Un bilan de décision}
\begin{questions}
\item Propose trois informations à réunir avant de remplacer un appareil encore fonctionnel par un autre.
\item Pourquoi la consommation à l’usage ne suffit-elle pas toujours ?
\end{questions}
\end{exercice}\end{blocLivre}
\pcPage{Défi · Prendre une initiative}{L’affiche d’une classe économe}
\phantomsection\label{defi-ex-16}
\begin{pcFiche}{defi-16}
Une classe veut réduire une consommation d’éclairage de 300 kWh par an. Une meilleure organisation permettrait d’éviter 60 kWh d’usage inutile ; un autre équipement, pour les mêmes heures restantes, permettrait encore 40 kWh d’économie. Les estimations sont fournies pour ce scénario précis.\par
\pcFigure{figures/a49852deeb58e915.pdf}{}
\begin{questions}
\item Calcule la consommation finale et l’économie totale.
\item Exprime l’économie comme une fraction simple de 300 kWh.
\item Rédige une affiche qui explique les hypothèses et n’annonce pas un impact nul.
\end{questions}
\begin{pcEncadre}{method}{Indice progressif}
Les deux économies du document concernent des étapes successives déjà définies ; ne pas additionner des pourcentages sans référence.\par
\end{pcEncadre}
\end{pcFiche}
\clearpage\pagestyle{otMain}
\begin{ouverture}[college4]{17}{Sources de lumière et vision}{Partie 4 · Des signaux pour observer et communiquer}{Que doit-il se passer pour que l’on voie un objet ?}
\objectif{Distinguer source primaire et objet diffusant.}
\objectif{Expliquer la vision par de la lumière qui entre dans l’œil.}
\objectif{Comparer matériaux transparents, translucides et opaques.}
\end{ouverture}
\pcPage{L’essentiel}{Chapitre 17 · Sources de lumière et vision}
\phantomsection\label{ch-17-cours-1}
\phantomsection\label{ess-17}
\begin{pcFiche}{cours-17-1}
\essSection{1}{Produire ou diffuser de la lumière}
Une source primaire émet de la lumière qu’elle produit : le Soleil et une lampe allumée en sont des exemples. Un objet diffusant renvoie dans plusieurs directions une partie de la lumière reçue. Une feuille éclairée, une planète ou la Lune sont visibles grâce à cette lumière diffusée. Un objet peut donc nous envoyer de la lumière sans être lui-même une source primaire.\par
\pcFigure{figures/dcf4bdddfb91f348.pdf}{Une source éclaire un objet ; une partie de la lumière diffusée par celui-ci peut atteindre l’œil.}
\begin{pcEncadre}{definition}{À retenir}
Un objet lumineux pour un observateur peut être une source primaire ou un objet diffusant. Préciser l’origine de la lumière évite l’ambiguïté.\par
\end{pcEncadre}
\begin{pcEncadre}{method}{Une méthode concrète}
Demander si l’objet produit la lumière ou reçoit celle d’une autre source. Ne pas confondre « brillant », « visible » et « source primaire ».\par
\end{pcEncadre}
\end{pcFiche}
\pcPage{L’essentiel}{Chapitre 17 · Sources de lumière et vision}
\phantomsection\label{ch-17-cours-2}
\begin{pcFiche}{cours-17-2}
\essSection{2}{La lumière doit entrer dans l’œil}
Pour voir un objet, de la lumière provenant de cet objet doit entrer dans l’œil. Un objet opaque non lumineux doit donc être éclairé et diffuser vers l’observateur une partie de la lumière reçue. L’œil n’envoie pas un rayon qui va chercher l’objet. Un obstacle opaque peut empêcher la lumière d’atteindre l’œil, même si l’objet reçoit de la lumière ailleurs dans la pièce.\par
\pcFigure{figures/dcf4bdddfb91f348.pdf}{Les flèches montrent un sens de propagation de la lumière, depuis la source jusqu’à l’œil via l’objet.}
\begin{pcEncadre}{definition}{À retenir}
Être éclairé ne garantit pas d’être visible pour tous les observateurs. Il faut aussi qu’un chemin de lumière relie l’objet à l’œil.\par
\end{pcEncadre}
\begin{pcEncadre}{method}{Une méthode concrète}
Tracer deux étapes : source \(\rightarrow\) objet, puis objet \(\rightarrow\) œil. Rechercher les obstacles éventuels sur chaque trajet.\par
\end{pcEncadre}
\end{pcFiche}
\pcPage{L’essentiel}{Chapitre 17 · Sources de lumière et vision}
\phantomsection\label{ch-17-cours-3}
\begin{pcFiche}{cours-17-3}
\essSection{3}{Traverser un matériau}
Un matériau transparent laisse passer de la lumière avec une transmission qui permet de voir nettement à travers dans les conditions étudiées. Un matériau translucide transmet mais diffuse la lumière : on ne voit pas nettement à travers. Un matériau opaque ne transmet pas la lumière visible de façon perceptible dans ces conditions ; il peut en absorber et en diffuser une partie. L’épaisseur et l’état du matériau influencent ces observations.\par
\pcFigure{figures/8cb7548c42defaeb.pdf}{Modèles simplifiés de la transmission de lumière visible à travers trois matériaux.}
\begin{pcEncadre}{definition}{À retenir}
Opaque ne veut pas dire noir ; une feuille blanche est opaque et diffuse de la lumière. Transparent ne veut pas dire invisible ni sans réflexion.\par
\end{pcEncadre}
\begin{pcEncadre}{method}{Une méthode concrète}
Comparer les matériaux devant le même objet et la même source adaptée. Ne jamais utiliser le Soleil ou un laser comme source d’observation directe.\par
\end{pcEncadre}
\end{pcFiche}
\pcPage{Activité expérimentale}{Voir à travers, voir l’objet}
\phantomsection\label{ch-17-activite}
\begin{pcFiche}{activite-17}
\begin{pcEncadre}{method}{Mission}
Distinguer transmission et diffusion avec une lampe adaptée.\par
\end{pcEncadre}
\begin{questions}
\item Le professeur fournit une source de faible puissance adaptée, un objet et des écrans transparent, translucide et opaque. Ne pas diriger une source intense vers les yeux.
\item Observer l’objet à travers chaque écran, avec un éclairage identique. Noter si sa forme est vue nettement, de façon floue ou pas à travers l’écran.
\item Dessiner le trajet de la lumière pour l’observation de l’objet éclairé, puis préciser quel écran interrompt ou modifie ce trajet.
\end{questions}
\pcFigure{figures/2158b9095894ca6c.pdf}{Observer une même scène à travers plusieurs matériaux ; les résultats dépendent aussi de leur épaisseur.}
\begin{questions}
\item Quel écran permet une image nette de l’objet ?
\item Un écran opaque blanc peut-il être lui-même visible ?
\item Pourquoi la flèche objet \(\rightarrow\) œil est-elle indispensable dans l’explication ?
\end{questions}
\end{pcFiche}
\pcPage{L’essentiel à retenir}{Chapitre 17 · Synthèse et autoévaluation}
\phantomsection\label{ch-17-bilan}
\begin{pcFiche}{synthese-17}
\begin{pcEncadre}{definition}{Repère 1}
Une source primaire produit la lumière qu’elle émet.\par
\end{pcEncadre}
\begin{pcEncadre}{method}{Repère 2}
Un objet diffusant renvoie une partie de la lumière qu’il reçoit.\par
\end{pcEncadre}
\begin{pcEncadre}{definition}{Repère 3}
Voir un objet exige que de la lumière provenant de lui entre dans l’œil.\par
\end{pcEncadre}
\begin{pcEncadre}{method}{Repère 4}
Transparent, translucide et opaque décrivent des comportements de transmission.\par
\end{pcEncadre}
\begin{pcEncadre}{definition}{Repère 5}
Les flèches d’un schéma lumineux suivent le sens de propagation de la lumière.\par
\end{pcEncadre}
\begin{pcEncadre}{attention}{Une erreur à éviter}
« L’œil éclaire l’objet pour le voir. » Dans ce modèle, la lumière entre dans l’œil ; elle n’en sort pas pour rechercher l’objet.\par
\end{pcEncadre}
\end{pcFiche}
\clearpage\markright{Exercices}
\begin{blocLivre}{exo-17-1}
\rubrique{Je démarre}
\begin{exercice}[1]{Deux types de sources}
\begin{questions}
\item Classe une lampe allumée et la Lune éclairée : source primaire ou objet diffusant.
\end{questions}
\end{exercice}\end{blocLivre}
\begin{blocLivre}{exo-17-2}
\begin{exercice}[1]{Le sens de la lumière}
\begin{questions}
\item Pour voir un livre éclairé, la lumière va-t-elle du livre vers l’œil ou de l’œil vers le livre ?
\item Que reçoit le livre au départ ?
\end{questions}
\end{exercice}\end{blocLivre}
\begin{blocLivre}{exo-17-3}
\begin{exercice}[1]{Un écran blanc}
\begin{questions}
\item Une feuille blanche qui bloque la vue à travers elle est-elle opaque ?
\item Peut-elle être visible quand elle est éclairée ?
\end{questions}
\end{exercice}\end{blocLivre}
\begin{blocLivre}{exo-17-4}
\begin{exercice}[1]{Nommer un matériau}
\begin{questions}
\item Un verre dépoli transmet de la lumière mais ne permet pas une vision nette à travers. Comment le qualifier ?
\item Comment qualifie-t-on une vitre permettant une vision nette ?
\end{questions}
\end{exercice}\end{blocLivre}
\begin{blocLivre}{exo-17-5}
\begin{exercice}[1]{Lire les flèches}
\pcFigure{figures/cd80bfc6ec09d32d.pdf}{}
\begin{questions}
\item Quel rôle joue l’objet central ?
\item Quelle flèche correspond à la lumière qui entre dans l’œil ?
\end{questions}
\end{exercice}\end{blocLivre}
\begin{blocLivre}{exo-17-6}
\rubrique{Je m’entraîne}
\begin{exercice}[2]{Dans l’obscurité}
\begin{questions}
\item Un objet non lumineux peut-il être vu dans une obscurité totale s’il n’émet aucune lumière ?
\item Quelle condition manque ?
\end{questions}
\end{exercice}\end{blocLivre}
\begin{blocLivre}{exo-17-7}
\begin{exercice}[2]{Un objet caché}
\pcFigure{figures/9eaef4335ed1b9e0.pdf}{Document de travail · modèle simplifié.}
\begin{questions}
\item Une lampe éclaire un livre, mais un écran opaque masque le livre à l’observateur. Pourquoi peut-il ne pas le voir ?
\item Quelle étape de la chaîne est interrompue ?
\end{questions}
\end{exercice}\end{blocLivre}
\begin{blocLivre}{exo-17-8}
\begin{exercice}[2]{Brillant ne veut pas dire primaire}
\begin{questions}
\item Une cuillère éclairée brille. Pourquoi cela ne prouve-t-il pas qu’elle produit sa propre lumière ?
\item Quelle origine de la lumière peut-on proposer ?
\end{questions}
\end{exercice}\end{blocLivre}
\begin{blocLivre}{exo-17-9}
\begin{exercice}[2]{Comparer les trois écrans}
\pcFigure{figures/2158b9095894ca6c.pdf}{}
\begin{questions}
\item Quel matériau interrompt la transmission représentée ?
\item Quel matériau disperse les directions transmises ?
\end{questions}
\end{exercice}\end{blocLivre}
\begin{blocLivre}{exo-17-10}
\begin{exercice}[2]{Éteindre la source}
\begin{questions}
\item Une lampe éclaire une balle non lumineuse. On éteint la seule source dans la pièce. Que devient la condition de visibilité ?
\item Précise l’hypothèse importante.
\end{questions}
\end{exercice}\end{blocLivre}
\begin{blocLivre}{exo-17-11}
\begin{exercice}[2]{Une planète visible}
\begin{questions}
\item Pourquoi une planète peut-elle être observée alors qu’elle n’est pas une étoile ?
\item Quel astre fournit principalement la lumière visible diffusée dans le Système solaire ?
\end{questions}
\end{exercice}\end{blocLivre}
\begin{blocLivre}{exo-17-12}
\rubrique{Je raisonne}
\begin{exercice}[3]{Un schéma à corriger}
\pcFigure{figures/86f8a5c97970f71a.pdf}{Document de travail · modèle simplifié.}
\begin{questions}
\item Un dessin montre uniquement une flèche œil \(\rightarrow\) objet pour expliquer la vision. Que faut-il modifier ?
\item Quelle autre information est nécessaire pour un objet diffusant ?
\end{questions}
\end{exercice}\end{blocLivre}
\begin{blocLivre}{exo-17-13}
\begin{exercice}[3]{Transparent et réflexion}
\begin{questions}
\item Peut-on voir le reflet d’une vitre qui laisse aussi voir à travers ?
\item Pourquoi cela ne contredit-il pas sa transparence ?
\end{questions}
\end{exercice}\end{blocLivre}
\begin{blocLivre}{exo-17-14}
\begin{exercice}[3]{L’épaisseur du matériau}
\begin{questions}
\item Une fine couche laisse passer de la lumière ; une grande épaisseur du même matériau en laisse très peu. Quelle précision faut-il ajouter au classement ?
\item Comment rendre l’essai comparable ?
\end{questions}
\end{exercice}\end{blocLivre}
\begin{blocLivre}{exo-17-15}
\begin{exercice}[3]{Préparer une démonstration}
\begin{questions}
\item Conçois une explication avec lampe, objet non lumineux et écran opaque pour montrer les deux conditions de vision.
\item Quel résultat attendre en masquant chacun des deux trajets ?
\end{questions}
\end{exercice}\end{blocLivre}
\pcPage{Défi · Prendre une initiative}{La boîte à regarder}
\phantomsection\label{defi-ex-17}
\begin{pcFiche}{defi-17}
Une boîte opaque contient une balle non lumineuse. Un trou sert à regarder ; un autre laisse entrer la lumière d’une lampe adaptée. On peut boucher chaque trou séparément. Prépare une explication des observations.\par
\pcFigure{figures/5d4e6a937a274a41.pdf}{}
\begin{questions}
\item Décris les deux trajets nécessaires pour voir la balle.
\item Prévois le résultat si le trou d’éclairage est bouché, sans autre source interne.
\item Explique pourquoi un trou de regard seul ne crée pas de lumière.
\end{questions}
\begin{pcEncadre}{method}{Indice progressif}
Identifier l’origine de la lumière et sa destination finale dans l’œil.\par
\end{pcEncadre}
\end{pcFiche}
\clearpage\pagestyle{otMain}
\begin{ouverture}[college4]{18}{Propagation de la lumière et ombres}{Partie 4 · Des signaux pour observer et communiquer}{Comment prévoir la position et la taille d’une ombre ?}
\objectif{Utiliser la propagation rectiligne dans un milieu transparent homogène.}
\objectif{Tracer des rayons orientés et reconnaître les limites du modèle.}
\objectif{Distinguer ombre propre, zone d’ombre et ombre portée.}
\end{ouverture}
\pcPage{L’essentiel}{Chapitre 18 · Propagation de la lumière et ombres}
\phantomsection\label{ch-18-cours-1}
\phantomsection\label{ess-18}
\begin{pcFiche}{cours-18-1}
\essSection{1}{Un modèle de propagation rectiligne}
Dans un milieu transparent homogène, la lumière se propage en ligne droite. On représente un trajet par un rayon : un trait rectiligne fléché dans le sens de propagation. Un faisceau est un ensemble de rayons ; le trait du schéma n’est pas un fil matériel. Pour laisser passer la lumière d’une petite source à travers plusieurs trous, les trous doivent être alignés avec la source dans le modèle.\par
\pcFigure{figures/70f381ff739f6c1a.pdf}{Trois trous alignés laissent passer un trajet lumineux direct. Source scolaire adaptée uniquement.}
\begin{pcEncadre}{definition}{À retenir}
La règle de ligne droite s’applique aux conditions annoncées. La lumière peut changer de direction à une interface ou sur un miroir ; ces situations ne sont pas modélisées ici.\par
\end{pcEncadre}
\begin{pcEncadre}{method}{Une méthode concrète}
Placer d’abord la source, les ouvertures et l’écran, puis tracer à la règle les trajets possibles avec des flèches.\par
\end{pcEncadre}
\end{pcFiche}
\pcPage{L’essentiel}{Chapitre 18 · Propagation de la lumière et ombres}
\phantomsection\label{ch-18-cours-2}
\begin{pcFiche}{cours-18-2}
\essSection{2}{Un objet opaque arrête certains trajets}
Un objet opaque intercepte une partie de la lumière d’une source. Derrière lui, certains points ne reçoivent pas la lumière directe de cette source : ils appartiennent à une zone d’ombre. Avec une petite source assimilée à un point, on peut tracer deux rayons limites passant au bord de l’objet. Leur prolongement délimite la zone d’ombre et sa rencontre avec un écran.\par
\pcFigure{figures/0705f5efefc3dc82.pdf}{Vue de côté avec une source ponctuelle S et un objet opaque ; les rayons limites sont prolongés jusqu’à l’écran.}
\begin{pcEncadre}{definition}{À retenir}
Un rayon limite est un outil géométrique. Il ne traverse pas l’objet opaque : il passe au bord et délimite les trajets bloqués.\par
\end{pcEncadre}
\begin{pcEncadre}{method}{Une méthode concrète}
Relier la source aux deux bords pertinents de l’objet, puis prolonger jusqu’à l’écran. Hachurer ou colorer seulement la zone privée de lumière directe.\par
\end{pcEncadre}
\end{pcFiche}
\pcPage{L’essentiel}{Chapitre 18 · Propagation de la lumière et ombres}
\phantomsection\label{ch-18-cours-3}
\begin{pcFiche}{cours-18-3}
\essSection{3}{Distinguer les différentes ombres}
L’ombre propre est la partie de la surface de l’objet non directement éclairée par la source. La zone d’ombre est la région de l’espace située derrière l’objet qui ne reçoit pas directement cette lumière. L’ombre portée est la tache obtenue lorsque cette zone rencontre un écran. À source et écran fixes, rapprocher l’objet de la source agrandit son ombre portée dans ce modèle de petite source.\par
\pcFigure{figures/83b67a7fd7986a89.pdf}{Une seule petite source est considérée. Les lumières diffusées par la pièce sont négligées dans ce modèle.}
\begin{pcEncadre}{definition}{À retenir}
Une ombre dépend de la source, de l’objet et de l’écran. Sa taille ne se déduit pas de la seule taille de l’objet.\par
\end{pcEncadre}
\begin{pcEncadre}{method}{Une méthode concrète}
Pour comparer deux essais, fixer deux positions et ne déplacer que le troisième élément. Préciser ce qui reste fixe avant d’annoncer une règle.\par
\end{pcEncadre}
\end{pcFiche}
\pcPage{Activité expérimentale}{Faire varier une ombre}
\phantomsection\label{ch-18-activite}
\begin{pcFiche}{activite-18}
\begin{pcEncadre}{method}{Mission}
Tester l’effet de la position d’un objet sur son ombre portée.\par
\end{pcEncadre}
\begin{questions}
\item Avec le professeur, utiliser une petite source adaptée, un objet opaque et un écran. Ne pas employer de laser ni regarder une source intense.
\item Garder source et écran fixes. Placer l’objet à deux positions entre eux, puis comparer la taille de l’ombre portée.
\item Tracer les deux constructions dans le cahier avec les mêmes échelles et le même objet.
\end{questions}
\pcFigure{figures/3f89b4be149ab996.pdf}{Le schéma est un modèle de côté ; il ne représente pas la lumière parasite de la pièce.}
\begin{questions}
\item Quel élément doit-on déplacer seul pour tester la règle annoncée ?
\item Quand l’objet se rapproche de la source, comment évolue l’ombre portée ?
\item Pourquoi faut-il garder la même source et le même écran ?
\end{questions}
\end{pcFiche}
\pcPage{L’essentiel à retenir}{Chapitre 18 · Synthèse et autoévaluation}
\phantomsection\label{ch-18-bilan}
\begin{pcFiche}{synthese-18}
\begin{pcEncadre}{definition}{Repère 1}
Dans un milieu transparent homogène, les trajets lumineux sont rectilignes.\par
\end{pcEncadre}
\begin{pcEncadre}{method}{Repère 2}
Un rayon est un trait fléché qui représente un trajet possible.\par
\end{pcEncadre}
\begin{pcEncadre}{definition}{Repère 3}
Un objet opaque intercepte des trajets lumineux.\par
\end{pcEncadre}
\begin{pcEncadre}{method}{Repère 4}
Ombre propre, zone d’ombre et ombre portée désignent des régions différentes.\par
\end{pcEncadre}
\begin{pcEncadre}{definition}{Repère 5}
La géométrie source-objet-écran détermine l’ombre dans le modèle choisi.\par
\end{pcEncadre}
\begin{pcEncadre}{attention}{Une erreur à éviter}
« L’ombre a toujours la taille de l’objet. » Sa taille dépend aussi des distances à la source et à l’écran.\par
\end{pcEncadre}
\end{pcFiche}
\clearpage\markright{Exercices}
\begin{blocLivre}{exo-18-1}
\rubrique{Je démarre}
\begin{exercice}[1]{Une condition de milieu}
\begin{questions}
\item Dans quelles conditions du cours représente-t-on la propagation par une droite ?
\item Quel symbole ajoute-t-on pour donner le sens ?
\end{questions}
\end{exercice}\end{blocLivre}
\begin{blocLivre}{exo-18-2}
\begin{exercice}[1]{Nommer un rayon}
\begin{questions}
\item Un rayon lumineux est-il un fil matériel ?
\item Quel est son rôle dans un schéma ?
\end{questions}
\end{exercice}\end{blocLivre}
\begin{blocLivre}{exo-18-3}
\begin{exercice}[1]{Les trous alignés}
\pcFigure{figures/3f9ef26636cab967.pdf}{}
\begin{questions}
\item Pourquoi l’alignement des trous est-il important ?
\item Que se passe-t-il si un écran masque le trajet direct ?
\end{questions}
\end{exercice}\end{blocLivre}
\begin{blocLivre}{exo-18-4}
\begin{exercice}[1]{L’ombre sur un mur}
\begin{questions}
\item Comment nomme-t-on la tache sombre obtenue sur un écran derrière un objet opaque ?
\item Comment nomme-t-on la partie non directement éclairée de l’objet ?
\end{questions}
\end{exercice}\end{blocLivre}
\begin{blocLivre}{exo-18-5}
\begin{exercice}[1]{Un objet opaque}
\begin{questions}
\item Un objet transparent donne-t-il nécessairement le même résultat qu’un objet opaque ?
\item Quelle propriété de l’objet opaque est utilisée dans le modèle ?
\end{questions}
\end{exercice}\end{blocLivre}
\begin{blocLivre}{exo-18-6}
\rubrique{Je m’entraîne}
\begin{exercice}[2]{Une ouverture déplacée}
\pcFigure{figures/62d9a8700e3bfecf.pdf}{}
\begin{questions}
\item Le trou de l’écran central est décalé. Pourquoi le trajet représenté s’arrête-t-il ?
\item Faut-il dessiner un rayon courbé pour qu’il rejoigne le trou ?
\end{questions}
\end{exercice}\end{blocLivre}
\begin{blocLivre}{exo-18-7}
\begin{exercice}[2]{Repérer les limites}
\pcFigure{figures/3f89b4be149ab996.pdf}{}
\begin{questions}
\item Où passent les deux rayons limites ?
\item À quoi sert leur prolongement ?
\end{questions}
\end{exercice}\end{blocLivre}
\begin{blocLivre}{exo-18-8}
\begin{exercice}[2]{Déplacer l’objet}
\begin{questions}
\item Source et écran sont fixes. L’objet se rapproche de la petite source. Comment évolue l’ombre portée ?
\item Quelle condition dois-tu conserver dans la réponse ?
\end{questions}
\end{exercice}\end{blocLivre}
\begin{blocLivre}{exo-18-9}
\begin{exercice}[2]{Un écran derrière l’objet}
\begin{questions}
\item Un écran est placé hors de la zone d’ombre. Reçoit-il nécessairement la même tache sombre ?
\item Pourquoi sa position compte-t-elle ?
\end{questions}
\end{exercice}\end{blocLivre}
\begin{blocLivre}{exo-18-10}
\begin{exercice}[2]{Deux essais comparables}
\begin{questions}
\item Un groupe change à la fois la lampe, l’objet et la distance de l’écran. Peut-il attribuer une variation d’ombre à la seule distance ?
\item Propose un essai mieux contrôlé.
\end{questions}
\end{exercice}\end{blocLivre}
\begin{blocLivre}{exo-18-11}
\begin{exercice}[2]{Le sens des flèches}
\begin{questions}
\item Pourquoi les flèches partent-elles de S vers l’écran dans la construction ?
\item Le sens de lecture de la page décide-t-il de leur orientation ?
\end{questions}
\end{exercice}\end{blocLivre}
\begin{blocLivre}{exo-18-12}
\rubrique{Je raisonne}
\begin{exercice}[3]{Une ombre moins nette}
\begin{questions}
\item L’expérience réelle montre des bords moins nets que le schéma. Donne deux explications compatibles avec les limites annoncées.
\end{questions}
\end{exercice}\end{blocLivre}
\clearpage
\begin{blocLivre}{exo-18-13}
\begin{exercice}[3]{Distinguer trois régions}
\begin{questions}
\item Explique la différence entre ombre propre, zone d’ombre et ombre portée sans employer trois fois seulement le mot « sombre ».
\end{questions}
\end{exercice}\end{blocLivre}
\begin{blocLivre}{exo-18-14}
\begin{exercice}[3]{Une construction impossible}
\pcFigure{figures/2973c974832181c2.pdf}{Document de travail · modèle simplifié.}
\begin{questions}
\item Un élève trace un rayon qui traverse le centre d’un objet opaque pour rejoindre l’écran. Quelle erreur commet-il ?
\item Comment construire les limites correctes ?
\end{questions}
\end{exercice}\end{blocLivre}
\begin{blocLivre}{exo-18-15}
\begin{exercice}[3]{Confronter modèle et mesure}
\begin{questions}
\item La construction prévoit une ombre plus grande lorsque l’objet approche de la source, mais la mesure paraît inchangée. Que vérifier avant de rejeter la règle ?
\item Quelle limite de mesure peut intervenir ?
\end{questions}
\end{exercice}\end{blocLivre}
\pcPage{Défi · Prendre une initiative}{Le théâtre d’ombres}
\phantomsection\label{defi-ex-18}
\begin{pcFiche}{defi-18}
Pour un petit théâtre d’ombres de classe, on souhaite agrandir la silhouette d’un même objet sans changer la source ni déplacer l’écran. La source est petite et les autres éclairages sont limités.\par
\pcFigure{figures/7a7475ea31ccba85.pdf}{}
\begin{questions}
\item Quel déplacement de l’objet proposer ?
\item Dessine les rayons limites avant et après.
\item Explique les limites de la prévision pour une lampe très étendue.
\end{questions}
\begin{pcEncadre}{method}{Indice progressif}
Comparer les angles sous lesquels les bords de l’objet sont vus depuis la petite source.\par
\end{pcEncadre}
\end{pcFiche}
\clearpage\pagestyle{otMain}
\begin{ouverture}[college4]{19}{Produire et transmettre un son}{Partie 4 · Des signaux pour observer et communiquer}{Pourquoi entend-on un objet qui vibre ?}
\objectif{Relier l’émission d’un son à une vibration.}
\objectif{Identifier le milieu de propagation et utiliser une vitesse fournie.}
\objectif{Distinguer hauteur d’un son, niveau sonore et conditions d’écoute.}
\end{ouverture}
\pcPage{L’essentiel}{Chapitre 19 · Produire et transmettre un son}
\phantomsection\label{ch-19-cours-1}
\phantomsection\label{ess-19}
\begin{pcFiche}{cours-19-1}
\essSection{1}{Une vibration à l’origine du son}
Une source sonore met en vibration un milieu matériel : corde, membrane, règle ou colonne d’air, par exemple. La perturbation se transmet de proche en proche dans le milieu et peut atteindre un récepteur comme l’oreille ou un microphone. L’air ne se déplace pas tout entier depuis la source jusqu’à l’oreille avec le son ; il transmet une perturbation. Une vibration peut rester peu visible tout en étant audible.\par
\pcFigure{figures/81a7d4712fdd1032.pdf}{Une règle vibre et met l’air en vibration ; la flèche indique la propagation, pas le transport de toute la matière.}
\begin{pcEncadre}{definition}{À retenir}
L’émetteur vibre, le milieu transmet et le récepteur reçoit. Ces trois rôles doivent être distingués dans l’explication.\par
\end{pcEncadre}
\begin{pcEncadre}{method}{Une méthode concrète}
Identifier ce qui vibre à la source et le milieu situé entre source et récepteur. Ne pas répondre seulement « le son va dans l’air ».\par
\end{pcEncadre}
\end{pcFiche}
\pcPage{L’essentiel}{Chapitre 19 · Produire et transmettre un son}
\phantomsection\label{ch-19-cours-2}
\begin{pcFiche}{cours-19-2}
\essSection{2}{Un milieu et une durée de propagation}
Le son se propage dans les gaz, les liquides et les solides, avec une vitesse qui dépend du milieu et de ses conditions. Il ne se propage pas dans le vide. Dans l’air proche de la température ambiante, on utilise ici la valeur arrondie de 340 m/s. Le son met donc un temps mesurable à parcourir une distance. Un écho correspond à un trajet aller puis retour après réflexion.\par
\pcFigure{figures/803676b91c818133.pdf}{Une sonnerie reste en fonctionnement sous une cloche ; raréfier l’air atténue la transmission par l’air. D’autres chemins par les supports peuvent subsister.}
\begin{pcEncadre}{definition}{À retenir}
Pour un écho, la distance parcourue par le son est deux fois la distance de l’obstacle dans le modèle d’aller-retour direct.\par
\end{pcEncadre}
\begin{pcEncadre}{method}{Une méthode concrète}
Utiliser distance = vitesse \(\times\) durée avec la valeur fournie. Dessiner le trajet avant de diviser par deux pour une distance à un obstacle.\par
\end{pcEncadre}
\end{pcFiche}
\pcPage{L’essentiel}{Chapitre 19 · Produire et transmettre un son}
\phantomsection\label{ch-19-cours-3}
\begin{pcFiche}{cours-19-3}
\essSection{3}{Aigu, grave, fort : des propriétés différentes}
La hauteur distingue un son aigu d’un son grave : une vibration plus rapide correspond à une fréquence plus élevée, exprimée en hertz (Hz). Le niveau sonore, exprimé en décibels (dB), renseigne sur une autre propriété liée à la force du son reçu. On peut donc avoir un son grave fort ou aigu faible. Les risques auditifs dépendent notamment du niveau et de la durée d’exposition, pas seulement du caractère agréable du son.\par
\pcFigure{figures/fefea017df4e810c.pdf}{Traces schématiques à mêmes échelles de temps et d’amplitude : B présente plus d’oscillations que A ; C a une amplitude plus grande que A.}
\begin{pcEncadre}{definition}{À retenir}
Réduire le niveau d’écoute, limiter la durée et s’éloigner des sources fortes protège l’audition. Une exposition gênante doit être interrompue ; ne pas attendre une douleur pour réduire le son.\par
\end{pcEncadre}
\begin{pcEncadre}{method}{Une méthode concrète}
Comparer d’abord le nombre d’oscillations sur une même durée, puis l’amplitude sur une même échelle. Ne pas additionner les décibels comme des longueurs.\par
\end{pcEncadre}
\end{pcFiche}
\pcPage{Activité expérimentale}{Observer et enregistrer une vibration}
\phantomsection\label{ch-19-activite}
\begin{pcFiche}{activite-19}
\begin{pcEncadre}{method}{Mission}
Mettre en relation une source vibrante et un signal enregistré à niveau modéré.\par
\end{pcEncadre}
\begin{questions}
\item Avec le professeur, faire vibrer doucement une règle maintenue sur une table ou utiliser une source scolaire adaptée. Garder les oreilles à distance raisonnable.
\item Observer le mouvement et, si le matériel le permet, enregistrer un signal avec un microphone. Conserver les réglages pour comparer deux essais.
\item Sans enregistrement, exploiter les trois traces schématiques proposées et leurs échelles identiques.
\end{questions}
\pcFigure{figures/ccf2d1287fdfb2dc.pdf}{A et C ont le même nombre d’oscillations ; B en présente davantage sur la même durée.}
\begin{questions}
\item Quelle trace correspond à une vibration plus rapide que A ?
\item Quelle trace garde le même rythme mais une plus grande amplitude ?
\item Pourquoi modifier le volume n’est-il pas la même chose que modifier la hauteur ?
\end{questions}
\end{pcFiche}
\pcPage{L’essentiel à retenir}{Chapitre 19 · Synthèse et autoévaluation}
\phantomsection\label{ch-19-bilan}
\begin{pcFiche}{synthese-19}
\begin{pcEncadre}{definition}{Repère 1}
Un son provient d’une source qui vibre.\par
\end{pcEncadre}
\begin{pcEncadre}{method}{Repère 2}
Sa propagation exige un milieu matériel ; elle n’a pas lieu dans le vide.\par
\end{pcEncadre}
\begin{pcEncadre}{definition}{Repère 3}
Dans l’air ambiant, on utilise souvent environ 340 m/s.\par
\end{pcEncadre}
\begin{pcEncadre}{method}{Repère 4}
Fréquence et hauteur diffèrent du niveau sonore.\par
\end{pcEncadre}
\begin{pcEncadre}{definition}{Repère 5}
Le niveau, la durée et la distance à la source comptent pour l’exposition sonore.\par
\end{pcEncadre}
\begin{pcEncadre}{attention}{Une erreur à éviter}
« Un son aigu est forcément plus fort. » La hauteur et le niveau sonore sont deux propriétés différentes.\par
\end{pcEncadre}
\end{pcFiche}
\clearpage\markright{Exercices}
\begin{blocLivre}{exo-19-1}
\rubrique{Je démarre}
\begin{exercice}[1]{À la source}
\begin{questions}
\item Quel mouvement d’une membrane de haut-parleur est à l’origine du son ?
\item Comment nomme-t-on un dispositif qui reçoit ce son pour l’enregistrer ?
\end{questions}
\end{exercice}\end{blocLivre}
\begin{blocLivre}{exo-19-2}
\begin{exercice}[1]{Le milieu nécessaire}
\begin{questions}
\item Le son peut-il se propager dans le vide ?
\item Donne deux types de milieux matériels où il peut se propager.
\end{questions}
\end{exercice}\end{blocLivre}
\begin{blocLivre}{exo-19-3}
\begin{exercice}[1]{Une vitesse avec son unité}
\begin{questions}
\item Quelle valeur arrondie utilise le cours pour le son dans l’air ambiant ?
\item Quelle unité accompagne 340 ?
\end{questions}
\end{exercice}\end{blocLivre}
\begin{blocLivre}{exo-19-4}
\begin{exercice}[1]{Deux propriétés}
\begin{questions}
\item Un son est grave et fort. Ces deux mots décrivent-ils la même propriété ?
\item Quelle propriété est associée à la fréquence ?
\end{questions}
\end{exercice}\end{blocLivre}
\begin{blocLivre}{exo-19-5}
\begin{exercice}[1]{Une protection raisonnée}
\begin{questions}
\item Cite deux moyens de réduire l’exposition à un son trop fort.
\item Faut-il attendre d’avoir mal pour agir ?
\end{questions}
\end{exercice}\end{blocLivre}
\begin{blocLivre}{exo-19-6}
\rubrique{Je m’entraîne}
\begin{exercice}[2]{Distance parcourue en une seconde}
\begin{questions}
\item Avec 340 m/s, quelle distance le son parcourt-il en 1 s dans le modèle ?
\item Et en 2 s ?
\end{questions}
\end{exercice}\end{blocLivre}
\begin{blocLivre}{exo-19-7}
\begin{exercice}[2]{Un trajet direct}
\begin{questions}
\item Deux points sont séparés de 170 m. Calcule la durée de propagation directe avec 340 m/s.
\item L’unité du résultat est-elle le mètre ?
\end{questions}
\end{exercice}\end{blocLivre}
\begin{blocLivre}{exo-19-8}
\begin{exercice}[2]{Comparer des traces}
\pcFigure{figures/ccf2d1287fdfb2dc.pdf}{}
\begin{questions}
\item Quelle trace a une fréquence plus élevée que A ?
\item Entre A et C, qu’est-ce qui change dans le modèle ?
\end{questions}
\end{exercice}\end{blocLivre}
\begin{blocLivre}{exo-19-9}
\begin{exercice}[2]{La cloche et le vide}
\pcFigure{figures/e500fed44c170644.pdf}{}
\begin{questions}
\item Pourquoi le son transmis par l’air s’atténue-t-il lorsqu’on raréfie fortement l’air ?
\item Pourquoi faut-il mentionner les supports du montage ?
\end{questions}
\end{exercice}\end{blocLivre}
\begin{blocLivre}{exo-19-10}
\begin{exercice}[2]{Un écho simple}
\pcFigure{figures/f7ff4aa87373e30c.pdf}{}
\begin{questions}
\item Un écho revient 1 s après un son bref. Quelle distance totale le son a-t-il parcourue ?
\item À quelle distance se situe l’obstacle dans le modèle aller-retour ?
\end{questions}
\end{exercice}\end{blocLivre}
\begin{blocLivre}{exo-19-11}
\begin{exercice}[2]{La matière ne voyage pas tout entière}
\begin{questions}
\item Pourquoi la flèche de propagation n’indique-t-elle pas que l’air de la source arrive entièrement dans l’oreille ?
\item Qu’est-ce qui se transmet ?
\end{questions}
\end{exercice}\end{blocLivre}
\clearpage
\begin{blocLivre}{exo-19-12}
\rubrique{Je raisonne}
\begin{exercice}[3]{Un délai de 3 secondes}
\begin{questions}
\item Un signal lumineux et un son sont émis simultanément à distance. Le son arrive environ 3 s plus tard ; on néglige ici le temps de la lumière. Estime la distance.
\item Quelle approximation utilises-tu ?
\end{questions}
\end{exercice}\end{blocLivre}
\begin{blocLivre}{exo-19-13}
\begin{exercice}[3]{Un écho plus rapide}
\begin{questions}
\item Le retour d’un écho prend 0,4 s. Calcule la distance de l’obstacle.
\item Quelle erreur donnerait une distance deux fois trop grande ?
\end{questions}
\end{exercice}\end{blocLivre}
\begin{blocLivre}{exo-19-14}
\begin{exercice}[3]{Comparer des enregistrements}
\pcFigure{figures/346715a5a3080bb8.pdf}{Document de travail · modèle simplifié.}
\begin{questions}
\item Deux traces ont le même nombre apparent d’oscillations, mais l’une représente 1 s et l’autre 2 s. Peut-on conclure à la même hauteur ?
\item Quel contrôle manque ?
\end{questions}
\end{exercice}\end{blocLivre}
\begin{blocLivre}{exo-19-15}
\begin{exercice}[3]{Une idée à discuter}
\begin{questions}
\item « Ma musique préférée ne peut pas abîmer l’audition. » Explique pourquoi cette affirmation est fausse.
\item Quelles données comptent davantage que le goût musical ?
\end{questions}
\end{exercice}\end{blocLivre}
\pcPage{Défi · Prendre une initiative}{Le mur repéré par un écho}
\phantomsection\label{defi-ex-19}
\begin{pcFiche}{defi-19}
Un document d’acquisition donne un écho 0,6 s après l’émission d’un son bref. Le trajet est direct jusqu’à un mur puis revient au récepteur situé près de la source. On utilise 340 m/s et on néglige les autres réflexions.\par
\pcFigure{figures/2b0eef694cb98dd7.pdf}{}
\begin{questions}
\item Dessine l’aller et le retour.
\item Calcule la distance parcourue par le son puis la distance au mur.
\item Explique pourquoi une mesure réelle peut différer légèrement du modèle.
\end{questions}
\begin{pcEncadre}{method}{Indice progressif}
La durée fournie concerne le trajet complet, pas seulement l’aller.\par
\end{pcEncadre}
\end{pcFiche}
\clearpage\pagestyle{otMain}
\begin{ouverture}[college4]{20}{Un signal pour transmettre une information}{Partie 4 · Des signaux pour observer et communiquer}{Comment un signal devient-il une information ?}
\objectif{Identifier émetteur, signal, milieu et récepteur.}
\objectif{Utiliser un code simple pour transmettre un message.}
\objectif{Distinguer signal et information, et exploiter un délai de transmission.}
\end{ouverture}
\pcPage{L’essentiel}{Chapitre 20 · Un signal pour transmettre une information}
\phantomsection\label{ch-20-cours-1}
\phantomsection\label{ess-20}
\begin{pcFiche}{cours-20-1}
\essSection{1}{De l’émetteur au récepteur}
Un signal est une variation physique que l’on peut émettre, transmettre et recevoir. Il peut être lumineux ou sonore dans les exemples du livre. L’émetteur le produit ; un milieu ou un trajet permet sa propagation ; un récepteur le détecte. L’information est le sens que l’on attribue à ce signal à l’aide de règles et d’un contexte. Le même signal peut porter des informations différentes selon le code partagé.\par
\pcFigure{figures/4fa5323551c9cd30.pdf}{Un signal physique et le message interprété ne sont pas la même chose.}
\begin{pcEncadre}{definition}{À retenir}
Un voyant rouge ne signifie pas à lui seul la même chose sur tous les appareils. Il faut connaître sa fonction ou sa légende.\par
\end{pcEncadre}
\begin{pcEncadre}{method}{Une méthode concrète}
Présenter la chaîne : émetteur \(\rightarrow\) nature du signal et trajet \(\rightarrow\) récepteur \(\rightarrow\) information interprétée. Vérifier que les participants partagent le même code.\par
\end{pcEncadre}
\end{pcFiche}
\pcPage{L’essentiel}{Chapitre 20 · Un signal pour transmettre une information}
\phantomsection\label{ch-20-cours-2}
\begin{pcFiche}{cours-20-2}
\essSection{2}{Coder avec deux états}
On peut coder un message avec deux états distinguables : lampe allumée ou éteinte, notés 1 et 0. Il faut aussi convenir des durées, du début et de la fin du message et de la façon de regrouper les valeurs. Avec des groupes de deux valeurs, les combinaisons 00, 01, 10 et 11 permettent de coder quatre symboles. Un état physique n’acquiert son sens qu’avec le code annoncé.\par
\pcFigure{figures/d157ef27a61ec9a4.pdf}{Huit intervalles identiques portent huit valeurs. Le découpage en groupes doit être convenu avant la transmission.}
\begin{pcEncadre}{definition}{À retenir}
Un zéro peut porter de l’information : il ne signifie pas nécessairement une absence de message. Un signal trop faible ou masqué peut entraîner une erreur de lecture.\par
\end{pcEncadre}
\begin{pcEncadre}{method}{Une méthode concrète}
Écrire le tableau du code, repérer les intervalles, puis regrouper les valeurs sans en sauter. Relire le résultat avec le récepteur.\par
\end{pcEncadre}
\end{pcFiche}
\pcPage{L’essentiel}{Chapitre 20 · Un signal pour transmettre une information}
\phantomsection\label{ch-20-cours-3}
\begin{pcFiche}{cours-20-3}
\essSection{3}{Mesurer et transmettre prend du temps}
Un capteur transforme une grandeur mesurée en un signal exploitable. Par exemple, une température peut donner un signal électrique ensuite utilisé pour un affichage. Une variation continue et une suite de valeurs codées ne se décrivent pas de la même manière. La propagation prend aussi un temps : la lumière voyage dans le vide à environ 300 000 km/s, bien plus vite que le son dans l’air, sans être instantanée.\par
\pcFigure{figures/62de05a7e9247aa7.pdf}{Le capteur relie une grandeur physique à un signal ; l’affichage utilise une conversion et un code.}
\begin{pcEncadre}{definition}{À retenir}
Une valeur affichée dépend du capteur, de son réglage et de sa précision. La lecture d’un écran ne dispense pas de vérifier l’unité ni les conditions de mesure.\par
\end{pcEncadre}
\begin{pcEncadre}{method}{Une méthode concrète}
Pour un délai de propagation, utiliser la distance et la vitesse dans des unités compatibles. Préciser si la durée concerne un aller simple ou un aller-retour.\par
\end{pcEncadre}
\end{pcFiche}
\pcPage{Activité expérimentale}{Transmettre quatre symboles}
\phantomsection\label{ch-20-activite}
\begin{pcFiche}{activite-20}
\begin{pcEncadre}{method}{Mission}
Mettre à l’épreuve un code lumineux simple à faible intensité.\par
\end{pcEncadre}
\begin{questions}
\item Convenir du code A = 10, B = 01, C = 11, D = 00 ; choisir des intervalles de 1 s et une longueur de message connue. Utiliser une lampe adaptée, sans viser les yeux.
\item Préparer le message ACDA en groupes de deux valeurs, puis transmettre avec un repère de départ partagé.
\item Faire noter les huit valeurs par le récepteur avant de traduire. Comparer les erreurs possibles si le début ou les intervalles sont mal repérés.
\end{questions}
\pcFigure{figures/36fef9ca654da749.pdf}{Code de l’activité : 10 | 11 | 00 | 10. Chaque valeur occupe ici un intervalle de 1 s.}
\begin{questions}
\item Quel message le récepteur doit-il retrouver ?
\item Quelle durée faut-il pour les huit intervalles, hors préparation ?
\item Pourquoi le début, la durée des intervalles et le tableau de code sont-ils indispensables ?
\end{questions}
\end{pcFiche}
\pcPage{L’essentiel à retenir}{Chapitre 20 · Synthèse et autoévaluation}
\phantomsection\label{ch-20-bilan}
\begin{pcFiche}{synthese-20}
\begin{pcEncadre}{definition}{Repère 1}
Un signal physique est produit, transmis et reçu.\par
\end{pcEncadre}
\begin{pcEncadre}{method}{Repère 2}
L’information dépend d’un code et d’un contexte partagés.\par
\end{pcEncadre}
\begin{pcEncadre}{definition}{Repère 3}
Un codage à deux états nécessite un repérage des durées et des groupes.\par
\end{pcEncadre}
\begin{pcEncadre}{method}{Repère 4}
Un capteur associe une grandeur mesurée à un signal exploitable.\par
\end{pcEncadre}
\begin{pcEncadre}{definition}{Repère 5}
La propagation lumineuse ou sonore a une vitesse finie et un délai.\par
\end{pcEncadre}
\begin{pcEncadre}{attention}{Une erreur à éviter}
« Une lampe éteinte ne transmet rien. » Dans un code convenu, un intervalle éteint peut représenter la valeur 0 et porter de l’information.\par
\end{pcEncadre}
\end{pcFiche}
\clearpage\markright{Exercices}
\begin{blocLivre}{exo-20-1}
\rubrique{Je démarre}
\begin{exercice}[1]{Nommer les rôles}
\begin{questions}
\item Une lampe clignote et un capteur la détecte. Qui est l’émetteur et qui est le récepteur ?
\item Quelle est la nature du signal ?
\end{questions}
\end{exercice}\end{blocLivre}
\begin{blocLivre}{exo-20-2}
\begin{exercice}[1]{Signal ou information ?}
\begin{questions}
\item Pour une sonnette, distingue la variation physique entendue et le sens habituel donné au son.
\end{questions}
\end{exercice}\end{blocLivre}
\begin{blocLivre}{exo-20-3}
\begin{exercice}[1]{Deux valeurs possibles}
\begin{questions}
\item Quels états d’une lampe peuvent coder 1 et 0 ?
\item L’état 0 peut-il faire partie d’un message ?
\end{questions}
\end{exercice}\end{blocLivre}
\begin{blocLivre}{exo-20-4}
\begin{exercice}[1]{Quatre combinaisons}
\begin{questions}
\item Liste tous les groupes différents de deux valeurs 0 ou 1.
\item Combien de symboles peuvent-ils représenter avec un code sans ambiguïté ?
\end{questions}
\end{exercice}\end{blocLivre}
\begin{blocLivre}{exo-20-5}
\begin{exercice}[1]{Lire les valeurs}
\pcFigure{figures/6994b66b13a29818.pdf}{}
\begin{questions}
\item Écris la suite des huit valeurs.
\item Combien de groupes de deux contient-elle ?
\end{questions}
\end{exercice}\end{blocLivre}
\begin{blocLivre}{exo-20-6}
\rubrique{Je m’entraîne}
\begin{exercice}[2]{Décoder un message}
\begin{questions}
\item Avec A = 10, B = 01, C = 11, D = 00, décode 10 | 01 | 11.
\item Pourquoi les séparations sont-elles utiles ?
\end{questions}
\end{exercice}\end{blocLivre}
\begin{blocLivre}{exo-20-7}
\begin{exercice}[2]{Coder un autre message}
\begin{questions}
\item Avec le même code, écris DAB en groupes.
\item Combien d’intervalles sont nécessaires ?
\end{questions}
\end{exercice}\end{blocLivre}
\begin{blocLivre}{exo-20-8}
\begin{exercice}[2]{Calculer la durée}
\begin{questions}
\item Huit valeurs occupent chacune 0,5 s. Quelle durée totale faut-il, hors préparation ?
\item Quelle donnée modifierait ce résultat ?
\end{questions}
\end{exercice}\end{blocLivre}
\begin{blocLivre}{exo-20-9}
\begin{exercice}[2]{Un capteur de température}
\pcFigure{figures/d7b34218d2642183.pdf}{}
\begin{questions}
\item Quelle grandeur est mesurée par le capteur ?
\item Pourquoi faut-il vérifier l’unité sur l’affichage ?
\end{questions}
\end{exercice}\end{blocLivre}
\begin{blocLivre}{exo-20-10}
\begin{exercice}[2]{Un délai lumineux fourni}
\begin{questions}
\item La lumière parcourt 600 000 km dans le vide. Avec 300 000 km/s, quelle durée de propagation obtient-on ?
\item La propagation est-elle instantanée ?
\end{questions}
\end{exercice}\end{blocLivre}
\begin{blocLivre}{exo-20-11}
\begin{exercice}[2]{Un code non partagé}
\begin{questions}
\item Deux groupes donnent des lettres différentes à 10. Pourquoi peuvent-ils lire des messages différents à partir du même signal ?
\item Que faut-il convenir avant ?
\end{questions}
\end{exercice}\end{blocLivre}
\begin{blocLivre}{exo-20-12}
\rubrique{Je raisonne}
\begin{exercice}[3]{La lumière du Soleil}
\begin{questions}
\item On utilise 150 millions de km et 300 000 km/s. Calcule la durée approximative du trajet Soleil-Terre.
\item Exprime-la en minutes et secondes.
\end{questions}
\end{exercice}\end{blocLivre}
\begin{blocLivre}{exo-20-13}
\begin{exercice}[3]{Un décalage de début}
\begin{questions}
\item Pourquoi commencer à lire une valeur trop tard peut-il changer tout un message codé par groupes de deux ?
\item Propose une règle de départ.
\end{questions}
\end{exercice}\end{blocLivre}
\begin{blocLivre}{exo-20-14}
\begin{exercice}[3]{Un signal masqué}
\begin{questions}
\item Une lumière ambiante empêche de distinguer lampe allumée et éteinte. Le tableau de code suffit-il à corriger le problème ?
\item Quel aspect de la chaîne faut-il améliorer ?
\end{questions}
\end{exercice}\end{blocLivre}
\begin{blocLivre}{exo-20-15}
\begin{exercice}[3]{Deux sens pour le même voyant}
\begin{questions}
\item Un voyant rouge signifie « enregistrement » sur un appareil et « défaut » sur un autre. Est-ce une contradiction physique ?
\item Quelle distinction explique la situation ?
\end{questions}
\end{exercice}\end{blocLivre}
\pcPage{Défi · Prendre une initiative}{Un message que l’on peut vérifier}
\phantomsection\label{defi-ex-20}
\begin{pcFiche}{defi-20}
Tu dois transmettre le mot BACA avec le code A = 10, B = 01, C = 11, D = 00. Chaque valeur dure 0,5 s. La longueur est connue du récepteur. Prévois ensuite une méthode simple pour vérifier qu’il a reçu le bon message.\par
\pcFigure{figures/4216f83c8cbc933f.pdf}{}
\begin{questions}
\item Écris les huit valeurs dans l’ordre.
\item Calcule la durée de transmission, hors préparation et vérification.
\item Propose une règle de confirmation et explique sa limite.
\end{questions}
\begin{pcEncadre}{method}{Indice progressif}
Une répétition ou un retour du message peut révéler une erreur, sans garantir à lui seul l’absence de toute erreur.\par
\end{pcEncadre}
\end{pcFiche}
\clearpage\gdef\otChapNum{0}\pagestyle{otFront}
\pcPage{Bilans des quatre thèmes}{Mobiliser ses connaissances dans des situations variées}
\phantomsection\label{bilans-debut}
\pdfbookmark[0]{Bilans des quatre thèmes}{bm-bilans-debut}
Ces quatre bilans proposent cinq missions chacun. Pour chaque mission, réponds aux deux questions et justifie ta réponse. Un barème indicatif est proposé dans les corrigés.\par
\otLigneRessource{bilan-A}{Bilan A · Matière et Système solaire}
\otLigneRessource{bilan-B}{Bilan B · Mouvement et interactions}
\otLigneRessource{bilan-C}{Bilan C · Énergie et circuits}
\otLigneRessource{bilan-D}{Bilan D · Lumière, sons et information}
\begin{pcEncadre}{method}{Relire avant de corriger}
Ai-je précisé les conditions, les unités, le sens des flèches et les étapes du raisonnement ? Une valeur juste sans explication ne répond pas toujours à toute la question.\par
\end{pcEncadre}
\pcPage{Bilan A}{Matière et Système solaire}
\phantomsection\label{bilan-A}
\begin{pcFiche}{bilan-A}
\begin{blocLivre}{mission-A-1}
\pcMission{A}{1}{Les passages de l’eau}
\begin{questions}
\item Nomme le passage solide \(\rightarrow\) liquide et le passage gaz \(\rightarrow\) liquide.
\item Dans un sachet fermé sans fuite, la masse de la glace est 40 g. Quelle masse d’eau obtient-on après fusion ?
\end{questions}
\end{blocLivre}
\begin{blocLivre}{mission-A-2}
\pcMission{A}{2}{Le contenu d’un flacon}
\begin{questions}
\item Un flacon contient eau, sel entièrement dissous et sable visible. Est-il globalement homogène ou hétérogène à l’œil nu ?
\item Décris deux étapes pour récupérer sable puis sel.
\end{questions}
\end{blocLivre}
\begin{blocLivre}{mission-A-3}
\pcMission{A}{3}{Deux volumes à lire}
\begin{questions}
\item Une éprouvette passe de 26 mL à 40 mL après immersion complète d’un solide. Calcule son volume.
\item Donne deux conditions nécessaires pour que la mesure corresponde au volume du solide.
\end{questions}
\end{blocLivre}
\begin{blocLivre}{mission-A-4}
\pcMission{A}{4}{Une conclusion de chimie}
\begin{questions}
\item Un gaz trouble l’eau de chaux initialement limpide dans le test fourni. Quelle espèce est présente ?
\item Pourquoi cela ne prouve-t-il pas que le gaz est pur ?
\end{questions}
\end{blocLivre}
\begin{blocLivre}{mission-A-5}
\pcMission{A}{5}{Une place dans l’espace}
\begin{questions}
\item Classe Soleil, Terre et Lune par catégorie d’astre.
\item À l’échelle 1 cm pour 50 millions de km, place la Terre dont la distance moyenne au Soleil vaut environ 150 millions de km.
\end{questions}
\end{blocLivre}
\end{pcFiche}
\pcPage{Bilan B}{Mouvement et interactions}
\phantomsection\label{bilan-B}
\begin{pcFiche}{bilan-B}
\begin{blocLivre}{mission-B-1}
\pcMission{B}{1}{Décrire avec un repère}
\begin{questions}
\item Un passager reste assis dans un train. Décris sa situation par rapport au siège puis au quai lorsque le train avance.
\item Pourquoi les deux descriptions ne se contredisent-elles pas ?
\end{questions}
\end{blocLivre}
\begin{blocLivre}{mission-B-2}
\pcMission{B}{2}{Lire des positions}
\pcFigure{figures/ff02b2b8cc9b0134.pdf}{}
\begin{questions}
\item Les positions sont prises à durées égales. Quelle ligne correspond à une valeur de vitesse constante ?
\item Quelle ligne montre une augmentation de cette valeur, et sur quel indice t’appuies-tu ?
\end{questions}
\end{blocLivre}
\begin{blocLivre}{mission-B-3}
\pcMission{B}{3}{Un déplacement chiffré}
\begin{questions}
\item Un mobile parcourt 90 m en 30 s. Calcule sa vitesse moyenne.
\item Exprime cette vitesse en km/h.
\end{questions}
\end{blocLivre}
\begin{blocLivre}{mission-B-4}
\pcMission{B}{4}{Un aller avec une pause}
\begin{questions}
\item Une sortie fait 2 km en 30 min, pause comprise. Calcule sa vitesse moyenne en km/h.
\item La moyenne permet-elle d’affirmer une vitesse constante pendant la sortie ?
\end{questions}
\end{blocLivre}
\begin{blocLivre}{mission-B-5}
\pcMission{B}{5}{Un livre au repos}
\begin{questions}
\item Cite deux objets qui interagissent avec un livre posé sur une table, en négligeant l’air.
\item Classe les interactions en contact ou à distance.
\end{questions}
\end{blocLivre}
\end{pcFiche}
\pcPage{Bilan C}{Énergie et circuits}
\phantomsection\label{bilan-C}
\begin{pcFiche}{bilan-C}
\begin{blocLivre}{mission-C-1}
\pcMission{C}{1}{Une chaîne à compléter}
\begin{questions}
\item Une pile alimente une lampe. Nomme la réserve initiale et l’effet lumineux recherché.
\item Indique un autre transfert à ajouter au bilan.
\end{questions}
\end{blocLivre}
\begin{blocLivre}{mission-C-2}
\pcMission{C}{2}{Une branche ouverte}
\pcFigure{figures/2ded9452e6318e61.pdf}{}
\begin{questions}
\item Quelle lampe conserve un chemin fermé avec le générateur ?
\item Où placer une commande qui coupe les deux lampes ?
\end{questions}
\end{blocLivre}
\begin{blocLivre}{mission-C-3}
\pcMission{C}{3}{Un test à interpréter}
\begin{questions}
\item Un métal peint n’allume pas la lampe d’un circuit de test. Donne une cause possible autre que « le métal est isolant ».
\item Cite un contrôle indispensable.
\end{questions}
\end{blocLivre}
\begin{blocLivre}{mission-C-4}
\pcMission{C}{4}{Un bilan de 120 J}
\begin{questions}
\item Un appareil sans stockage reçoit 120 J et transfère 45 J utilement. Quelle autre énergie faut-il compter ?
\item Cette énergie a-t-elle disparu ?
\end{questions}
\end{blocLivre}
\begin{blocLivre}{mission-C-5}
\pcMission{C}{5}{Une économie documentée}
\begin{questions}
\item Deux solutions comparables utilisent 180 kWh et 120 kWh sur une année. Calcule l’économie.
\item Pourquoi cette comparaison ne constitue-t-elle pas tout le bilan environnemental ?
\end{questions}
\end{blocLivre}
\end{pcFiche}
\pcPage{Bilan D}{Lumière, sons et information}
\phantomsection\label{bilan-D}
\begin{pcFiche}{bilan-D}
\begin{blocLivre}{mission-D-1}
\pcMission{D}{1}{Expliquer la vision}
\pcFigure{figures/9c3300c051b4f703.pdf}{}
\begin{questions}
\item Nomme les deux étapes nécessaires pour voir un objet non lumineux éclairé.
\item Quel sens doit avoir la flèche qui rejoint l’œil ?
\end{questions}
\end{blocLivre}
\begin{blocLivre}{mission-D-2}
\pcMission{D}{2}{Une ombre à prévoir}
\begin{questions}
\item Dans quel type de milieu utilise-t-on une propagation rectiligne ?
\item À source et écran fixes, comment varie l’ombre portée si l’objet se rapproche de la petite source ?
\end{questions}
\end{blocLivre}
\begin{blocLivre}{mission-D-3}
\pcMission{D}{3}{La distance d’un écho}
\begin{questions}
\item Un écho revient après 0,5 s. Avec 340 m/s, calcule la distance totale parcourue.
\item Déduis la distance à l’obstacle dans le modèle aller-retour direct.
\end{questions}
\end{blocLivre}
\begin{blocLivre}{mission-D-4}
\pcMission{D}{4}{Un code partagé}
\begin{questions}
\item Avec A = 10, B = 01, C = 11 et D = 00, décode 11 | 10 | 01.
\item Chaque valeur dure 1 s. Quelle durée occupent les six valeurs ?
\end{questions}
\end{blocLivre}
\begin{blocLivre}{mission-D-5}
\pcMission{D}{5}{Écouter en préservant son audition}
\begin{questions}
\item Pourquoi un son apprécié peut-il présenter un risque ?
\item Cite deux façons de réduire l’exposition.
\end{questions}
\end{blocLivre}
\end{pcFiche}
\partieCorriges
\clearpage\corrigesChapitre{1}{La matière et ses états}
\begin{blocLivre}{cor-1-1}
\begin{corrige}{1}{Associer une matière à son état}
L’acier de la règle est solide ; le lait et l’huile sont liquides ; l’air est gazeux. Le glaçon est de l’eau solide. La vapeur d’eau est de l’eau à l’état gazeux, invisible. L’état dépend des conditions, notamment de la température.\par
\end{corrige}\end{blocLivre}
\begin{blocLivre}{cor-1-2}
\begin{corrige}{2}{Les bons mots}
Forme ; volume. Un liquide possède un volume propre mais prend la forme de la partie du récipient qu’il occupe. Un gaz occupe tout l’espace qui lui est accessible. On compare ici les récipients sans changer la température ni comprimer la matière.\par
\end{corrige}\end{blocLivre}
\begin{blocLivre}{cor-1-3}
\begin{corrige}{3}{Dessiner une surface libre}
Dans les deux récipients, la surface libre étendue est horizontale. Elle ne reste pas parallèle au bord du récipient incliné. On trace un trait horizontal et on le légende « surface libre de l’eau au repos ». La courbure près des parois n’est pas étudiée ici.\par
\begin{pcEncadre}{method}{Pour comprendre}
Repère d’abord l’horizontale de la feuille. Le récipient peut être incliné, mais pas la grande surface libre d’un liquide au repos.\par
\end{pcEncadre}
\pcFigure{figures/df269ec6e47083ab.pdf}{Modèle de correction · les proportions sont indicatives sauf mention contraire.}
\end{corrige}\end{blocLivre}
\begin{blocLivre}{cor-1-4}
\begin{corrige}{4}{De l’air dans une seringue}
Le volume occupé par l’air passe de 50 mL à 30 mL. L’embout est fermé et ne laisse pas fuir l’air : la même quantité d’air occupe moins de place. Le gaz est compressible. La diminution de volume ne signifie pas que l’air a disparu.\par
\pcFigure{figures/dbbdc1cf32213dde.pdf}{Modèle de correction · les proportions sont indicatives sauf mention contraire.}
\end{corrige}\end{blocLivre}
\begin{blocLivre}{cor-1-5}
\begin{corrige}{5}{Vrai ou faux ? Justifier}
Faux : l’air et la vapeur d’eau sont invisibles. Faux : la vapeur d’eau est gazeuse ; le brouillard visible contient de fines gouttelettes liquides. Vrai : une surface libre suffisamment étendue est plane et horizontale au repos, loin des parois.\par
\end{corrige}\end{blocLivre}
\begin{blocLivre}{cor-1-6}
\begin{corrige}{6}{Peut-on reconnaître un liquide à sa forme ?}
Non. Un ensemble de grains peut s’écouler et épouser la forme du récipient. Chaque grain de sable est pourtant solide et conserve sa propre forme. Il faut préciser l’objet observé et combiner plusieurs critères.\par
\end{corrige}\end{blocLivre}
\begin{blocLivre}{cor-1-7}
\begin{corrige}{7}{Le papier reste sec}
De l’air est piégé dans le verre. Il occupe une partie du volume intérieur et empêche l’eau de remplir immédiatement tout le verre. Si l’on incline suffisamment le verre, l’air peut sortir sous forme de bulles ; l’eau prend alors davantage de place dans le verre et peut mouiller le papier.\par
\begin{pcEncadre}{method}{Pour comprendre}
Suis les places occupées par l’air et par l’eau. Une matière invisible peut empêcher une autre matière d’occuper le même volume.\par
\end{pcEncadre}
\pcFigure{figures/992a247d3beea186.pdf}{Modèle de correction · les proportions sont indicatives sauf mention contraire.}
\end{corrige}\end{blocLivre}
\begin{blocLivre}{cor-1-8}
\begin{corrige}{8}{Changer de récipient}
Le volume final est de 120 mL. Dans le récipient large, l’eau s’étale davantage et la hauteur peut être plus faible. La hauteur seule ne mesure pas le volume : la forme et les dimensions du récipient interviennent.\par
\pcFigure{figures/d84809871d7bd12e.pdf}{Modèle de correction · les proportions sont indicatives sauf mention contraire.}
\end{corrige}\end{blocLivre}
\begin{blocLivre}{cor-1-9}
\begin{corrige}{9}{Observation ou interprétation ?}
« Des bulles quittent le verre » et « Le papier est sec » décrivent des observations. « L’air occupait une partie du verre » est une interprétation. Une conclusion possible : l’air est une matière invisible qui occupe une place ; il peut être déplacé et remplacé par de l’eau.\par
\end{corrige}\end{blocLivre}
\begin{blocLivre}{cor-1-10}
\begin{corrige}{10}{Lire un modèle}
L’état liquide. Dans le modèle du gaz, les particules sont plus espacées ; leur taille n’a pas augmenté. Ce n’est pas une photographie, mais une représentation simplifiée qui aide à expliquer les propriétés observées. Les points ne représentent pas des grains de sable visibles.\par
\pcFigure{figures/281e473284f774b5.pdf}{Modèle de correction · les proportions sont indicatives sauf mention contraire.}
\end{corrige}\end{blocLivre}
\begin{blocLivre}{cor-1-11}
\begin{corrige}{11}{Tester une idée correctement}
Le test n’est pas convaincant : dans la seringue ouverte, l’air peut sortir. Utiliser deux seringues adaptées, sans aiguille, de même taille et fermées, avec une poussée douce comparable. L’eau doit être sans bulle d’air, sinon on comprimerait surtout la bulle. L’expérience est réalisée avec le professeur, sans forcer.\par
\end{corrige}\end{blocLivre}
\clearpage
\begin{blocLivre}{cor-1-12}
\begin{corrige}{12}{L’air possède-t-il une masse ?}
239,6 g - 238,4 g = 1,2 g. La masse a augmenté après l’admission d’air. Si seul le contenu en air a changé, cette augmentation montre que l’air possède une masse. Être invisible ne signifie pas être dépourvu de matière.\par
\end{corrige}\end{blocLivre}
\begin{blocLivre}{cor-1-13}
\begin{corrige}{13}{Une légende scientifique}
La vapeur d’eau est invisible. Le brouillard visible contient des gouttelettes d’eau liquide en suspension dans l’air. Une légende correcte serait : « De fines gouttelettes d’eau liquide rendent le brouillard visible. L’air contient aussi de la vapeur d’eau invisible. »\par
\end{corrige}\end{blocLivre}
\begin{blocLivre}{cor-1-14}
\begin{corrige}{14}{Le sable met la règle à l’épreuve}
Le sable est un contre-exemple : on peut le verser, pourtant chaque grain est solide. L’ensemble change de forme quand les grains se déplacent les uns par rapport aux autres. Une description rigoureuse précise donc « un ensemble de grains solides » plutôt que « un liquide ».\par
\pcFigure{figures/e2ca89e45732b69c.pdf}{Modèle de correction · les proportions sont indicatives sauf mention contraire.}
\end{corrige}\end{blocLivre}
\begin{blocLivre}{cor-1-15}
\begin{corrige}{15}{Identifier trois échantillons}
A est un solide compact : forme et volume propres. B est un liquide : il change de forme, conserve son volume et se comprime très peu dans les conditions du test. C est un gaz : il occupe tout l’espace accessible et son volume diminue lorsqu’on le comprime. Un ensemble de grains montre qu’un seul critère, comme prendre la forme d’un récipient, peut tromper.\par
\end{corrige}\end{blocLivre}
\pcPage{Activité et défi}{Chapitre 01 · Éléments de réponse}
\phantomsection\label{solutions-1}
\begin{pcFiche}{solutions-1}
\begin{pcEncadre}{method}{CHAPITRE 1 · LE VERRE PROTÈGE-T-IL LE PAPIER ?}
Le papier peut rester sec lorsque le gobelet est descendu verticalement : l’air piégé empêche l’eau d’occuper tout l’intérieur. Si le gobelet est assez incliné, des bulles s’échappent et l’eau entre davantage. Ces observations mettent en évidence une matière invisible qui occupe un volume. Un protocole valide doit distinguer les observations de l’explication.\par
\end{pcEncadre}
\pcRetour{ch-1-activite}{Revenir à l’activité}
\phantomsection\label{defi-cor-1}
\pcSousTitre{Le trajet de l’air}
Une solution consiste à placer l’ouverture de B au-dessus de celle de A, puis à incliner doucement A pour libérer des bulles. Les bulles remontent et entrent dans B ; elles y remplacent une partie de l’eau.\par
\pcFigure{figures/6473550b13c99cd2.pdf}{État après un transfert suffisant : A contient davantage d’eau et B contient de l’air.}
\begin{pcEncadre}{definition}{CE QU’IL FALLAIT JUSTIFIER}
L’air est une matière qui occupe un volume et peut se déplacer. L’eau peut être remplacée par de l’air dans un récipient. La réussite dépend de la position des ouvertures : les bulles doivent pouvoir atteindre l’intérieur de B.\par
\end{pcEncadre}
Accepter tout montage équivalent, correctement légendé et compatible avec le trajet ascendant des bulles.\par
\pcRetour{defi-ex-1}{Revenir au défi}
\end{pcFiche}
\clearpage\corrigesChapitre{2}{Mesurer une masse et un volume}
\begin{blocLivre}{cor-2-1}
\begin{corrige}{1}{Choisir la grandeur et l’instrument}
Une balance, avec une masse en g ou en kg. Pour un peu d’eau, une éprouvette graduée et une lecture en mL conviennent. Le cm\textsuperscript{3} est une unité de volume : il ne mesure pas une masse.\par
\end{corrige}\end{blocLivre}
\begin{blocLivre}{cor-2-2}
\begin{corrige}{2}{Lire une pesée}
La masse du liquide est de 37,6 g. La tare a remis l’affichage à zéro avec le récipient déjà posé : sa masse a été prise en compte par la balance. L’affichage représente donc uniquement la masse ajoutée.\par
\end{corrige}\end{blocLivre}
\begin{blocLivre}{cor-2-3}
\begin{corrige}{3}{Lire une éprouvette}
Entre 20 mL et 30 mL, il y a 5 intervalles égaux : une petite division vaut 2 mL. Le bas du ménisque atteint 24 mL. On écrit V = 24 mL, et non 24 g.\par
\begin{pcEncadre}{method}{Pour comprendre}
Compte les intervalles, pas les traits : cinq intervalles séparent 20 mL et 30 mL, donc chaque intervalle vaut (30 − 20) \(\div\) 5 = 2 mL.\par
\end{pcEncadre}
\pcFigure{figures/98a329ea090aad5f.pdf}{Modèle de correction · les proportions sont indicatives sauf mention contraire.}
\end{corrige}\end{blocLivre}
\begin{blocLivre}{cor-2-4}
\begin{corrige}{4}{Convertir sans confondre}
1,25 kg = 1 250 g ; 750 g = 0,750 kg. 2,4 L = 2 400 mL ; 350 mL = 0,350 L. 18 mL = 18 cm\textsuperscript{3} ; 3 L = 3 dm\textsuperscript{3}. La grandeur ne change pas pendant une conversion.\par
\end{corrige}\end{blocLivre}
\begin{blocLivre}{cor-2-5}
\begin{corrige}{5}{Corriger l’unité}
L’éprouvette fournit un volume, pas une masse. On peut écrire : « Le volume du liquide est V = 35 mL. » Pour sa masse, il faut une pesée. Une égalité numérique entre mL et g n’est pas une règle générale.\par
\end{corrige}\end{blocLivre}
\begin{blocLivre}{cor-2-6}
\begin{corrige}{6}{Peser sans utiliser la tare}
La masse totale réunit le récipient et le liquide. Masse du liquide = masse totale - masse du récipient = 187,3 g - 62,8 g = 124,5 g. On soustrait des masses exprimées dans la même unité.\par
\end{corrige}\end{blocLivre}
\begin{blocLivre}{cor-2-7}
\begin{corrige}{7}{Deux méthodes, un même résultat}
Méthode A : 118,7 g - 45,2 g = 73,5 g. Méthode B : environ 73,5 g, si la même quantité de liquide est versée et si les mesures sont réalisées dans les mêmes conditions. Les deux méthodes retirent l’effet de la masse du récipient.\par
\end{corrige}\end{blocLivre}
\begin{blocLivre}{cor-2-8}
\begin{corrige}{8}{Mesurer un petit solide}
Avant : 32 mL. Après : 46 mL. Le solide déplace 46 mL - 32 mL = 14 mL d’eau. Son volume vaut donc 14 cm\textsuperscript{3}. L’eau et l’objet occupent ensemble le volume final ; ce volume final n’est pas celui du solide seul.\par
\begin{pcEncadre}{method}{Pour comprendre}
L’augmentation du niveau correspond seulement au volume ajouté par le solide entièrement immergé. On soustrait la lecture avant immersion à celle après immersion.\par
\end{pcEncadre}
\pcFigure{figures/07e0c4aee7da0a60.pdf}{Modèle de correction · les proportions sont indicatives sauf mention contraire.}
\end{corrige}\end{blocLivre}
\begin{blocLivre}{cor-2-9}
\begin{corrige}{9}{Un repère pour l’eau}
Environ 250 g, car pour l’eau usuelle à température ambiante, 1 mL correspond approximativement à 1 g. Ce repère ne s’applique pas à tous les liquides. La masse de 250 mL d’huile doit être mesurée ou déterminée à partir de données propres à cette huile.\par
\end{corrige}\end{blocLivre}
\begin{blocLivre}{cor-2-10}
\begin{corrige}{10}{Exploiter une proportionnalité}
Les masses sont approximativement 60 g et 80 g. Le volume est multiplié par 4 entre 20 mL et 80 mL ; la masse l’est aussi : 20 g \(\times\) 4 = 80 g. Il faut garder la même substance, ici l’eau, et la même température. Ces valeurs pédagogiques sont arrondies.\par
\end{corrige}\end{blocLivre}
\begin{blocLivre}{cor-2-11}
\begin{corrige}{11}{Choisir une verrerie adaptée}
L’éprouvette de 50 mL convient : sa capacité suffit pour 38 mL et ses divisions de 1 mL permettent une lecture plus fine. Celle de 1 L possède une capacité inutilement grande ici et des divisions trop larges pour distinguer facilement 38 mL.\par
\end{corrige}\end{blocLivre}
\clearpage
\begin{blocLivre}{cor-2-12}
\begin{corrige}{12}{Comparer sans mélanger les grandeurs}
A a le plus grand volume : 45 mL contre 30 mL. B a la plus grande masse : 60 g contre 40 g. Les substances sont différentes ; un plus grand volume ne garantit donc pas une plus grande masse. Le mot « lourd » doit ici être remplacé par une comparaison précise des masses.\par
\end{corrige}\end{blocLivre}
\begin{blocLivre}{cor-2-13}
\begin{corrige}{13}{Un protocole à discuter}
Non. Une partie du bois dépasse de l’eau ; seule la partie immergée déplace de l’eau. Le protocole simple exige un objet entièrement immergé. On peut choisir un solide compact qui s’immerge complètement, ou faire adapter le protocole par le professeur. On ne compte pas le volume d’un doigt ajouté pour maintenir l’objet.\par
\pcFigure{figures/a4a062060ae3acf4.pdf}{Modèle de correction · les proportions sont indicatives sauf mention contraire.}
\end{corrige}\end{blocLivre}
\begin{blocLivre}{cor-2-14}
\begin{corrige}{14}{Une identification prudente}
Oui. Doubler 25 mL donne 50 mL et doubler 22,5 g donne 45 g : les données sont compatibles avec la même proportionnalité. Cela ne prouve pas l’identité du liquide : plusieurs substances peuvent présenter des valeurs voisines et toute mesure possède une précision limitée. D’autres observations ou tests seraient nécessaires.\par
\end{corrige}\end{blocLivre}
\begin{blocLivre}{cor-2-15}
\begin{corrige}{15}{Réunir deux mesures indépendantes}
Masse du solide : 125,6 g - 91,4 g = 34,2 g. Volume : 52 mL - 38 mL = 14 mL, soit 14 cm\textsuperscript{3}. Les valeurs 34,2 et 14 sont différentes : il n’existe pas de correspondance universelle 1 mL = 1 g. On a mesuré deux grandeurs distinctes sur le même échantillon.\par
\begin{pcEncadre}{method}{Pour comprendre}
Fais deux calculs indépendants : une différence de masses donne une masse ; une différence de volumes donne un volume. Ne mélange pas les unités.\par
\end{pcEncadre}
\pcFigure{figures/db09f4b30b65117c.pdf}{Modèle de correction · les proportions sont indicatives sauf mention contraire.}
\end{corrige}\end{blocLivre}
\pcPage{Activité et défi}{Chapitre 02 · Éléments de réponse}
\phantomsection\label{solutions-2}
\begin{pcFiche}{solutions-2}
\begin{pcEncadre}{method}{CHAPITRE 2 · PRÉPARER 50 mL D’EAU}
Une masse voisine de 50 g est attendue pour 50 mL d’eau à température ambiante. Les valeurs réelles dépendent de la précision, du ménisque, de la tare et des pertes au transfert. Une table de résultats doit contenir les mesures obtenues, pas des valeurs inventées pour correspondre à la prévision. Comparer les essais et expliquer les écarts avec le protocole.\par
\end{pcEncadre}
\pcRetour{ch-2-activite}{Revenir à l’activité}
\phantomsection\label{defi-cor-2}
\pcSousTitre{Passer du lot à une bille}
\begin{pcEncadre}{definition}{1 · MESURER UN VOLUME TOTAL}
Le lot déplace 50 mL - 40 mL = 10 mL d’eau. Ce volume total est plus facile à distinguer sur les graduations que le volume d’une bille isolée.\par
\end{pcEncadre}
\begin{pcEncadre}{method}{2 · TROUVER UN VOLUME MOYEN}
Les billes sont supposées identiques : 10 mL \(\div\) 20 = 0,5 mL par bille, soit 0,5 cm\textsuperscript{3}.\par
\end{pcEncadre}
\begin{pcEncadre}{definition}{3 · TROUVER UNE MASSE MOYENNE}
78,0 g \(\div\) 20 = 3,9 g par bille. Cette division concerne la masse ; elle ne convertit pas un volume en masse.\par
\end{pcEncadre}
\begin{pcEncadre}{attention}{4 · DISCUTER LA MESURE}
Les lectures restent approximatives. Il faut éviter les bulles, immerger toutes les billes et ne pas perdre d’eau. Le calcul donne une moyenne et n’invente pas une graduation plus fine. Si les billes diffèrent, on ne connaît pas exactement chaque bille.\par
\end{pcEncadre}
\pcRetour{defi-ex-2}{Revenir au défi}
\end{pcFiche}
\clearpage\corrigesChapitre{3}{Les changements d’état de l’eau}
\begin{blocLivre}{cor-3-1}
\begin{corrige}{1}{Nommer les passages}
La glace subit une fusion. Le passage de l’eau liquide à la glace est une solidification. On donne le nom du passage, et non seulement le nom de l’état final.\par
\pcFigure{figures/da279b44dcb5e183.pdf}{Modèle de correction · les proportions sont indicatives sauf mention contraire.}
\end{corrige}\end{blocLivre}
\begin{blocLivre}{cor-3-2}
\begin{corrige}{2}{La buée sur une vitre}
La vapeur d’eau se liquéfie : c’est une liquéfaction. Les gouttelettes sont liquides. La vapeur invisible présente dans l’air ne doit pas être confondue avec la buée visible.\par
\end{corrige}\end{blocLivre}
\begin{blocLivre}{cor-3-3}
\begin{corrige}{3}{Lire une température}
L’axe vertical porte la température \(\theta\), en °C. Le palier est situé à 0 °C. On lit la valeur sur l’axe vertical, et non une durée sur l’axe horizontal.\par
\end{corrige}\end{blocLivre}
\begin{blocLivre}{cor-3-4}
\begin{corrige}{4}{Évaporation ou ébullition ?}
Le linge sèche par évaporation ; l’eau n’a pas besoin d’atteindre 100 °C. Une vaporisation avec des bulles de vapeur dans tout le liquide est une ébullition.\par
\end{corrige}\end{blocLivre}
\begin{blocLivre}{cor-3-5}
\begin{corrige}{5}{Choisir la bonne unité}
La température vaut 0 °C. La minute est une unité de durée et ne mesure pas une température. Même lorsque la valeur numérique est nulle, l’unité est indispensable.\par
\end{corrige}\end{blocLivre}
\begin{blocLivre}{cor-3-6}
\begin{corrige}{6}{Peser pendant la fusion}
On doit retrouver 180,0 g pour l’ensemble fermé. Il ne faut ni fuite, ni ajout, ni perte de matière. La masse de l’ensemble se conserve pendant cette transformation physique.\par
\end{corrige}\end{blocLivre}
\begin{blocLivre}{cor-3-7}
\begin{corrige}{7}{Deux grandeurs à distinguer}
Non. L’eau solide occupe davantage de volume que l’eau liquide de même masse. Le niveau renseigne sur le volume occupé ; il ne démontre pas une augmentation de masse.\par
\pcFigure{figures/4bf3cc5c97b47074.pdf}{Modèle de correction · les proportions sont indicatives sauf mention contraire.}
\end{corrige}\end{blocLivre}
\begin{blocLivre}{cor-3-8}
\begin{corrige}{8}{Durée du palier}
La durée est 8 - 3 = 5 min. Pour de l’eau pure en fusion, on peut observer simultanément de la glace et de l’eau liquide. La durée n’est pas la température du changement.\par
\end{corrige}\end{blocLivre}
\begin{blocLivre}{cor-3-9}
\begin{corrige}{9}{Dessiner une courbe}
Placer les points (0 ; -5), (1 ; 0), (2 ; 0), (3 ; 0), (4 ; 5), puis les relier dans l’ordre. La portion horizontale de 1 à 3 min est compatible avec une fusion à 0 °C si glace et eau liquide coexistent.\par
\begin{pcEncadre}{method}{Pour comprendre}
Écris le temps sur l’axe horizontal et la température sur l’axe vertical. Choisis des graduations régulières, place les cinq points puis relie-les dans l’ordre chronologique.\par
\end{pcEncadre}
\pcFigure{figures/107e6d91513d1085.pdf}{Modèle de correction · les proportions sont indicatives sauf mention contraire.}
\end{corrige}\end{blocLivre}
\begin{blocLivre}{cor-3-10}
\begin{corrige}{10}{Le récipient ouvert}
250 - 242 = 8 g ont quitté l’ensemble pesé, si aucune autre perte n’a eu lieu. La vapeur d’eau s’est dispersée dans l’air. Le récipient ouvert n’est pas un système fermé : il faut compter aussi l’eau partie dans l’environnement.\par
\end{corrige}\end{blocLivre}
\begin{blocLivre}{cor-3-11}
\begin{corrige}{11}{Le bon placement du thermomètre}
Le fond peut être plus chaud que le contenu et fausser la température recherchée. Le capteur doit mesurer le milieu étudié selon les consignes du matériel. Noter aussi la date et la présence des différents états, par exemple glace et eau liquide.\par
\end{corrige}\end{blocLivre}
\begin{blocLivre}{cor-3-12}
\begin{corrige}{12}{Un glaçon dans un sachet}
Peser un sachet fermé et étanche contenant un glaçon, le laisser fondre, sécher son extérieur et repeser le même ensemble. Une fuite ou de l’eau ajoutée à l’extérieur invalide la comparaison. On compare les masses avec la précision de la balance.\par
\end{corrige}\end{blocLivre}
\begin{blocLivre}{cor-3-13}
\begin{corrige}{13}{Une eau qui bout à 94 °C}
Le thermomètre n’est pas forcément défectueux : la température d’ébullition dépend de la pression. Il faut préciser la pression, notamment la pression atmosphérique normale pour la valeur voisine de 100 °C. Une vérification de l’instrument reste possible.\par
\end{corrige}\end{blocLivre}
\begin{blocLivre}{cor-3-14}
\begin{corrige}{14}{Un palier trop vite annoncé}
Un point ne suffit pas à montrer une température stable pendant une durée. Il faut plusieurs mesures successives à intervalles réguliers et l’observation des états présents. Un palier de fusion est étayé par une température presque constante et la coexistence du solide et du liquide.\par
\end{corrige}\end{blocLivre}
\begin{blocLivre}{cor-3-15}
\begin{corrige}{15}{Interpréter le modèle}
On ne fait pas grossir les particules. Dans un modèle simple de changement d’état, la nature et le nombre des particules d’eau sont conservés. Leur organisation et leurs possibilités de mouvement changent. Le modèle n’est pas une photographie.\par
\pcFigure{figures/281e473284f774b5.pdf}{Modèle de correction · les proportions sont indicatives sauf mention contraire.}
\end{corrige}\end{blocLivre}
\pcPage{Activité et défi}{Chapitre 03 · Éléments de réponse}
\phantomsection\label{solutions-3}
\begin{pcFiche}{solutions-3}
\begin{pcEncadre}{method}{Suivre la fusion de la glace}
Le palier s’étend de 1 à 4 min à 0 °C. Si l’on observe simultanément de la glace et de l’eau liquide, il correspond à la fusion. La courbe seule ne décrit pas tout le contenu du récipient. Après disparition de la glace, l’eau liquide peut se réchauffer : la température augmente. Les mesures expérimentales peuvent osciller légèrement autour de 0 °C.\par
\end{pcEncadre}
\pcRetour{ch-3-activite}{Revenir à l’activité}
\phantomsection\label{defi-cor-3}
\begin{pcEncadre}{definition}{Enquêter sur une bouteille déformée}
Le camarade confond masse et volume. À quantité d’eau conservée, la masse reste identique, mais la glace occupe davantage de volume. La place disponible dans une bouteille presque pleine devient insuffisante. Pour vérifier la masse, utiliser une expérience de fusion dans un récipient adapté, fermé et étanche ; ne pas congeler une bouteille pleine et rigide.\par
\end{pcEncadre}
\pcRetour{defi-ex-3}{Revenir au défi}
\end{pcFiche}
\clearpage\corrigesChapitre{4}{Corps purs, mélanges et composition de l’air}
\begin{blocLivre}{cor-4-1}
\begin{corrige}{1}{Corps pur ou mélange ?}
L’eau seule est un corps pur dans le modèle ; l’eau salée et l’air sec sont des mélanges. Le critère est le nombre d’espèces chimiques présentes, pas le nombre d’états physiques.\par
\end{corrige}\end{blocLivre}
\begin{blocLivre}{cor-4-2}
\begin{corrige}{2}{Nommer l’aspect}
Le jus avec pulpe est hétérogène à l’œil nu. L’eau salée sans cristaux visibles est homogène à cette échelle. Dans les deux cas, plusieurs espèces sont présentes.\par
\end{corrige}\end{blocLivre}
\begin{blocLivre}{cor-4-3}
\begin{corrige}{3}{Les deux gaz principaux}
Le diazote représente environ 78 \% et le dioxygène environ 21 \% du volume de l’air sec. Les autres gaz représentent ensemble environ 1 \%.\par
\end{corrige}\end{blocLivre}
\begin{blocLivre}{cor-4-4}
\begin{corrige}{4}{Lire un modèle}
Le dessin avec une seule couleur représente un corps pur ; celui à deux couleurs représente un mélange. Chaque couleur code une espèce chimique, et non la couleur visible réelle du gaz ou du liquide.\par
\pcFigure{figures/2b735c49a64f52ea.pdf}{Modèle de correction · les proportions sont indicatives sauf mention contraire.}
\end{corrige}\end{blocLivre}
\begin{blocLivre}{cor-4-5}
\begin{corrige}{5}{Déjouer une affirmation}
Faux. De l’eau contenant du sel entièrement dissous peut être incolore et transparente tout en contenant plusieurs espèces. L’aspect ne suffit pas pour établir la pureté.\par
\end{corrige}\end{blocLivre}
\begin{blocLivre}{cor-4-6}
\begin{corrige}{6}{Lire une étiquette}
Oui : la composition mentionne plusieurs espèces. Une eau peut respecter les critères de potabilité tout en étant un mélange ; la potabilité et la pureté chimique répondent à des questions différentes.\par
\end{corrige}\end{blocLivre}
\begin{blocLivre}{cor-4-7}
\begin{corrige}{7}{Air sec dans une seringue}
Environ 21 mL correspondent au dioxygène et 78 mL au diazote, à conditions de mesure identiques. Il s’agit d’une description en proportions de volumes ; les gaz ne forment pas des couches distinctes.\par
\end{corrige}\end{blocLivre}
\begin{blocLivre}{cor-4-8}
\begin{corrige}{8}{Changer de quantité}
Les volumes correspondants sont environ 42 mL de dioxygène et 156 mL de diazote. Les pourcentages restent 21 \% et 78 \% : c’est la quantité totale qui a doublé.\par
\end{corrige}\end{blocLivre}
\begin{blocLivre}{cor-4-9}
\begin{corrige}{9}{Agiter puis laisser reposer}
L’apparition durable de deux phases après repos est décisive. Le mélange est alors hétérogène. Une dispersion temporaire de petites gouttes ne prouve pas que les deux liquides sont miscibles.\par
\pcFigure{figures/9a252fdceeeb3df9.pdf}{Modèle de correction · les proportions sont indicatives sauf mention contraire.}
\end{corrige}\end{blocLivre}
\begin{blocLivre}{cor-4-10}
\begin{corrige}{10}{Faire le total}
100 - 78 - 21 = 1 \%. Ce reste comprend plusieurs gaz, dont l’argon et le dioxyde de carbone. Le dioxyde de carbone n’en représente qu’une partie.\par
\end{corrige}\end{blocLivre}
\begin{blocLivre}{cor-4-11}
\begin{corrige}{11}{Préciser l’échelle}
Il faut préciser l’échelle d’observation. Un aspect uniforme à l’œil nu peut masquer une dispersion visible à une autre échelle. Les deux descriptions ne sont pas contradictoires si leurs conditions sont annoncées.\par
\end{corrige}\end{blocLivre}
\begin{blocLivre}{cor-4-12}
\begin{corrige}{12}{Deux flacons identiques}
Sur de petits échantillons, le professeur peut faire évaporer l’eau : un résidu solide est attendu pour l’eau salée. L’absence de résidu visible ne prouve pas à elle seule une pureté parfaite ; la quantité et la sensibilité de la méthode comptent.\par
\end{corrige}\end{blocLivre}
\begin{blocLivre}{cor-4-13}
\begin{corrige}{13}{Un modèle à corriger}
Le dessin ne représente pas sa composition de mélange s’il annonce une seule espèce. Utiliser plusieurs sortes de points avec une légende, en respectant approximativement les proportions pour un échantillon assez grand. Ne pas confondre la taille des points et la proportion.\par
\begin{pcEncadre}{method}{Pour comprendre}
Une couleur code ici une espèce chimique. Changer la quantité de points d’une seule couleur ne transforme pas un corps pur en mélange.\par
\end{pcEncadre}
\pcFigure{figures/2a10d8ca7b9c06ed.pdf}{Modèle de correction · les proportions sont indicatives sauf mention contraire.}
\end{corrige}\end{blocLivre}
\begin{blocLivre}{cor-4-14}
\begin{corrige}{14}{Passer à un litre}
1 L = 1 000 mL, soit dix fois 100 mL. Environ 780 mL correspondent au diazote et 210 mL au dioxygène. Les 10 mL restants correspondent à d’autres gaz, selon les valeurs arrondies retenues.\par
\end{corrige}\end{blocLivre}
\begin{blocLivre}{cor-4-15}
\begin{corrige}{15}{Argumenter avec deux indices}
Le liquide contenait au moins une substance non volatile dans les conditions de l’expérience : il n’était pas constitué seulement d’eau. On ne peut pas identifier précisément le dépôt par sa seule présence. Il faut d’autres informations ou des tests adaptés.\par
\end{corrige}\end{blocLivre}
\pcPage{Activité et défi}{Chapitre 04 · Éléments de réponse}
\phantomsection\label{solutions-4}
\begin{pcFiche}{solutions-4}
\begin{pcEncadre}{method}{Classer sans se fier à la couleur}
L’eau distillée est le corps pur dans le modèle annoncé. L’eau salée est un mélange homogène ; l’eau et l’huile forment un mélange hétérogène après repos. L’eau distillée et l’eau salée peuvent toutes deux être transparentes : la transparence ne révèle pas le nombre d’espèces. Il faut utiliser la composition connue ou réaliser des tests adaptés.\par
\end{pcEncadre}
\pcRetour{ch-4-activite}{Revenir à l’activité}
\phantomsection\label{defi-cor-4}
\begin{pcEncadre}{definition}{Préparer un modèle de l’air}
Utiliser 78 jetons pour le diazote, 21 pour le dioxygène et 1 pour les autres gaz. Mélanger les jetons dans la représentation, avec une légende. Le modèle code des proportions arrondies en volume ; il ne montre ni les tailles réelles ni les mouvements. Avec 200 jetons : 156, 42 et 2.\par
\end{pcEncadre}
\pcFigure{figures/06dd6c9bc410a401.pdf}{Construction ou modèle de correction.}
\pcRetour{defi-ex-4}{Revenir au défi}
\end{pcFiche}
\clearpage\corrigesChapitre{5}{Dissolution et miscibilité}
\begin{blocLivre}{cor-5-1}
\begin{corrige}{1}{Le vocabulaire de la solution}
L’eau est le solvant ; le sucre est le soluté. On obtient une solution aqueuse de sucre. Le mot aqueuse précise que le solvant est l’eau.\par
\end{corrige}\end{blocLivre}
\begin{blocLivre}{cor-5-2}
\begin{corrige}{2}{Dissoudre ou fondre ?}
Le glaçon subit une fusion, changement d’état de l’eau. Le sel subit une dissolution dans l’eau, qui joue le rôle de solvant. Une dissolution n’est pas définie par le passage du solide au liquide.\par
\end{corrige}\end{blocLivre}
\begin{blocLivre}{cor-5-3}
\begin{corrige}{3}{Lire les récipients}
Le dernier récipient montre un dépôt. Une solution saturée contient la quantité maximale de soluté qui peut être dissoute dans les conditions étudiées, après un temps et une agitation suffisants.\par
\pcFigure{figures/42794c855063a88c.pdf}{Modèle de correction · les proportions sont indicatives sauf mention contraire.}
\end{corrige}\end{blocLivre}
\begin{blocLivre}{cor-5-4}
\begin{corrige}{4}{Deux liquides au repos}
Elles sont non miscibles dans ces conditions. Le mélange est hétérogène, car on distingue deux phases liquides.\par
\pcFigure{figures/9a252fdceeeb3df9.pdf}{Modèle de correction · les proportions sont indicatives sauf mention contraire.}
\end{corrige}\end{blocLivre}
\begin{blocLivre}{cor-5-5}
\begin{corrige}{5}{La substance invisible}
Le sel reste présent sous forme dissoute, même si ses cristaux ne sont plus visibles. La solution contient de l’eau et du sel : c’est un mélange homogène.\par
\end{corrige}\end{blocLivre}
\begin{blocLivre}{cor-5-6}
\begin{corrige}{6}{Conserver la masse}
La masse de solution est 80 + 12 = 92 g. On additionne des masses exprimées dans la même unité, pas des volumes.\par
\end{corrige}\end{blocLivre}
\begin{blocLivre}{cor-5-7}
\begin{corrige}{7}{Retrouver la masse du soluté}
125 - 110 = 15 g de soluté ont été ajoutés. On suppose qu’il n’y a eu aucune perte ni aucun autre ajout et que le mélange ne contient que ces constituants.\par
\end{corrige}\end{blocLivre}
\begin{blocLivre}{cor-5-8}
\begin{corrige}{8}{Ne changer qu’un facteur}
Non : volume et température changent en plus de l’agitation. Utiliser la même substance, la même masse de solide, le même volume d’eau, la même température et des récipients comparables ; seul le mode d’agitation diffère.\par
\end{corrige}\end{blocLivre}
\begin{blocLivre}{cor-5-9}
\begin{corrige}{9}{Encadrer une limite}
La limite est au moins 12 g et inférieure à 16 g pour 40 g d’eau, à cette température. Tester par exemple 14 g, avec les mêmes conditions, affine l’encadrement.\par
\end{corrige}\end{blocLivre}
\begin{blocLivre}{cor-5-10}
\begin{corrige}{10}{Une masse totale avec un dépôt}
La masse totale du contenu est 85 g. Elle inclut la solution et le dépôt. La masse de la seule phase liquide ne peut pas être obtenue en affirmant que les 25 g se sont tous dissous.\par
\end{corrige}\end{blocLivre}
\begin{blocLivre}{cor-5-11}
\begin{corrige}{11}{Comparer deux volumes de solvant}
100 g d’eau représentent deux fois 50 g : la masse maximale correspondante est 36 g. On suppose que les conditions sont identiques et que la limite dissoute est proportionnelle à la quantité de solvant, pour cette substance.\par
\end{corrige}\end{blocLivre}
\begin{blocLivre}{cor-5-12}
\begin{corrige}{12}{Le dépôt disparaît lentement}
Non. Le dépôt pouvait seulement traduire une dissolution inachevée. Il faut observer après une durée et une agitation suffisantes. On écrit « un dépôt est encore visible à cet instant », puis on conclut avec les observations finales.\par
\end{corrige}\end{blocLivre}
\begin{blocLivre}{cor-5-13}
\begin{corrige}{13}{Choisir le mot précis}
On dit que l’eau et l’huile sont non miscibles dans les conditions de l’essai. La propriété concerne un couple de substances, pas une qualité générale de l’eau. Un même liquide peut se mélanger différemment avec des liquides différents.\par
\end{corrige}\end{blocLivre}
\begin{blocLivre}{cor-5-14}
\begin{corrige}{14}{Calculer la masse dissoute}
22 - 4 = 18 g ont été dissous. La phase liquide contient 50 + 18 = 68 g. Le système entier contenait 72 g avant séparation ; les 4 g restants appartiennent au dépôt.\par
\begin{pcEncadre}{method}{Pour comprendre}
Sur les 22 g ajoutés, 4 g restent hors de la solution : 22 − 4 = 18 g sont dissous. La solution contient alors 50 + 18 = 68 g, et l’ensemble solution + dépôt contient 72 g.\par
\end{pcEncadre}
\pcFigure{figures/8161903911a85a92.pdf}{Modèle de correction · les proportions sont indicatives sauf mention contraire.}
\end{corrige}\end{blocLivre}
\begin{blocLivre}{cor-5-15}
\begin{corrige}{15}{Le volume ne s’additionne pas toujours}
Non : le volume du mélange n’est pas toujours la somme des volumes initiaux. Une pesée avant et après mélange dans un système adapté permet de contrôler la conservation de la masse. On ne déduit pas une perte de matière de la seule baisse de volume.\par
\end{corrige}\end{blocLivre}
\pcPage{Activité et défi}{Chapitre 05 · Éléments de réponse}
\phantomsection\label{solutions-5}
\begin{pcFiche}{solutions-5}
\begin{pcEncadre}{method}{Chercher une limite expérimentale}
La masse maximale dissoute est au moins égale à 15 g et inférieure à 20 g, dans les conditions données. Tester par exemple 17 g ou 18 g avec 50 g d’eau à la même température précise l’encadrement. Changer la température modifierait une autre variable susceptible de changer la limite ; la comparaison ne porterait plus seulement sur la masse ajoutée.\par
\end{pcEncadre}
\pcRetour{ch-5-activite}{Revenir à l’activité}
\phantomsection\label{defi-cor-5}
\begin{pcEncadre}{definition}{Obtenir une solution sans dépôt}
Il faut 40 \(\times\) 1,5 = 60 g d’eau au total. Ajouter au minimum 20 g d’eau aux 40 g initiaux. Agiter suffisamment à température constante et vérifier la disparition du dépôt. La valeur théorique est une limite ; les mesures ont une précision finie.\par
\end{pcEncadre}
\pcRetour{defi-ex-5}{Revenir au défi}
\end{pcFiche}
\clearpage\corrigesChapitre{6}{Séparer les constituants d’un mélange}
\begin{blocLivre}{cor-6-1}
\begin{corrige}{1}{Choisir une technique}
On réalise une filtration. Le liquide qui traverse le filtre est le filtrat ; le sable retenu forme le résidu.\par
\end{corrige}\end{blocLivre}
\begin{blocLivre}{cor-6-2}
\begin{corrige}{2}{Nommer les produits}
Le résidu solide est sur le papier-filtre ; le filtrat est dans le récipient inférieur. Le papier retient les particules assez grosses pour ne pas traverser ses pores.\par
\pcFigure{figures/3bc7ac9680ed49fa.pdf}{Modèle de correction · les proportions sont indicatives sauf mention contraire.}
\end{corrige}\end{blocLivre}
\begin{blocLivre}{cor-6-3}
\begin{corrige}{3}{Laisser reposer}
On utilise une décantation. Le dépôt provient de particules qui se sont rassemblées au fond ; il ne prouve pas une dissolution.\par
\end{corrige}\end{blocLivre}
\begin{blocLivre}{cor-6-4}
\begin{corrige}{4}{Récupérer du sel}
L’évaporation de l’eau laisse le sel non volatil. L’eau est passée dans l’air à l’état gazeux ; elle n’a pas disparu.\par
\end{corrige}\end{blocLivre}
\begin{blocLivre}{cor-6-5}
\begin{corrige}{5}{Deux changements successifs}
Il s’agit d’abord d’une vaporisation, puis d’une liquéfaction. Le réfrigérant favorise le refroidissement et la liquéfaction de la vapeur.\par
\end{corrige}\end{blocLivre}
\begin{blocLivre}{cor-6-6}
\begin{corrige}{6}{Un filtre face au sel}
Oui. Le sel dissous est dispersé dans l’eau et traverse le papier-filtre ordinaire avec elle. La filtration scolaire ne sépare pas les espèces dissoutes de cette manière.\par
\begin{pcEncadre}{method}{Pour comprendre}
Le filtre retient des particules solides assez grandes. Le sel dissous traverse avec l’eau : filtrer n’est pas dessaler.\par
\end{pcEncadre}
\pcFigure{figures/7073ce8fbf40e320.pdf}{Modèle de correction · les proportions sont indicatives sauf mention contraire.}
\end{corrige}\end{blocLivre}
\begin{blocLivre}{cor-6-7}
\begin{corrige}{7}{Comparer deux objectifs}
Pour le sel, une évaporation adaptée convient. Pour récupérer l’eau liquide, choisir une distillation, avec liquéfaction de la vapeur et collecte du distillat.\par
\end{corrige}\end{blocLivre}
\begin{blocLivre}{cor-6-8}
\begin{corrige}{8}{Lire un montage}
Le liquide se vaporise dans le récipient chauffé à gauche. Le réfrigérant retire de l’énergie à la vapeur et permet sa liquéfaction. Le distillat est recueilli à droite.\par
\pcFigure{figures/9a0a5539496a7e8b.pdf}{Modèle de correction · les proportions sont indicatives sauf mention contraire.}
\end{corrige}\end{blocLivre}
\begin{blocLivre}{cor-6-9}
\begin{corrige}{9}{Une séparation en deux étapes}
Filtrer d’abord pour retenir le sable. Le filtrat contient l’eau et le sel dissous ; il faut une autre étape, comme l’évaporation ou la distillation, selon le constituant recherché.\par
\end{corrige}\end{blocLivre}
\begin{blocLivre}{cor-6-10}
\begin{corrige}{10}{Le résultat de l’évaporation}
Il y avait au moins 7,5 g de substance récupérée. Pour assimiler ce résultat à toute la masse initialement dissoute, il faut supposer une récupération complète, sans perte ni transformation du solide et sans contamination.\par
\end{corrige}\end{blocLivre}
\begin{blocLivre}{cor-6-11}
\begin{corrige}{11}{Des volumes de distillat}
60 - 45 = 15 mL n’ont pas été recueillis, en comparant les volumes à conditions équivalentes. Ils peuvent rester dans le montage ou être perdus sous forme de vapeur ou de liquide. Rien ne permet d’affirmer une destruction de l’eau.\par
\end{corrige}\end{blocLivre}
\begin{blocLivre}{cor-6-12}
\begin{corrige}{12}{Filtrer avant ou après ?}
Une décantation peut retirer une grande partie du dépôt avant filtration du liquide surnageant. Des particules fines et des substances dissoutes peuvent néanmoins rester ; le résultat dépend de l’objectif et du filtre.\par
\end{corrige}\end{blocLivre}
\begin{blocLivre}{cor-6-13}
\begin{corrige}{13}{Un mot trompeur}
Demander quelles substances et quelles particules ont réellement été retirées, avec quel dispositif et quels contrôles. On peut écrire « eau filtrée avec tel dispositif » sans affirmer une pureté chimique totale.\par
\end{corrige}\end{blocLivre}
\begin{blocLivre}{cor-6-14}
\begin{corrige}{14}{Éviter une mauvaise conclusion}
Il peut contenir des espèces dissoutes et des particules ou micro-organismes non retenus par le dispositif. L’aspect ne garantit pas l’absence de danger ; un échantillon de laboratoire ne se goûte jamais.\par
\end{corrige}\end{blocLivre}
\begin{blocLivre}{cor-6-15}
\begin{corrige}{15}{Construire un protocole complet}
Filtrer ; récupérer le sable retenu et le sécher selon le protocole. Faire évaporer le filtrat pour obtenir le sel, puis le sécher si nécessaire. Une partie du sel peut rester sur le sable humide, et des pertes peuvent avoir lieu lors des transferts. Les produits ne sont pas déclarés purs sans contrôle.\par
\pcFigure{figures/41ce2a3b8f6f60c1.pdf}{Modèle de correction · les proportions sont indicatives sauf mention contraire.}
\end{corrige}\end{blocLivre}
\pcPage{Activité et défi}{Chapitre 06 · Éléments de réponse}
\phantomsection\label{solutions-6}
\begin{pcFiche}{solutions-6}
\begin{pcEncadre}{method}{Séparer un mélange à trois constituants}
Le sable est principalement retenu dans le résidu ; le filtrat est de l’eau salée. L’évaporation du filtrat laisse le sel non volatil. Pour recueillir l’eau liquide, le professeur peut remplacer l’évaporation simple par une distillation avec refroidissement et liquéfaction de la vapeur. Les pertes et les contaminations doivent être prises en compte.\par
\end{pcEncadre}
\pcRetour{ch-6-activite}{Revenir à l’activité}
\phantomsection\label{defi-cor-6}
\begin{pcEncadre}{definition}{Une mini-chaîne de séparation}
Retirer le liège par prélèvement adapté. Filtrer pour retenir le sable, puis sécher les solides selon les consignes. Le filtrat est une solution salée ; son évaporation permet de récupérer le sel. Pour récupérer l’eau, utiliser une distillation plutôt qu’une évaporation simple. Signaler les pertes et les limites de pureté.\par
\end{pcEncadre}
\pcFigure{figures/41ce2a3b8f6f60c1.pdf}{Construction ou modèle de correction.}
\pcRetour{defi-ex-6}{Revenir au défi}
\end{pcFiche}
\clearpage\corrigesChapitre{7}{Reconnaître une transformation chimique}
\begin{blocLivre}{cor-7-1}
\begin{corrige}{1}{Nommer réactif et produit}
Une espèce consommée est un réactif. Une espèce formée est un produit. Une espèce présente avant et après ne doit pas être déclarée consommée sans information supplémentaire.\par
\end{corrige}\end{blocLivre}
\begin{blocLivre}{cor-7-2}
\begin{corrige}{2}{Trois situations}
La fusion est un changement d’état ; sable et eau forment un mélange ; la formation de rouille correspond à une transformation chimique, car de nouvelles espèces se forment.\par
\end{corrige}\end{blocLivre}
\begin{blocLivre}{cor-7-3}
\begin{corrige}{3}{Un test fourni}
Le test indique la présence de dioxyde de carbone dans les conditions fournies. Il ne prouve pas que le gaz est pur : d’autres gaz peuvent être présents.\par
\end{corrige}\end{blocLivre}
\begin{blocLivre}{cor-7-4}
\begin{corrige}{4}{Des bulles sans réaction}
L’ébullition de l’eau est un exemple : les bulles contiennent de la vapeur d’eau. L’espèce eau est conservée ; son état physique change.\par
\end{corrige}\end{blocLivre}
\begin{blocLivre}{cor-7-5}
\begin{corrige}{5}{Rendre compte d’un test}
J’observe que l’eau de chaux initialement limpide se trouble. D’après la banque de tests, ce résultat indique que le gaz testé contient du dioxyde de carbone.\par
\pcFigure{figures/b594f708d644c622.pdf}{Modèle de correction · les proportions sont indicatives sauf mention contraire.}
\end{corrige}\end{blocLivre}
\begin{blocLivre}{cor-7-6}
\begin{corrige}{6}{Peser un ensemble fermé}
On attend environ 152,0 g, dans la précision de la balance. Il faut peser tout le même ensemble, y compris le récipient et le dispositif qui retient le gaz.\par
\end{corrige}\end{blocLivre}
\begin{blocLivre}{cor-7-7}
\begin{corrige}{7}{Un gaz qui part}
98,4 - 97,8 = 0,6 g ont quitté le système sous forme de gaz. Si une éclaboussure ou une autre perte se produit, la baisse ne peut pas être attribuée au seul gaz.\par
\end{corrige}\end{blocLivre}
\begin{blocLivre}{cor-7-8}
\begin{corrige}{8}{Interpréter un test de l’eau}
Le test fourni met en évidence la présence d’eau. Il ne prouve pas que le liquide est uniquement de l’eau ; une solution contenant d’autres espèces peut aussi donner ce résultat.\par
\pcFigure{figures/ceaa2beaf10e86f3.pdf}{Modèle de correction · les proportions sont indicatives sauf mention contraire.}
\end{corrige}\end{blocLivre}
\begin{blocLivre}{cor-7-9}
\begin{corrige}{9}{Un ballon qui gonfle}
Le ballon peut recueillir le gaz et permettre d’inclure sa masse dans l’ensemble pesé. Éviter de produire un gaz dans un récipient rigide hermétiquement bouché : la pression pourrait augmenter dangereusement. Le professeur choisit le montage.\par
\pcFigure{figures/420b2fff6466bf8c.pdf}{Modèle de correction · les proportions sont indicatives sauf mention contraire.}
\end{corrige}\end{blocLivre}
\begin{blocLivre}{cor-7-10}
\begin{corrige}{10}{Le rôle de l’air}
Le dioxygène de l’air est un réactif de la combustion. Sa masse fait partie de la matière initiale consommée ; ignorer les échanges avec l’air conduit à un bilan incomplet.\par
\end{corrige}\end{blocLivre}
\begin{blocLivre}{cor-7-11}
\begin{corrige}{11}{Présence ou absence ?}
Pas nécessairement : l’espèce peut être trop peu abondante pour la sensibilité du test, ou le protocole peut être inadapté. Vérifier le réactif, les quantités et les conditions, par exemple avec un témoin positif adapté fourni par le professeur.\par
\end{corrige}\end{blocLivre}
\begin{blocLivre}{cor-7-12}
\begin{corrige}{12}{Un changement de couleur}
Non. La dilution ou la dispersion du colorant peut expliquer l’aspect rose sans nouvelle espèce. Il faut rechercher d’autres indices ou des tests qui montrent une transformation de la composition chimique.\par
\end{corrige}\end{blocLivre}
\begin{blocLivre}{cor-7-13}
\begin{corrige}{13}{Une conclusion trop large}
On conclut que le gaz contient du dioxyde de carbone dans les conditions du test. Il faut avoir observé le passage d’une eau de chaux initialement limpide à un trouble, selon le protocole fourni. La pureté n’est pas démontrée.\par
\end{corrige}\end{blocLivre}
\begin{blocLivre}{cor-7-14}
\begin{corrige}{14}{Comparer ouvert et fermé}
La masse totale de toute la matière se conserve dans les deux situations. Dans A, la balance ne pèse plus le gaz parti dans l’air ; dans B, elle pèse encore l’ensemble fermé. La différence concerne les limites du système mesuré.\par
\end{corrige}\end{blocLivre}
\begin{blocLivre}{cor-7-15}
\begin{corrige}{15}{Une preuve construite}
Décrire les substances initiales, recueillir ou analyser le gaz par les tests fournis et comparer les espèces finales. L’eau vaporisée peut redonner de l’eau par liquéfaction ; un gaz contenant du dioxyde de carbone trouble l’eau de chaux. Les bulles indiquent seulement une phase gazeuse, pas son identité ni son mode de formation.\par
\end{corrige}\end{blocLivre}
\pcPage{Activité et défi}{Chapitre 07 · Éléments de réponse}
\phantomsection\label{solutions-7}
\begin{pcFiche}{solutions-7}
\begin{pcEncadre}{method}{Une enquête à partir d’un gaz}
Dans B, le trouble de l’eau de chaux indique la présence de dioxyde de carbone dans le gaz testé. L’apparition de cette espèce, dans la réaction donnée, étaye une transformation chimique. Dans A, l’eau change d’état et redevient liquide au refroidissement : la vaporisation n’est pas une transformation chimique. Des bulles seules ne suffisent donc pas pour trancher.\par
\end{pcEncadre}
\pcRetour{ch-7-activite}{Revenir à l’activité}
\phantomsection\label{defi-cor-7}
\begin{pcEncadre}{definition}{Le mystère des 1,2 gramme}
Les deux résultats peuvent être cohérents : dans le système ouvert, 1,2 g de gaz ont quitté l’ensemble pesé ; dans le système fermé, le gaz est encore inclus. Contrôler l’absence de fuite ou d’éclaboussure, utiliser le même ensemble avant et après, vérifier la stabilité et la précision de la balance. La matière n’a pas été détruite.\par
\end{pcEncadre}
\pcRetour{defi-ex-7}{Revenir au défi}
\end{pcFiche}
\clearpage\corrigesChapitre{8}{La Terre dans le Système solaire}
\begin{blocLivre}{cor-8-1}
\begin{corrige}{1}{Classer les astres}
Le Soleil est une étoile ; la Terre une planète ; la Lune un satellite naturel de la Terre. Ces catégories n’indiquent pas simplement la taille des objets.\par
\end{corrige}\end{blocLivre}
\begin{blocLivre}{cor-8-2}
\begin{corrige}{2}{Retrouver l’ordre}
Mercure, Vénus, Terre et Mars. Jupiter vient ensuite. L’ordre est celui des distances moyennes au Soleil.\par
\end{corrige}\end{blocLivre}
\begin{blocLivre}{cor-8-3}
\begin{corrige}{3}{Identifier la source de lumière}
La Lune diffuse une partie de la lumière solaire reçue. Elle n’est pas une source primaire de lumière visible dans ce modèle ; c’est un objet diffusant.\par
\end{corrige}\end{blocLivre}
\begin{blocLivre}{cor-8-4}
\begin{corrige}{4}{Rotation ou révolution ?}
L’alternance jour-nuit est liée à la rotation terrestre. La révolution de la Terre autour du Soleil dure environ une année.\par
\end{corrige}\end{blocLivre}
\begin{blocLivre}{cor-8-5}
\begin{corrige}{5}{Lire une représentation}
Non. Le dessin indique un ordre et annonce des tailles et distances non à l’échelle. Sans échelle de distance, une mesure sur la page n’est pas convertible en distance réelle.\par
\end{corrige}\end{blocLivre}
\begin{blocLivre}{cor-8-6}
\begin{corrige}{6}{Une unité de distance}
150 millions de km = 150 000 000 km. C’est une distance, pas une durée. Elle sert ici de repère approximatif pour la distance moyenne Terre-Soleil.\par
\end{corrige}\end{blocLivre}
\begin{blocLivre}{cor-8-7}
\begin{corrige}{7}{Placer la Terre}
150 \(\div\) 50 = 3 cm sur la feuille. La valeur 3 n’est pas exprimée en millions de km : elle décrit la maquette.\par
\end{corrige}\end{blocLivre}
\begin{blocLivre}{cor-8-8}
\begin{corrige}{8}{Placer Mars}
228 \(\div\) 50 = 4,56 cm. Oui : 4,56 cm dépassent 4 cm, qui représentent 200 millions de km.\par
\end{corrige}\end{blocLivre}
\begin{blocLivre}{cor-8-9}
\begin{corrige}{9}{Jour et nuit}
À un même instant, les lieux peuvent être sur les moitiés éclairée et non éclairée de la Terre. Le Soleil reste lumineux ; c’est l’orientation du lieu terrestre par rapport à lui qui change avec la rotation.\par
\pcFigure{figures/756e7a4628904c01.pdf}{Modèle de correction · les proportions sont indicatives sauf mention contraire.}
\end{corrige}\end{blocLivre}
\begin{blocLivre}{cor-8-10}
\begin{corrige}{10}{Comparer des distances}
780 \(\div\) 150 = 5,2 fois. Non : la séparation entre les deux planètes dépend de leurs positions sur leurs orbites. Les distances données sont chacune mesurées depuis le Soleil.\par
\end{corrige}\end{blocLivre}
\begin{blocLivre}{cor-8-11}
\begin{corrige}{11}{Lire un repère en ua}
3 \(\times\) 150 = 450 millions de km environ. L’unité astronomique permet d’exprimer plus simplement certaines distances dans le Système solaire.\par
\end{corrige}\end{blocLivre}
\begin{blocLivre}{cor-8-12}
\begin{corrige}{12}{Le modèle trop petit}
4 500 \(\div\) 50 = 90 cm. Une feuille de 21 cm de large ne suffit pas ; il faut une bande plus longue ou une autre échelle.\par
\end{corrige}\end{blocLivre}
\begin{blocLivre}{cor-8-13}
\begin{corrige}{13}{Une échelle incohérente}
Il n’applique pas le même rapport aux distances. Avec 1 cm pour 50 millions de km, Mars doit être à 4,56 cm du Soleil. Une même longueur sur la maquette doit représenter la même distance réelle partout.\par
\begin{pcEncadre}{method}{Pour comprendre}
Calcule séparément chaque longueur avec le même diviseur 50. Les longueurs calculées sont celles de la maquette ; le schéma de correction sert à comparer les positions.\par
\end{pcEncadre}
\pcFigure{figures/025dab5ad291e881.pdf}{Modèle de correction · les proportions sont indicatives sauf mention contraire.}
\end{corrige}\end{blocLivre}
\begin{blocLivre}{cor-8-14}
\begin{corrige}{14}{Relier plusieurs mouvements}
Elle accompagne la rotation terrestre et la révolution de la Terre autour du Soleil. Son repos est relatif au jardin ; un autre observateur peut décrire d’autres mouvements. La description dépend du repère choisi.\par
\end{corrige}\end{blocLivre}
\begin{blocLivre}{cor-8-15}
\begin{corrige}{15}{Évaluer une source}
Vérifier les catégories, l’auteur, la date, les unités et les sources des données ; comparer avec une source reconnue. Le CNES ou un observatoire astronomique peut fournir des repères fiables. Une image attractive ne garantit pas l’exactitude scientifique.\par
\end{corrige}\end{blocLivre}
\pcPage{Activité et défi}{Chapitre 08 · Éléments de réponse}
\phantomsection\label{solutions-8}
\begin{pcFiche}{solutions-8}
\begin{pcEncadre}{method}{Une maquette sur une bande de papier}
Les positions sont 1,16 cm ; 2,16 cm ; 3,00 cm ; 4,56 cm. Jupiter serait à 15,6 cm et tient sur une bande de 20 cm. Les diamètres réels, à cette échelle des distances, seraient trop petits pour être lisibles : les disques sont agrandis pour repérer les planètes. Le modèle ne prétend donc pas respecter simultanément l’échelle des tailles.\par
\end{pcEncadre}
\pcRetour{ch-8-activite}{Revenir à l’activité}
\phantomsection\label{defi-cor-8}
\begin{pcEncadre}{definition}{Le Système solaire dans la cour}
1 m représente 100 millions de km. La Terre est à 1,50 m, Mars à 2,28 m et Jupiter à 7,80 m du Soleil. Neptune est à 45 m. Les diamètres des planètes à cette échelle seraient minuscules ; employer des repères agrandis clairement annoncés, sans prétendre respecter l’échelle des tailles.\par
\end{pcEncadre}
\pcFigure{figures/e7785657822d78af.pdf}{Construction ou modèle de correction.}
\pcRetour{defi-ex-8}{Revenir au défi}
\end{pcFiche}
\clearpage\corrigesChapitre{9}{Décrire un mouvement}
\begin{blocLivre}{cor-9-1}
\begin{corrige}{1}{Préciser le repère}
Elle est immobile par rapport à son siège si elle reste assise. Elle est en mouvement par rapport à l’arrêt quand le bus roule. Il faut annoncer le repère pour éviter une contradiction apparente.\par
\end{corrige}\end{blocLivre}
\begin{blocLivre}{cor-9-2}
\begin{corrige}{2}{Nommer un chemin}
Les trajectoires sont respectivement rectiligne et circulaire. Ces noms ne précisent pas si la valeur de la vitesse est constante.\par
\end{corrige}\end{blocLivre}
\begin{blocLivre}{cor-9-3}
\begin{corrige}{3}{Lire le sens}
Oui, la forme peut être rectiligne pour les deux. Il faut préciser le sens du mouvement, par exemple de A vers B et de B vers A.\par
\pcFigure{figures/07bfff6c23b887fa.pdf}{Modèle de correction · les proportions sont indicatives sauf mention contraire.}
\end{corrige}\end{blocLivre}
\begin{blocLivre}{cor-9-4}
\begin{corrige}{4}{Un passager observé}
On peut choisir le train ou le quai. Le voyageur assis est au repos par rapport au train et en mouvement par rapport au quai lorsque le train se déplace.\par
\pcFigure{figures/5f58c69c5265b1b2.pdf}{Modèle de correction · les proportions sont indicatives sauf mention contraire.}
\end{corrige}\end{blocLivre}
\begin{blocLivre}{cor-9-5}
\begin{corrige}{5}{Des mots différents}
Rectiligne décrit la forme de la trajectoire. Uniforme indique une valeur de vitesse constante. On peut aller en ligne droite en accélérant ou en ralentissant.\par
\end{corrige}\end{blocLivre}
\begin{blocLivre}{cor-9-6}
\begin{corrige}{6}{Des positions régulières}
La même distance est parcourue pendant chaque durée égale : la valeur de la vitesse est constante sur ces intervalles. L’égalité des intervalles de temps est essentielle.\par
\begin{pcEncadre}{method}{Pour comprendre}
Compare des distances sur des durées égales. Sans connaître les dates des images, les seuls espacements ne suffisent pas à conclure.\par
\end{pcEncadre}
\pcFigure{figures/2f5fc91e576092f8.pdf}{Modèle de correction · les proportions sont indicatives sauf mention contraire.}
\end{corrige}\end{blocLivre}
\begin{blocLivre}{cor-9-7}
\begin{corrige}{7}{Trois lignes à comparer}
A : valeur constante ; B : valeur croissante ; C : valeur décroissante. On utilise les espacements successifs, avec les mêmes durées entre positions dans chaque ligne.\par
\pcFigure{figures/b9d2213e035c1cd3.pdf}{Modèle de correction · les proportions sont indicatives sauf mention contraire.}
\end{corrige}\end{blocLivre}
\begin{blocLivre}{cor-9-8}
\begin{corrige}{8}{Un tour de manège}
La trajectoire du siège est circulaire dans le modèle. La direction du déplacement change au cours du tour, même si la valeur de la vitesse est constante.\par
\end{corrige}\end{blocLivre}
\begin{blocLivre}{cor-9-9}
\begin{corrige}{9}{Le bon point à suivre}
Changer de point peut modifier la trajectoire relevée et fausser la comparaison. Le centre du ballon constitue un choix simple pour décrire son déplacement global, en laissant de côté sa rotation.\par
\end{corrige}\end{blocLivre}
\begin{blocLivre}{cor-9-10}
\begin{corrige}{10}{Des durées inégales}
Non. Les deux intervalles correspondent à 1 m par seconde en moyenne. La plus grande distance du second intervalle est accompagnée d’une durée deux fois plus longue ; elle ne prouve pas une augmentation de vitesse.\par
\end{corrige}\end{blocLivre}
\begin{blocLivre}{cor-9-11}
\begin{corrige}{11}{Une description complète}
Par rapport au repère fixe choisi, le point étudié suit une trajectoire rectiligne vers la droite et la valeur de sa vitesse augmente sur les intervalles observés. Les trois informations sont repère, trajectoire et évolution de vitesse.\par
\pcFigure{figures/0ee4546ede86e6f7.pdf}{Modèle de correction · les proportions sont indicatives sauf mention contraire.}
\end{corrige}\end{blocLivre}
\begin{blocLivre}{cor-9-12}
\begin{corrige}{12}{L’appareil se déplace}
L’appareil change de position par rapport à l’élève filmé. Le mouvement apparent dépend donc du repère de l’appareil. Par rapport au sol ou à un repère fixé au sol, l’élève filmé est au repos dans la situation donnée.\par
\end{corrige}\end{blocLivre}
\begin{blocLivre}{cor-9-13}
\begin{corrige}{13}{Une photographie unique}
Non. Il faut des positions à plusieurs dates connues ou une mesure de vitesse fournie par un instrument. Une photo seule indique une position ou un aspect à un instant, pas l’évolution temporelle.\par
\end{corrige}\end{blocLivre}
\begin{blocLivre}{cor-9-14}
\begin{corrige}{14}{Une hypothèse à tester}
Pas encore. Il faut vérifier que les durées entre images sont égales. Des espacements égaux parcourus en durées différentes ne correspondent pas à une même vitesse moyenne sur les intervalles.\par
\end{corrige}\end{blocLivre}
\begin{blocLivre}{cor-9-15}
\begin{corrige}{15}{Comparer deux observateurs}
Les deux descriptions peuvent être correctes. « Nous sommes au repos par rapport au train » et « Les voyageurs sont en mouvement par rapport au quai » précisent chacune le repère. Une description scientifique n’omet pas ce choix.\par
\end{corrige}\end{blocLivre}
\pcPage{Activité et défi}{Chapitre 09 · Éléments de réponse}
\phantomsection\label{solutions-9}
\begin{pcFiche}{solutions-9}
\begin{pcEncadre}{method}{Photographier un petit mobile}
A montre des espacements constants : mouvement rectiligne de valeur de vitesse constante. En B, les espacements augmentent ; en C, ils diminuent. Un appareil fixe conserve le même repère ; s’il se déplace, les positions apparentes changent aussi à cause de l’observateur. La comparaison suppose des intervalles de temps égaux et une échelle de longueur commune.\par
\end{pcEncadre}
\pcRetour{ch-9-activite}{Revenir à l’activité}
\phantomsection\label{defi-cor-9}
\begin{pcEncadre}{definition}{Le reportage de deux témoins}
Pour la balle tenue, son centre est immobile par rapport au train et se déplace en ligne droite par rapport au quai, dans le modèle donné. Après lâcher, sa position change aussi par rapport au train ; il faut alors relever de nouvelles positions pour décrire les trajectoires dans chaque repère. Les commentaires doivent annoncer le repère au lieu de chercher un mouvement unique absolu.\par
\end{pcEncadre}
\pcRetour{defi-ex-9}{Revenir au défi}
\end{pcFiche}
\clearpage\corrigesChapitre{10}{Mesurer et calculer une vitesse}
\begin{blocLivre}{cor-10-1}
\begin{corrige}{1}{Choisir une unité}
On obtient respectivement des m/s et des km/h. Le résultat doit inclure à la fois une longueur et une durée dans son unité.\par
\end{corrige}\end{blocLivre}
\begin{blocLivre}{cor-10-2}
\begin{corrige}{2}{Une première vitesse}
20 \(\div\) 5 = 4 m/s. Si le mouvement est uniforme, le mobile parcourt 4 m chaque seconde.\par
\end{corrige}\end{blocLivre}
\begin{blocLivre}{cor-10-3}
\begin{corrige}{3}{Convertir une durée}
2 min = 120 s. 1 min 30 s = 60 + 30 = 90 s. Le système des durées utilise 60 secondes par minute.\par
\end{corrige}\end{blocLivre}
\begin{blocLivre}{cor-10-4}
\begin{corrige}{4}{Lire une droite}
Le graphique donne 12 m en 4 s. La vitesse vaut 12 \(\div\) 4 = 3 m/s. La droite passant par l’origine indique une distance proportionnelle au temps dans le modèle.\par
\pcFigure{figures/c12d6c191831498b.pdf}{Modèle de correction · les proportions sont indicatives sauf mention contraire.}
\end{corrige}\end{blocLivre}
\begin{blocLivre}{cor-10-5}
\begin{corrige}{5}{Le bon trajet}
Il faut utiliser les 120 m réellement parcourus sur le chemin. Les 80 m correspondent à la séparation en ligne droite, qui n’est pas la distance du trajet.\par
\pcFigure{figures/38c5864ac59a2ed8.pdf}{Modèle de correction · les proportions sont indicatives sauf mention contraire.}
\end{corrige}\end{blocLivre}
\begin{blocLivre}{cor-10-6}
\begin{corrige}{6}{Un trajet à vélo}
12 \(\div\) 0,5 = 24 km/h. La durée 0,5 h correspond à 30 min, et non à 50 min.\par
\end{corrige}\end{blocLivre}
\begin{blocLivre}{cor-10-7}
\begin{corrige}{7}{Un petit robot}
150 cm = 1,50 m. La vitesse moyenne vaut 1,50 \(\div\) 5 = 0,30 m/s. Sans conversion, on obtiendrait 30 cm/s, pas 30 m/s.\par
\end{corrige}\end{blocLivre}
\begin{blocLivre}{cor-10-8}
\begin{corrige}{8}{Changer d’unité de vitesse}
5 \(\times\) 3,6 = 18 km/h. 36 \(\div\) 3,6 = 10 m/s. On contrôle le sens de conversion avec le repère 1 m/s = 3,6 km/h.\par
\end{corrige}\end{blocLivre}
\begin{blocLivre}{cor-10-9}
\begin{corrige}{9}{Prévoir une distance}
La distance vaut 2 \(\times\) 15 = 30 m. On suppose que la valeur de la vitesse reste égale à 2 m/s pendant les 15 s.\par
\end{corrige}\end{blocLivre}
\begin{blocLivre}{cor-10-10}
\begin{corrige}{10}{Prévoir une durée}
La durée vaut 100 \(\div\) 4 = 25 s. On utilise durée = distance \(\div\) vitesse, avec des unités cohérentes.\par
\end{corrige}\end{blocLivre}
\begin{blocLivre}{cor-10-11}
\begin{corrige}{11}{Inclure un arrêt}
La moyenne de la sortie entière vaut 3 km/h. La durée annoncée inclut l’arrêt ; utiliser 45 min calculerait une moyenne sur le seul temps de déplacement, qui répond à une autre question.\par
\begin{pcEncadre}{method}{Pour comprendre}
La durée demandée est celle de la sortie complète : les 15 min d’arrêt font partie de l’heure. La vitesse moyenne vaut donc 3 km \(\div\) 1 h = 3 km/h.\par
\end{pcEncadre}
\pcFigure{figures/14814a5dcde7089b.pdf}{Modèle de correction · les proportions sont indicatives sauf mention contraire.}
\end{corrige}\end{blocLivre}
\begin{blocLivre}{cor-10-12}
\begin{corrige}{12}{Deux étapes différentes}
La distance totale est 6 km et la durée totale 60 min = 1 h : moyenne de 6 km/h. Chaque étape a aussi une moyenne de 6 km/h ; les distances sont proportionnelles aux durées.\par
\end{corrige}\end{blocLivre}
\begin{blocLivre}{cor-10-13}
\begin{corrige}{13}{Une moyenne trompeuse}
Distances : 10 km et 8 km, soit 18 km. Durée : 3 h. La moyenne vaut 18 \(\div\) 3 = 6 km/h, et non 7 km/h. Les deux vitesses n’ont pas été maintenues pendant des durées égales.\par
\begin{pcEncadre}{method}{Pour comprendre}
Calcule chaque distance avant de faire le total : 10 \(\times\) 1 = 10 km et 4 \(\times\) 2 = 8 km. Puis 18 km \(\div\) 3 h = 6 km/h. La moyenne simple des deux vitesses ne convient pas pour des durées différentes.\par
\end{pcEncadre}
\pcFigure{figures/1c7cb8757e8d04d5.pdf}{Modèle de correction · les proportions sont indicatives sauf mention contraire.}
\end{corrige}\end{blocLivre}
\begin{blocLivre}{cor-10-14}
\begin{corrige}{14}{Comparer des performances}
A : 5 m/s. B : 6 m/s. B est plus rapide en moyenne. Les distances diffèrent : une durée plus courte ne suffit pas pour classer les vitesses.\par
\end{corrige}\end{blocLivre}
\begin{blocLivre}{cor-10-15}
\begin{corrige}{15}{Contrôler un résultat}
Le résultat correct est 100 \(\div\) 200 = 0,5 m/s. Une confusion entre multiplication, division ou unités a probablement eu lieu. La valeur 500 m/s est incohérente avec une marche ; l’ordre de grandeur aide à détecter l’erreur.\par
\end{corrige}\end{blocLivre}
\pcPage{Activité et défi}{Chapitre 10 · Éléments de réponse}
\phantomsection\label{solutions-10}
\begin{pcFiche}{solutions-10}
\begin{pcEncadre}{method}{Mesurer une vitesse de marche}
Les vitesses moyennes sont 10 \(\div\) 8 = 1,25 m/s ; 10 \(\div\) 10 = 1 m/s ; 10 \(\div\) 12 \(\approx\) 0,83 m/s. Le premier essai est le plus rapide en moyenne. Le rythme de marche, le déclenchement du chronomètre, le repérage des lignes et la précision de la distance contribuent aux variations.\par
\end{pcEncadre}
\pcRetour{ch-10-activite}{Revenir à l’activité}
\phantomsection\label{defi-cor-10}
\begin{pcEncadre}{definition}{Le trajet avec une pause}
Distance totale : 1 000 m. Durée totale : 8 + 2 + 5 = 15 min, soit 0,25 h. Moyenne : environ 66,7 m/min, ou 4 km/h. La moyenne arithmétique des deux vitesses ne tient pas compte des durées différentes ni de la pause. Utiliser les totaux du même trajet.\par
\end{pcEncadre}
\pcRetour{defi-ex-10}{Revenir au défi}
\end{pcFiche}
\clearpage\corrigesChapitre{11}{Décrire des actions et des interactions}
\begin{blocLivre}{cor-11-1}
\begin{corrige}{1}{Un effet visible}
On observe une déformation de l’élastique. La main et l’élastique sont en interaction de contact.\par
\end{corrige}\end{blocLivre}
\begin{blocLivre}{cor-11-2}
\begin{corrige}{2}{Classer les interactions}
Pied-ballon : contact. Terre-ballon : interaction à distance. Aimant-trombone sans contact : interaction à distance.\par
\end{corrige}\end{blocLivre}
\begin{blocLivre}{cor-11-3}
\begin{corrige}{3}{Le livre immobile}
Oui. Le livre interagit aussi avec la table qui le soutient. Son immobilité n’efface pas ces interactions.\par
\pcFigure{figures/852f07059882b685.pdf}{Modèle de correction · les proportions sont indicatives sauf mention contraire.}
\end{corrige}\end{blocLivre}
\begin{blocLivre}{cor-11-4}
\begin{corrige}{4}{L’action de l’air}
C’est une interaction de contact entre l’air et la voile. L’air est une matière ; son invisibilité ne l’empêche pas d’exercer une action.\par
\end{corrige}\end{blocLivre}
\begin{blocLivre}{cor-11-5}
\begin{corrige}{5}{Lire un diagramme}
Le système est l’objet central. Le lien pointillé avec la Terre représente une interaction à distance, conformément à la légende. Les liens pleins indiquent des contacts.\par
\pcFigure{figures/b0f4a40eb7b49549.pdf}{Modèle de correction · les proportions sont indicatives sauf mention contraire.}
\end{corrige}\end{blocLivre}
\begin{blocLivre}{cor-11-6}
\begin{corrige}{6}{Après le tir}
Il n’y a plus de contact entre le pied et le ballon. L’interaction avec la Terre reste présente. D’autres interactions, comme celle avec l’air, peuvent être retenues selon le modèle.\par
\end{corrige}\end{blocLivre}
\begin{blocLivre}{cor-11-7}
\begin{corrige}{7}{Une bille déviée}
Une interaction magnétique à distance est mise en évidence dans la situation donnée. Son effet observé est une modification de trajectoire.\par
\end{corrige}\end{blocLivre}
\begin{blocLivre}{cor-11-8}
\begin{corrige}{8}{Préciser les objets}
Il manque l’objet qui reçoit l’action. Une description précise identifie les deux objets en interaction ; le seul nom de la table ne suffit pas.\par
\end{corrige}\end{blocLivre}
\begin{blocLivre}{cor-11-9}
\begin{corrige}{9}{Changer le système}
Non : la table devient le système central. On liste les objets qui interagissent avec elle. Le système central dépend de la question étudiée, et non d’un choix définitif pour toute la scène.\par
\end{corrige}\end{blocLivre}
\begin{blocLivre}{cor-11-10}
\begin{corrige}{10}{Main retirée}
Le lien avec la main disparaît. Le lien avec la table reste puisque le livre est toujours en contact avec elle. Le lien avec la Terre reste également.\par
\begin{pcEncadre}{method}{Pour comprendre}
Conserve le même système, le livre. Supprime seulement le lien avec la main retirée ; la table et la Terre restent en interaction avec lui.\par
\end{pcEncadre}
\pcFigure{figures/9e9f7fd040937727.pdf}{Modèle de correction · les proportions sont indicatives sauf mention contraire.}
\end{corrige}\end{blocLivre}
\begin{blocLivre}{cor-11-11}
\begin{corrige}{11}{Ne pas tout déduire}
Non. Le diagramme inventorie les interactions retenues dans la situation. Une valeur de vitesse exige une mesure de distance et de durée ou une autre information adaptée.\par
\end{corrige}\end{blocLivre}
\begin{blocLivre}{cor-11-12}
\begin{corrige}{12}{Une chute dans l’air}
La Terre interagit à distance avec la balle ; l’air interagit par contact. Ici, on décrit une chute dans l’air ; le modèle simplifié qui néglige l’air doit être annoncé explicitement.\par
\pcFigure{figures/7c1b3bfa354b787d.pdf}{Modèle de correction · les proportions sont indicatives sauf mention contraire.}
\end{corrige}\end{blocLivre}
\begin{blocLivre}{cor-11-13}
\begin{corrige}{13}{Choisir un modèle simple}
Elle indique que le diagramme représente un modèle choisi, pas une liste exhaustive de toutes les influences possibles. On garde l’interaction de contact avec la table et l’interaction à distance avec la Terre.\par
\pcFigure{figures/852f07059882b685.pdf}{Modèle de correction · les proportions sont indicatives sauf mention contraire.}
\end{corrige}\end{blocLivre}
\begin{blocLivre}{cor-11-14}
\begin{corrige}{14}{Une généralisation à discuter}
Les matériaux métalliques n’ont pas tous le même comportement magnétique. Comparer plusieurs échantillons identifiés avec le même aimant, à la même distance et sans contact initial, en notant les résultats. Ne pas étendre la propriété d’un trombone en acier à tous les métaux.\par
\end{corrige}\end{blocLivre}
\begin{blocLivre}{cor-11-15}
\begin{corrige}{15}{Une action sans mouvement visible}
Non. La main ou le corps peut exercer une action même sans déplacement visible du meuble. Le contact observé et l’effort exercé font partie des informations de la situation. Il faut distinguer existence d’une interaction et effet visible sur le mouvement.\par
\end{corrige}\end{blocLivre}
\pcPage{Activité et défi}{Chapitre 11 · Éléments de réponse}
\phantomsection\label{solutions-11}
\begin{pcFiche}{solutions-11}
\begin{pcEncadre}{method}{Enquêter autour d’un livre}
Retirer le lien avec la main dans A ; garder la table et la Terre. L’attraction terrestre existe que le livre soit au repos ou en mouvement. Le diagramme liste les interactions retenues, mais ne donne ni leur valeur, ni leur direction détaillée, ni la vitesse de l’objet. Ces descriptions seront approfondies aux niveaux suivants.\par
\end{pcEncadre}
\pcRetour{ch-11-activite}{Revenir à l’activité}
\phantomsection\label{defi-cor-11}
\begin{pcEncadre}{definition}{Trois scènes, trois diagrammes}
Ballon tenu : contact avec la main et interaction à distance avec la Terre. Ballon lâché : Terre à distance et air par contact si on le prend en compte. Ballon posé : contact avec le sol et interaction avec la Terre. L’air entoure aussi les scènes tenues et posées ; il peut être négligé dans un modèle simple annoncé. L’interaction terrestre reste présente dans les trois situations.\par
\end{pcEncadre}
\pcFigure{figures/d05a150d20f841c3.pdf}{Construction ou modèle de correction.}
\pcRetour{defi-ex-11}{Revenir au défi}
\end{pcFiche}
\clearpage\corrigesChapitre{12}{Sources, formes et conversions d’énergie}
\begin{blocLivre}{cor-12-1}
\begin{corrige}{1}{Ressource ou appareil ?}
Vent et Soleil sont des ressources ; éolienne et panneau photovoltaïque sont des convertisseurs. Une même chaîne doit distinguer ce qui fournit l’énergie et l’objet technique qui l’exploite.\par
\end{corrige}\end{blocLivre}
\begin{blocLivre}{cor-12-2}
\begin{corrige}{2}{Une réserve dans la pile}
La pile possède une réserve d’énergie chimique. Cette réserve diminue au cours du fonctionnement ; de l’énergie est transférée au circuit et à l’environnement.\par
\end{corrige}\end{blocLivre}
\begin{blocLivre}{cor-12-3}
\begin{corrige}{3}{Nommer une conversion}
Le moteur sert à produire un mouvement ; la lampe sert à éclairer. Les deux dispositifs peuvent également transférer de l’énergie thermiquement à l’environnement.\par
\end{corrige}\end{blocLivre}
\begin{blocLivre}{cor-12-4}
\begin{corrige}{4}{L’unité d’énergie}
Le joule, symbole J, est une unité d’énergie. Il ne s’agit pas d’une unité de masse, de durée ou de vitesse.\par
\end{corrige}\end{blocLivre}
\begin{blocLivre}{cor-12-5}
\begin{corrige}{5}{Lire une chaîne}
Le convertisseur est la lampe. Les sorties représentent la lumière et un transfert thermique. La pile fournit l’énergie par transfert électrique.\par
\pcFigure{figures/6d0e39ac014c4fa7.pdf}{Modèle de correction · les proportions sont indicatives sauf mention contraire.}
\end{corrige}\end{blocLivre}
\begin{blocLivre}{cor-12-6}
\begin{corrige}{6}{Compléter un bilan}
80 - 25 = 55 J sont transférés autrement, par exemple thermiquement selon l’appareil. Cette énergie n’est pas détruite ; elle doit apparaître dans le bilan.\par
\begin{pcEncadre}{method}{Pour comprendre}
À durée d’étude identique et sans stockage, énergie reçue = totalité des énergies transférées. L’énergie non utile vaut 80 − 25 = 55 J.\par
\end{pcEncadre}
\pcFigure{figures/58e5095fca5761c3.pdf}{Modèle de correction · les proportions sont indicatives sauf mention contraire.}
\end{corrige}\end{blocLivre}
\begin{blocLivre}{cor-12-7}
\begin{corrige}{7}{Changer d’objectif}
Non. Pour éclairer, le transfert thermique peut être un effet secondaire ; pour chauffer, il peut constituer l’effet recherché. Le mot utile dépend de la fonction choisie.\par
\end{corrige}\end{blocLivre}
\begin{blocLivre}{cor-12-8}
\begin{corrige}{8}{Une chaîne solaire}
Le Soleil fournit le rayonnement reçu ; le panneau photovoltaïque convertit une partie de l’énergie reçue en transfert électrique ; le moteur produit notamment un mouvement. D’autres transferts doivent aussi être pris en compte.\par
\end{corrige}\end{blocLivre}
\begin{blocLivre}{cor-12-9}
\begin{corrige}{9}{Un bilan à contrôler}
Oui : 30 + 70 = 100 J. Les valeurs doivent concerner la même durée d’étude et la même unité. Un bilan pris à des durées différentes serait incohérent.\par
\pcFigure{figures/21eaae42816f98e4.pdf}{Modèle de correction · les proportions sont indicatives sauf mention contraire.}
\end{corrige}\end{blocLivre}
\begin{blocLivre}{cor-12-10}
\begin{corrige}{10}{Une réserve qui augmente}
La réserve augmente de 50 - 35 = 15 J. L’écart correspond au stockage dans le système ; ce n’est pas une disparition d’énergie.\par
\begin{pcEncadre}{method}{Pour comprendre}
Le bilan comprend aussi le stockage : entrée = sortie + augmentation de la réserve. On trouve 50 − 35 = 15 J stockés.\par
\end{pcEncadre}
\pcFigure{figures/5b40d7c84961bce3.pdf}{Modèle de correction · les proportions sont indicatives sauf mention contraire.}
\end{corrige}\end{blocLivre}
\begin{blocLivre}{cor-12-11}
\begin{corrige}{11}{Des mots précis}
Le vent met l’éolienne en mouvement ; son dispositif convertit une partie de l’énergie reçue et la transfère notamment sous forme électrique. L’énergie est transférée et convertie, pas créée à partir de rien.\par
\end{corrige}\end{blocLivre}
\begin{blocLivre}{cor-12-12}
\begin{corrige}{12}{Comparer deux appareils}
B fournit davantage d’énergie utile. Les autres transferts valent 80 J pour A et 65 J pour B. La comparaison suppose le même objectif et des valeurs mesurées sur des conditions comparables.\par
\end{corrige}\end{blocLivre}
\begin{blocLivre}{cor-12-13}
\begin{corrige}{13}{Une chaîne incomplète}
Il manque le convertisseur moteur, la nature du transfert entrant et les autres transferts sortants. Ajouter pile \(\rightarrow\) transfert électrique \(\rightarrow\) moteur, puis les sorties mouvement et transfert thermique, et éventuellement son selon le modèle.\par
\pcFigure{figures/4884783ed01a165e.pdf}{Modèle de correction · les proportions sont indicatives sauf mention contraire.}
\end{corrige}\end{blocLivre}
\begin{blocLivre}{cor-12-14}
\begin{corrige}{14}{Ne pas additionner des grandeurs différentes}
Il additionne une énergie et une durée, grandeurs différentes. On n’additionne que des grandeurs compatibles exprimées dans la même unité ; l’unité ne se choisit pas après le calcul.\par
\end{corrige}\end{blocLivre}
\begin{blocLivre}{cor-12-15}
\begin{corrige}{15}{Un effet secondaire observé}
Ajouter des transferts thermiques et sonores vers l’environnement. Ils peuvent être moins utiles pour faire tourner le moteur, mais l’énergie existe toujours dans le bilan. Le schéma initial était seulement incomplet.\par
\end{corrige}\end{blocLivre}
\pcPage{Activité et défi}{Chapitre 12 · Éléments de réponse}
\phantomsection\label{solutions-12}
\begin{pcFiche}{solutions-12}
\begin{pcEncadre}{method}{Construire une chaîne énergétique}
Le convertisseur et l’effet utile changent : la lampe émet de la lumière, le moteur met un objet en mouvement. Les deux appareils peuvent aussi transférer de l’énergie thermiquement à leur environnement. La réserve d’énergie chimique de la pile diminue. Le schéma doit montrer cette origine et éviter de faire disparaître les transferts secondaires.\par
\end{pcEncadre}
\pcRetour{ch-12-activite}{Revenir à l’activité}
\phantomsection\label{defi-cor-12}
\begin{pcEncadre}{definition}{L’éclairage à améliorer}
A transfère 160 J autrement et B 130 J. B fournit plus d’énergie lumineuse utile pour 200 J reçus. Un choix réel demande aussi la répartition de la lumière, le confort visuel, la durée de vie, le prix et les conditions d’utilisation. Les données de l’essai ne suffisent pas à classer tous les produits.\par
\end{pcEncadre}
\pcRetour{defi-ex-12}{Revenir au défi}
\end{pcFiche}
\clearpage\corrigesChapitre{13}{Réaliser et schématiser un circuit}
\begin{blocLivre}{cor-13-1}
\begin{corrige}{1}{Nommer les composants}
Une pile peut être un générateur. Une lampe et un moteur sont des récepteurs. Leur compatibilité avec le générateur doit être vérifiée.\par
\end{corrige}\end{blocLivre}
\begin{blocLivre}{cor-13-2}
\begin{corrige}{2}{Deux bornes}
Un dipôle possède deux bornes de connexion. Le symbole d’une lampe comporte deux connexions vers les fils du circuit.\par
\end{corrige}\end{blocLivre}
\begin{blocLivre}{cor-13-3}
\begin{corrige}{3}{Ouvrir le circuit}
L’interrupteur ouvert interrompt la boucle ; la lampe ne fonctionne pas dans ce montage. La présence de la pile seule ne suffit pas : la boucle doit être fermée et les composants adaptés.\par
\pcFigure{figures/a7491f435a5e5e7a.pdf}{Modèle de correction · les proportions sont indicatives sauf mention contraire.}
\end{corrige}\end{blocLivre}
\begin{blocLivre}{cor-13-4}
\begin{corrige}{4}{Lire les symboles}
La pile est représentée par deux traits parallèles de longueurs différentes ; la lampe par un cercle avec une croix ; le moteur par un cercle avec M. Le trait long est la borne positive.\par
\pcFigure{figures/d31c9e690caff648.pdf}{Modèle de correction · les proportions sont indicatives sauf mention contraire.}
\end{corrige}\end{blocLivre}
\begin{blocLivre}{cor-13-5}
\begin{corrige}{5}{Le sens conventionnel}
On le représente du pôle positif vers le pôle négatif dans le circuit extérieur fermé. Il s’agit d’une convention de description.\par
\pcFigure{figures/c6a7cbcb169788d4.pdf}{Modèle de correction · les proportions sont indicatives sauf mention contraire.}
\end{corrige}\end{blocLivre}
\begin{blocLivre}{cor-13-6}
\begin{corrige}{6}{Lire une coupure}
L’interrupteur est ouvert et coupe le trajet conducteur. Le fermer permet de compléter la boucle, si les autres connexions et les composants sont corrects.\par
\begin{pcEncadre}{method}{Pour comprendre}
Pars d’une borne de la pile et suis tout le trajet. Une coupure suffit à empêcher le courant dans cette boucle unique.\par
\end{pcEncadre}
\pcFigure{figures/8a70d74bfce1a098.pdf}{Modèle de correction · les proportions sont indicatives sauf mention contraire.}
\end{corrige}\end{blocLivre}
\begin{blocLivre}{cor-13-7}
\begin{corrige}{7}{Un dessin ou un schéma ?}
Le schéma code le circuit par des symboles. Il doit préserver les composants et leurs connexions, pas leurs dimensions ou leurs positions réelles sur la table.\par
\end{corrige}\end{blocLivre}
\begin{blocLivre}{cor-13-8}
\begin{corrige}{8}{Une panne à rechercher}
Causes possibles : mauvais contact, lampe défectueuse, générateur déchargé ou inadapté. Une observation identique peut avoir plusieurs causes ; il faut les tester successivement selon le protocole de classe.\par
\end{corrige}\end{blocLivre}
\begin{blocLivre}{cor-13-9}
\begin{corrige}{9}{Le moteur inversé}
Le moteur peut tourner dans le sens opposé. Le sens de rotation est une propriété du mouvement du moteur ; la forme rectangulaire des fils sur le schéma n’en constitue pas une trajectoire.\par
\end{corrige}\end{blocLivre}
\begin{blocLivre}{cor-13-10}
\begin{corrige}{10}{Suivre une boucle}
On suit un fil, la lampe, l’interrupteur fermé et les autres liaisons jusqu’à la borne négative. L’ordre dépend du schéma, mais le trajet doit relier les deux bornes du générateur par les récepteurs.\par
\pcFigure{figures/c6a7cbcb169788d4.pdf}{Modèle de correction · les proportions sont indicatives sauf mention contraire.}
\end{corrige}\end{blocLivre}
\begin{blocLivre}{cor-13-11}
\begin{corrige}{11}{Courant et énergie}
La lampe reçoit de l’énergie par transfert électrique et en transfère sous forme lumineuse et thermique. Dans une boucle fermée en fonctionnement, on ne décrit pas un courant qui disparaît progressivement dans la lampe.\par
\end{corrige}\end{blocLivre}
\begin{blocLivre}{cor-13-12}
\begin{corrige}{12}{Le bon témoin}
Un montage témoin valide la source et les connexions. Remplacer seulement la lampe connue par celle à tester permet de comparer sans multiplier les causes possibles. Le professeur choisit des composants compatibles.\par
\end{corrige}\end{blocLivre}
\begin{blocLivre}{cor-13-13}
\begin{corrige}{13}{Une connexion trompeuse}
Le fil ne touche pas nécessairement une partie conductrice reliée à la borne. Il faut identifier les vraies bornes de connexion et contrôler le contact électrique, pas seulement le contact visuel avec l’objet.\par
\end{corrige}\end{blocLivre}
\begin{blocLivre}{cor-13-14}
\begin{corrige}{14}{Deux schémas équivalents}
Non. Comparer les connexions entre les bornes des mêmes composants. Des dessins de formes différentes peuvent représenter le même circuit si tous les liens sont conservés.\par
\begin{pcEncadre}{method}{Pour comprendre}
Compare les bornes reliées entre elles : la forme des fils et la place du symbole sur la page ne changent pas les connexions.\par
\end{pcEncadre}
\pcFigure{figures/23cb91a1b45aefe5.pdf}{Modèle de correction · les proportions sont indicatives sauf mention contraire.}
\end{corrige}\end{blocLivre}
\begin{blocLivre}{cor-13-15}
\begin{corrige}{15}{Répondre à un cahier des charges}
Utiliser un générateur adapté, une lampe, des fils et un interrupteur dans une boucle unique. Le générateur fournit l’énergie ; la lampe éclaire ; les fils réalisent les liaisons ; l’interrupteur ouvre ou ferme la boucle. Le montage est vérifié par le professeur.\par
\pcFigure{figures/da232bc04f83598f.pdf}{Modèle de correction · les proportions sont indicatives sauf mention contraire.}
\end{corrige}\end{blocLivre}
\pcPage{Activité et défi}{Chapitre 13 · Éléments de réponse}
\phantomsection\label{solutions-13}
\begin{pcFiche}{solutions-13}
\begin{pcEncadre}{method}{Allumer, éteindre et dépanner}
La boucle relie une borne du générateur à l’autre en traversant la lampe et l’interrupteur fermé. L’ouverture interrompt ce trajet. Modifier une seule cause à la fois permet de relier un changement de fonctionnement à un élément précis ; changer simultanément pile, fils et lampe empêcherait d’identifier la panne.\par
\end{pcEncadre}
\pcRetour{ch-13-activite}{Revenir à l’activité}
\phantomsection\label{defi-cor-13}
\begin{pcEncadre}{definition}{Le dossier de dépannage}
Contrôler les bornes, les contacts et l’interrupteur. Tester ensuite la lampe de référence dans le même circuit ; si elle fonctionne, la lampe initiale est suspecte. Sinon, contrôler les fils un par un puis le générateur avec le professeur. Chaque test compare une seule modification. Une panne peut aussi résulter de plusieurs défauts : ne pas généraliser trop vite.\par
\end{pcEncadre}
\pcRetour{defi-ex-13}{Revenir au défi}
\end{pcFiche}
\clearpage\corrigesChapitre{14}{Conducteurs, isolants et sécurité électrique}
\begin{blocLivre}{cor-14-1}
\begin{corrige}{1}{Conducteur ou isolant ?}
Le cuivre est conducteur et la gaine plastique sèche est isolante dans ces conditions. Le classement décrit leur comportement dans le contexte étudié.\par
\end{corrige}\end{blocLivre}
\begin{blocLivre}{cor-14-2}
\begin{corrige}{2}{Le rôle de la gaine}
Le cœur métallique permet la liaison conductrice. Le plastique limite le contact avec le conducteur et sépare les fils. Une gaine abîmée ne doit pas être ignorée.\par
\end{corrige}\end{blocLivre}
\begin{blocLivre}{cor-14-3}
\begin{corrige}{3}{Lire le test}
L’échantillon complète la liaison prévue sur le côté du schéma. La lampe sert de témoin de fonctionnement et de récepteur ; on ne relie pas directement les bornes du générateur avec l’échantillon.\par
\pcFigure{figures/4855a50c85e1943a.pdf}{Modèle de correction · les proportions sont indicatives sauf mention contraire.}
\end{corrige}\end{blocLivre}
\begin{blocLivre}{cor-14-4}
\begin{corrige}{4}{Les expériences autorisées}
Non. On utilise uniquement le matériel scolaire de très basse tension prévu et validé par le professeur. Le secteur ne sert jamais à ces expériences.\par
\end{corrige}\end{blocLivre}
\begin{blocLivre}{cor-14-5}
\begin{corrige}{5}{Un chemin dangereux}
C’est un court-circuit du générateur. Il peut entraîner un courant important, un échauffement ou une détérioration ; on ne le réalise pas pour vérifier le risque.\par
\pcFigure{figures/e82b4cba3fbea2a7.pdf}{Modèle de correction · les proportions sont indicatives sauf mention contraire.}
\end{corrige}\end{blocLivre}
\begin{blocLivre}{cor-14-6}
\begin{corrige}{6}{Un témoin utile}
Non : le montage n’est pas validé. Rechercher la cause de panne ou de mauvais contact avec le professeur, puis refaire le test après avoir retrouvé le fonctionnement du témoin.\par
\end{corrige}\end{blocLivre}
\begin{blocLivre}{cor-14-7}
\begin{corrige}{7}{Un métal recouvert}
La peinture peut empêcher le contact avec le métal. Il faut identifier la surface réellement touchée par les connexions et contrôler le contact, selon le protocole de classe.\par
\end{corrige}\end{blocLivre}
\begin{blocLivre}{cor-14-8}
\begin{corrige}{8}{Repérer le fil ajouté}
Le fil contourne l’unique récepteur et laisse un trajet presque direct entre les bornes de la pile. C’est une situation de court-circuit à reconnaître sur document, sans la réaliser.\par
\pcFigure{figures/e82b4cba3fbea2a7.pdf}{Modèle de correction · les proportions sont indicatives sauf mention contraire.}
\end{corrige}\end{blocLivre}
\begin{blocLivre}{cor-14-9}
\begin{corrige}{9}{Le corps humain}
Le corps peut conduire le courant et ne constitue pas un échantillon de test. L’absence de sensation dans une situation ne prouve pas la sécurité d’autres tensions, sources ou conditions.\par
\end{corrige}\end{blocLivre}
\begin{blocLivre}{cor-14-10}
\begin{corrige}{10}{L’eau et les risques}
Les eaux usuelles contiennent souvent des espèces dissoutes et ne sont pas des isolants fiables. Respecter les consignes d’usage et tenir l’eau à l’écart des montages et appareils non prévus pour cela ; ne pas improviser de test.\par
\end{corrige}\end{blocLivre}
\begin{blocLivre}{cor-14-11}
\begin{corrige}{11}{Deux contrôles de contact}
Une surface peinte ou oxydée et une connexion mal placée peuvent empêcher un bon contact. Contrôler les bornes et utiliser un témoin dans les mêmes conditions avant de classer l’échantillon.\par
\end{corrige}\end{blocLivre}
\begin{blocLivre}{cor-14-12}
\begin{corrige}{12}{Rédiger un résultat limité}
« Avec ce morceau de bois sec, la lampe ne s’allume pas dans les conditions du test validé. » L’humidité et la composition peuvent modifier le comportement ; le résultat ne décrit pas toutes les situations possibles.\par
\end{corrige}\end{blocLivre}
\begin{blocLivre}{cor-14-13}
\begin{corrige}{13}{La pile chauffe}
Arrêter l’activité selon les consignes, ne pas toucher la partie échauffée et prévenir le professeur. Un échauffement peut signaler un mauvais branchement ou un défaut ; poursuivre peut aggraver le risque.\par
\end{corrige}\end{blocLivre}
\begin{blocLivre}{cor-14-14}
\begin{corrige}{14}{La variable étudiée}
Le fonctionnement peut changer à cause de la source ou du matériau. Conserver le même montage validé et modifier uniquement l’échantillon, avec des contacts comparables et des témoins.\par
\end{corrige}\end{blocLivre}
\begin{blocLivre}{cor-14-15}
\begin{corrige}{15}{Concevoir une fiche de test}
Indiquer le matériau et son état, le schéma et le matériel, les témoins, les contacts, l’observation et l’interprétation. Ajouter les limites : sensibilité du dispositif, humidité, revêtement et éventuels défauts de contact. Ne pas présenter un résultat unique comme universel.\par
\end{corrige}\end{blocLivre}
\pcPage{Activité et défi}{Chapitre 14 · Éléments de réponse}
\phantomsection\label{solutions-14}
\begin{pcFiche}{solutions-14}
\begin{pcEncadre}{method}{Comparer trois matériaux}
Le témoin vérifie que le montage reste fonctionnel avant et après la série. Un résultat négatif s’écrit : « la lampe ne s’allume pas avec cet échantillon dans les conditions du test », puis s’interprète avec les contrôles. Le métal réalise la liaison conductrice ; le revêtement limite les contacts involontaires avec le conducteur.\par
\end{pcEncadre}
\pcRetour{ch-14-activite}{Revenir à l’activité}
\phantomsection\label{defi-cor-14}
\begin{pcEncadre}{definition}{Le métal déclaré isolant}
Les groupes peuvent avoir touché des zones différentes ou utilisé un montage défectueux. Valider d’abord le circuit témoin ; comparer ensuite la zone peinte et la zone métallique accessible sans modifier le reste. La tige peut être conductrice tandis que sa peinture limite le contact. Toute préparation de l’échantillon relève du professeur.\par
\end{pcEncadre}
\pcRetour{defi-ex-14}{Revenir au défi}
\end{pcFiche}
\clearpage\corrigesChapitre{15}{Circuits en série et en dérivation}
\begin{blocLivre}{cor-15-1}
\begin{corrige}{1}{Une boucle unique}
Il existe un seul chemin dans la boucle ; on traverse les deux lampes successivement. Aucune jonction ne permet de choisir une autre branche.\par
\end{corrige}\end{blocLivre}
\begin{blocLivre}{cor-15-2}
\begin{corrige}{2}{Nommer l’association}
Les deux lampes sont en série. Elles appartiennent à une seule boucle et sont traversées successivement en suivant le trajet.\par
\pcFigure{figures/1b916fb9a3d436c2.pdf}{Modèle de correction · les proportions sont indicatives sauf mention contraire.}
\end{corrige}\end{blocLivre}
\begin{blocLivre}{cor-15-3}
\begin{corrige}{3}{Une ouverture en série}
L’autre lampe s’éteint également : la boucle unique est interrompue. Il n’existe pas de chemin alternatif dans ce montage.\par
\pcFigure{figures/32535d3f51866fd7.pdf}{Modèle de correction · les proportions sont indicatives sauf mention contraire.}
\end{corrige}\end{blocLivre}
\begin{blocLivre}{cor-15-4}
\begin{corrige}{4}{Une dérivation}
C’est une association en dérivation. La propriété vient des connexions ; les fils peuvent être dessinés autrement tout en conservant les mêmes jonctions.\par
\pcFigure{figures/5fb4215f70151542.pdf}{Modèle de correction · les proportions sont indicatives sauf mention contraire.}
\end{corrige}\end{blocLivre}
\begin{blocLivre}{cor-15-5}
\begin{corrige}{5}{Une commande indépendante}
Le placer dans cette branche, sans interrompre les portions communes aux autres chemins. Vérifier sur le schéma que les autres récepteurs conservent une boucle fermée avec le générateur.\par
\pcFigure{figures/71e1cac1257b2ab4.pdf}{Modèle de correction · les proportions sont indicatives sauf mention contraire.}
\end{corrige}\end{blocLivre}
\begin{blocLivre}{cor-15-6}
\begin{corrige}{6}{Suivre les chemins}
Chaque chemin relie les deux bornes du générateur en passant par une branche de lampe. Ils partagent le générateur et certaines liaisons principales ; les lampes occupent des branches différentes.\par
\pcFigure{figures/5fb4215f70151542.pdf}{Modèle de correction · les proportions sont indicatives sauf mention contraire.}
\end{corrige}\end{blocLivre}
\begin{blocLivre}{cor-15-7}
\begin{corrige}{7}{La branche ouverte}
L\textsubscript{2} reste sur une boucle fermée et peut fonctionner. La branche de L₁ est ouverte par son interrupteur et ne permet pas son fonctionnement.\par
\begin{pcEncadre}{method}{Pour comprendre}
Cherche séparément une boucle fermée pour chaque lampe. Une branche ouverte n’interrompt pas l’autre branche si celle-ci relie encore les deux bornes du générateur.\par
\end{pcEncadre}
\pcFigure{figures/413870db94cfb67c.pdf}{Modèle de correction · les proportions sont indicatives sauf mention contraire.}
\end{corrige}\end{blocLivre}
\begin{blocLivre}{cor-15-8}
\begin{corrige}{8}{Commander les deux lampes}
Il réalise une commande générale. Son ouverture interrompt tous les chemins qui ont besoin de cette liaison ; les deux lampes s’éteignent.\par
\pcFigure{figures/c132ea2c6e94fe88.pdf}{Modèle de correction · les proportions sont indicatives sauf mention contraire.}
\end{corrige}\end{blocLivre}
\begin{blocLivre}{cor-15-9}
\begin{corrige}{9}{Changer l’ordre}
Non. Le circuit conserve une seule boucle et les deux lampes restent successives. Le nombre de branches et la structure des connexions restent identiques.\par
\end{corrige}\end{blocLivre}
\begin{blocLivre}{cor-15-10}
\begin{corrige}{10}{Les conditions d’éclat}
L’éclat dépend des caractéristiques des lampes et de la source. Contrôler les modèles de lampes, leur compatibilité, le générateur et ses conditions de fonctionnement. Un exemple ne justifie pas une loi pour tous les composants.\par
\end{corrige}\end{blocLivre}
\begin{blocLivre}{cor-15-11}
\begin{corrige}{11}{Reconnaître une jonction}
Les points indiquent une connexion conductrice. Un croisement sans légende peut laisser hésiter entre fils connectés et fils qui se croisent sans contact ; le schéma doit utiliser une convention explicite.\par
\begin{pcEncadre}{method}{Pour comprendre}
Un point noir est une jonction : les conducteurs y communiquent. Un pont de fil indique ici un croisement sans contact. La légende fixe la convention.\par
\end{pcEncadre}
\pcFigure{figures/8787152f2a635390.pdf}{Modèle de correction · les proportions sont indicatives sauf mention contraire.}
\end{corrige}\end{blocLivre}
\begin{blocLivre}{cor-15-12}
\begin{corrige}{12}{Deux lampes pour un couloir}
Choisir la dérivation : chaque lampe possède sa branche, formant un chemin avec le générateur. L’ouverture d’une seule branche n’interrompt pas forcément le chemin de l’autre. Il s’agit d’un modèle scolaire, pas d’une instruction de câblage domestique.\par
\end{corrige}\end{blocLivre}
\begin{blocLivre}{cor-15-13}
\begin{corrige}{13}{Une fausse dérivation}
Non. Le fil direct peut court-circuiter le générateur. Il faut contrôler tous les chemins, notamment ceux qui évitent les récepteurs. Le nombre de branches ne garantit pas la sécurité.\par
\pcFigure{figures/c3df0e174684c6a7.pdf}{Modèle de correction · les proportions sont indicatives sauf mention contraire.}
\end{corrige}\end{blocLivre}
\begin{blocLivre}{cor-15-14}
\begin{corrige}{14}{Un symptôme, plusieurs causes}
Vérifier que le montage réel correspond au schéma et que l’interrupteur n’est pas dans une portion commune. Vérifier aussi le générateur et les connexions. L’observation seule ne permet pas d’identifier sans contrôle la cause exacte.\par
\end{corrige}\end{blocLivre}
\begin{blocLivre}{cor-15-15}
\begin{corrige}{15}{Trois fonctions à réaliser}
Placer le général dans une portion commune ; placer chacun des deux autres dans sa branche de lampe. Général fermé, ouvrir successivement un interrupteur de branche et vérifier que seule la lampe correspondante s’éteint, dans les conditions du montage validé.\par
\end{corrige}\end{blocLivre}
\pcPage{Activité et défi}{Chapitre 15 · Éléments de réponse}
\phantomsection\label{solutions-15}
\begin{pcFiche}{solutions-15}
\begin{pcEncadre}{method}{Choisir une commande d’éclairage}
Le montage en dérivation permet à L\textsubscript{2} de conserver sa boucle avec le générateur. Une commande générale se place dans une portion commune à tous les trajets, par exemple sur une liaison principale du générateur. Un fil qui contourne les récepteurs peut créer un court-circuit ; il ne constitue pas une commande d’éclairage correcte.\par
\end{pcEncadre}
\pcRetour{ch-15-activite}{Revenir à l’activité}
\phantomsection\label{defi-cor-15}
\begin{pcEncadre}{definition}{Le tableau de commande}
Choisir deux branches en dérivation, chacune avec son interrupteur, et un interrupteur général sur une liaison commune. Général fermé : ouvrir les deux branches, fermer seulement celle de L₁, seulement celle de L\textsubscript{2}, ou les deux. Général ouvert : toutes les lampes s’éteignent quel que soit l’état des commandes de branche.\par
\end{pcEncadre}
\pcFigure{figures/187096fd29805e74.pdf}{Construction ou modèle de correction.}
\pcRetour{defi-ex-15}{Revenir au défi}
\end{pcFiche}
\clearpage\corrigesChapitre{16}{Énergie et choix responsables}
\begin{blocLivre}{cor-16-1}
\begin{corrige}{1}{Classer des ressources}
Vent et Soleil sont des ressources renouvelables à cette échelle ; le pétrole provient d’un stock non renouvelable. Cette classification ne résume pas tous les impacts des équipements.\par
\end{corrige}\end{blocLivre}
\begin{blocLivre}{cor-16-2}
\begin{corrige}{2}{Lire les unités}
W est une unité de puissance. Wh et kWh sont des unités d’énergie. Les symboles ne sont pas interchangeables.\par
\end{corrige}\end{blocLivre}
\begin{blocLivre}{cor-16-3}
\begin{corrige}{3}{Une conversion d’énergie}
2 kWh = 2 000 Wh. 500 Wh = 0,5 kWh. On garde une grandeur énergie pendant la conversion.\par
\end{corrige}\end{blocLivre}
\begin{blocLivre}{cor-16-4}
\begin{corrige}{4}{Le bon critère}
On peut comparer leurs nombres, mais pas conclure à une meilleure efficacité pour le même besoin sans service équivalent. Il faut demander un éclairage utile et des conditions d’usage comparables.\par
\end{corrige}\end{blocLivre}
\begin{blocLivre}{cor-16-5}
\begin{corrige}{5}{Un renouvelable sans impact ?}
Le vent est une ressource renouvelable, mais la fabrication, l’installation, l’entretien et la fin de vie de l’éolienne mobilisent des matériaux et de l’énergie. Renouvelable ne signifie pas sans ressources ni sans impact.\par
\end{corrige}\end{blocLivre}
\begin{blocLivre}{cor-16-6}
\begin{corrige}{6}{Lire deux barres}
120 - 36 = 84 kWh. L’écart concerne une année selon le document. Il ne faut pas le présenter comme une économie quotidienne.\par
\end{corrige}\end{blocLivre}
\begin{blocLivre}{cor-16-7}
\begin{corrige}{7}{La durée de fonctionnement}
À puissance constante, 50 W pendant 2 h correspondent à 100 Wh ; pendant 1 h, à 50 Wh. L’énergie est proportionnelle à la durée dans ce cas.\par
\end{corrige}\end{blocLivre}
\begin{blocLivre}{cor-16-8}
\begin{corrige}{8}{Passer à plusieurs jours}
5 \(\times\) 40 = 200 Wh, soit 0,2 kWh. Le calcul suppose que l’usage quotidien reste identique.\par
\end{corrige}\end{blocLivre}
\begin{blocLivre}{cor-16-9}
\begin{corrige}{9}{Réduire un usage inutile}
Éteindre l’éclairage lorsque le service n’est pas nécessaire peut réduire l’énergie d’usage. La décision doit respecter les besoins de sécurité, de confort et les conditions propres au lieu.\par
\end{corrige}\end{blocLivre}
\begin{blocLivre}{cor-16-10}
\begin{corrige}{10}{Un stock non renouvelable}
Non : le stock d’uranium n’est pas renouvelable à l’échelle humaine. Distinguer renouvellement de la ressource et émissions ou impacts du dispositif. Une seule propriété ne remplace pas un bilan complet.\par
\end{corrige}\end{blocLivre}
\begin{blocLivre}{cor-16-11}
\begin{corrige}{11}{Un même service sur deux ans}
L’économie annuelle est 84 kWh ; sur deux ans, 168 kWh. On suppose le même usage, les mêmes performances et le même service chaque année.\par
\end{corrige}\end{blocLivre}
\begin{blocLivre}{cor-16-12}
\begin{corrige}{12}{Moins puissant, mais plus longtemps}
A utilise 100 Wh et B 80 Wh. L’appareil moins puissant peut utiliser davantage d’énergie s’il fonctionne plus longtemps. Il faut encore vérifier que le service comparé est pertinent.\par
\end{corrige}\end{blocLivre}
\begin{blocLivre}{cor-16-13}
\begin{corrige}{13}{Une donnée publicitaire}
Demander 30 \% de quelle grandeur, par rapport à quel appareil, pour quel service et quelles conditions. Des périodes ou des durées différentes changent l’énergie ; une comparaison honnête précise sa référence.\par
\end{corrige}\end{blocLivre}
\begin{blocLivre}{cor-16-14}
\begin{corrige}{14}{Calculer une part économisée}
L’économie est de 50 kWh. 50 sur 200 représente un quart, soit 25 \%. Le pourcentage est calculé par rapport à la consommation de départ.\par
\end{corrige}\end{blocLivre}
\begin{blocLivre}{cor-16-15}
\begin{corrige}{15}{Un bilan de décision}
Rechercher le service réellement fourni, la consommation selon l’usage, la durée de vie et la réparabilité, ainsi que les impacts de fabrication et de fin de vie. Le remplacement mobilise aussi des ressources ; un choix raisonné considère davantage que le seul chiffre de consommation.\par
\end{corrige}\end{blocLivre}
\pcPage{Activité et défi}{Chapitre 16 · Éléments de réponse}
\phantomsection\label{solutions-16}
\begin{pcFiche}{solutions-16}
\begin{pcEncadre}{method}{Comparer deux éclairages fictifs}
L’économie vaut 120 - 36 = 84 kWh, soit 70 \% de la consommation de référence A ; c’est plus de la moitié. Les seules consommations à l’usage ne décrivent pas fabrication, durée de vie, transport, réparabilité ni fin de vie. La conclusion porte sur l’énergie d’usage pour le service et la période annoncés.\par
\end{pcEncadre}
\pcRetour{ch-16-activite}{Revenir à l’activité}
\phantomsection\label{defi-cor-16}
\begin{pcEncadre}{definition}{L’affiche d’une classe économe}
Consommation finale : 300 - 60 - 40 = 200 kWh ; économie : 100 kWh, soit un tiers de la référence. L’affiche précise le même service, les durées supposées et le caractère estimé des chiffres. Elle distingue économies à l’usage et impacts du matériel. Elle ne prétend pas que l’énergie ou les impacts deviennent nuls.\par
\end{pcEncadre}
\pcRetour{defi-ex-16}{Revenir au défi}
\end{pcFiche}
\clearpage\corrigesChapitre{17}{Sources de lumière et vision}
\begin{blocLivre}{cor-17-1}
\begin{corrige}{1}{Deux types de sources}
La lampe allumée est une source primaire ; la Lune éclairée diffuse la lumière du Soleil. Elle ne produit pas cette lumière visible.\par
\end{corrige}\end{blocLivre}
\begin{blocLivre}{cor-17-2}
\begin{corrige}{2}{Le sens de la lumière}
La lumière diffusée va du livre vers l’œil. Le livre reçoit auparavant de la lumière d’une source ; l’œil ne l’éclaire pas pour le voir.\par
\end{corrige}\end{blocLivre}
\begin{blocLivre}{cor-17-3}
\begin{corrige}{3}{Un écran blanc}
Oui, elle peut être opaque tout en étant visible grâce à la lumière qu’elle diffuse. Opaque et invisible ne sont pas synonymes.\par
\end{corrige}\end{blocLivre}
\begin{blocLivre}{cor-17-4}
\begin{corrige}{4}{Nommer un matériau}
Le verre dépoli est translucide dans ces conditions. Une vitre permettant une vision nette est transparente.\par
\end{corrige}\end{blocLivre}
\begin{blocLivre}{cor-17-5}
\begin{corrige}{5}{Lire les flèches}
L’objet central diffuse une partie de la lumière reçue. La flèche allant de l’objet à l’œil représente la lumière reçue par l’observateur.\par
\pcFigure{figures/cd80bfc6ec09d32d.pdf}{Modèle de correction · les proportions sont indicatives sauf mention contraire.}
\end{corrige}\end{blocLivre}
\begin{blocLivre}{cor-17-6}
\begin{corrige}{6}{Dans l’obscurité}
Non dans le modèle étudié : aucune lumière provenant de lui n’entre dans l’œil. Il faut un éclairage et un trajet de lumière vers l’observateur.\par
\end{corrige}\end{blocLivre}
\begin{blocLivre}{cor-17-7}
\begin{corrige}{7}{Un objet caché}
L’écran interrompt le trajet de lumière entre le livre et l’œil. Le trajet source-livre peut rester possible ; cela ne suffit pas pour assurer la vision de l’objet.\par
\pcFigure{figures/9eaef4335ed1b9e0.pdf}{Modèle de correction · les proportions sont indicatives sauf mention contraire.}
\end{corrige}\end{blocLivre}
\begin{blocLivre}{cor-17-8}
\begin{corrige}{8}{Brillant ne veut pas dire primaire}
Elle peut renvoyer la lumière d’une source extérieure. Le fait de paraître brillante ne démontre pas une production de lumière ; il faut rechercher l’origine du rayonnement reçu par l’œil.\par
\end{corrige}\end{blocLivre}
\begin{blocLivre}{cor-17-9}
\begin{corrige}{9}{Comparer les trois écrans}
Le matériau opaque interrompt la transmission visible du schéma. Le translucide transmet dans plusieurs directions, ce qui empêche ici une observation nette à travers.\par
\pcFigure{figures/2158b9095894ca6c.pdf}{Modèle de correction · les proportions sont indicatives sauf mention contraire.}
\end{corrige}\end{blocLivre}
\begin{blocLivre}{cor-17-10}
\begin{corrige}{10}{Éteindre la source}
La balle ne reçoit plus l’éclairage nécessaire et ne peut plus diffuser cette lumière vers l’œil. On suppose qu’aucune autre source ni lumière extérieure n’éclaire la scène.\par
\end{corrige}\end{blocLivre}
\begin{blocLivre}{cor-17-11}
\begin{corrige}{11}{Une planète visible}
Elle diffuse une partie de la lumière solaire reçue. Le Soleil est la source primaire principale considérée pour sa visibilité.\par
\end{corrige}\end{blocLivre}
\begin{blocLivre}{cor-17-12}
\begin{corrige}{12}{Un schéma à corriger}
Inverser la flèche pour représenter la lumière objet \(\rightarrow\) œil. Ajouter une source qui éclaire l’objet, puis le trajet source \(\rightarrow\) objet. La légende doit préciser qu’on suit la lumière.\par
\begin{pcEncadre}{method}{Pour comprendre}
Suis la lumière : une source éclaire l’objet, puis l’objet diffuse vers l’œil. La flèche à corriger doit donc finir du côté de l’œil.\par
\end{pcEncadre}
\pcFigure{figures/cd80bfc6ec09d32d.pdf}{Modèle de correction · les proportions sont indicatives sauf mention contraire.}
\end{corrige}\end{blocLivre}
\begin{blocLivre}{cor-17-13}
\begin{corrige}{13}{Transparent et réflexion}
Oui. Une vitre peut transmettre une partie de la lumière et en réfléchir une autre. Transparent ne signifie pas que toute la lumière traverse sans autre effet.\par
\end{corrige}\end{blocLivre}
\begin{blocLivre}{cor-17-14}
\begin{corrige}{14}{L’épaisseur du matériau}
Préciser l’épaisseur et les conditions d’observation. Comparer avec la même source, le même objet, la même géométrie et des épaisseurs annoncées. Le mot matériau seul ne décrit pas toutes les situations.\par
\end{corrige}\end{blocLivre}
\begin{blocLivre}{cor-17-15}
\begin{corrige}{15}{Préparer une démonstration}
Éclairer l’objet et observer la lumière diffusée vers l’œil. Masquer source-objet empêche l’éclairage ; masquer objet-œil empêche la lumière diffusée d’atteindre l’observateur. Utiliser une source scolaire adaptée et ne pas diriger une lumière intense vers les yeux.\par
\pcFigure{figures/c9d32114c6d4de65.pdf}{Modèle de correction · les proportions sont indicatives sauf mention contraire.}
\end{corrige}\end{blocLivre}
\pcPage{Activité et défi}{Chapitre 17 · Éléments de réponse}
\phantomsection\label{solutions-17}
\begin{pcFiche}{solutions-17}
\begin{pcEncadre}{method}{Voir à travers, voir l’objet}
Le matériau transparent permet une observation nette dans les conditions prévues. L’écran opaque blanc peut être visible en diffusant la lumière reçue, tout en empêchant de voir l’objet à travers. La flèche objet \(\rightarrow\) œil représente la lumière qui doit entrer dans l’œil ; l’éclairage de l’objet seul ne suffit pas à expliquer la vision.\par
\end{pcEncadre}
\pcRetour{ch-17-activite}{Revenir à l’activité}
\phantomsection\label{defi-cor-17}
\begin{pcEncadre}{definition}{La boîte à regarder}
La lumière doit entrer dans la boîte, atteindre la balle puis être diffusée vers l’œil par le trou de regard. Si l’unique éclairage est bloqué, la balle ne reçoit plus la lumière requise. Un trou de regard ne produit pas de lumière ; il permet seulement un trajet. Une autre lumière parasite doit être exclue pour interpréter le test.\par
\end{pcEncadre}
\pcFigure{figures/cd80bfc6ec09d32d.pdf}{Construction ou modèle de correction.}
\pcRetour{defi-ex-17}{Revenir au défi}
\end{pcFiche}
\clearpage\corrigesChapitre{18}{Propagation de la lumière et ombres}
\begin{blocLivre}{cor-18-1}
\begin{corrige}{1}{Une condition de milieu}
Dans un milieu transparent homogène. Une flèche indique le sens de propagation de la lumière sur le rayon représenté.\par
\end{corrige}\end{blocLivre}
\begin{blocLivre}{cor-18-2}
\begin{corrige}{2}{Nommer un rayon}
Non : c’est une représentation géométrique d’un trajet de lumière. Le trait fléché aide à prévoir les chemins possibles et ne décrit pas un objet matériel tendu dans l’espace.\par
\end{corrige}\end{blocLivre}
\begin{blocLivre}{cor-18-3}
\begin{corrige}{3}{Les trous alignés}
La propagation rectiligne exige un trajet droit qui passe par les ouvertures. Si un écran opaque masque ce trajet, la lumière directe ne poursuit pas ce chemin.\par
\pcFigure{figures/3f9ef26636cab967.pdf}{Modèle de correction · les proportions sont indicatives sauf mention contraire.}
\end{corrige}\end{blocLivre}
\begin{blocLivre}{cor-18-4}
\begin{corrige}{4}{L’ombre sur un mur}
La tache sur l’écran est l’ombre portée. La partie de la surface de l’objet non directement éclairée est l’ombre propre.\par
\pcFigure{figures/3f89b4be149ab996.pdf}{Modèle de correction · les proportions sont indicatives sauf mention contraire.}
\end{corrige}\end{blocLivre}
\begin{blocLivre}{cor-18-5}
\begin{corrige}{5}{Un objet opaque}
Non. Le modèle d’ombre utilisé suppose que l’objet intercepte les trajets lumineux. Un matériau transparent transmet une partie de la lumière et demande une analyse différente.\par
\end{corrige}\end{blocLivre}
\begin{blocLivre}{cor-18-6}
\begin{corrige}{6}{Une ouverture déplacée}
Il rencontre la partie opaque de l’écran central. On ne courbe pas artificiellement le rayon pour rejoindre l’ouverture : dans le milieu homogène du modèle, le trajet reste rectiligne.\par
\pcFigure{figures/62d9a8700e3bfecf.pdf}{Modèle de correction · les proportions sont indicatives sauf mention contraire.}
\end{corrige}\end{blocLivre}
\begin{blocLivre}{cor-18-7}
\begin{corrige}{7}{Repérer les limites}
Ils partent de la source ponctuelle et passent aux bords de l’objet opaque. Leur prolongement délimite la zone d’ombre et la taille de l’ombre portée sur l’écran.\par
\pcFigure{figures/3f89b4be149ab996.pdf}{Modèle de correction · les proportions sont indicatives sauf mention contraire.}
\end{corrige}\end{blocLivre}
\begin{blocLivre}{cor-18-8}
\begin{corrige}{8}{Déplacer l’objet}
Elle grandit dans le modèle étudié. Il faut préciser que source et écran restent fixes et que l’on utilise la même petite source et le même objet entre les deux.\par
\end{corrige}\end{blocLivre}
\begin{blocLivre}{cor-18-9}
\begin{corrige}{9}{Un écran derrière l’objet}
Non. L’ombre portée correspond à la rencontre de l’écran avec la zone d’ombre. Déplacer l’écran peut modifier la tache ou le placer en dehors de cette région.\par
\end{corrige}\end{blocLivre}
\begin{blocLivre}{cor-18-10}
\begin{corrige}{10}{Deux essais comparables}
Non. Conserver la même source, le même objet, le même écran et ne faire varier qu’une distance définie. Noter toutes les positions pour pouvoir comparer.\par
\end{corrige}\end{blocLivre}
\begin{blocLivre}{cor-18-11}
\begin{corrige}{11}{Le sens des flèches}
Les flèches suivent la propagation depuis la source. Leur sens n’est pas choisi selon la lecture gauche-droite : si la source est à droite, les flèches peuvent pointer vers la gauche.\par
\end{corrige}\end{blocLivre}
\begin{blocLivre}{cor-18-12}
\begin{corrige}{12}{Une ombre moins nette}
La source peut être étendue plutôt que ponctuelle, et d’autres lumières peuvent éclairer la zone. L’objet et l’écran peuvent aussi être mal positionnés. Le modèle simplifié ne prétend pas reproduire tous les détails.\par
\end{corrige}\end{blocLivre}
\begin{blocLivre}{cor-18-13}
\begin{corrige}{13}{Distinguer trois régions}
L’ombre propre appartient à la surface de l’objet ; la zone d’ombre est un volume de l’espace qui ne reçoit pas directement la source ; l’ombre portée est l’intersection de cette zone avec un écran. La localisation distingue les trois notions.\par
\pcFigure{figures/3f89b4be149ab996.pdf}{Modèle de correction · les proportions sont indicatives sauf mention contraire.}
\end{corrige}\end{blocLivre}
\begin{blocLivre}{cor-18-14}
\begin{corrige}{14}{Une construction impossible}
Il fait traverser la lumière à travers l’objet supposé opaque. Les rayons limites passent par les bords de l’objet depuis la source ; on les prolonge jusqu’à l’écran sans les faire traverser le solide.\par
\begin{pcEncadre}{method}{Pour comprendre}
Le solide opaque arrête les rayons qui le rencontrent. Pour délimiter l’ombre, utilise les rayons qui rasent les deux bords, puis prolonge-les jusqu’à l’écran.\par
\end{pcEncadre}
\pcFigure{figures/3f89b4be149ab996.pdf}{Modèle de correction · les proportions sont indicatives sauf mention contraire.}
\end{corrige}\end{blocLivre}
\begin{blocLivre}{cor-18-15}
\begin{corrige}{15}{Confronter modèle et mesure}
Vérifier les positions réellement fixes, la direction du déplacement, la taille de la source et la précision de mesure. Un déplacement trop faible peut produire une variation inférieure à la résolution de la mesure. Refaire une comparaison contrôlée et conserver les résultats observés.\par
\end{corrige}\end{blocLivre}
\pcPage{Activité et défi}{Chapitre 18 · Éléments de réponse}
\phantomsection\label{solutions-18}
\begin{pcFiche}{solutions-18}
\begin{pcEncadre}{method}{Faire varier une ombre}
On déplace seulement l’objet, entre source et écran fixes. Dans le modèle de petite source, son ombre portée grandit lorsqu’il se rapproche de la source. Garder la source, l’écran, l’objet et les autres conditions identiques permet d’attribuer la variation à la seule position. Une source étendue ou d’autres éclairages peuvent rendre les limites moins nettes.\par
\end{pcEncadre}
\pcRetour{ch-18-activite}{Revenir à l’activité}
\phantomsection\label{defi-cor-18}
\begin{pcEncadre}{definition}{Le théâtre d’ombres}
Rapprocher l’objet de la source en le maintenant entre elle et l’écran. Les rayons limites s’écartent davantage avant d’atteindre l’écran : la silhouette grandit. Une source étendue peut donner des bords moins nets ; il faut alors compléter le modèle ou tester le résultat. Utiliser uniquement la lampe adaptée fournie pour l’activité.\par
\end{pcEncadre}
\pcFigure{figures/fb6a0bd783e84e52.pdf}{Construction ou modèle de correction.}
\pcRetour{defi-ex-18}{Revenir au défi}
\end{pcFiche}
\clearpage\corrigesChapitre{19}{Produire et transmettre un son}
\begin{blocLivre}{cor-19-1}
\begin{corrige}{1}{À la source}
La membrane vibre. Un microphone peut recevoir le signal sonore et permettre son enregistrement. L’émetteur et le récepteur ont des rôles distincts.\par
\end{corrige}\end{blocLivre}
\begin{blocLivre}{cor-19-2}
\begin{corrige}{2}{Le milieu nécessaire}
Le son ne se propage pas dans le vide. Il peut se propager dans un gaz comme l’air, dans un liquide ou dans un solide, avec des vitesses différentes.\par
\end{corrige}\end{blocLivre}
\begin{blocLivre}{cor-19-3}
\begin{corrige}{3}{Une vitesse avec son unité}
On utilise environ 340 m/s. La vitesse est exprimée en mètres par seconde ; 340 seul ne constitue pas un résultat complet.\par
\end{corrige}\end{blocLivre}
\begin{blocLivre}{cor-19-4}
\begin{corrige}{4}{Deux propriétés}
Non. Grave décrit la hauteur ; fort décrit le niveau reçu. La fréquence est liée à la hauteur : une fréquence plus élevée correspond à un son plus aigu.\par
\end{corrige}\end{blocLivre}
\begin{blocLivre}{cor-19-5}
\begin{corrige}{5}{Une protection raisonnée}
Réduire le niveau, limiter la durée et s’éloigner de la source sont des moyens de réduire l’exposition. Il ne faut pas attendre la douleur pour diminuer un son gênant ou trop fort.\par
\end{corrige}\end{blocLivre}
\begin{blocLivre}{cor-19-6}
\begin{corrige}{6}{Distance parcourue en une seconde}
Il parcourt environ 340 m en 1 s et 680 m en 2 s, si les conditions de propagation restent celles du modèle.\par
\end{corrige}\end{blocLivre}
\begin{blocLivre}{cor-19-7}
\begin{corrige}{7}{Un trajet direct}
170 \(\div\) 340 = 0,5 s. Le résultat est une durée en secondes, pas une distance en mètres.\par
\end{corrige}\end{blocLivre}
\begin{blocLivre}{cor-19-8}
\begin{corrige}{8}{Comparer des traces}
B présente plus d’oscillations dans la même durée : sa fréquence est plus élevée. Entre A et C, l’amplitude change tandis que le rythme des oscillations reste identique.\par
\pcFigure{figures/ccf2d1287fdfb2dc.pdf}{Modèle de correction · les proportions sont indicatives sauf mention contraire.}
\end{corrige}\end{blocLivre}
\begin{blocLivre}{cor-19-9}
\begin{corrige}{9}{La cloche et le vide}
Le milieu gazeux qui transmet la perturbation est fortement raréfié. Les supports solides peuvent encore transmettre des vibrations ; l’expérience réelle ne se réduit pas à un trajet uniquement dans l’air.\par
\end{corrige}\end{blocLivre}
\begin{blocLivre}{cor-19-10}
\begin{corrige}{10}{Un écho simple}
Distance totale : 340 \(\times\) 1 = 340 m. L’obstacle est à 340 \(\div\) 2 = 170 m, car le son parcourt l’aller et le retour.\par
\begin{pcEncadre}{method}{Pour comprendre}
Le temps mesuré inclut deux trajets de même longueur : l’aller vers l’obstacle et le retour. Après le produit vitesse \(\times\) durée, partage la distance totale en deux.\par
\end{pcEncadre}
\pcFigure{figures/2ffed6ca376994a6.pdf}{Modèle de correction · les proportions sont indicatives sauf mention contraire.}
\end{corrige}\end{blocLivre}
\begin{blocLivre}{cor-19-11}
\begin{corrige}{11}{La matière ne voyage pas tout entière}
Une perturbation vibratoire se transmet de proche en proche dans le milieu. Cela ne signifie pas que toute la masse d’air parcourt le trajet depuis la source jusqu’au récepteur.\par
\end{corrige}\end{blocLivre}
\begin{blocLivre}{cor-19-12}
\begin{corrige}{12}{Un délai de 3 secondes}
La distance est environ 340 \(\times\) 3 = 1 020 m. On utilise une propagation directe dans l’air avec 340 m/s et on néglige le temps de propagation de la lumière sur cette distance.\par
\end{corrige}\end{blocLivre}
\begin{blocLivre}{cor-19-13}
\begin{corrige}{13}{Un écho plus rapide}
Distance parcourue : 340 \(\times\) 0,4 = 136 m. Distance à l’obstacle : 68 m. Oublier de diviser le trajet aller-retour par deux donnerait 136 m.\par
\begin{pcEncadre}{method}{Pour comprendre}
Écris d’abord d’aller-retour = 340 \(\times\) 0,4 = 136 m. La distance recherchée est la moitié : 136 \(\div\) 2 = 68 m.\par
\end{pcEncadre}
\pcFigure{figures/018e6252910da6e9.pdf}{Modèle de correction · les proportions sont indicatives sauf mention contraire.}
\end{corrige}\end{blocLivre}
\begin{blocLivre}{cor-19-14}
\begin{corrige}{14}{Comparer des enregistrements}
Non. Le nombre d’oscillations doit être comparé sur une même durée ; à nombre égal sur deux fois plus de temps, le rythme est plus faible. Il faut vérifier les échelles temporelles avant la comparaison.\par
\begin{pcEncadre}{method}{Pour comprendre}
Compare un nombre d’oscillations par unité de temps : deux oscillations en 1 s et deux en 2 s ne correspondent pas au même rythme.\par
\end{pcEncadre}
\pcFigure{figures/346715a5a3080bb8.pdf}{Modèle de correction · les proportions sont indicatives sauf mention contraire.}
\end{corrige}\end{blocLivre}
\begin{blocLivre}{cor-19-15}
\begin{corrige}{15}{Une idée à discuter}
Le caractère agréable ne supprime pas le risque lié à une exposition trop intense ou trop longue. Le niveau reçu, la durée, la distance et les protections comptent. Une gêne ou des signes inhabituels doivent conduire à arrêter l’exposition et à en parler à un adulte.\par
\end{corrige}\end{blocLivre}
\pcPage{Activité et défi}{Chapitre 19 · Éléments de réponse}
\phantomsection\label{solutions-19}
\begin{pcFiche}{solutions-19}
\begin{pcEncadre}{method}{Observer et enregistrer une vibration}
B présente davantage d’oscillations que A pendant la même durée : sa fréquence est plus élevée. C conserve le rythme de A mais avec une amplitude plus grande dans ce modèle. Hauteur et niveau sonore ne sont donc pas la même propriété. La comparaison suppose les mêmes échelles et les mêmes réglages d’acquisition ; l’activité ne demande pas de produire des sons forts.\par
\end{pcEncadre}
\pcRetour{ch-19-activite}{Revenir à l’activité}
\phantomsection\label{defi-cor-19}
\begin{pcEncadre}{definition}{Le mur repéré par un écho}
Le son parcourt environ 340 \(\times\) 0,6 = 204 m au total ; le mur est à environ 102 m. Les deux trajets doivent être comptés. La température, la précision de mesure, l’emplacement du récepteur et d’autres réflexions peuvent produire des écarts. Le résultat est une estimation dans les conditions du document.\par
\end{pcEncadre}
\pcFigure{figures/9b1c11bdeee83145.pdf}{Construction ou modèle de correction.}
\pcRetour{defi-ex-19}{Revenir au défi}
\end{pcFiche}
\clearpage\corrigesChapitre{20}{Un signal pour transmettre une information}
\begin{blocLivre}{cor-20-1}
\begin{corrige}{1}{Nommer les rôles}
La lampe est l’émetteur et le capteur le récepteur. Le signal est lumineux. L’information dépend du code attribué au clignotement.\par
\end{corrige}\end{blocLivre}
\begin{blocLivre}{cor-20-2}
\begin{corrige}{2}{Signal ou information ?}
Le signal est sonore. L’information habituelle est qu’une personne appelle à la porte, selon le contexte. Le son et son interprétation sont deux éléments distincts.\par
\end{corrige}\end{blocLivre}
\begin{blocLivre}{cor-20-3}
\begin{corrige}{3}{Deux valeurs possibles}
On peut convenir que 1 signifie allumée et 0 éteinte pendant un intervalle. L’état 0 fait partie du message si le code et les durées sont partagés.\par
\end{corrige}\end{blocLivre}
\begin{blocLivre}{cor-20-4}
\begin{corrige}{4}{Quatre combinaisons}
00, 01, 10 et 11 : quatre groupes différents. On peut leur attribuer quatre symboles distincts, à condition de connaître le découpage en groupes.\par
\end{corrige}\end{blocLivre}
\begin{blocLivre}{cor-20-5}
\begin{corrige}{5}{Lire les valeurs}
La suite est 10110010. Elle contient quatre groupes de deux : 10 | 11 | 00 | 10. La lecture respecte les intervalles et leur ordre.\par
\begin{pcEncadre}{method}{Pour comprendre}
Lis les intervalles du début à la fin, puis regroupe deux valeurs à la fois. Les séparations des groupes font partie de la règle de décodage.\par
\end{pcEncadre}
\pcFigure{figures/94bffc6bdd8aed1c.pdf}{Modèle de correction · les proportions sont indicatives sauf mention contraire.}
\end{corrige}\end{blocLivre}
\begin{blocLivre}{cor-20-6}
\begin{corrige}{6}{Décoder un message}
Le message est ABC. Les séparations indiquent les groupes à traduire et évitent un décalage de lecture.\par
\end{corrige}\end{blocLivre}
\begin{blocLivre}{cor-20-7}
\begin{corrige}{7}{Coder un autre message}
DAB donne 00 | 10 | 01, soit six intervalles. La longueur du message doit être connue ou sa fin signalée par une règle convenue.\par
\begin{pcEncadre}{method}{Pour comprendre}
Remplace chaque lettre par son groupe : D devient 00, A devient 10 et B devient 01. Trois lettres occupent six intervalles.\par
\end{pcEncadre}
\pcFigure{figures/dbf04db4052baf24.pdf}{Modèle de correction · les proportions sont indicatives sauf mention contraire.}
\end{corrige}\end{blocLivre}
\begin{blocLivre}{cor-20-8}
\begin{corrige}{8}{Calculer la durée}
8 \(\times\) 0,5 = 4 s. Une durée différente d’intervalle ou l’ajout de pauses de séparation modifierait la durée totale.\par
\end{corrige}\end{blocLivre}
\begin{blocLivre}{cor-20-9}
\begin{corrige}{9}{Un capteur de température}
Il mesure une température. Une valeur affichée sans unité ne permet pas de savoir précisément ce qui est annoncé ; l’unité et les conditions de mesure restent nécessaires.\par
\end{corrige}\end{blocLivre}
\begin{blocLivre}{cor-20-10}
\begin{corrige}{10}{Un délai lumineux fourni}
600 000 \(\div\) 300 000 = 2 s environ. La durée est courte mais non nulle : la propagation n’est pas instantanée.\par
\end{corrige}\end{blocLivre}
\begin{blocLivre}{cor-20-11}
\begin{corrige}{11}{Un code non partagé}
L’interprétation dépend du tableau de code. Les groupes doivent partager les mêmes associations, le découpage et les durées avant la transmission pour retrouver la même information.\par
\end{corrige}\end{blocLivre}
\begin{blocLivre}{cor-20-12}
\begin{corrige}{12}{La lumière du Soleil}
150 000 000 \(\div\) 300 000 = 500 s, soit 8 min 20 s. Ce calcul emploie une distance moyenne et une vitesse arrondie ; le résultat est un ordre de grandeur.\par
\pcFigure{figures/fa85bf8daf0c26fc.pdf}{Modèle de correction · les proportions sont indicatives sauf mention contraire.}
\end{corrige}\end{blocLivre}
\begin{blocLivre}{cor-20-13}
\begin{corrige}{13}{Un décalage de début}
Le regroupement peut être décalé et attribuer de mauvais couples aux symboles. Convenir d’un signal de départ distinct ou d’un compte à rebours et d’une longueur de message connue permet de synchroniser la lecture.\par
\end{corrige}\end{blocLivre}
\begin{blocLivre}{cor-20-14}
\begin{corrige}{14}{Un signal masqué}
Non. Il faut d’abord que les deux états physiques soient détectables. Améliorer le contraste et les conditions de réception avec le matériel adapté, sans augmenter dangereusement l’intensité vers les yeux.\par
\end{corrige}\end{blocLivre}
\begin{blocLivre}{cor-20-15}
\begin{corrige}{15}{Deux sens pour le même voyant}
Non. Le signal physique peut être semblable tandis que l’information change avec le code et le contexte de l’appareil. Lire la légende ou la notice permet d’interpréter le voyant correctement.\par
\end{corrige}\end{blocLivre}
\pcPage{Activité et défi}{Chapitre 20 · Éléments de réponse}
\phantomsection\label{solutions-20}
\begin{pcFiche}{solutions-20}
\begin{pcEncadre}{method}{Transmettre quatre symboles}
Le message est ACDA. Les huit intervalles de 1 s durent 8 s. Sans synchronisation de départ ou sans durée commune, le récepteur peut décaler le regroupement. Sans le tableau partagé, les mêmes valeurs ne donnent pas nécessairement le même message. La longueur connue aide aussi à distinguer les zéros de fin et une transmission terminée.\par
\end{pcEncadre}
\pcRetour{ch-20-activite}{Revenir à l’activité}
\phantomsection\label{defi-cor-20}
\begin{pcEncadre}{definition}{Un message que l’on peut vérifier}
BACA donne 01 | 10 | 11 | 10, soit 01101110. Huit valeurs de 0,5 s durent 4 s. Le récepteur peut relire ou retransmettre sa suite au premier groupe pour comparaison. Cette confirmation aide à repérer des erreurs, mais une même erreur répétée reste possible ; le code, la synchronisation et les conditions de réception doivent rester contrôlés.\par
\end{pcEncadre}
\pcFigure{figures/6c1555715472135a.pdf}{Construction ou modèle de correction.}
\pcRetour{defi-ex-20}{Revenir au défi}
\end{pcFiche}
\clearpage\gdef\otChapNum{0}\pagestyle{otFront}
\pcPage{Corrigé du bilan A}{Matière et Système solaire}
\phantomsection\label{cor-bilan-A}
Barème indicatif : 2 points pour chaque réponse correcte et justifiée. Pour un calcul, valoriser la méthode et l’unité ; pour une explication, les idées exactes et reliées.\par
\begin{blocLivre}{cor-mission-A-1}
\pcSolutionMission{A}{1}{Les passages de l’eau}
Les passages sont la fusion et la liquéfaction. On obtient 40 g d’eau après fusion, si l’on considère la même eau sans perte. Le changement d’état conserve la masse dans le système fermé.\par
\pcFigure{figures/da279b44dcb5e183.pdf}{Modèle de correction.}
\end{blocLivre}
\begin{blocLivre}{cor-mission-A-2}
\pcSolutionMission{A}{2}{Le contenu d’un flacon}
Le mélange est hétérogène car le sable est visible comme phase distincte. Filtrer pour retenir le sable ; évaporer l’eau du filtrat pour récupérer le sel non volatil. Le filtrat contient encore le sel dissous.\par
\pcFigure{figures/41ce2a3b8f6f60c1.pdf}{Modèle de correction.}
\end{blocLivre}
\begin{blocLivre}{cor-mission-A-3}
\pcSolutionMission{A}{3}{Deux volumes à lire}
Le solide occupe 40 - 26 = 14 mL, soit 14 cm\textsuperscript{3}. Il doit être complètement immergé, non absorbant, insoluble et sans réaction avec l’eau ; éviter les bulles et tout débordement.\par
\pcFigure{figures/8be5ac2b5ba2901b.pdf}{Modèle de correction.}
\end{blocLivre}
\begin{blocLivre}{cor-mission-A-4}
\pcSolutionMission{A}{4}{Une conclusion de chimie}
Le test indique la présence de dioxyde de carbone dans les conditions fournies. D’autres gaz peuvent être présents dans l’échantillon ; une présence ne démontre pas une pureté.\par
\pcFigure{figures/b594f708d644c622.pdf}{Modèle de correction.}
\end{blocLivre}
\begin{blocLivre}{cor-mission-A-5}
\pcSolutionMission{A}{5}{Une place dans l’espace}
Soleil : étoile ; Terre : planète ; Lune : satellite naturel de la Terre. La Terre se place à 150 \(\div\) 50 = 3 cm du Soleil. Les diamètres ne sont pas automatiquement à la même échelle.\par
\begin{pcEncadre}{method}{Utiliser le résultat}
Repère la connaissance ou la méthode à reprendre. Reviens au chapitre concerné, puis refais une mission sans regarder la solution avant de comparer.\par
\end{pcEncadre}
\end{blocLivre}
\pcPage{Corrigé du bilan B}{Mouvement et interactions}
\phantomsection\label{cor-bilan-B}
Barème indicatif : 2 points pour chaque réponse correcte et justifiée. Pour un calcul, valoriser la méthode et l’unité ; pour une explication, les idées exactes et reliées.\par
\begin{blocLivre}{cor-mission-B-1}
\pcSolutionMission{B}{1}{Décrire avec un repère}
Il est au repos par rapport au siège et en mouvement par rapport au quai. Les descriptions portent sur des repères différents ; le mouvement doit toujours être associé à un repère.\par
\pcFigure{figures/5f58c69c5265b1b2.pdf}{Modèle de correction.}
\end{blocLivre}
\begin{blocLivre}{cor-mission-B-2}
\pcSolutionMission{B}{2}{Lire des positions}
A présente des espacements constants. B montre des espacements croissants : davantage de distance est parcourue pendant chaque intervalle de même durée.\par
\pcFigure{figures/b9d2213e035c1cd3.pdf}{Modèle de correction.}
\end{blocLivre}
\begin{blocLivre}{cor-mission-B-3}
\pcSolutionMission{B}{3}{Un déplacement chiffré}
90 \(\div\) 30 = 3 m/s. 3 \(\times\) 3,6 = 10,8 km/h. Le calcul utilise la distance réellement parcourue et la durée du même trajet.\par
\end{blocLivre}
\begin{blocLivre}{cor-mission-B-4}
\pcSolutionMission{B}{4}{Un aller avec une pause}
30 min = 0,5 h ; 2 \(\div\) 0,5 = 4 km/h. La moyenne ne prouve pas une vitesse constante ; une pause est même annoncée.\par
\end{blocLivre}
\begin{blocLivre}{cor-mission-B-5}
\pcSolutionMission{B}{5}{Un livre au repos}
La table interagit par contact et la Terre à distance. Le repos du livre ne signifie pas absence d’interactions.\par
\pcFigure{figures/852f07059882b685.pdf}{Modèle de correction.}
\begin{pcEncadre}{method}{Utiliser le résultat}
Repère la connaissance ou la méthode à reprendre. Reviens au chapitre concerné, puis refais une mission sans regarder la solution avant de comparer.\par
\end{pcEncadre}
\end{blocLivre}
\pcPage{Corrigé du bilan C}{Énergie et circuits}
\phantomsection\label{cor-bilan-C}
Barème indicatif : 2 points pour chaque réponse correcte et justifiée. Pour un calcul, valoriser la méthode et l’unité ; pour une explication, les idées exactes et reliées.\par
\begin{blocLivre}{cor-mission-C-1}
\pcSolutionMission{C}{1}{Une chaîne à compléter}
La pile possède une réserve d’énergie chimique ; la lampe reçoit un transfert électrique et émet de la lumière. Un transfert thermique vers l’environnement doit aussi être considéré.\par
\pcFigure{figures/6d0e39ac014c4fa7.pdf}{Modèle de correction.}
\end{blocLivre}
\begin{blocLivre}{cor-mission-C-2}
\pcSolutionMission{C}{2}{Une branche ouverte}
L\textsubscript{2} conserve une boucle fermée ; L₁ est dans la branche ouverte. Une commande générale se place dans une portion commune à tous les chemins, par exemple sur une liaison principale du générateur.\par
\pcFigure{figures/0e679a159d2530f6.pdf}{Modèle de correction.}
\end{blocLivre}
\begin{blocLivre}{cor-mission-C-3}
\pcSolutionMission{C}{3}{Un test à interpréter}
La peinture peut empêcher le contact électrique ; un mauvais contact ou un défaut du montage est aussi possible. Vérifier avec un conducteur témoin dans un montage validé.\par
\end{blocLivre}
\begin{blocLivre}{cor-mission-C-4}
\pcSolutionMission{C}{4}{Un bilan de 120 J}
120 - 45 = 75 J sont transférés autrement. L’énergie n’a pas disparu ; elle peut notamment être transférée thermiquement à l’environnement selon l’appareil.\par
\pcFigure{figures/1b3b7cd46e676e3a.pdf}{Modèle de correction.}
\end{blocLivre}
\begin{blocLivre}{cor-mission-C-5}
\pcSolutionMission{C}{5}{Une économie documentée}
L’économie est 60 kWh sur l’année. Il reste à considérer fabrication, durée de vie, transport, réparation et fin de vie, ainsi que la validité du service comparé.\par
\begin{pcEncadre}{method}{Utiliser le résultat}
Repère la connaissance ou la méthode à reprendre. Reviens au chapitre concerné, puis refais une mission sans regarder la solution avant de comparer.\par
\end{pcEncadre}
\end{blocLivre}
\pcPage{Corrigé du bilan D}{Lumière, sons et information}
\phantomsection\label{cor-bilan-D}
Barème indicatif : 2 points pour chaque réponse correcte et justifiée. Pour un calcul, valoriser la méthode et l’unité ; pour une explication, les idées exactes et reliées.\par
\begin{blocLivre}{cor-mission-D-1}
\pcSolutionMission{D}{1}{Expliquer la vision}
La source éclaire l’objet ; l’objet diffuse une partie de la lumière vers l’œil. La dernière flèche va de l’objet vers l’œil, car la lumière doit y entrer.\par
\pcFigure{figures/cd80bfc6ec09d32d.pdf}{Modèle de correction.}
\end{blocLivre}
\begin{blocLivre}{cor-mission-D-2}
\pcSolutionMission{D}{2}{Une ombre à prévoir}
On utilise le modèle dans un milieu transparent homogène. L’ombre portée grandit dans cette géométrie à petite source, avec le même objet entre source et écran.\par
\pcFigure{figures/fb6a0bd783e84e52.pdf}{Modèle de correction.}
\end{blocLivre}
\begin{blocLivre}{cor-mission-D-3}
\pcSolutionMission{D}{3}{La distance d’un écho}
Le trajet total vaut 340 \(\times\) 0,5 = 170 m. L’obstacle se situe à 85 m, car le son a effectué l’aller et le retour.\par
\pcFigure{figures/70c340791eaa0ac1.pdf}{Modèle de correction.}
\end{blocLivre}
\begin{blocLivre}{cor-mission-D-4}
\pcSolutionMission{D}{4}{Un code partagé}
Le message est CAB. Six valeurs de 1 s durent 6 s, hors préparation. Le code et le repérage des intervalles doivent être partagés.\par
\end{blocLivre}
\begin{blocLivre}{cor-mission-D-5}
\pcSolutionMission{D}{5}{Écouter en préservant son audition}
Le risque dépend notamment du niveau et de la durée, pas seulement du plaisir d’écoute. Réduire le niveau, limiter la durée, s’éloigner des sources fortes et utiliser les protections adaptées sont des moyens de prévention.\par
\begin{pcEncadre}{method}{Utiliser le résultat}
Repère la connaissance ou la méthode à reprendre. Reviens au chapitre concerné, puis refais une mission sans regarder la solution avant de comparer.\par
\end{pcEncadre}
\end{blocLivre}
\clearpage\gdef\otChapNum{0}\gdef\otAccent{otNavy}\pagestyle{otFront}\phantomsection\label{annexes-debut}
\pcPage{Mesurer et contrôler un résultat}{FICHE MÉTHODE 1}
\phantomsection\label{methode-mesure}
\pdfbookmark[0]{Mesurer et contrôler un résultat}{bm-methode-mesure}
\begin{pcFiche}{annexe-methode-mesure}
\begin{pcEncadre}{definition}{1 · CHOISIR}
Identifier la grandeur, l’instrument, sa capacité et sa graduation. Une balance mesure une masse ; une éprouvette, un volume ; un thermomètre, une température.\par
\end{pcEncadre}
\begin{pcEncadre}{method}{2 · RELEVER}
Effectuer la tare si nécessaire, placer l’œil correctement, attendre la stabilisation et noter toutes les unités. Garder les résultats réellement mesurés.\par
\end{pcEncadre}
\begin{pcEncadre}{definition}{3 · COMPARER}
Répéter la mesure dans des conditions comparables. Deux valeurs proches peuvent différer à cause de la précision des instruments ou du protocole.\par
\end{pcEncadre}
\begin{pcEncadre}{method}{4 · CONCLURE}
Écrire une phrase adaptée à la précision disponible. Ne pas donner une longue suite de décimales qui ne correspond pas à la mesure. Un calcul ne rend pas une graduation plus fine.\par
\end{pcEncadre}
Exemple : si le pas de lecture est 2 mL, un résultat calculé comme 24,000 000 mL ne décrit pas une lecture aussi précise.\par
\end{pcFiche}
\pcPage{Construire et lire un graphique}{FICHE MÉTHODE 2}
\phantomsection\label{methode-graphique}
\begin{pcFiche}{annexe-methode-graphique}
\begin{questions}
\item Identifier la grandeur placée sur chaque axe et son unité.
\item Choisir des graduations régulières et des échelles lisibles.
\item Placer les points à partir du tableau ; ne pas inverser les coordonnées.
\item Décrire la tendance avant de proposer une interprétation.
\end{questions}
\pcFigure{figures/1448bdea1799d68c.pdf}{Exemple : 0, 3, 6, 9 et 12 m aux dates 0, 1, 2, 3 et 4 s.}
\begin{pcEncadre}{method}{UNE LECTURE ARGUMENTÉE}
À 2 s, la distance lue est 6 m. La droite passe par l’origine et indique ici une distance proportionnelle à la durée. Le modèle correspond à 3 m parcourus par seconde.\par
\end{pcEncadre}
\end{pcFiche}
\pcPage{Calculer sans perdre les unités}{FICHE MÉTHODE 3}
\phantomsection\label{methode-calcul}
\begin{pcFiche}{annexe-methode-calcul}
\begin{pcEncadre}{definition}{UN EXEMPLE COMPLET}
Question : vitesse moyenne pour 240 m en 2 min.\par Données : distance 240 m ; durée 2 min = 120 s.\par Méthode : vitesse moyenne = distance \(\div\) durée.\par Calcul : 240 \(\div\) 120 = 2.\par Conclusion : la vitesse moyenne est de 2 m/s.\par
\end{pcEncadre}
\begin{questions}
\item Choisir les unités avant le calcul.
\item Garder une durée et une distance qui concernent le même trajet.
\item Contrôler l’ordre de grandeur et l’unité du résultat.
\item Pour un trajet par étapes : additionner les distances et les durées avant de calculer la moyenne.
\end{questions}
\begin{pcEncadre}{method}{UN CONTRÔLE RAPIDE}
À 2 m/s pendant 120 s, on retrouve 240 m. Ce retour aux données aide à repérer une inversion de quotient ou une erreur de conversion.\par
\end{pcEncadre}
\end{pcFiche}
\pcPage{Grandeurs et unités}{AIDE-MÉMOIRE}
\phantomsection\label{unites}
\pdfbookmark[0]{Grandeurs et unités}{bm-unites}
\begin{pcFiche}{annexe-unites}
\begin{center}
{\renewcommand{\arraystretch}{1.35}\fontsize{11}{14}\selectfont
\setlength{\tabcolsep}{2mm}\begin{tabular}{@{}>{\raggedright\arraybackslash}p{41.20mm}>{\raggedright\arraybackslash}p{41.20mm}>{\raggedright\arraybackslash}p{69.59mm}@{}}\toprule
\otTete{Grandeur} & \otTete{Unité utilisée} & \otTete{Instrument ou document}\\
\midrule
Masse & kg, g & Balance\\
Volume & m\textsuperscript{3}, L, mL & Éprouvette, dimensions\\
Température & °C & Thermomètre\\
Durée & s, min, h & Chronomètre, horloge\\
Distance & m, km & Règle, mètre, carte\\
Vitesse moyenne & m/s, km/h & Distance et durée\\
Énergie & J, Wh, kWh & Bilan ou mesure d’énergie\\
Puissance & W & Document ou mesure adaptée\\
Fréquence & Hz & Acquisition ou document\\
Niveau sonore & dB & Sonomètre ou document\\
\bottomrule\end{tabular}}\end{center}
\begin{pcEncadre}{method}{NE PAS CONFONDRE}
Masse et volume ; durée et date ; vitesse et distance ; énergie et puissance ; fréquence et niveau sonore. Une unité permet souvent de retrouver la grandeur avant tout calcul.\par
\end{pcEncadre}
\end{pcFiche}
\pcPage{Les équivalences à retrouver}{AIDE-MÉMOIRE}
\phantomsection\label{conversions}
\begin{pcFiche}{annexe-conversions}
\begin{center}
{\renewcommand{\arraystretch}{1.35}\fontsize{11}{14}\selectfont
\setlength{\tabcolsep}{2mm}\begin{tabular}{@{}>{\raggedright\arraybackslash}p{50.36mm}>{\raggedright\arraybackslash}p{51.28mm}>{\raggedright\arraybackslash}p{50.36mm}@{}}\toprule
\otTete{Masses} & \otTete{Volumes} & \otTete{Durées}\\
\midrule
1 kg = 1 000 g & 1 L = 1 000 mL & 1 h = 60 min\\
1 g = 1 000 mg & 1 L = 100 cL & 1 min = 60 s\\
0,5 kg = 500 g & 1 L = 10 dL & 1 h = 3 600 s\\
250 g = 0,25 kg & 1 L = 1 dm\textsuperscript{3} & 0,5 h = 30 min\\
 & 1 mL = 1 cm\textsuperscript{3} & 1 min 30 s = 90 s\\
\bottomrule\end{tabular}}\end{center}
\begin{pcEncadre}{definition}{VITESSES}
1 m/s = 3,6 km/h. Pour convertir de m/s en km/h, multiplier par 3,6 ; pour le trajet inverse, diviser par 3,6.\par
\end{pcEncadre}
\begin{pcEncadre}{method}{ÉNERGIES}
1 kWh = 1 000 Wh. Le symbole kWh désigne une énergie ; kW désigne une puissance. Le repère 1 mL d’eau \(\approx\) 1 g à température ambiante n’est pas une conversion universelle entre volume et masse.\par
\end{pcEncadre}
\end{pcFiche}
\pcPage{Faire le point sur tes acquis}{BILAN PERSONNEL}
\phantomsection\label{competences}
\pdfbookmark[0]{Faire le point sur tes acquis}{bm-competences}
\begin{pcFiche}{annexe-competences}
Reprends les exercices et les bilans après correction. Pour chaque compétence, retrouve un exemple que tu sais expliquer sans regarder la solution.\par
\begin{pcEncadre}{definition}{DÉCRIRE}
Nommer les objets, les substances, les états, les composants et les conditions.\par
\end{pcEncadre}
\begin{pcEncadre}{definition}{EXPÉRIMENTER}
Proposer une hypothèse, choisir les témoins et ne changer qu’un facteur à la fois.\par
\end{pcEncadre}
\begin{pcEncadre}{definition}{MESURER}
Choisir un instrument, lire ses graduations et exprimer une valeur avec son unité.\par
\end{pcEncadre}
\begin{pcEncadre}{definition}{MODÉLISER}
Utiliser des particules, un circuit, une chaîne ou des rayons avec une légende et des limites.\par
\end{pcEncadre}
\begin{pcEncadre}{definition}{ARGUMENTER}
Relier une observation à une connaissance, faire un calcul cohérent et nuancer la conclusion.\par
\end{pcEncadre}
\end{pcFiche}
\pcPage{Lexique scientifique}{80 mots pour comprendre et expliquer}
\phantomsection\label{lexique}
\pdfbookmark[0]{Lexique scientifique}{bm-lexique}
\begin{pcLexique}{Action mécanique}{11}
Action pouvant déformer un objet ou modifier son mouvement.\par
\end{pcLexique}
\begin{pcLexique}{Air}{4}
Mélange de gaz ; dans l’air sec, environ 78 \% de diazote et 21 \% de dioxygène en volume.\par
\end{pcLexique}
\begin{pcLexique}{Boucle}{13}
Chemin fermé passant par le générateur et des récepteurs dans un circuit.\par
\end{pcLexique}
\begin{pcLexique}{Changement d’état}{3}
Transformation physique qui conserve la nature de la substance et modifie son état.\par
\end{pcLexique}
\begin{pcLexique}{Circuit ouvert}{13}
Circuit dont un trajet est interrompu ; une boucle ouverte ne permet pas le courant.\par
\end{pcLexique}
\begin{pcLexique}{Code}{20}
Règles partagées qui associent des signaux ou des suites de valeurs à une information.\par
\end{pcLexique}
\begin{pcLexique}{Conducteur}{14}
Matériau permettant le passage du courant dans les conditions considérées.\par
\end{pcLexique}
\begin{pcLexique}{Conservation}{3}
Maintien d’une grandeur totale dans les conditions et le système précisés.\par
\end{pcLexique}
\begin{pcLexique}{Convertisseur}{12}
Dispositif qui reçoit de l’énergie et la transfère sous d’autres formes ou modes.\par
\end{pcLexique}
\begin{pcLexique}{Corps pur}{4}
Échantillon constitué d’une seule espèce chimique dans le modèle étudié.\par
\end{pcLexique}
\begin{pcLexique}{Court-circuit}{14}
Liaison de très faible résistance qui contourne un dipôle ; celle du générateur présente un risque.\par
\end{pcLexique}
\begin{pcLexique}{Décantation}{6}
Séparation de phases d’un mélange hétérogène après repos.\par
\end{pcLexique}
\begin{pcLexique}{Dérivation}{15}
Association avec des branches reliées aux mêmes deux points de connexion.\par
\end{pcLexique}
\begin{pcLexique}{Diffusion lumineuse}{17}
Renvoi d’une partie de la lumière reçue dans plusieurs directions.\par
\end{pcLexique}
\begin{pcLexique}{Dipôle}{13}
Composant électrique possédant deux bornes de connexion.\par
\end{pcLexique}
\begin{pcLexique}{Dissolution}{5}
Dispersion d’un soluté dans un solvant donnant une solution, sans être une fusion.\par
\end{pcLexique}
\begin{pcLexique}{Distillation}{6}
Séparation avec vaporisation puis liquéfaction d’un constituant recueilli comme distillat.\par
\end{pcLexique}
\begin{pcLexique}{Durée}{10}
Temps écoulé entre deux dates ; elle peut s’exprimer en secondes.\par
\end{pcLexique}
\begin{pcLexique}{Écho}{19}
Signal sonore reçu après réflexion ; sa durée peut inclure un aller-retour.\par
\end{pcLexique}
\begin{pcLexique}{Émetteur}{20}
Objet ou dispositif qui produit le signal transmis.\par
\end{pcLexique}
\begin{pcLexique}{Énergie}{12}
Grandeur transférée et convertie dans les systèmes ; son unité SI est le joule.\par
\end{pcLexique}
\begin{pcLexique}{Espèce chimique}{4}
Catégorie de constituants identiques du point de vue chimique.\par
\end{pcLexique}
\begin{pcLexique}{Étoile}{8}
Astre qui produit la lumière qu’il émet, comme le Soleil.\par
\end{pcLexique}
\begin{pcLexique}{Évaporation}{3}
Vaporisation à la surface d’un liquide, possible sans ébullition.\par
\end{pcLexique}
\begin{pcLexique}{Filtrat}{6}
Liquide qui traverse le filtre pendant une filtration.\par
\end{pcLexique}
\begin{pcLexique}{Filtration}{6}
Séparation par une barrière poreuse retenant certaines particules.\par
\end{pcLexique}
\begin{pcLexique}{Fréquence}{19}
Nombre de vibrations par seconde pour une vibration périodique ; unité : hertz.\par
\end{pcLexique}
\begin{pcLexique}{Fusion}{3}
Passage de l’état solide à l’état liquide.\par
\end{pcLexique}
\begin{pcLexique}{Générateur}{13}
Dipôle qui permet de fournir de l’énergie au circuit, par exemple une pile.\par
\end{pcLexique}
\begin{pcLexique}{Hétérogène}{4}
Se dit d’un mélange dont on distingue plusieurs phases ou constituants.\par
\end{pcLexique}
\begin{pcLexique}{Homogène}{4}
Se dit d’un mélange présentant une seule phase visible à l’échelle considérée.\par
\end{pcLexique}
\begin{pcLexique}{Hypothèse}{1}
Proposition d’explication ou de réponse à mettre à l’épreuve.\par
\end{pcLexique}
\begin{pcLexique}{Information}{20}
Sens attribué à un signal selon un code et un contexte.\par
\end{pcLexique}
\begin{pcLexique}{Interaction}{11}
Action mutuelle entre deux objets, de contact ou à distance.\par
\end{pcLexique}
\begin{pcLexique}{Isolant}{14}
Matériau limitant fortement le passage du courant dans les conditions considérées.\par
\end{pcLexique}
\begin{pcLexique}{Liquéfaction}{3}
Passage de l’état gazeux à l’état liquide.\par
\end{pcLexique}
\begin{pcLexique}{Masse}{2}
Grandeur mesurée avec une balance ; unité SI : kilogramme.\par
\end{pcLexique}
\begin{pcLexique}{Mélange}{4}
Échantillon contenant plusieurs espèces chimiques.\par
\end{pcLexique}
\begin{pcLexique}{Ménisque}{2}
Courbure de la surface d’un liquide près des parois ; pour l’eau usuelle, lire le bas.\par
\end{pcLexique}
\begin{pcLexique}{Miscibilité}{5}
Aptitude de deux liquides à former une phase homogène dans les conditions étudiées.\par
\end{pcLexique}
\begin{pcLexique}{Modèle}{1}
Représentation simplifiée choisie pour expliquer ou prévoir certains phénomènes.\par
\end{pcLexique}
\begin{pcLexique}{Niveau sonore}{19}
Grandeur liée au son reçu, exprimée en décibels ; différente de la hauteur.\par
\end{pcLexique}
\begin{pcLexique}{Ombre portée}{18}
Tache obtenue sur un écran à l’intersection avec une zone d’ombre.\par
\end{pcLexique}
\begin{pcLexique}{Ombre propre}{18}
Partie de la surface d’un objet non directement éclairée par la source considérée.\par
\end{pcLexique}
\begin{pcLexique}{Opaque}{17}
Se dit d’un matériau qui ne transmet pas perceptiblement la lumière visible dans les conditions étudiées.\par
\end{pcLexique}
\begin{pcLexique}{Palier}{3}
Portion de courbe pendant laquelle la grandeur reste presque constante.\par
\end{pcLexique}
\begin{pcLexique}{Planète}{8}
Astre en révolution autour d’une étoile ; il diffuse la lumière reçue.\par
\end{pcLexique}
\begin{pcLexique}{Produit}{7}
Espèce formée au cours d’une transformation chimique.\par
\end{pcLexique}
\begin{pcLexique}{Propagation}{19}
Transmission d’une perturbation ou d’un signal depuis une source.\par
\end{pcLexique}
\begin{pcLexique}{Puissance}{16}
Rythme du transfert d’énergie ; unité SI : watt.\par
\end{pcLexique}
\begin{pcLexique}{Rayon lumineux}{18}
Trait rectiligne fléché représentant un trajet de lumière dans le modèle.\par
\end{pcLexique}
\begin{pcLexique}{Réactif}{7}
Espèce consommée au cours d’une transformation chimique.\par
\end{pcLexique}
\begin{pcLexique}{Récepteur}{13}
Dispositif qui reçoit un signal ; en électricité, dipôle qui reçoit l’énergie transférée.\par
\end{pcLexique}
\begin{pcLexique}{Renouvelable}{16}
Se dit d’une ressource disponible sous forme de flux ou renouvelée à l’échelle humaine.\par
\end{pcLexique}
\begin{pcLexique}{Repère}{9}
Référence choisie pour décrire des positions et un mouvement.\par
\end{pcLexique}
\begin{pcLexique}{Résidu}{6}
Matière retenue sur un filtre ou restant après une opération de séparation.\par
\end{pcLexique}
\begin{pcLexique}{Révolution}{8}
Mouvement orbital d’un astre autour d’un autre.\par
\end{pcLexique}
\begin{pcLexique}{Rotation}{8}
Mouvement d’un objet sur lui-même autour d’un axe.\par
\end{pcLexique}
\begin{pcLexique}{Satellite naturel}{8}
Astre en révolution autour d’une planète ; la Lune est celui de la Terre.\par
\end{pcLexique}
\begin{pcLexique}{Saturation}{5}
État d’une solution qui ne dissout plus de soluté supplémentaire dans les conditions données.\par
\end{pcLexique}
\begin{pcLexique}{Sens conventionnel du courant}{13}
Sens de description du courant dans le circuit extérieur : du + vers le - de la pile.\par
\end{pcLexique}
\begin{pcLexique}{Série}{15}
Association de dipôles successifs dans une même branche.\par
\end{pcLexique}
\begin{pcLexique}{Signal}{20}
Variation physique émise, transmise et reçue ; à distinguer de son interprétation.\par
\end{pcLexique}
\begin{pcLexique}{Solidification}{3}
Passage de l’état liquide à l’état solide.\par
\end{pcLexique}
\begin{pcLexique}{Solubilité}{5}
Limite de quantité dissoute pour un soluté, un solvant et des conditions précisés.\par
\end{pcLexique}
\begin{pcLexique}{Soluté}{5}
Substance dissoute dans un solvant.\par
\end{pcLexique}
\begin{pcLexique}{Solution}{5}
Mélange homogène contenant un solvant et un ou plusieurs solutés.\par
\end{pcLexique}
\begin{pcLexique}{Solvant}{5}
Liquide dans lequel le soluté se dissout dans les situations du livre.\par
\end{pcLexique}
\begin{pcLexique}{Source primaire}{17}
Objet qui produit la lumière qu’il émet.\par
\end{pcLexique}
\begin{pcLexique}{Système}{11}
Objet ou ensemble choisi comme sujet d’étude, avec des limites précisées.\par
\end{pcLexique}
\begin{pcLexique}{Tare}{2}
Remise à zéro d’une balance avec le récipient afin de mesurer la masse ajoutée.\par
\end{pcLexique}
\begin{pcLexique}{Température}{3}
Grandeur mesurée avec un thermomètre ; elle est ici exprimée en degré Celsius.\par
\end{pcLexique}
\begin{pcLexique}{Trajectoire}{9}
Ensemble des positions successives d’un point dans le repère choisi.\par
\end{pcLexique}
\begin{pcLexique}{Transformation chimique}{7}
Transformation qui consomme des espèces et en forme de nouvelles.\par
\end{pcLexique}
\begin{pcLexique}{Translucide}{17}
Se dit d’un matériau transmettant de la lumière sans permettre une vision nette à travers.\par
\end{pcLexique}
\begin{pcLexique}{Transparent}{17}
Se dit d’un matériau permettant une vision nette à travers dans les conditions étudiées.\par
\end{pcLexique}
\begin{pcLexique}{Unité}{2}
Référence utilisée pour exprimer la mesure d’une grandeur.\par
\end{pcLexique}
\begin{pcLexique}{Vaporisation}{3}
Passage de l’état liquide à l’état gazeux, par évaporation ou ébullition.\par
\end{pcLexique}
\begin{pcLexique}{Vitesse moyenne}{10}
Distance totale parcourue divisée par la durée correspondante.\par
\end{pcLexique}
\begin{pcLexique}{Volume}{2}
Grandeur qui mesure l’espace occupé ; unité SI : mètre cube.\par
\end{pcLexique}
\pcPage{Programme et progression}{Les choix du manuel}
\phantomsection\label{programme}
\pdfbookmark[0]{Programme et progression}{bm-programme}
Le programme national définit les attendus sur les trois années du cycle 4. Ce manuel propose une progression complète pour la 5e, construite à partir des repères de début de cycle. L’équipe pédagogique peut organiser les apprentissages autrement.\par
\begin{center}
{\renewcommand{\arraystretch}{1.35}\fontsize{11}{14}\selectfont
\setlength{\tabcolsep}{2mm}\begin{tabular}{@{}>{\raggedright\arraybackslash}p{122.17mm}>{\raggedright\arraybackslash}p{33.83mm}@{}}\toprule
\otTete{Thème du programme} & \otTete{Chapitres}\\
\midrule
Organisation et transformations de la matière & 1 à 8\\
Mouvement et interactions & 9 à 11\\
Énergie, transferts et conversions & 12 à 16\\
Signaux pour observer et communiquer & 17 à 20\\
\bottomrule\end{tabular}}\end{center}
\begin{pcEncadre}{method}{Une progression sur trois années}
Les modèles atomiques détaillés, les équations de réaction, la grandeur quotient masse volumique, la modélisation formelle des forces et les lois électriques sont approfondis aux niveaux suivants. Ici, les circuits, les interactions et les transferts d’énergie sont étudiés surtout par observations et raisonnements simples.\par
\end{pcEncadre}
Le titre « manuel complet » désigne cette progression de 5e ; il ne prétend pas traiter à lui seul tous les attendus de fin de 3e.\par
\pcPage{Références}{Sources consultées le 23 octobre 2026}
\phantomsection\label{references}
\pdfbookmark[0]{Références}{bm-references}
\begin{blocLivre}{source-1}
\pcSousTitre{1. Programme du cycle 4}
BO n° 31 du 30 juillet 2020, rubrique physique-chimie, pages 96 à 105 de la version Éduscol.\par
{\otLight\fontsize{9}{12}\selectfont\color{otMuted}\url{https://eduscol.education.fr/document/621/download}}\par\end{blocLivre}
\begin{blocLivre}{source-2}
\pcSousTitre{2. Programmes applicables en 2026-2027}
Éduscol, « Enseigner au cycle 4 » : confirmation du texte de référence pour la 5e.\par
{\otLight\fontsize{9}{12}\selectfont\color{otMuted}\url{https://eduscol.education.gouv.fr/4362/enseigner-au-cycle-4}}\par\end{blocLivre}
\begin{blocLivre}{source-3}
\pcSousTitre{3. Ressources de physique-chimie au cycle 4}
Éduscol, page actualisée en mars 2026 : programme, progressions et démarches pédagogiques.\par
{\otLight\fontsize{9}{12}\selectfont\color{otMuted}\url{https://eduscol.education.gouv.fr/5739/ressources-d-accompagnement-du-programme-de-physique-chimie-au-cycle-4}}\par\end{blocLivre}
\begin{blocLivre}{source-4}
\pcSousTitre{4. Construire une progression au cycle 4}
Éduscol, 2016 : repères de progressivité et exemples non prescriptifs.\par
{\otLight\fontsize{9}{12}\selectfont\color{otMuted}\url{https://eduscol.education.fr/document/17731/download}}\par\end{blocLivre}
\begin{blocLivre}{source-5}
\pcSousTitre{5. Le Système international d’unités}
BIPM, 9e édition, mise à jour 2026 : kilogramme, unités de volume et usage du litre.\par
{\otLight\fontsize{9}{12}\selectfont\color{otMuted}\url{https://www.bipm.org/fr/publications/si-brochure/}}\par\end{blocLivre}
\begin{blocLivre}{source-6}
\pcSousTitre{6. Lire un volume}
Collège Carnot, ENT Occitanie : position de l’œil et lecture du bas du ménisque.\par
{\otLight\fontsize{9}{12}\selectfont\color{otMuted}\url{https://carnot.mon-ent-occitanie.fr/espaces-pedagogiques/socle-commun/activites-d-accompagnement/comment-mesurer-un-volume--14816.htm}}\par\end{blocLivre}
\begin{blocLivre}{source-7}
\pcSousTitre{7. Explorer le Système solaire}
CNES : organisation du Système solaire et diversité des huit planètes.\par
{\otLight\fontsize{9}{12}\selectfont\color{otMuted}\url{https://cnes.fr/dossiers/systeme-solaire}}\par\end{blocLivre}
\begin{blocLivre}{source-8}
\pcSousTitre{8. Se repérer dans l’espace}
CNES, glossaire spatial : unités et vocabulaire astronomiques.\par
{\otLight\fontsize{9}{12}\selectfont\color{otMuted}\url{https://cnes.fr/dossiers/glossaire-spatial}}\par\end{blocLivre}
\begin{blocLivre}{source-9}
\pcSousTitre{9. Comprendre le risque électrique}
INRS : risque de contact, court-circuit et prévention. Les montages du livre utilisent uniquement le matériel scolaire de très basse tension.\par
{\otLight\fontsize{9}{12}\selectfont\color{otMuted}\url{https://www.inrs.fr/risques/electriques/ce-qu-il-faut-retenir.html}}\par\end{blocLivre}
\begin{blocLivre}{source-10}
\pcSousTitre{10. Préserver son audition}
Assurance Maladie : effets des bruits intenses et des expositions prolongées.\par
{\otLight\fontsize{9}{12}\selectfont\color{otMuted}\url{https://www.ameli.fr/assure/sante/themes/bruit-sante/effets-du-bruit-sur-la-sante}}\par\end{blocLivre}
\begin{blocLivre}{source-11}
\pcSousTitre{11. Identifier un gaz}
Éduscol : exemple d’exploitation du test du dioxyde de carbone à l’eau de chaux.\par
{\otLight\fontsize{9}{12}\selectfont\color{otMuted}\url{https://eduscol.education.fr/document/29578/download}}\par\end{blocLivre}
\begin{blocLivre}{source-12}
\pcSousTitre{12. Ressources renouvelables}
Ministère de la Transition écologique : familles de ressources et conversions énergétiques.\par
{\otLight\fontsize{9}{12}\selectfont\color{otMuted}\url{https://www.ecologie.gouv.fr/politiques-publiques/energies-renouvelables}}\par\end{blocLivre}
\begin{blocLivre}{source-13}
\pcSousTitre{13. Sons et effets du bruit}
Ressource Éduscol sur le son, les milieux de propagation et la santé auditive.\par
{\otLight\fontsize{9}{12}\selectfont\color{otMuted}\url{https://eduscol.education.fr/document/30202/download}}\par\end{blocLivre}
\begin{pcEncadre}{definition}{Créations originales}
Textes, exercices, corrigés et schémas sont des créations pour cet ouvrage. Les documents officiels servent de références de programme et de démarche. Les séries annoncées comme pédagogiques sont des exemples à exploiter, pas des mesures effectuées par l’élève.\par
\end{pcEncadre}
\end{document}
